%============================================================
%  Bd. 23 - Halbgruppen
%============================================================

\documentclass[main.tex]{subfiles}

\ifSubfilesClassLoaded{
  \BandSetupStandalone{B23}
}{
}

\title{Bd. 23 - Halbgruppen}
\author{Martin Kunze}
\date{}
\setcounter{file}{23}

\hypersetup{
  pdftitle={Bd. 23 - Halbgruppen},
  pdfauthor={Martin Kunze}
}

\begin{document}

\title{Bd. 23 - Halbgruppen}
\author{Martin Kunze}
\date{}
\hypersetup{pageanchor=false}
\maketitle
\hypersetup{pageanchor=true}
\setcounter{tocdepth}{3}
\tableofcontents
\setcounter{file}{23}
\renewcommand{\FormulaBandID}{23}

\chapter{Einführung}

Dieser Band entwickelt Halbgruppen aus einer assoziativen binären Operation.
Behandelt werden Beispiele, Teilbarkeit, Homomorphismen, Isomorphismen und
die funktorielle Konstruktion der Potenzhalbgruppe. Zusätzliche
Strukturaxiome werden in den folgenden Bänden getrennt untersucht.

\chapter{Halbgruppen}

\section{Definition der Halbgruppe}

\FormulaDefDeltaK[Begriff der Halbgruppe]{%
  \SemigroupAlg{A}{\star}%
}{SemigroupDef}{%
  \DeltaRow{Mengen}{A\dsep x\dsep y\dsep z}
  \DeltaRow{\textbf{Axiome}}{}
  \DeltaRow{Abgeschlossenheit}{%
    x,y\in A\vdash x\star y\in A}
  \DeltaRow{Assoziativität}{%
    x,y,z\in A\vdash (x\star y)\star z=x\star(y\star z)}
  \DeltaRow{\textbf{Neues Symbol}}{}
  \DeltaPrem{Halbgruppe}{\SemigroupAlg{A}{\star}}
}

\FormulaAxiomDeltaK[Zugriff auf die Abgeschlossenheit einer Halbgruppe]{%
  \SemigroupAlg{A}{\star}\dsep x,y\in A
  \vdash x\star y\in A%
}{SemigroupClosureAxiom}{%
  \DeltaRow{Mengen}{A\dsep x\dsep y}
}

\FormulaAxiomDeltaK[Zugriff auf die Assoziativität einer Halbgruppe]{%
  \SemigroupAlg{A}{\star}\dsep x,y,z\in A
  \vdash (x\star y)\star z=x\star(y\star z)%
}{SemigroupAssociativityAxiom}{%
  \DeltaRow{Mengen}{A\dsep x\dsep y\dsep z}
}

\FormulaAxiomDeltaK[Zugriff auf die binäre Operation einer Halbgruppe]{%
  \SemigroupAlg{A}{\star}\vdash
  \mathsf{BinOp}\!\left(A,\star\right)%
}{SemigroupBinaryOperationAxiom}{%
  \DeltaRow{Mengen}{A}
}

\FormulaThmDeltaK[Fünfgliedrige Umklammerung in Halbgruppen]{%
  \begin{aligned}[t]
    &\SemigroupAlg{A}{\star}\dsep a,b,c,d,e\in A\vdash{}\\[-2pt]
    &(a\star b)\star\bigl((c\star d)\star e\bigr)
      =\bigl((a\star(b\star c))\star d\bigr)\star e
  \end{aligned}%
}{SemigroupFiveFactorReassociation}{%
  \DeltaRow{Mengen}{A\dsep a\dsep b\dsep c\dsep d\dsep e}
  \DeltaRow{Funktionen}{\star}
}
\begin{tabproofwide}
  \proofstepwidestar[1]{\SemigroupAlg{A}{\star}}{\rA}
  \proofstepwidestar[2]{a,b,c,d,e\in A}{\rA}
  \proofstepwidestar[1,2]{a\star b\in A}{%
    \FormulaRefAuto{SemigroupClosureAxiom}[ax]{1,2}}
  \proofstepwidestar[1,2]{d\star e\in A}{%
    \FormulaRefAuto{SemigroupClosureAxiom}[ax]{1,2}}
  \proofstepwidestar[1,2]{(a\star b)\star c\in A}{%
    \FormulaRefAuto{SemigroupClosureAxiom}[ax]{1,3,2}}
  \proofstepwidestar[1,2]{%
    (c\star d)\star e=c\star(d\star e)}{%
    \FormulaRefAuto{SemigroupAssociativityAxiom}[ax]{1,2}}
  \proofstepwidestar[1,2]{%
    \bigl((a\star b)\star c\bigr)\star(d\star e)
      =(a\star b)\star\bigl(c\star(d\star e)\bigr)}{%
    \hspace*{.75em}%
    \FormulaRefAuto{SemigroupAssociativityAxiom}[ax]{1,3,2,4}}
  \proofstepwidestarbreakkeep[1,2]{%
    \bigl(((a\star b)\star c)\star d\bigr)\star e
      =\bigl((a\star b)\star c\bigr)\star(d\star e)}{%
    \FormulaRefAuto{SemigroupAssociativityAxiom}[ax]{1,5,2}}
  \proofstepwidestar[1,2]{%
    (a\star b)\star c=a\star(b\star c)}{%
    \FormulaRefAuto{SemigroupAssociativityAxiom}[ax]{1,2}}
  \proofstepwide[1,2]{%
    (a\star b)\star\bigl((c\star d)\star e\bigr)}{=}{%
    (a\star b)\star\bigl(c\star(d\star e)\bigr)}{%
    \hspace*{1em}%
    \FormulaRefAuto{BinaryFunctionRightArgumentCongruence}[thm]{6}}
  \proofstepwide[1,2]{}{=}{%
    \bigl((a\star b)\star c\bigr)\star(d\star e)}{%
    \hspace*{1em}%
    \FormulaRefAuto{a=b\vdash b=a}[thm]{7}}
  \proofstepwide[1,2]{}{=}{%
    \bigl(((a\star b)\star c)\star d\bigr)\star e}{%
    \hspace*{1em}%
    \FormulaRefAuto{a=b\vdash b=a}[thm]{8}}
  \proofstepwide[1,2]{}{=}{%
    \bigl((a\star(b\star c))\star d\bigr)\star e}{%
    \hspace*{1em}%
    \FormulaRefAuto{BinaryFunctionLeftArgumentCongruence}[thm]{%
      \FormulaRefAuto{BinaryFunctionLeftArgumentCongruence}[thm]{9}}}
  \proofstepwidestar[1,2]{%
    (a\star b)\star\bigl((c\star d)\star e\bigr)
      =\bigl((a\star(b\star c))\star d\bigr)\star e}{%
    \rChain{10,13}}
\end{tabproofwide}

\section{Produkte endlicher Wörter}

Der fünfgliedrige Satz aus dem vorigen Abschnitt zeigt die Kernidee an einem
konkreten Produkt. Für beliebig lange Produkte werden nun die in Band~22
konstruierten Wortwerkzeuge verwendet.

Für eine Halbgruppe \(\SemigroupAlg{A}{\star}\) bezeichnet
\(\WordFold{\star}{u}\) die in Band~22 rekursiv definierte Linksfaltung des
nichtleeren Wortes \(u\in\NonemptyWords{A}\). Das folgende Blockgesetz ist
der einzige Punkt des Abschnitts, an dem die Assoziativität der
Halbgruppenoperation benötigt wird.

\FormulaThmDeltaK[Induktionsanfang des Blockgesetzes]{%
  \SemigroupAlg{A}{\star}\dsep u\in\NonemptyWords{A}\dsep a\in A
  \vdash
  \WordFold{\star}{\WordConcat{u}{\LetterWord{a}}}
  =\WordFold{\star}{u}\star\WordFold{\star}{\LetterWord{a}}%
}{SemigroupWordBlockBase}{%
  \DeltaRow{Mengen}{A\dsep a\dsep u}
  \DeltaRow{Funktionen}{\star}
}
\begin{tabproofwide}
  \proofstepwidestar[1]{\SemigroupAlg{A}{\star}}{\rA}
  \proofstepwidestar[2]{u\in\NonemptyWords{A}}{\rA}
  \proofstepwidestar[3]{a\in A}{\rA}
  \proofstepwidestar[1]{\mathsf{BinOp}(A,\star)}{%
    \FormulaRefAuto{SemigroupBinaryOperationAxiom}[ax]{1}}
  \proofstepwidestar[1,3]{%
    \WordFold{\star}{\LetterWord{a}}=a}{%
    \FormulaRefAuto{LeftWordFoldEquations}[thm]{4,3}}
  \proofstepwidestar[1,3]{%
    a=\WordFold{\star}{\LetterWord{a}}}{%
    \FormulaRefAuto{a=b\vdash b=a}[thm]{5}}
  \proofstepwide[1,2,3]{%
    \WordFold{\star}{\WordConcat{u}{\LetterWord{a}}}}{=}{%
    \WordFold{\star}{u}\star a}{%
    \FormulaRefAuto{LeftWordFoldEquations}[thm]{4,2,3}}
  \proofstepwide[1,2,3]{}{=}{%
    \WordFold{\star}{u}\star\WordFold{\star}{\LetterWord{a}}}{%
    \FormulaRefAuto{BinaryFunctionRightArgumentCongruence}[thm]{6}}
  \proofstepwidestar[1,2,3]{%
    \WordFold{\star}{\WordConcat{u}{\LetterWord{a}}}
    =\WordFold{\star}{u}\star\WordFold{\star}{\LetterWord{a}}}{%
    \rChain{7,8}}
\end{tabproofwide}

\FormulaThmDeltaK[Induktionsschritt des Blockgesetzes]{%
  \begin{aligned}[t]
    &\SemigroupAlg{A}{\star}\dsep
      u,v\in\NonemptyWords{A}\dsep a\in A\dsep
      \WordFold{\star}{\WordConcat{u}{v}}
        =\WordFold{\star}{u}\star\WordFold{\star}{v}\vdash{}\\[-2pt]
    &\WordFold{\star}{\WordConcat{u}{\WordConcat{v}{\LetterWord{a}}}}
      =\WordFold{\star}{u}\star
        \WordFold{\star}{\WordConcat{v}{\LetterWord{a}}}
  \end{aligned}%
}{SemigroupWordBlockStep}{%
  \DeltaRow{Mengen}{A\dsep a\dsep u\dsep v}
  \DeltaRow{Funktionen}{\star}
}
\begin{tabproofwide}
  \proofstepwidestar[1]{\SemigroupAlg{A}{\star}}{\rA}
  \proofstepwidestar[2]{u\in\NonemptyWords{A}}{\rA}
  \proofstepwidestar[3]{v\in\NonemptyWords{A}}{\rA}
  \proofstepwidestar[4]{a\in A}{\rA}
  \proofstepwidestar[5]{%
    \WordFold{\star}{\WordConcat{u}{v}}
      =\WordFold{\star}{u}\star\WordFold{\star}{v}}{\rA}
  \proofstepwidestar[1]{\mathsf{BinOp}(A,\star)}{%
    \FormulaRefAuto{SemigroupBinaryOperationAxiom}[ax]{1}}
  \proofstepwidestar[1]{%
    \Pi_{\star}\colon\NonemptyWords{A}\to A}{%
    \FormulaRefAuto{LeftWordFoldFunction}[thm]{6}}
  \proofstepwidestar[1,2]{\WordFold{\star}{u}\in A}{%
    \FormulaRefAuto{F\colon A\to B\dsep x\in A\vdash F(x)\in B}[thm]{7,2}}
  \proofstepwidestar[1,3]{\WordFold{\star}{v}\in A}{%
    \FormulaRefAuto{F\colon A\to B\dsep x\in A\vdash F(x)\in B}[thm]{7,3}}
  \proofstepwidestar[2,3]{\WordConcat{u}{v}\in\NonemptyWords{A}}{%
    \FormulaRefAuto{NonemptyWordConcatenationClosure}[thm]{2,3}}
  \proofstepwidestar[4]{\LetterWord{a}\in\NonemptyWords{A}}{%
    \FormulaRefAuto{LetterWordFacts}[thm]{4}}
  \proofstepwidestar[2]{u\in\FiniteWords{A}}{%
    \rAEa{\FormulaRefAuto{NonemptyWordMembership}[thm]{2}}}
  \proofstepwidestar[3]{v\in\FiniteWords{A}}{%
    \rAEa{\FormulaRefAuto{NonemptyWordMembership}[thm]{3}}}
  \proofstepwidestar[4]{\LetterWord{a}\in\FiniteWords{A}}{%
    \rAEa{\FormulaRefAuto{NonemptyWordMembership}[thm]{11}}}
  \proofstepwidestar[2,3,4]{%
    \WordConcat{\WordConcat{u}{v}}{\LetterWord{a}}
      =\WordConcat{u}{\WordConcat{v}{\LetterWord{a}}}}{%
    \FormulaRefAuto{WordConcatenationAssociative}[thm]{12,13,14}}
  \proofstepwidestar[2,3,4]{%
    \WordConcat{u}{\WordConcat{v}{\LetterWord{a}}}
      =\WordConcat{\WordConcat{u}{v}}{\LetterWord{a}}}{%
    \FormulaRefAuto{a=b\vdash b=a}[thm]{15}}
  \proofstepwidestar[1,3,4]{%
    \WordFold{\star}{\WordConcat{v}{\LetterWord{a}}}
      =\WordFold{\star}{v}\star a}{%
    \FormulaRefAuto{LeftWordFoldEquations}[thm]{6,3,4}}
  \proofstepwidestar[1,3,4]{%
    \WordFold{\star}{v}\star a
      =\WordFold{\star}{\WordConcat{v}{\LetterWord{a}}}}{%
    \FormulaRefAuto{a=b\vdash b=a}[thm]{17}}
  \proofstepwide[1,2,3,4,5]{%
    \WordFold{\star}{\WordConcat{u}{\WordConcat{v}{\LetterWord{a}}}}}{=}{%
    \WordFold{\star}{\WordConcat{\WordConcat{u}{v}}{\LetterWord{a}}}}{%
    \rIE{16}}
  \proofstepwide[1,2,3,4,5]{}{=}{%
    \WordFold{\star}{\WordConcat{u}{v}}\star a}{%
    \FormulaRefAuto{LeftWordFoldEquations}[thm]{6,10,4}}
  \proofstepwide[1,2,3,4,5]{}{=}{%
    \bigl(\WordFold{\star}{u}\star\WordFold{\star}{v}\bigr)\star a}{%
    \FormulaRefAuto{BinaryFunctionLeftArgumentCongruence}[thm]{5}}
  \proofstepwide[1,2,3,4,5]{}{=}{%
    \WordFold{\star}{u}\star
      \bigl(\WordFold{\star}{v}\star a\bigr)}{%
    \FormulaRefAuto{SemigroupAssociativityAxiom}[ax]{1,8,9,4}}
  \proofstepwide[1,2,3,4,5]{}{=}{%
    \WordFold{\star}{u}\star
      \WordFold{\star}{\WordConcat{v}{\LetterWord{a}}}}{%
    \FormulaRefAuto{BinaryFunctionRightArgumentCongruence}[thm]{18}}
  \proofstepwidestar[1,2,3,4,5]{%
    \WordFold{\star}{\WordConcat{u}{\WordConcat{v}{\LetterWord{a}}}}
      =\WordFold{\star}{u}\star
        \WordFold{\star}{\WordConcat{v}{\LetterWord{a}}}}{%
    \rChain{19,23}}
\end{tabproofwide}

\FormulaThmDeltaK[Blockgesetz für Halbgruppenprodukte]{%
  \SemigroupAlg{A}{\star}\dsep u,v\in\NonemptyWords{A}
  \vdash
  \WordFold{\star}{\WordConcat{u}{v}}
    =\WordFold{\star}{u}\star\WordFold{\star}{v}%
}{SemigroupWordBlockLaw}{%
  \DeltaRow{Mengen}{A\dsep u\dsep v}
  \DeltaRow{Funktionen}{\star}
}
\begin{tabproofwide}
  \proofstepwidestar[1]{\SemigroupAlg{A}{\star}}{\rA}
  \proofstepwidestar[2]{u\in\NonemptyWords{A}}{\rA}
  \proofstepwidestar[3]{v\in\NonemptyWords{A}}{\rA}
  \proofstepwidestarbreakkeep[1,2,3]{%
    \WordFold{\star}{\WordConcat{u}{v}}
      =\WordFold{\star}{u}\star\WordFold{\star}{v}}{%
    \FormulaRefAuto{NonemptyWordInduction}[thm]{%
      \FormulaRefAuto{SemigroupWordBlockBase}[thm]{1,2},
      \FormulaRefAuto{SemigroupWordBlockStep}[thm]{1,2},3}}
\end{tabproofwide}

\section{Klammerungsunabhängigkeit}

\FormulaThmDeltaK[Blattfall der Baum-Normalform]{%
  \SemigroupAlg{A}{\star}\dsep a\in A\vdash
  \TreeEvaluation{\star}{\LeafTree{A}{a}}
  =\WordFold{\star}{\TreeLeafWord{A}{\LeafTree{A}{a}}}%
}{SemigroupTreeNormalFormLeaf}{%
  \DeltaRow{Mengen}{A\dsep a}
  \DeltaRow{Funktionen}{\star}
}
\begin{tabproofwide}
  \proofstepwidestar[1]{\SemigroupAlg{A}{\star}}{\rA}
  \proofstepwidestar[2]{a\in A}{\rA}
  \proofstepwidestar[1]{\mathsf{BinOp}(A,\star)}{%
    \FormulaRefAuto{SemigroupBinaryOperationAxiom}[ax]{1}}
  \proofstepwidestar[1,2]{%
    \WordFold{\star}{\LetterWord{a}}=a}{%
    \FormulaRefAuto{LeftWordFoldEquations}[thm]{3,2}}
  \proofstepwidestar[1,2]{%
    a=\WordFold{\star}{\LetterWord{a}}}{%
    \FormulaRefAuto{a=b\vdash b=a}[thm]{4}}
  \proofstepwidestar[2]{%
    \TreeLeafWord{A}{\LeafTree{A}{a}}=\LetterWord{a}}{%
    \FormulaRefAuto{TreeLeafWordEquations}[thm]{2}}
  \proofstepwidestar[2]{%
    \LetterWord{a}=\TreeLeafWord{A}{\LeafTree{A}{a}}}{%
    \FormulaRefAuto{a=b\vdash b=a}[thm]{6}}
  \proofstepwide[1,2]{%
    \TreeEvaluation{\star}{\LeafTree{A}{a}}}{=}{a}{%
    \FormulaRefAuto{TreeEvaluationEquations}[thm]{3,2}}
  \proofstepwide[1,2]{}{=}{\WordFold{\star}{\LetterWord{a}}}{%
    \rIE{5}}
  \proofstepwide[1,2]{}{=}{%
    \WordFold{\star}{\TreeLeafWord{A}{\LeafTree{A}{a}}}}{%
    \rIE{7}}
  \proofstepwidestar[1,2]{%
    \TreeEvaluation{\star}{\LeafTree{A}{a}}
    =\WordFold{\star}{\TreeLeafWord{A}{\LeafTree{A}{a}}}}{%
    \rChain{8,10}}
\end{tabproofwide}

\FormulaThmDeltaK[Knotenschritt der Baum-Normalform]{%
  \begin{aligned}[t]
    &\SemigroupAlg{A}{\star}\dsep
      S,T\in\BracketTreeSet{A}\dsep
      \TreeEvaluation{\star}{S}=\WordFold{\star}{\TreeLeafWord{A}{S}}\dsep{}\\[-2pt]
    &\TreeEvaluation{\star}{T}=\WordFold{\star}{\TreeLeafWord{A}{T}}
      \vdash{}\\[-2pt]
    &\TreeEvaluation{\star}{\NodeTree{A}{S}{T}}
      =\WordFold{\star}{\TreeLeafWord{A}{\NodeTree{A}{S}{T}}}
  \end{aligned}%
}{SemigroupTreeNormalFormNode}{%
  \DeltaRow{Mengen}{A\dsep S\dsep T}
  \DeltaRow{Funktionen}{\star}
}
\begin{tabproofwide}
  \proofstepwidestar[1]{\SemigroupAlg{A}{\star}}{\rA}
  \proofstepwidestar[2]{S\in\BracketTreeSet{A}}{\rA}
  \proofstepwidestar[3]{T\in\BracketTreeSet{A}}{\rA}
  \proofstepwidestar[4]{%
    \TreeEvaluation{\star}{S}=\WordFold{\star}{\TreeLeafWord{A}{S}}}{\rA}
  \proofstepwidestar[5]{%
    \TreeEvaluation{\star}{T}=\WordFold{\star}{\TreeLeafWord{A}{T}}}{\rA}
  \proofstepwidestar[1]{\mathsf{BinOp}(A,\star)}{%
    \FormulaRefAuto{SemigroupBinaryOperationAxiom}[ax]{1}}
  \proofstepwidestar[]{%
    \operatorname{wort}_A\colon
      \BracketTreeSet{A}\to\NonemptyWords{A}}{%
    \FormulaRefAuto{TreeLeafWordEquations}[thm]}
  \proofstepwidestar[2]{\TreeLeafWord{A}{S}\in\NonemptyWords{A}}{%
    \FormulaRefAuto{F\colon A\to B\dsep x\in A\vdash F(x)\in B}[thm]{%
      7,2}}
  \proofstepwidestar[3]{\TreeLeafWord{A}{T}\in\NonemptyWords{A}}{%
    \FormulaRefAuto{F\colon A\to B\dsep x\in A\vdash F(x)\in B}[thm]{%
      7,3}}
  \proofstepwidestar[1,2,3]{%
    \begin{aligned}[t]
      &\WordFold{\star}{%
        \WordConcat{\TreeLeafWord{A}{S}}{\TreeLeafWord{A}{T}}}\\[-2pt]
      &\qquad
        =\WordFold{\star}{\TreeLeafWord{A}{S}}\star
          \WordFold{\star}{\TreeLeafWord{A}{T}}
    \end{aligned}}{%
    \FormulaRefAuto{SemigroupWordBlockLaw}[thm]{1,8,9}}
  \proofstepwidestar[1,2,3]{%
    \begin{aligned}[t]
      &\WordFold{\star}{\TreeLeafWord{A}{S}}\star
        \WordFold{\star}{\TreeLeafWord{A}{T}}\\[-2pt]
      &\qquad
        =\WordFold{\star}{%
          \WordConcat{\TreeLeafWord{A}{S}}{\TreeLeafWord{A}{T}}}
    \end{aligned}}{%
    \FormulaRefAuto{a=b\vdash b=a}[thm]{10}}
  \proofstepwidestar[2,3]{%
    \begin{aligned}[t]
      &\TreeLeafWord{A}{\NodeTree{A}{S}{T}}\\[-2pt]
      &\qquad
        =\WordConcat{\TreeLeafWord{A}{S}}{\TreeLeafWord{A}{T}}
    \end{aligned}}{%
    \FormulaRefAuto{TreeLeafWordEquations}[thm]{2,3}}
  \proofstepwidestar[2,3]{%
    \begin{aligned}[t]
      &\WordConcat{\TreeLeafWord{A}{S}}{\TreeLeafWord{A}{T}}\\[-2pt]
      &\qquad
        =\TreeLeafWord{A}{\NodeTree{A}{S}{T}}
    \end{aligned}}{%
    \FormulaRefAuto{a=b\vdash b=a}[thm]{12}}
  \proofstepwide[1,2,3,4,5]{%
    \TreeEvaluation{\star}{\NodeTree{A}{S}{T}}}{=}{%
    \TreeEvaluation{\star}{S}\star\TreeEvaluation{\star}{T}}{%
    \FormulaRefAuto{TreeEvaluationEquations}[thm]{6,2,3}}
  \proofstepwide[1,2,3,4,5]{}{=}{%
    \WordFold{\star}{\TreeLeafWord{A}{S}}\star
      \TreeEvaluation{\star}{T}}{%
    \FormulaRefAuto{BinaryFunctionLeftArgumentCongruence}[thm]{4}}
  \proofstepwide[1,2,3,4,5]{}{=}{%
    \begin{aligned}[t]
      &\WordFold{\star}{\TreeLeafWord{A}{S}}\\[-2pt]
      &\quad\star\WordFold{\star}{\TreeLeafWord{A}{T}}
    \end{aligned}}{%
    \FormulaRefAuto{BinaryFunctionRightArgumentCongruence}[thm]{5}}
  \proofstepwide[1,2,3,4,5]{}{=}{%
    \WordFold{\star}{%
      \WordConcat{\TreeLeafWord{A}{S}}{\TreeLeafWord{A}{T}}}}{%
    \rIE{11}}
  \proofstepwide[1,2,3,4,5]{}{=}{%
    \WordFold{\star}{\TreeLeafWord{A}{\NodeTree{A}{S}{T}}}}{%
    \rIE{13}}
  \proofstepwidestar[1,2,3,4,5]{%
    \TreeEvaluation{\star}{\NodeTree{A}{S}{T}}
      =\WordFold{\star}{\TreeLeafWord{A}{\NodeTree{A}{S}{T}}}}{%
    \rChain{14,18}}
\end{tabproofwide}

\FormulaThmDeltaK[Baum-Normalform für Halbgruppenprodukte]{%
  \SemigroupAlg{A}{\star}\dsep T\in\BracketTreeSet{A}\vdash
  \TreeEvaluation{\star}{T}=\WordFold{\star}{\TreeLeafWord{A}{T}}%
}{SemigroupTreeNormalForm}{%
  \DeltaRow{Mengen}{A\dsep T}
  \DeltaRow{Funktionen}{\star}
}
\begin{tabproofwide}
  \proofstepwidestar[1]{\SemigroupAlg{A}{\star}}{\rA}
  \proofstepwidestar[2]{T\in\BracketTreeSet{A}}{\rA}
  \proofstepwidestarbreakkeep[1,2]{%
    \TreeEvaluation{\star}{T}=\WordFold{\star}{\TreeLeafWord{A}{T}}}{%
    \FormulaRefAuto{FullPlanarBinaryTreeStructuralInduction}[thm]{%
      \FormulaRefAuto{SemigroupTreeNormalFormLeaf}[thm]{1},
      \FormulaRefAuto{SemigroupTreeNormalFormNode}[thm]{1},2}}
\end{tabproofwide}

\FormulaThmDeltaK[Unabhängigkeit von der Klammerung]{%
  \begin{aligned}[t]
    &\SemigroupAlg{A}{\star}\dsep w\in\NonemptyWords{A}\dsep
      S,T\in\WordBracketings{A}{w}\vdash{}\\[-2pt]
    &\TreeEvaluation{\star}{S}=\TreeEvaluation{\star}{T}
  \end{aligned}%
}{SemigroupBracketingIndependence}{%
  \DeltaRow{Mengen}{A\dsep S\dsep T\dsep w}
  \DeltaRow{Funktionen}{\star}
}
\begin{tabproofwide}
  \proofstepwidestar[1]{\SemigroupAlg{A}{\star}}{\rA}
  \proofstepwidestar[2]{w\in\NonemptyWords{A}}{\rA}
  \proofstepwidestar[3]{S\in\WordBracketings{A}{w}}{\rA}
  \proofstepwidestar[4]{T\in\WordBracketings{A}{w}}{\rA}
  \proofstepwidestar[2,3]{%
    S\in\BracketTreeSet{A}\land\TreeLeafWord{A}{S}=w}{%
    \FormulaRefAuto{WordBracketingMembership}[thm]{2,3}}
  \proofstepwidestar[2,4]{%
    T\in\BracketTreeSet{A}\land\TreeLeafWord{A}{T}=w}{%
    \FormulaRefAuto{WordBracketingMembership}[thm]{2,4}}
  \proofstepwidestar[2,3]{S\in\BracketTreeSet{A}}{\rAEa{5}}
  \proofstepwidestar[2,3]{\TreeLeafWord{A}{S}=w}{\rAEb{5}}
  \proofstepwidestar[2,4]{T\in\BracketTreeSet{A}}{\rAEa{6}}
  \proofstepwidestar[2,4]{\TreeLeafWord{A}{T}=w}{\rAEb{6}}
  \proofstepwide[1,2,3,4]{\TreeEvaluation{\star}{S}}{=}{%
    \WordFold{\star}{\TreeLeafWord{A}{S}}}{%
    \FormulaRefAuto{SemigroupTreeNormalForm}[thm]{1,7}}
  \proofstepwide[1,2,3,4]{}{=}{\WordFold{\star}{w}}{\rIE{8}}
  \proofstepwide[1,2,3,4]{}{=}{%
    \WordFold{\star}{\TreeLeafWord{A}{T}}}{%
    \rIE{\FormulaRefAuto{a=b\vdash b=a}[thm]{10}}}
  \proofstepwide[1,2,3,4]{}{=}{\TreeEvaluation{\star}{T}}{%
    \FormulaRefAuto{a=b\vdash b=a}[thm]{%
      \FormulaRefAuto{SemigroupTreeNormalForm}[thm]{1,9}}}
  \proofstepwidestar[1,2,3,4]{%
    \TreeEvaluation{\star}{S}=\TreeEvaluation{\star}{T}}{%
    \rChain{11,14}}
\end{tabproofwide}

\section{Beispiele von Halbgruppen}

\subsection{Die Potenzhalbgruppe}

Aus jeder Halbgruppe entsteht eine weitere Halbgruppe, deren Elemente die
nichtleeren Teilmengen des ursprünglichen Trägers sind. Die Menge dieser
Teilmengen schreiben wir als \(\PowNE{A}\).

Für \(X,Y\in\PowNE{A}\) wird die ursprüngliche Operation mengenweise
fortgesetzt. Der Index \(\mathcal P\) unterscheidet diese Operation auf
Teilmengen von der Operation \(\star\) auf einzelnen Elementen.

\FormulaDefDeltaK[Mengenweise Fortsetzung einer Halbgruppenoperation]{%
  \begin{aligned}[t]
    &\SemigroupAlg{A}{\star}\dsep X,Y\in\PowNE{A}\vdash{}\\[-2pt]
    &X\star_{\mathcal P}Y\coloneqq
      \bigl\{z\in A\bigm|
        \exists x\in X\,\exists y\in Y\,(z=x\star y)\bigr\}
  \end{aligned}
}{PowerSemigroupProductDef}{%
  \DeltaRow{Mengen}{A\dsep X\dsep Y\dsep x\dsep y\dsep z}
  \DeltaRow{Funktionen}{\star\dsep\star_{\mathcal P}}
  \DeltaRow{\textbf{Neues Symbol}}{}
  \DeltaPrem{Mengenweise Operation}{\star_{\mathcal P}}
}

\FormulaThmDeltaK[Elementkriterium für das Mengenprodukt]{%
  \begin{aligned}[t]
    &\SemigroupAlg{A}{\star}\dsep X,Y\in\PowNE{A}\vdash{}\\[-2pt]
    &z\in X\star_{\mathcal P}Y
      \leftrightarrow
      \Bigl(z\in A\land
        \exists x\in X\,\exists y\in Y\,(z=x\star y)\Bigr)
  \end{aligned}
}{PowerSemigroupProductMembership}{%
  \DeltaRow{Mengen}{A\dsep X\dsep Y\dsep x\dsep y\dsep z}
  \DeltaRow{Funktionen}{\star\dsep\star_{\mathcal P}}
}
\begin{tabproofwide}
  \proofstepwidestar[1]{\SemigroupAlg{A}{\star}}{\rA}
  \proofstepwidestar[2]{X\in\PowNE{A}}{\rA}
  \proofstepwidestar[3]{Y\in\PowNE{A}}{\rA}
  \proofstepwidestar[1,2,3]{%
    X\star_{\mathcal P}Y=
    \bigl\{z\in A\bigm|
      \exists x\in X\,\exists y\in Y\,(z=x\star y)\bigr\}}{%
    \FormulaRefAuto{PowerSemigroupProductDef}[def]{1,2,3}}
  \proofstepwidestar[1,2,3]{%
    \begin{aligned}[t]
      z\in X\star_{\mathcal P}Y
      \leftrightarrow
      \Bigl(z\in A\land
        \exists x\in X\,\exists y\in Y\,(z=x\star y)\Bigr)
    \end{aligned}}{%
    \FormulaRefAuto{x\in\{u\in A\mid P(u)\}\eqvdash
      x\in A\land P(x)}[thm]{4}}
\end{tabproofwide}

Im folgenden Beweis fassen \(\exists I^{*}\) und \(\exists E^{*}\) endlich
viele geschachtelte Anwendungen von \(\exists I\) beziehungsweise
\(\exists E\) zusammen. Bei beschränkten Quantoren schließen diese
Kurzschreibweisen die erforderlichen \(\land\)-Schritte ein.

\FormulaThmDeltaK[Die nichtleeren Teilmengen bilden eine Halbgruppe]{%
  \SemigroupAlg{A}{\star}\vdash
  \SemigroupAlg{\PowNE{A}}{\star_{\mathcal P}}
}{NonemptyPowerSemigroup}{%
  \DeltaRow{Mengen}{%
    A\dsep X\dsep Y\dsep Z\dsep u\dsep v\dsep w\dsep x\dsep y\dsep z}
  \DeltaRow{Funktionen}{\star\dsep\star_{\mathcal P}}
}
\begin{tabproofsplitwide}
  \proofpartwideindR[Trägerelement einer nichtleeren Teilmenge]{%
    X\in\PowNE{A}\dsep x\in X\vdash x\in A
  }{%
    X\in\PowNE{A}\dsep x\in X\vdash x\in A
  }[PowerSemigroupCarrierElement]
    \proofstepwidestar[1]{X\in\PowNE{A}}{\rA}
    \proofstepwidestar[2]{x\in X}{\rA}
    \proofstepwidestar[1]{X\subseteq A}{%
      \rAEa{\FormulaRefAuto{NonemptyPowersetMembership}[thm]{1}}}
    \proofstepwidestar[1,2]{x\in A}{%
      \FormulaRefAuto{A\subseteq B,\,x\in A\vdash x\in B}[thm]{3,2}}
  \closeproofpartwideindR

  \proofpartwideindR[Elementare Produkte liegen im Mengenprodukt]{%
    \begin{aligned}[t]
      &\SemigroupAlg{A}{\star}\dsep X,Y\in\PowNE{A}\dsep
        x\in X\dsep y\in Y\vdash{}\\[-2pt]
      &x\star y\in X\star_{\mathcal P}Y
    \end{aligned}
  }{%
    \begin{aligned}[t]
      &\SemigroupAlg{A}{\star}\dsep X,Y\in\PowNE{A}\dsep
        x\in X\dsep y\in Y\vdash{}\\[-2pt]
      &x\star y\in X\star_{\mathcal P}Y
    \end{aligned}
  }[PowerSemigroupProductContains]
    \proofstepwidestar[1]{\SemigroupAlg{A}{\star}}{\rA}
    \proofstepwidestar[2]{X\in\PowNE{A}}{\rA}
    \proofstepwidestar[3]{Y\in\PowNE{A}}{\rA}
    \proofstepwidestar[4]{x\in X}{\rA}
    \proofstepwidestar[5]{y\in Y}{\rA}
    \proofstepwidestar[2,4]{x\in A}{%
      \FormulaRefAuto{PowerSemigroupCarrierElement}[thm]{2,4}}
    \proofstepwidestar[3,5]{y\in A}{%
      \FormulaRefAuto{PowerSemigroupCarrierElement}[thm]{3,5}}
    \proofstepwidestar[1,2,3,4,5]{x\star y\in A}{%
      \FormulaRefAuto{SemigroupClosureAxiom}[ax]{1,6,7}}
    \proofstepwidestar[4,5]{%
      \exists u\in X\,\exists v\in Y\,(x\star y=u\star v)}{%
      \rEI{4,\rEI{5,\rII}}}
    \proofstepwidestar[1,2,3,4,5]{%
      x\star y\in A\land
      \exists u\in X\,\exists v\in Y\,(x\star y=u\star v)}{%
      \rAI{8,9}}
    \proofstepwidestar[1,2,3,4,5]{x\star y\in X\star_{\mathcal P}Y}{%
      \FormulaRefAuto{PowerSemigroupProductMembership}[thm]{1,2,3,10}}
  \closeproofpartwideindR

  \proofpartwideindR[Das Mengenprodukt bleibt im ursprünglichen Träger]{%
    \SemigroupAlg{A}{\star}\dsep X,Y\in\PowNE{A}
    \vdash X\star_{\mathcal P}Y\subseteq A
  }{%
    \SemigroupAlg{A}{\star}\dsep X,Y\in\PowNE{A}
    \vdash X\star_{\mathcal P}Y\subseteq A
  }[PowerSemigroupProductSubset]
    \proofstepwidestar[1]{\SemigroupAlg{A}{\star}}{\rA}
    \proofstepwidestar[2]{X\in\PowNE{A}}{\rA}
    \proofstepwidestar[3]{Y\in\PowNE{A}}{\rA}
    \proofstepwidestar[4]{w\in X\star_{\mathcal P}Y}{\rA}
    \proofstepwidestar[1,2,3,4]{%
      w\in A\land
      \exists x\in X\,\exists y\in Y\,(w=x\star y)}{%
      \FormulaRefAuto{PowerSemigroupProductMembership}[thm]{1,2,3,4}}
    \proofstepwidestar[1,2,3,4]{w\in A}{\rAEa{5}}
    \proofstepwidestar[1,2,3]{X\star_{\mathcal P}Y\subseteq A}{%
      \FormulaRefAuto{A\subseteq B\coloneq
        \forall x\,(x\in A\rightarrow x\in B)}[def]{\rUI{\rRI{4,6}}}}
  \closeproofpartwideindR

  \proofpartwideindR[Das Mengenprodukt ist nichtleer]{%
    \SemigroupAlg{A}{\star}\dsep X,Y\in\PowNE{A}
    \vdash X\star_{\mathcal P}Y\neq\varnothing
  }{%
    \SemigroupAlg{A}{\star}\dsep X,Y\in\PowNE{A}
    \vdash X\star_{\mathcal P}Y\neq\varnothing
  }[PowerSemigroupProductNonempty]
    \proofstepwidestar[1]{\SemigroupAlg{A}{\star}}{\rA}
    \proofstepwidestar[2]{X\in\PowNE{A}}{\rA}
    \proofstepwidestar[3]{Y\in\PowNE{A}}{\rA}
    \proofstepwidestar[2]{X\neq\varnothing}{%
      \rAEb{\FormulaRefAuto{NonemptyPowersetMembership}[thm]{2}}}
    \proofstepwidestar[3]{Y\neq\varnothing}{%
      \rAEb{\FormulaRefAuto{NonemptyPowersetMembership}[thm]{3}}}
    \proofstepwidestar[2]{\exists x\,(x\in X)}{%
      \FormulaRefAuto{A\neq\varnothing\eqvdash
        \exists x\,(x\in A)}[thm]{4}}
    \proofstepwidestar[3]{\exists y\,(y\in Y)}{%
      \FormulaRefAuto{A\neq\varnothing\eqvdash
        \exists x\,(x\in A)}[thm]{5}}
    \proofstepwidestar[6]{x\in X}{\rA}
    \proofstepwidestar[7]{y\in Y}{\rA}
    \proofstepwidestar[1,2,3,8,9]{x\star y\in X\star_{\mathcal P}Y}{%
      \FormulaRefAuto{PowerSemigroupProductContains}[thm]{1,2,3,8,9}}
    \proofstepwidestar[1,2,3,8,9]{X\star_{\mathcal P}Y\neq\varnothing}{%
      \FormulaRefAuto{a\in A\vdash A\neq\varnothing}[thm]{10}}
    \proofstepwidestar[1,2,3,8]{X\star_{\mathcal P}Y\neq\varnothing}{%
      \rEE{7,9,11}}
    \proofstepwidestar[1,2,3]{X\star_{\mathcal P}Y\neq\varnothing}{%
      \rEE{6,8,12}}
  \closeproofpartwideindR

  \proofpartwideindR[Abgeschlossenheit der mengenweisen Operation]{%
    \SemigroupAlg{A}{\star}\dsep X,Y\in\PowNE{A}
    \vdash X\star_{\mathcal P}Y\in\PowNE{A}
  }{%
    \SemigroupAlg{A}{\star}\dsep X,Y\in\PowNE{A}
    \vdash X\star_{\mathcal P}Y\in\PowNE{A}
  }[PowerSemigroupProductClosure]
    \proofstepwidestar[1]{\SemigroupAlg{A}{\star}}{\rA}
    \proofstepwidestar[2]{X\in\PowNE{A}}{\rA}
    \proofstepwidestar[3]{Y\in\PowNE{A}}{\rA}
    \proofstepwidestar[1,2,3]{X\star_{\mathcal P}Y\subseteq A}{%
      \FormulaRefAuto{PowerSemigroupProductSubset}[thm]{1,2,3}}
    \proofstepwidestar[1,2,3]{X\star_{\mathcal P}Y\neq\varnothing}{%
      \FormulaRefAuto{PowerSemigroupProductNonempty}[thm]{1,2,3}}
    \proofstepwidestar[1,2,3]{%
      X\star_{\mathcal P}Y\subseteq A\land
      X\star_{\mathcal P}Y\neq\varnothing}{\rAI{4,5}}
    \proofstepwidestar[1,2,3]{X\star_{\mathcal P}Y\in\PowNE{A}}{%
      \FormulaRefAuto{NonemptyPowersetMembership}[thm]{6}}
  \closeproofpartwideindR

  \proofpartwideindR[Zeugen eines Mengenprodukts]{%
    \begin{aligned}[t]
      &\SemigroupAlg{A}{\star}\dsep X,Y\in\PowNE{A}\dsep
        w\in X\star_{\mathcal P}Y\vdash{}\\[-2pt]
      &\exists x\in X\,\exists y\in Y\,(w=x\star y)
    \end{aligned}
  }{%
    \begin{aligned}[t]
      &\SemigroupAlg{A}{\star}\dsep X,Y\in\PowNE{A}\dsep
        w\in X\star_{\mathcal P}Y\vdash{}\\[-2pt]
      &\exists x\in X\,\exists y\in Y\,(w=x\star y)
    \end{aligned}
  }[PowerSemigroupProductWitnesses]
    \proofstepwidestar[1]{\SemigroupAlg{A}{\star}}{\rA}
    \proofstepwidestar[2]{X\in\PowNE{A}}{\rA}
    \proofstepwidestar[3]{Y\in\PowNE{A}}{\rA}
    \proofstepwidestar[4]{w\in X\star_{\mathcal P}Y}{\rA}
    \proofstepwidestar[1,2,3,4]{%
      w\in A\land
      \exists x\in X\,\exists y\in Y\,(w=x\star y)}{%
      \FormulaRefAuto{PowerSemigroupProductMembership}[thm]{1,2,3,4}}
    \proofstepwidestar[1,2,3,4]{%
      \exists x\in X\,\exists y\in Y\,(w=x\star y)}{\rAEb{5}}
  \closeproofpartwideindR

  \proofpartwideindR[Assoziativität auf Teilmengenelementen]{%
    \begin{aligned}[t]
      &\SemigroupAlg{A}{\star}\dsep X,Y,Z\in\PowNE{A}\dsep
        x\in X\dsep y\in Y\dsep z\in Z\vdash{}\\[-2pt]
      &(x\star y)\star z=x\star(y\star z)
    \end{aligned}
  }{%
    \begin{aligned}[t]
      &\SemigroupAlg{A}{\star}\dsep X,Y,Z\in\PowNE{A}\dsep
        x\in X\dsep y\in Y\dsep z\in Z\vdash{}\\[-2pt]
      &(x\star y)\star z=x\star(y\star z)
    \end{aligned}
  }[PowerSemigroupElementAssociativity]
    \proofstepwidestar[1]{\SemigroupAlg{A}{\star}}{\rA}
    \proofstepwidestar[2]{X\in\PowNE{A}}{\rA}
    \proofstepwidestar[3]{Y\in\PowNE{A}}{\rA}
    \proofstepwidestar[4]{Z\in\PowNE{A}}{\rA}
    \proofstepwidestar[5]{x\in X}{\rA}
    \proofstepwidestar[6]{y\in Y}{\rA}
    \proofstepwidestar[7]{z\in Z}{\rA}
    \proofstepwidestar[2,5]{x\in A}{%
      \FormulaRefAuto{PowerSemigroupCarrierElement}[thm]{2,5}}
    \proofstepwidestar[3,6]{y\in A}{%
      \FormulaRefAuto{PowerSemigroupCarrierElement}[thm]{3,6}}
    \proofstepwidestar[4,7]{z\in A}{%
      \FormulaRefAuto{PowerSemigroupCarrierElement}[thm]{4,7}}
    \proofstepwidestar[1,2,3,4,5,6,7]{%
      (x\star y)\star z=x\star(y\star z)}{%
      \FormulaRefAuto{SemigroupAssociativityAxiom}[ax]{1,8,9,10}}
  \closeproofpartwideindR

  \proofpartwideindR[Dreifachzeugen der links geklammerten Seite]{%
    \begin{aligned}[t]
      &\SemigroupAlg{A}{\star}\dsep X,Y,Z\in\PowNE{A}\dsep
        w\in (X\star_{\mathcal P}Y)\star_{\mathcal P}Z\vdash{}\\[-2pt]
      &\exists x\in X\,\exists y\in Y\,\exists z\in Z\,
        \bigl(w=(x\star y)\star z\bigr)
    \end{aligned}
  }{%
    \begin{aligned}[t]
      &\SemigroupAlg{A}{\star}\dsep X,Y,Z\in\PowNE{A}\dsep
        w\in (X\star_{\mathcal P}Y)\star_{\mathcal P}Z\vdash{}\\[-2pt]
      &\exists x\in X\,\exists y\in Y\,\exists z\in Z\,
        \bigl(w=(x\star y)\star z\bigr)
    \end{aligned}
  }[PowerSemigroupLeftTripleMembershipForward]
    \proofstepwidestar[1]{\SemigroupAlg{A}{\star}}{\rA}
    \proofstepwidestar[2]{X\in\PowNE{A}}{\rA}
    \proofstepwidestar[3]{Y\in\PowNE{A}}{\rA}
    \proofstepwidestar[4]{Z\in\PowNE{A}}{\rA}
    \proofstepwidestar[5]{%
      w\in (X\star_{\mathcal P}Y)\star_{\mathcal P}Z}{\rA}
    \proofstepwidestar[1,2,3]{X\star_{\mathcal P}Y\in\PowNE{A}}{%
      \FormulaRefAuto{PowerSemigroupProductClosure}[thm]{1,2,3}}
    \proofstepwidestar[1,2,3,4,5]{%
      \exists u\in X\star_{\mathcal P}Y\,
      \exists z\in Z\,(w=u\star z)}{%
      \FormulaRefAuto{PowerSemigroupProductWitnesses}[thm]{1,6,4,5}}
    \proofstepwidestar[7]{u\in X\star_{\mathcal P}Y}{\rA}
    \proofstepwidestar[7]{z\in Z}{\rA}
    \proofstepwidestar[7]{w=u\star z}{\rA}
    \proofstepwidestar[1,2,3,8]{%
      \exists x\in X\,\exists y\in Y\,(u=x\star y)}{%
      \FormulaRefAuto{PowerSemigroupProductWitnesses}[thm]{1,2,3,8}}
    \proofstepwidestar[11]{x\in X}{\rA}
    \proofstepwidestar[11]{y\in Y}{\rA}
    \proofstepwidestar[11]{u=x\star y}{\rA}
    \proofstepwidestar[7,11]{w=(x\star y)\star z}{\rIE{14,10}}
    \proofstepwidestar[7,11]{%
      \exists x\in X\,\exists y\in Y\,\exists z\in Z\,
      \bigl(w=(x\star y)\star z\bigr)}{%
      \rEIStar{12,13,9,15}}
    \proofstepwidestar[1,2,3,4,5]{%
      \exists x\in X\,\exists y\in Y\,\exists z\in Z\,
      \bigl(w=(x\star y)\star z\bigr)}{%
      \rEEStar{7\text{--}16}}
  \closeproofpartwideindR

  \proofpartwideindR[Aufbau der links geklammerten Seite]{%
    \begin{aligned}[t]
      &\SemigroupAlg{A}{\star}\dsep X,Y,Z\in\PowNE{A}\dsep{}\\[-2pt]
      &\exists x\in X\,\exists y\in Y\,\exists z\in Z\,
        \bigl(w=(x\star y)\star z\bigr)\vdash{}\\[-2pt]
      &w\in (X\star_{\mathcal P}Y)\star_{\mathcal P}Z
    \end{aligned}
  }{%
    \begin{aligned}[t]
      &\SemigroupAlg{A}{\star}\dsep X,Y,Z\in\PowNE{A}\dsep{}\\[-2pt]
      &\exists x\in X\,\exists y\in Y\,\exists z\in Z\,
        \bigl(w=(x\star y)\star z\bigr)\vdash{}\\[-2pt]
      &w\in (X\star_{\mathcal P}Y)\star_{\mathcal P}Z
    \end{aligned}
  }[PowerSemigroupLeftTripleMembershipBackward]
    \proofstepwidestar[1]{\SemigroupAlg{A}{\star}}{\rA}
    \proofstepwidestar[2]{X\in\PowNE{A}}{\rA}
    \proofstepwidestar[3]{Y\in\PowNE{A}}{\rA}
    \proofstepwidestar[4]{Z\in\PowNE{A}}{\rA}
    \proofstepwidestar[5]{%
      \exists x\in X\,\exists y\in Y\,\exists z\in Z\,
      \bigl(w=(x\star y)\star z\bigr)}{\rA}
    \proofstepwidestar[5]{x\in X}{\rA}
    \proofstepwidestar[5]{y\in Y}{\rA}
    \proofstepwidestar[5]{z\in Z}{\rA}
    \proofstepwidestar[5]{w=(x\star y)\star z}{\rA}
    \proofstepwidestar[1,2,3,6,7]{%
      x\star y\in X\star_{\mathcal P}Y}{%
      \FormulaRefAuto{PowerSemigroupProductContains}[thm]{1,2,3,6,7}}
    \proofstepwidestar[1,2,3]{X\star_{\mathcal P}Y\in\PowNE{A}}{%
      \FormulaRefAuto{PowerSemigroupProductClosure}[thm]{1,2,3}}
    \proofstepwidestar[1,2,3,4,6,7,8]{%
      (x\star y)\star z
      \in (X\star_{\mathcal P}Y)\star_{\mathcal P}Z}{%
      \FormulaRefAuto{PowerSemigroupProductContains}[thm]{1,11,4,10,8}}
    \proofstepwidestar[9]{(x\star y)\star z=w}{%
      \FormulaRefAuto{a=b\vdash b=a}[thm]{9}}
    \proofstepwidestar[1,2,3,4,5]{%
      w\in (X\star_{\mathcal P}Y)\star_{\mathcal P}Z}{\rIE{13,12}}
    \proofstepwidestar[1,2,3,4,5]{%
      w\in (X\star_{\mathcal P}Y)\star_{\mathcal P}Z}{%
      \rEEStar{5\text{--}14}}
  \closeproofpartwideindR

  \proofpartwideindR[Aufbau der rechts geklammerten Seite]{%
    \begin{aligned}[t]
      &\SemigroupAlg{A}{\star}\dsep X,Y,Z\in\PowNE{A}\dsep{}\\[-2pt]
      &\exists x\in X\,\exists y\in Y\,\exists z\in Z\,
        \bigl(w=x\star(y\star z)\bigr)\vdash{}\\[-2pt]
      &w\in X\star_{\mathcal P}(Y\star_{\mathcal P}Z)
    \end{aligned}
  }{%
    \begin{aligned}[t]
      &\SemigroupAlg{A}{\star}\dsep X,Y,Z\in\PowNE{A}\dsep{}\\[-2pt]
      &\exists x\in X\,\exists y\in Y\,\exists z\in Z\,
        \bigl(w=x\star(y\star z)\bigr)\vdash{}\\[-2pt]
      &w\in X\star_{\mathcal P}(Y\star_{\mathcal P}Z)
    \end{aligned}
  }[PowerSemigroupRightTripleMembershipForward]
    \proofstepwidestar[1]{\SemigroupAlg{A}{\star}}{\rA}
    \proofstepwidestar[2]{X\in\PowNE{A}}{\rA}
    \proofstepwidestar[3]{Y\in\PowNE{A}}{\rA}
    \proofstepwidestar[4]{Z\in\PowNE{A}}{\rA}
    \proofstepwidestar[5]{%
      \exists x\in X\,\exists y\in Y\,\exists z\in Z\,
      \bigl(w=x\star(y\star z)\bigr)}{\rA}
    \proofstepwidestar[5]{x\in X}{\rA}
    \proofstepwidestar[5]{y\in Y}{\rA}
    \proofstepwidestar[5]{z\in Z}{\rA}
    \proofstepwidestar[5]{w=x\star(y\star z)}{\rA}
    \proofstepwidestar[1,3,4,7,8]{%
      y\star z\in Y\star_{\mathcal P}Z}{%
      \FormulaRefAuto{PowerSemigroupProductContains}[thm]{1,3,4,7,8}}
    \proofstepwidestar[1,3,4]{Y\star_{\mathcal P}Z\in\PowNE{A}}{%
      \FormulaRefAuto{PowerSemigroupProductClosure}[thm]{1,3,4}}
    \proofstepwidestar[1,2,3,4,6,7,8]{%
      x\star(y\star z)
      \in X\star_{\mathcal P}(Y\star_{\mathcal P}Z)}{%
      \FormulaRefAuto{PowerSemigroupProductContains}[thm]{1,2,11,6,10}}
    \proofstepwidestar[9]{x\star(y\star z)=w}{%
      \FormulaRefAuto{a=b\vdash b=a}[thm]{9}}
    \proofstepwidestar[1,2,3,4,5]{%
      w\in X\star_{\mathcal P}(Y\star_{\mathcal P}Z)}{\rIE{13,12}}
    \proofstepwidestar[1,2,3,4,5]{%
      w\in X\star_{\mathcal P}(Y\star_{\mathcal P}Z)}{%
      \rEEStar{5\text{--}14}}
  \closeproofpartwideindR

  \proofpartwideindR[Dreifachzeugen der rechts geklammerten Seite]{%
    \begin{aligned}[t]
      &\SemigroupAlg{A}{\star}\dsep X,Y,Z\in\PowNE{A}\dsep
        w\in X\star_{\mathcal P}(Y\star_{\mathcal P}Z)\vdash{}\\[-2pt]
      &\exists x\in X\,\exists y\in Y\,\exists z\in Z\,
        \bigl(w=x\star(y\star z)\bigr)
    \end{aligned}
  }{%
    \begin{aligned}[t]
      &\SemigroupAlg{A}{\star}\dsep X,Y,Z\in\PowNE{A}\dsep
        w\in X\star_{\mathcal P}(Y\star_{\mathcal P}Z)\vdash{}\\[-2pt]
      &\exists x\in X\,\exists y\in Y\,\exists z\in Z\,
        \bigl(w=x\star(y\star z)\bigr)
    \end{aligned}
  }[PowerSemigroupRightTripleMembershipBackward]
    \proofstepwidestar[1]{\SemigroupAlg{A}{\star}}{\rA}
    \proofstepwidestar[2]{X\in\PowNE{A}}{\rA}
    \proofstepwidestar[3]{Y\in\PowNE{A}}{\rA}
    \proofstepwidestar[4]{Z\in\PowNE{A}}{\rA}
    \proofstepwidestar[5]{%
      w\in X\star_{\mathcal P}(Y\star_{\mathcal P}Z)}{\rA}
    \proofstepwidestar[1,3,4]{Y\star_{\mathcal P}Z\in\PowNE{A}}{%
      \FormulaRefAuto{PowerSemigroupProductClosure}[thm]{1,3,4}}
    \proofstepwidestar[1,2,3,4,5]{%
      \exists x\in X\,
      \exists v\in Y\star_{\mathcal P}Z\,(w=x\star v)}{%
      \FormulaRefAuto{PowerSemigroupProductWitnesses}[thm]{1,2,6,5}}
    \proofstepwidestar[7]{x\in X}{\rA}
    \proofstepwidestar[7]{v\in Y\star_{\mathcal P}Z}{\rA}
    \proofstepwidestar[7]{w=x\star v}{\rA}
    \proofstepwidestar[1,3,4,9]{%
      \exists y\in Y\,\exists z\in Z\,(v=y\star z)}{%
      \FormulaRefAuto{PowerSemigroupProductWitnesses}[thm]{1,3,4,9}}
    \proofstepwidestar[11]{y\in Y}{\rA}
    \proofstepwidestar[11]{z\in Z}{\rA}
    \proofstepwidestar[11]{v=y\star z}{\rA}
    \proofstepwidestar[7,11]{w=x\star(y\star z)}{\rIE{14,10}}
    \proofstepwidestar[7,11]{%
      \exists x\in X\,\exists y\in Y\,\exists z\in Z\,
      \bigl(w=x\star(y\star z)\bigr)}{%
      \rEIStar{8,12,13,15}}
    \proofstepwidestar[1,2,3,4,5]{%
      \exists x\in X\,\exists y\in Y\,\exists z\in Z\,
      \bigl(w=x\star(y\star z)\bigr)}{%
      \rEEStar{7\text{--}16}}
  \closeproofpartwideindR

  \proofpartwideindR[Umklammern der Dreifachzeugen nach rechts]{%
    \begin{aligned}[t]
      &\SemigroupAlg{A}{\star}\dsep X,Y,Z\in\PowNE{A}\dsep{}\\[-2pt]
      &\exists x\in X\,\exists y\in Y\,\exists z\in Z\,
        \bigl(w=(x\star y)\star z\bigr)\vdash{}\\[-2pt]
      &\exists x\in X\,\exists y\in Y\,\exists z\in Z\,
        \bigl(w=x\star(y\star z)\bigr)
    \end{aligned}
  }{%
    \begin{aligned}[t]
      &\SemigroupAlg{A}{\star}\dsep X,Y,Z\in\PowNE{A}\dsep{}\\[-2pt]
      &\exists x\in X\,\exists y\in Y\,\exists z\in Z\,
        \bigl(w=(x\star y)\star z\bigr)\vdash{}\\[-2pt]
      &\exists x\in X\,\exists y\in Y\,\exists z\in Z\,
        \bigl(w=x\star(y\star z)\bigr)
    \end{aligned}
  }[PowerSemigroupTripleWitnessAssociativityForward]
    \proofstepwidestar[1]{\SemigroupAlg{A}{\star}}{\rA}
    \proofstepwidestar[2]{X\in\PowNE{A}}{\rA}
    \proofstepwidestar[3]{Y\in\PowNE{A}}{\rA}
    \proofstepwidestar[4]{Z\in\PowNE{A}}{\rA}
    \proofstepwidestar[5]{%
      \exists x\in X\,\exists y\in Y\,\exists z\in Z\,
      \bigl(w=(x\star y)\star z\bigr)}{\rA}
    \proofstepwidestar[5]{x\in X}{\rA}
    \proofstepwidestar[5]{y\in Y}{\rA}
    \proofstepwidestar[5]{z\in Z}{\rA}
    \proofstepwide[5]{w}{=}{(x\star y)\star z}{\rA}
    \proofstepwidestarbreak[1,2,3,4,6,7,8]{%
      (x\star y)\star z=x\star(y\star z)}{%
      \FormulaRefAuto{PowerSemigroupElementAssociativity}[thm]{%
        1,2,3,4,6,7,8}}
    \proofstepwidestar[1,2,3,4,5]{%
      w=x\star(y\star z)}{\rChain{9,10}}
    \proofstepwidestar[1,2,3,4,5]{%
      \exists x\in X\,\exists y\in Y\,\exists z\in Z\,
      \bigl(w=x\star(y\star z)\bigr)}{%
      \rEIStar{6,7,8,11}}
    \proofstepwidestar[1,2,3,4,5]{%
      \exists x\in X\,\exists y\in Y\,\exists z\in Z\,
      \bigl(w=x\star(y\star z)\bigr)}{%
      \rEEStar{5\text{--}12}}
  \closeproofpartwideindR

  \proofpartwideindR[Umklammern der Dreifachzeugen nach links]{%
    \begin{aligned}[t]
      &\SemigroupAlg{A}{\star}\dsep X,Y,Z\in\PowNE{A}\dsep{}\\[-2pt]
      &\exists x\in X\,\exists y\in Y\,\exists z\in Z\,
        \bigl(w=x\star(y\star z)\bigr)\vdash{}\\[-2pt]
      &\exists x\in X\,\exists y\in Y\,\exists z\in Z\,
        \bigl(w=(x\star y)\star z\bigr)
    \end{aligned}
  }{%
    \begin{aligned}[t]
      &\SemigroupAlg{A}{\star}\dsep X,Y,Z\in\PowNE{A}\dsep{}\\[-2pt]
      &\exists x\in X\,\exists y\in Y\,\exists z\in Z\,
        \bigl(w=x\star(y\star z)\bigr)\vdash{}\\[-2pt]
      &\exists x\in X\,\exists y\in Y\,\exists z\in Z\,
        \bigl(w=(x\star y)\star z\bigr)
    \end{aligned}
  }[PowerSemigroupTripleWitnessAssociativityBackward]
    \proofstepwidestar[1]{\SemigroupAlg{A}{\star}}{\rA}
    \proofstepwidestar[2]{X\in\PowNE{A}}{\rA}
    \proofstepwidestar[3]{Y\in\PowNE{A}}{\rA}
    \proofstepwidestar[4]{Z\in\PowNE{A}}{\rA}
    \proofstepwidestar[5]{%
      \exists x\in X\,\exists y\in Y\,\exists z\in Z\,
      \bigl(w=x\star(y\star z)\bigr)}{\rA}
    \proofstepwidestar[5]{x\in X}{\rA}
    \proofstepwidestar[5]{y\in Y}{\rA}
    \proofstepwidestar[5]{z\in Z}{\rA}
    \proofstepwidestarbreak[1,2,3,4,6,7,8]{%
      (x\star y)\star z=x\star(y\star z)}{%
      \FormulaRefAuto{PowerSemigroupElementAssociativity}[thm]{%
        1,2,3,4,6,7,8}}
    \proofstepwide[5]{w}{=}{x\star(y\star z)}{\rA}
    \proofstepwide[1,2,3,4,6,7,8]{}{=}{(x\star y)\star z}{%
      \FormulaRefAuto{a=b\vdash b=a}[thm]{9}}
    \proofstepwidestar[1,2,3,4,5]{%
      w=(x\star y)\star z}{\rChain{10,11}}
    \proofstepwidestar[1,2,3,4,5]{%
      \exists x\in X\,\exists y\in Y\,\exists z\in Z\,
      \bigl(w=(x\star y)\star z\bigr)}{%
      \rEIStar{6,7,8,12}}
    \proofstepwidestar[1,2,3,4,5]{%
      \exists x\in X\,\exists y\in Y\,\exists z\in Z\,
      \bigl(w=(x\star y)\star z\bigr)}{%
      \rEEStar{5\text{--}13}}
  \closeproofpartwideindR

  \proofpartwideindR[Zeugenform der links geklammerten Seite]{%
    \begin{aligned}[t]
      &\SemigroupAlg{A}{\star}\dsep X,Y,Z\in\PowNE{A}\vdash{}\\[-2pt]
      &w\in (X\star_{\mathcal P}Y)\star_{\mathcal P}Z
        \leftrightarrow
        \exists x\in X\,\exists y\in Y\,\exists z\in Z\,
        \bigl(w=(x\star y)\star z\bigr)
    \end{aligned}
  }{%
    \begin{aligned}[t]
      &\SemigroupAlg{A}{\star}\dsep X,Y,Z\in\PowNE{A}\vdash{}\\[-2pt]
      &w\in (X\star_{\mathcal P}Y)\star_{\mathcal P}Z
        \leftrightarrow
        \exists x\in X\,\exists y\in Y\,\exists z\in Z\,
        \bigl(w=(x\star y)\star z\bigr)
    \end{aligned}
  }[PowerSemigroupLeftTripleMembershipEquivalence]
    \proofstepwidestar[1]{\SemigroupAlg{A}{\star}}{\rA}
    \proofstepwidestar[2]{X\in\PowNE{A}}{\rA}
    \proofstepwidestar[3]{Y\in\PowNE{A}}{\rA}
    \proofstepwidestar[4]{Z\in\PowNE{A}}{\rA}
    \proofstepwidestar[1,2,3,4]{%
      \begin{aligned}[t]
        &w\in (X\star_{\mathcal P}Y)\star_{\mathcal P}Z
          \leftrightarrow{}\\[-2pt]
        &\exists x\in X\,\exists y\in Y\\[-2pt]
        &\quad\exists z\in Z\,\bigl(w=(x\star y)\star z\bigr)
      \end{aligned}}{%
      \rLRI{%
        \FormulaRefAuto{PowerSemigroupLeftTripleMembershipForward}[thm],
        \FormulaRefAuto{PowerSemigroupLeftTripleMembershipBackward}[thm]}}
  \closeproofpartwideindR

  \proofpartwideindR[Umklammern der Zeugen]{%
    \begin{aligned}[t]
      &\SemigroupAlg{A}{\star}\dsep X,Y,Z\in\PowNE{A}\vdash{}\\[-2pt]
      &\exists x\in X\,\exists y\in Y\,\exists z\in Z\,
        \bigl(w=(x\star y)\star z\bigr)
        \leftrightarrow
        \exists x\in X\,\exists y\in Y\,\exists z\in Z\,
        \bigl(w=x\star(y\star z)\bigr)
    \end{aligned}
  }{%
    \begin{aligned}[t]
      &\SemigroupAlg{A}{\star}\dsep X,Y,Z\in\PowNE{A}\vdash{}\\[-2pt]
      &\exists x\in X\,\exists y\in Y\,\exists z\in Z\,
        \bigl(w=(x\star y)\star z\bigr)
        \leftrightarrow
        \exists x\in X\,\exists y\in Y\,\exists z\in Z\,
        \bigl(w=x\star(y\star z)\bigr)
    \end{aligned}
  }[PowerSemigroupTripleWitnessAssociativityEquivalence]
    \proofstepwidestar[1]{\SemigroupAlg{A}{\star}}{\rA}
    \proofstepwidestar[2]{X\in\PowNE{A}}{\rA}
    \proofstepwidestar[3]{Y\in\PowNE{A}}{\rA}
    \proofstepwidestar[4]{Z\in\PowNE{A}}{\rA}
    \proofstepwidestar[1,2,3,4]{%
      \begin{aligned}[t]
        &\exists x\in X\,\exists y\in Y\\[-2pt]
        &\quad\exists z\in Z\,\bigl(w=(x\star y)\star z\bigr)\\[-2pt]
        &\leftrightarrow \exists x\in X\,\exists y\in Y\\[-2pt]
        &\quad\exists z\in Z\,\bigl(w=x\star(y\star z)\bigr)
      \end{aligned}}{%
      \rLRI{%
        \FormulaRefAuto{PowerSemigroupTripleWitnessAssociativityForward}[thm],
        \FormulaRefAuto{PowerSemigroupTripleWitnessAssociativityBackward}[thm]}}
  \closeproofpartwideindR

  \proofpartwideindR[Zeugenform der rechts geklammerten Seite]{%
    \begin{aligned}[t]
      &\SemigroupAlg{A}{\star}\dsep X,Y,Z\in\PowNE{A}\vdash{}\\[-2pt]
      &\exists x\in X\,\exists y\in Y\,\exists z\in Z\,
        \bigl(w=x\star(y\star z)\bigr)
        \leftrightarrow
        w\in X\star_{\mathcal P}(Y\star_{\mathcal P}Z)
    \end{aligned}
  }{%
    \begin{aligned}[t]
      &\SemigroupAlg{A}{\star}\dsep X,Y,Z\in\PowNE{A}\vdash{}\\[-2pt]
      &\exists x\in X\,\exists y\in Y\,\exists z\in Z\,
        \bigl(w=x\star(y\star z)\bigr)
        \leftrightarrow
        w\in X\star_{\mathcal P}(Y\star_{\mathcal P}Z)
    \end{aligned}
  }[PowerSemigroupRightTripleMembershipEquivalence]
    \proofstepwidestar[1]{\SemigroupAlg{A}{\star}}{\rA}
    \proofstepwidestar[2]{X\in\PowNE{A}}{\rA}
    \proofstepwidestar[3]{Y\in\PowNE{A}}{\rA}
    \proofstepwidestar[4]{Z\in\PowNE{A}}{\rA}
    \proofstepwidestar[1,2,3,4]{%
      \begin{aligned}[t]
        &\exists x\in X\,\exists y\in Y\\[-2pt]
        &\quad\exists z\in Z\,\bigl(w=x\star(y\star z)\bigr)
          \leftrightarrow{}\\[-2pt]
        &w\in X\star_{\mathcal P}(Y\star_{\mathcal P}Z)
      \end{aligned}}{%
      \rLRI{%
        \FormulaRefAuto{PowerSemigroupRightTripleMembershipForward}[thm],
        \FormulaRefAuto{PowerSemigroupRightTripleMembershipBackward}[thm]}}
  \closeproofpartwideindR

  \proofpartwideindR[Assoziativität der mengenweisen Operation]{%
    \begin{aligned}[t]
      &\SemigroupAlg{A}{\star}\dsep X,Y,Z\in\PowNE{A}\vdash{}\\[-2pt]
      &(X\star_{\mathcal P}Y)\star_{\mathcal P}Z
        =X\star_{\mathcal P}(Y\star_{\mathcal P}Z)
    \end{aligned}
  }{%
    \begin{aligned}[t]
      &\SemigroupAlg{A}{\star}\dsep X,Y,Z\in\PowNE{A}\vdash{}\\[-2pt]
      &(X\star_{\mathcal P}Y)\star_{\mathcal P}Z
        =X\star_{\mathcal P}(Y\star_{\mathcal P}Z)
    \end{aligned}
  }[PowerSemigroupAssociativity]
    \proofstepwidestar[1]{\SemigroupAlg{A}{\star}}{\rA}
    \proofstepwidestar[2]{X\in\PowNE{A}}{\rA}
    \proofstepwidestar[3]{Y\in\PowNE{A}}{\rA}
    \proofstepwidestar[4]{Z\in\PowNE{A}}{\rA}
    \proofstepwide[1,2,3,4]{%
      w\in (X\star_{\mathcal P}Y)\star_{\mathcal P}Z}{\leftrightarrow}{%
      \begin{aligned}[t]
        &\exists x\in X\,\exists y\in Y\\[-2pt]
        &\quad\exists z\in Z\,\bigl(w=(x\star y)\star z\bigr)
      \end{aligned}}{%
      \FormulaRefAuto{PowerSemigroupLeftTripleMembershipEquivalence}[thm]}
    \proofstepwide[1,2,3,4]{}{\leftrightarrow}{%
      \begin{aligned}[t]
        &\exists x\in X\,\exists y\in Y\\[-2pt]
        &\quad\exists z\in Z\,\bigl(w=x\star(y\star z)\bigr)
      \end{aligned}}{%
      \FormulaRefAuto{PowerSemigroupTripleWitnessAssociativityEquivalence}[thm]}
    \proofstepwide[1,2,3,4]{}{\leftrightarrow}{%
      w\in X\star_{\mathcal P}(Y\star_{\mathcal P}Z)}{%
      \FormulaRefAuto{PowerSemigroupRightTripleMembershipEquivalence}[thm]}
    \proofstepwide[1,2,3,4]{%
      w\in (X\star_{\mathcal P}Y)\star_{\mathcal P}Z}{\leftrightarrow}{%
      w\in X\star_{\mathcal P}(Y\star_{\mathcal P}Z)}{\rChain{5,7}}
    \proofstepwidestar[1,2,3,4]{%
      \forall w\,\Bigl(
        w\in (X\star_{\mathcal P}Y)\star_{\mathcal P}Z
        \leftrightarrow
        w\in X\star_{\mathcal P}(Y\star_{\mathcal P}Z)
      \Bigr)}{\rUI{8}}
    \proofstepwidestar[1,2,3,4]{%
      (X\star_{\mathcal P}Y)\star_{\mathcal P}Z
      =X\star_{\mathcal P}(Y\star_{\mathcal P}Z)}{%
      \FormulaRefAuto{\forall x\,(x\in A\leftrightarrow x\in B)
        \vdash A=B}[ax]{9}}
  \closeproofpartwideindR

  \proofpartwide{Schluss}
    \proofstepwidestar[1]{\SemigroupAlg{A}{\star}}{\rA}
    \proofstepwidestar[2]{X\in\PowNE{A}}{\rA}
    \proofstepwidestar[3]{Y\in\PowNE{A}}{\rA}
    \proofstepwidestar[1,2,3]{X\star_{\mathcal P}Y\in\PowNE{A}}{%
      \FormulaRefAuto{PowerSemigroupProductClosure}[thm]{1,2,3}}
    \proofstepwidestar[5]{Z\in\PowNE{A}}{\rA}
    \proofstepwidestar[1,2,3,5]{%
      (X\star_{\mathcal P}Y)\star_{\mathcal P}Z
      =X\star_{\mathcal P}(Y\star_{\mathcal P}Z)}{%
      \FormulaRefAuto{PowerSemigroupAssociativity}[thm]{1,2,3,5}}
    \proofstepwidestar[1]{\SemigroupAlg{\PowNE{A}}{\star_{\mathcal P}}}{%
      \FormulaRefAuto{SemigroupDef}[def]{4,6}}
  \closeproofpartwide
\end{tabproofsplitwide}

Die Halbgruppe \((\PowNE{A},\star_{\mathcal P})\) heißt die
\emph{Potenzhalbgruppe} von \((A,\star)\). Ihre Konstruktion setzt keine
Endlichkeit von \(A\) voraus.

\FormulaDefDeltaK[Einermengenkopie einer Halbgruppe]{%
  \begin{aligned}[t]
    &\SemigroupAlg{A}{\star}\vdash{}\\[-2pt]
    &\operatorname{Sing}(A)\coloneqq
      \bigl\{\{a\}\bigm|a\in A\bigr\},\\[-2pt]
    &\eta_A\colon A\to\operatorname{Sing}(A),
      \qquad \eta_A(a)\coloneqq\{a\}
  \end{aligned}%
}{SemigroupSingletonCopyDef}{%
  \DeltaRow{Mengen}{A\dsep a}
  \DeltaRow{Funktionen}{\star\dsep\eta_A}
  \DeltaRow{\textbf{Neue Symbole}}{}
  \DeltaPrem{Einermengenschicht}{\operatorname{Sing}(A)}
  \DeltaPrem{Einermengenabbildung}{\eta_A}
}

\section{Weitere Beispiele}

Auf jeder total geordneten Menge liefert Band~15 die binären Operationen
\(\min_{A,\leq}\) und \(\max_{A,\leq}\); siehe
\FormulaRefAuto{TotalOrderBinaryMinimum}[def] und
\FormulaRefAuto{TotalOrderBinaryMaximum}[def]. Ihre dort bewiesene
Funktionstypisierung und Assoziativität geben unmittelbar zwei Halbgruppen.

\FormulaThmDeltaK[Die Minimumsoperation bildet eine Halbgruppe]{%
  \TotOrd{A,\leq}\vdash\SemigroupAlg{A}{\min_{A,\leq}}%
}{TotalOrderMinimumSemigroup}{%
  \DeltaRow{Mengen}{A\dsep x\dsep y\dsep z}
  \DeltaRow{Relationsoperatoren}{\leq}
  \DeltaRow{Funktionen}{\min_{A,\leq}}
}
\begin{tabproofwide}
  \proofstepwidestar[1]{\TotOrd{A,\leq}}{\rA}
  \proofstepwidestar[1]{\min_{A,\leq}\colon A\times A\to A}{%
    \FormulaRefAuto{TotalOrderBinaryMinimumFunctionType}[thm]{1}}
  \proofstepwidestar[3]{x\in A}{\rA}
  \proofstepwidestar[4]{y\in A}{\rA}
  \proofstepwidestar[3,4]{(x,y)\in A\times A}{%
    \FormulaRefAuto{a\in A\dsep b\in B
      \vdash(a,b)\in A\times B}[thm]{3,4}}
  \proofstepwidestar[1,3,4]{\min_{A,\leq}(x,y)\in A}{%
    \FormulaRefAuto{F\colon A\to B\dsep x\in A
      \vdash F(x)\in B}[thm]{2,5}}
  \proofstepwidestar[7]{z\in A}{\rA}
  \proofstepwidestarbreak[1,3,4,7]{%
    \begin{aligned}[t]
      &\min_{A,\leq}(\min_{A,\leq}(x,y),z)\\[-2pt]
      &\qquad=\min_{A,\leq}(x,\min_{A,\leq}(y,z))
    \end{aligned}}{%
    \FormulaRefAuto{TotalOrderBinaryMinimumAssociative}[thm]{1,3,4,7}}
  \proofstepwidestar[1]{\SemigroupAlg{A}{\min_{A,\leq}}}{%
    \FormulaRefAuto{SemigroupDef}[def]{6,8}}
\end{tabproofwide}

\newpage

\FormulaThmDeltaK[Die Maximumsoperation bildet eine Halbgruppe]{%
  \TotOrd{A,\leq}\vdash\SemigroupAlg{A}{\max_{A,\leq}}%
}{TotalOrderMaximumSemigroup}{%
  \DeltaRow{Mengen}{A\dsep x\dsep y\dsep z}
  \DeltaRow{Relationsoperatoren}{\leq}
  \DeltaRow{Funktionen}{\max_{A,\leq}}
}
\begin{tabproofwide}
  \proofstepwidestar[1]{\TotOrd{A,\leq}}{\rA}
  \proofstepwidestar[1]{\max_{A,\leq}\colon A\times A\to A}{%
    \FormulaRefAuto{TotalOrderBinaryMaximumFunctionType}[thm]{1}}
  \proofstepwidestar[3]{x\in A}{\rA}
  \proofstepwidestar[4]{y\in A}{\rA}
  \proofstepwidestar[3,4]{(x,y)\in A\times A}{%
    \FormulaRefAuto{a\in A\dsep b\in B
      \vdash(a,b)\in A\times B}[thm]{3,4}}
  \proofstepwidestar[1,3,4]{\max_{A,\leq}(x,y)\in A}{%
    \FormulaRefAuto{F\colon A\to B\dsep x\in A
      \vdash F(x)\in B}[thm]{2,5}}
  \proofstepwidestar[7]{z\in A}{\rA}
  \proofstepwidestarbreak[1,3,4,7]{%
    \begin{aligned}[t]
      &\max_{A,\leq}(\max_{A,\leq}(x,y),z)\\[-2pt]
      &\qquad=\max_{A,\leq}(x,\max_{A,\leq}(y,z))
    \end{aligned}}{%
    \FormulaRefAuto{TotalOrderBinaryMaximumAssociative}[thm]{1,3,4,7}}
  \proofstepwidestar[1]{\SemigroupAlg{A}{\max_{A,\leq}}}{%
    \FormulaRefAuto{SemigroupDef}[def]{6,8}}
\end{tabproofwide}

\FormulaThmDeltaK[Die Addition bildet eine Halbgruppe]{%
  \SemigroupAlg{\mathbb{N}}{+}%
}{NaturalAdditionSemigroup}{%
  \DeltaRow{Mengen}{\mathbb{N}\dsep a\dsep b\dsep c}
}
\begin{tabproof}
  \proofstep{1}{a\in\mathbb{N}}{\rA}
  \proofstep{2}{b\in\mathbb{N}}{\rA}
  \proofstep{1,2}{a+b\in\mathbb{N}}{%
    \FormulaRefAuto{PeanoAddClosure}[thm]{1,2}}
  \proofstep{4}{c\in\mathbb{N}}{\rA}
  \proofstep{1,2,4}{(a+b)+c=a+(b+c)}{%
    \FormulaRefAuto{PeanoAddAssociative}[thm]{1,2,4}}
  \proofstep{}{\SemigroupAlg{\mathbb{N}}{+}}{%
    \FormulaRefAuto{SemigroupDef}[def]{3,5}}
\end{tabproof}

\newpage
\FormulaThmDeltaK[Die Multiplikation bildet eine Halbgruppe]{%
  \SemigroupAlg{\mathbb{N}}{\cdot}%
}{NaturalMultiplicationSemigroup}{%
  \DeltaRow{Mengen}{\mathbb{N}\dsep a\dsep b\dsep c}
}
\begin{tabproof}
  \proofstep{1}{a\in\mathbb{N}}{\rA}
  \proofstep{2}{b\in\mathbb{N}}{\rA}
  \proofstep{1,2}{a\cdot b\in\mathbb{N}}{%
    \FormulaRefAuto{PeanoMulClosure}[thm]{1,2}}
  \proofstep{4}{c\in\mathbb{N}}{\rA}
  \proofstep{1,2,4}{(a\cdot b)\cdot c=a\cdot(b\cdot c)}{%
    \FormulaRefAuto{PeanoMulAssociative}[thm]{1,2,4}}
  \proofstep{}{\SemigroupAlg{\mathbb{N}}{\cdot}}{%
    \FormulaRefAuto{SemigroupDef}[def]{3,5}}
\end{tabproof}

\section{Potenzen in Halbgruppen}

In einer Halbgruppe gibt es im Allgemeinen kein neutrales Element. Daher
werden zun\"achst nur \emph{positive} Potenzen eingef\"uhrt. Die Konstruktion
wird auf die Rekursionsabbildung aus dem Band \emph{Nat\"urliche Zahlen}
zur\"uckgef\"uhrt. Auf diese Weise ist die Potenzfolge eine wirkliche
Abbildung und nicht lediglich eine informelle Rekursionsvorschrift.

\subsection{Rechtstranslation und Potenzfolge}

\FormulaDefDeltaK[Rechtstranslation]{%
  \SemigroupAlg{A}{\star}\dsep a\in A\vdash
  \begin{aligned}[t]
    &\rho_{A,\star,a}\colon A\to A,\\[-2pt]
    &\forall x\in A\,
      \bigl(\rho_{A,\star,a}(x)\coloneqq x\star a\bigr)
  \end{aligned}%
}{SemigroupRightTranslationDef}{%
  \DeltaRow{Mengen}{A\dsep a\dsep x}
  \DeltaRow{Funktionen}{\star\dsep\rho_{A,\star,a}}
  \DeltaRow{\textbf{Neues Symbol}}{}
  \DeltaPrem{Rechtstranslation}{\rho_{A,\star,a}}
}
\begin{tabproof}
  \proofstep{1}{\SemigroupAlg{A}{\star}}{\rA}
  \proofstep{2}{a\in A}{\rA}
  \proofstep{3}{x\in A}{\rA}
  \proofstep{1,2,3}{x\star a\in A}{%
    \FormulaRefAuto{SemigroupClosureAxiom}[ax]{1,3,2}}
  \proofstep{1,2}{\forall x\in A\,x\star a\in A}{%
    \rUI{\rRI{3,4}}}
\end{tabproof}

\FormulaThmDeltaK[Die Rechtstranslation ist eine Abbildung]{%
  \SemigroupAlg{A}{\star}\dsep a\in A\vdash
  \rho_{A,\star,a}\colon A\to A%
}{SemigroupRightTranslationFunction}{%
  \DeltaRow{Mengen}{A\dsep a}
  \DeltaRow{Funktionen}{\star\dsep\rho_{A,\star,a}}
}
\begin{tabproofwide}
  \proofstepwidestar[1]{\SemigroupAlg{A}{\star}}{\rA}
  \proofstepwidestar[2]{a\in A}{\rA}
  \proofstepwidestar[1,2]{\rho_{A,\star,a}\colon A\to A}{%
    \hspace{0.8em}\FormulaRefAuto{SemigroupRightTranslationDef}[def]{1,2}}
\end{tabproofwide}

\FormulaThmDeltaK[Auswertung der Rechtstranslation]{%
  \SemigroupAlg{A}{\star}\dsep a,x\in A\vdash
  \rho_{A,\star,a}(x)=x\star a%
}{SemigroupRightTranslationEvaluation}{%
  \DeltaRow{Mengen}{A\dsep a\dsep x}
  \DeltaRow{Funktionen}{\star\dsep\rho_{A,\star,a}}
}
\begin{tabproofwide}
  \proofstepwidestar[1]{\SemigroupAlg{A}{\star}}{\rA}
  \proofstepwidestar[2]{a\in A}{\rA}
  \proofstepwidestar[3]{x\in A}{\rA}
  \proofstepwidestar[1,2]{%
    \forall u\in A\,\rho_{A,\star,a}(u)=u\star a}{%
    \hspace{0.8em}\FormulaRefAuto{SemigroupRightTranslationDef}[def]{1,2}}
  \proofstepwidestar[1,2,3]{\rho_{A,\star,a}(x)=x\star a}{%
    \rRE{\rUE{4},3}}
\end{tabproofwide}

\FormulaDefDeltaK[Potenzfolge eines Halbgruppenelements]{%
  \SemigroupAlg{A}{\star}\dsep a\in A\vdash
  \operatorname{pow}_{A,\star}(a)\coloneqq
  \RecFun{a}{\rho_{A,\star,a}}%
}{SemigroupPowerSequenceDef}{%
  \DeltaRow{Mengen}{A\dsep a}
  \DeltaRow{Funktionen}{\star\dsep\rho_{A,\star,a}\dsep
    \operatorname{pow}_{A,\star}}
  \DeltaRow{\textbf{Neues Symbol}}{}
  \DeltaPrem{Potenzfolge}{\operatorname{pow}_{A,\star}(a)}
}

\FormulaThmDeltaK[Die Potenzfolge ist eine Abbildung]{%
  \SemigroupAlg{A}{\star}\dsep a\in A\vdash
  \operatorname{pow}_{A,\star}(a)\colon\mathbb N\to A%
}{SemigroupPowerSequenceFunction}{%
  \DeltaRow{Mengen}{A\dsep a}
  \DeltaRow{Funktionen}{\star\dsep\operatorname{pow}_{A,\star}}
}
\begin{tabproofwide}
  \proofstepwidestar[1]{\SemigroupAlg{A}{\star}}{\rA}
  \proofstepwidestar[2]{a\in A}{\rA}
  \proofstepwidestar[1,2]{\rho_{A,\star,a}\colon A\to A}{%
    \FormulaRefAuto{SemigroupRightTranslationFunction}[thm]{1,2}}
  \proofstepwidestar[1,2]{%
    \RecFun{a}{\rho_{A,\star,a}}\colon\mathbb N\to A}{%
    \FormulaRefAuto{m_0\in M\dsep f\colon M\to M\vdash
      \RecFun{m_0}{f}\colon\mathbb{N}\to M}[thm]{2,3}}
  \proofstepwidestarbreakkeep[1,2]{%
    \operatorname{pow}_{A,\star}(a)=
    \RecFun{a}{\rho_{A,\star,a}}}{%
    \FormulaRefAuto{SemigroupPowerSequenceDef}[def]{1,2}}
  \proofstepwidestar[1,2]{%
    \operatorname{pow}_{A,\star}(a)\colon\mathbb N\to A}{\rIE{5,4}}
\end{tabproofwide}

\subsection{Positive Potenzen}

Ist \(n\in\mathbb N\) und \(n\neq0\), so ist
\(n=\Succ(\Pred(n))\). Der Vorgänger dient daher nur dazu, die bei \(0\)
beginnende Potenzfolge mit den bei \(1\) beginnenden Exponenten zu
synchronisieren.

\FormulaDefDeltaK[Positive Potenz eines Halbgruppenelements]{%
  \begin{aligned}[t]
    &\SemigroupAlg{A}{\star}\dsep a\in A\dsep
      n\in\mathbb N\dsep n\neq0\vdash{}\\[-2pt]
    &a^n\coloneqq
      \operatorname{pow}_{A,\star}(a)(\Pred(n))
  \end{aligned}
}{SemigroupPositivePowerDef}{%
  \DeltaRow{Mengen}{A\dsep a\dsep n}
  \DeltaRow{Funktionen}{\star\dsep\operatorname{pow}_{A,\star}}
  \DeltaRow{\textbf{Neues Symbol}}{}
  \DeltaPrem{Positive Potenz}{a^n}
}

Wegen \(n+1=\Succ(n)\) formulieren wir die Nachfolgeraussagen im Folgenden
mit der Addition auf den natürlichen Zahlen.  Diese Schreibweise passt
unmittelbar zu den anschließenden Potenzgesetzen.

\FormulaThmDeltaK[Nachfolgerindex der Potenzfolge]{%
  \SemigroupAlg{A}{\star}\dsep a\in A\dsep n\in\mathbb N\vdash
  a^{n+1}=\operatorname{pow}_{A,\star}(a)(n)%
}{SemigroupPowerSuccessorIndex}{%
  \DeltaRow{Mengen}{A\dsep a\dsep n}
  \DeltaRow{Funktionen}{\star\dsep\operatorname{pow}_{A,\star}}
}
\begin{tabproofwide}
  \proofstepwidestar[1]{\SemigroupAlg{A}{\star}}{\rA}
  \proofstepwidestar[2]{a\in A}{\rA}
  \proofstepwidestar[3]{n\in\mathbb N}{\rA}
  \proofstepwidestar[3]{n+1=\Succ(n)}{%
    \FormulaRefAuto{PeanoAddOneSucc}[thm]{3}}
  \proofstepwidestar[3]{n+1\in\mathbb N}{%
    \rIE{4,\FormulaRefAuto{PeanoSuccClosure}[ax]{3}}}
  \proofstepwidestar[3]{n+1\neq0}{%
    \rIE{4,\FormulaRefAuto{PeanoSuccNonzero}[ax]{3}}}
  \proofstepwidestar[1,2,3]{%
    a^{n+1}=\operatorname{pow}_{A,\star}(a)(\Pred(n+1))}{%
    \FormulaRefAuto{SemigroupPositivePowerDef}[def]{1,2,5,6}}
  \proofstepwidestar[3]{\Pred(n+1)=n}{%
    \rIE{4,\FormulaRefAuto{PeanoPredSucc}[thm]{3}}}
  \proofstepwidestarbreakkeep[1,2,3]{%
    \operatorname{pow}_{A,\star}(a)(\Pred(n+1))=
    \operatorname{pow}_{A,\star}(a)(n)}{\rIE{8}}
  \proofstepwidestar[1,2,3]{%
    a^{n+1}=\operatorname{pow}_{A,\star}(a)(n)}{%
    \FormulaRefAuto{a = b,\, b = c \vdash a = c}[thm]{7,9}}
\end{tabproofwide}

\FormulaThmDeltaK[Erste positive Potenz]{%
  \SemigroupAlg{A}{\star}\dsep a\in A\vdash a^1=a%
}{SemigroupPowerOne}{%
  \DeltaRow{Mengen}{A\dsep a}
  \DeltaRow{Funktionen}{\star\dsep\operatorname{pow}_{A,\star}}
}
\begin{tabproofwide}
  \proofstepwidestar[1]{\SemigroupAlg{A}{\star}}{\rA}
  \proofstepwidestar[2]{a\in A}{\rA}
  \proofstepwidestar[]{0\in\mathbb N}{\FormulaRefAuto{PeanoZero}[ax]}
  \proofstepwidestar[1,2]{%
    a^{0+1}=\operatorname{pow}_{A,\star}(a)(0)}{%
    \FormulaRefAuto{SemigroupPowerSuccessorIndex}[thm]{1,2,3}}
  \proofstepwidestar[1,2]{%
    \operatorname{pow}_{A,\star}(a)(0)=
    \RecFun{a}{\rho_{A,\star,a}}(0)}{%
    \FormulaRefAuto{SemigroupPowerSequenceDef}[def]{1,2}}
  \proofstepwidestar[1,2]{\rho_{A,\star,a}\colon A\to A}{%
    \FormulaRefAuto{SemigroupRightTranslationFunction}[thm]{1,2}}
  \proofstepwidestar[1,2]{\RecFun{a}{\rho_{A,\star,a}}(0)=a}{%
    \FormulaRefAuto{m_0\in M\dsep f\colon M\to M\vdash
      \RecFun{m_0}{f}(0)=m_0}[thm]{2,6}}
  \proofstepwidestar[1,2]{a^{0+1}=a}{%
    \FormulaRefAuto{a = b,\, b = c \vdash a = c}[thm]{%
      4,\FormulaRefAuto{a = b,\, b = c \vdash a = c}[thm]{5,7}}}
  \proofstepwidestar[]{0+1=1}{%
    \FormulaRefAuto{PeanoAddLeftZero}[thm]{%
      \FormulaRefAuto{PeanoOneInN}[thm]}}
  \proofstepwidestar[1,2]{a^1=a}{\rIE{9,8}}
\end{tabproofwide}

\FormulaThmDeltaK[Nachfolgerregel positiver Potenzen]{%
  \begin{aligned}[t]
    &\SemigroupAlg{A}{\star}\dsep a\in A\dsep
      n\in\mathbb N\dsep n\neq0\vdash{}\\[-2pt]
    &a^{n+1}=a^n\star a
  \end{aligned}
}{SemigroupPowerSuccessor}{%
  \DeltaRow{Mengen}{A\dsep a\dsep n}
  \DeltaRow{Funktionen}{\star\dsep\operatorname{pow}_{A,\star}}
}
\begin{tabproofwide}
  \proofstepwidestar[1]{\SemigroupAlg{A}{\star}}{\rA}
  \proofstepwidestar[2]{a\in A}{\rA}
  \proofstepwidestar[3]{n\in\mathbb N}{\rA}
  \proofstepwidestar[4]{n\neq0}{\rA}
  \proofstepwidestar[1,2,3]{%
    a^{n+1}=\operatorname{pow}_{A,\star}(a)(n)}{%
    \FormulaRefAuto{SemigroupPowerSuccessorIndex}[thm]{1,2,3}}
  \proofstepwidestar[1,2]{\rho_{A,\star,a}\colon A\to A}{%
    \FormulaRefAuto{SemigroupRightTranslationFunction}[thm]{1,2}}
  \proofstepwidestarbreak[1,2,3,4]{%
    \operatorname{pow}_{A,\star}(a)(n)=
    \RecFun{a}{\rho_{A,\star,a}}(n)}{%
    \FormulaRefAuto{SemigroupPowerSequenceDef}[def]{1,2}}
  \proofstepwidestarbreak[1,2,3,4]{%
    \RecFun{a}{\rho_{A,\star,a}}(n)=
    \rho_{A,\star,a}
      (\RecFun{a}{\rho_{A,\star,a}}(\Pred(n)))}{%
    \FormulaRefAuto{m_0\in M\dsep f\colon M\to M\dsep
      n\in\mathbb{N}\dsep n\neq0\vdash
      \RecFun{m_0}{f}(n)=f(\RecFun{m_0}{f}(\Pred(n)))}[thm]{2,6,3,4}}
  \proofstepwidestarbreak[1,2,3]{%
    \RecFun{a}{\rho_{A,\star,a}}(\Pred(n))\in A}{%
    \FormulaRefAuto{m_0\in M\dsep f\colon M\to M\dsep k\in\mathbb{N}
      \vdash\RecFun{m_0}{f}(k)\in M}[thm]{%
      2,6,\FormulaRefAuto{PeanoPredNatural}[thm]{3,4}}}
  \proofstepwidestarbreak[1,2,3]{%
    \rho_{A,\star,a}
      (\RecFun{a}{\rho_{A,\star,a}}(\Pred(n)))
    =\RecFun{a}{\rho_{A,\star,a}}(\Pred(n))\star a}{%
    \FormulaRefAuto{SemigroupRightTranslationEvaluation}[thm]{1,2,9}}
  \proofstepwidestarbreak[1,2,3,4]{%
    a^n=\operatorname{pow}_{A,\star}(a)(\Pred(n))}{%
    \FormulaRefAuto{SemigroupPositivePowerDef}[def]{1,2,3,4}}
  \proofstepwidestarbreak[1,2,3,4]{%
    a^n=\RecFun{a}{\rho_{A,\star,a}}(\Pred(n))}{%
    \rIE{\FormulaRefAuto{SemigroupPowerSequenceDef}[def]{1,2},11}}
  \proofstepwidestarbreak[1,2,3,4]{%
    \RecFun{a}{\rho_{A,\star,a}}(\Pred(n))\star a=a^n\star a}{%
    \rIE{\FormulaRefAuto{a=b\vdash b=a}[thm]{12}}}
  \proofstepwidestarbreak[1,2,3,4]{a^{n+1}=a^n\star a}{%
    \FormulaRefAuto{a = b,\, b = c \vdash a = c}[thm]{%
      5,\FormulaRefAuto{a = b,\, b = c \vdash a = c}[thm]{%
        7,\FormulaRefAuto{a = b,\, b = c \vdash a = c}[thm]{%
          8,\FormulaRefAuto{a = b,\, b = c \vdash a = c}[thm]{10,13}}}}}
\end{tabproofwide}

\FormulaThmDeltaK[Abgeschlossenheit positiver Potenzen]{%
  \SemigroupAlg{A}{\star}\dsep a\in A\dsep
  n\in\mathbb N\dsep n\neq0\vdash a^n\in A%
}{SemigroupPowerClosure}{%
  \DeltaRow{Mengen}{A\dsep a\dsep n}
  \DeltaRow{Funktionen}{\star\dsep\operatorname{pow}_{A,\star}}
}
\begin{tabproofwide}
  \proofstepwidestar[1]{\SemigroupAlg{A}{\star}}{\rA}
  \proofstepwidestar[2]{a\in A}{\rA}
  \proofstepwidestar[3]{n\in\mathbb N}{\rA}
  \proofstepwidestar[4]{n\neq0}{\rA}
  \proofstepwidestar[1,2]{%
    \operatorname{pow}_{A,\star}(a)\colon\mathbb N\to A}{%
    \FormulaRefAuto{SemigroupPowerSequenceFunction}[thm]{1,2}}
  \proofstepwidestar[3,4]{\Pred(n)\in\mathbb N}{%
    \FormulaRefAuto{PeanoPredNatural}[thm]{3,4}}
  \proofstepwidestar[1,2,3]{%
    \operatorname{pow}_{A,\star}(a)(\Pred(n))\in A}{%
    \FormulaRefAuto{F\colon A\to B\dsep x\in A\vdash F(x)\in B}[thm]{5,6}}
  \proofstepwidestar[1,2,3,4]{%
    a^n=\operatorname{pow}_{A,\star}(a)(\Pred(n))}{%
    \FormulaRefAuto{SemigroupPositivePowerDef}[def]{1,2,3,4}}
  \proofstepwidestar[1,2,3,4]{a^n\in A}{\rIE{8,7}}
\end{tabproofwide}

\subsection{Potenzgesetze}

Die folgenden Gesetze zeigen, dass die Rekursionsklammerung wegen der
Assoziativit\"at keine zus\"atzliche Information tr\"agt.

\FormulaThmDeltaK[Additionsgesetz positiver Potenzen]{%
  \begin{aligned}[t]
    &\SemigroupAlg{A}{\star}\dsep a\in A\dsep
      m,n\in\mathbb N\dsep m\neq0\dsep n\neq0\vdash{}\\[-2pt]
    &a^{m+n}=a^m\star a^n
  \end{aligned}
}{SemigroupPowerAdditionLaw}{%
  \DeltaRow{Mengen}{A\dsep a\dsep m\dsep n\dsep k}
  \DeltaRow{Funktionen}{\star}
}
\begin{tabproofsplitwide}
  \proofpartwideindR[Induktionsanfang]{%
    \begin{aligned}[t]
      &\SemigroupAlg{A}{\star}\dsep a\in A\dsep
        m\in\mathbb N\dsep m\neq0\vdash{}\\[-2pt]
      &a^{m+\Succ(0)}=a^m\star a^{\Succ(0)}
    \end{aligned}
  }{%
    \SemigroupAlg{A}{\star}\dsep a\in A\dsep
    m\in\mathbb N\dsep m\neq0\vdash
    a^{m+\Succ(0)}=a^m\star a^{\Succ(0)}
  }[SemigroupPowerAdditionBase]
    \proofstepwidestar[1]{\SemigroupAlg{A}{\star}}{\rA}
    \proofstepwidestar[2]{a\in A}{\rA}
    \proofstepwidestar[3]{m\in\mathbb N}{\rA}
    \proofstepwidestar[4]{m\neq0}{\rA}
    \proofstepwidestar[]{1=\Succ(0)}{\FormulaRefAuto{PeanoOne}[def]}
    \proofstepwidestar[1,2,3,4]{a^{m+1}=a^m\star a}{%
      \FormulaRefAuto{SemigroupPowerSuccessor}[thm]{1,2,3,4}}
    \proofstepwidestar[1,2]{a^1=a}{%
      \FormulaRefAuto{SemigroupPowerOne}[thm]{1,2}}
    \proofstepwidestar[1,2,3,4]{%
      a^m\star a=a^m\star a^1}{%
      \rIE{\FormulaRefAuto{a=b\vdash b=a}[thm]{7}}}
    \proofstepwidestar[1,2,3,4]{%
      a^{m+\Succ(0)}=a^m\star a^{\Succ(0)}}{%
      \rIE{5,\FormulaRefAuto{a = b,\, b = c \vdash a = c}[thm]{6,8}}}
  \closeproofpartwideindR

  \proofpartwideindR[Induktionsschritt]{%
    \begin{aligned}[t]
      &\SemigroupAlg{A}{\star}\dsep a\in A\dsep
        m,k\in\mathbb N\dsep m\neq0\dsep
        a^{m+\Succ(k)}=a^m\star a^{\Succ(k)}\vdash{}\\[-2pt]
      &a^{m+\Succ(\Succ(k))}=a^m\star a^{\Succ(\Succ(k))}
    \end{aligned}
  }{%
    \begin{aligned}[t]
      &\SemigroupAlg{A}{\star}\dsep a\in A\dsep
        m,k\in\mathbb N\dsep m\neq0\dsep
        a^{m+\Succ(k)}=a^m\star a^{\Succ(k)}\vdash{}\\[-2pt]
      &a^{m+\Succ(\Succ(k))}=a^m\star a^{\Succ(\Succ(k))}
    \end{aligned}
  }[SemigroupPowerAdditionStep]
    \proofstepwidestar[1]{\SemigroupAlg{A}{\star}}{\rA}
    \proofstepwidestar[2]{a\in A}{\rA}
    \proofstepwidestar[3]{m\in\mathbb N}{\rA}
    \proofstepwidestar[4]{k\in\mathbb N}{\rA}
    \proofstepwidestar[5]{m\neq0}{\rA}
    \proofstepwidestar[6]{a^{m+\Succ(k)}=a^m\star a^{\Succ(k)}}{\rA}
    \proofstepwidestar[4]{\Succ(k)\in\mathbb N}{%
      \FormulaRefAuto{PeanoSuccClosure}[ax]{4}}
    \proofstepwidestar[3,4]{m+\Succ(k)\in\mathbb N}{%
      \FormulaRefAuto{PeanoAddClosure}[thm]{3,7}}
    \proofstepwidestar[3,4]{m+\Succ(k)=\Succ(m+k)}{%
      \FormulaRefAuto{PeanoAddSuccRight}[thm]{3,4}}
    \proofstepwidestarbreak[3,4]{m+\Succ(k)\neq0}{%
      \rIE{\FormulaRefAuto{a=b\vdash b=a}[thm]{9},
        \FormulaRefAuto{PeanoSuccNonzero}[ax]{%
          \FormulaRefAuto{PeanoAddClosure}[thm]{3,4}}}}
    \proofstepwidestarbreak[3,4]{%
      m+\Succ(\Succ(k))=(m+\Succ(k))+1}{%
      \FormulaRefAuto{a = b,\, b = c \vdash a = c}[thm]{%
        \FormulaRefAuto{PeanoAddSuccRight}[thm]{3,7},
        \FormulaRefAuto{a=b\vdash b=a}[thm]{%
          \FormulaRefAuto{PeanoAddOneSucc}[thm]{8}}}}
    \proofstepwidestar[1,2,3,4]{%
      a^{(m+\Succ(k))+1}=a^{m+\Succ(k)}\star a}{%
      \FormulaRefAuto{SemigroupPowerSuccessor}[thm]{1,2,8,10}}
    \proofstepwidestar[1,2,3,4,5,6]{%
      a^{m+\Succ(k)}\star a
      =(a^m\star a^{\Succ(k)})\star a}{\rIE{6}}
    \proofstepwidestarbreak[1,2,3,4,5]{%
      (a^m\star a^{\Succ(k)})\star a
      =a^m\star(a^{\Succ(k)}\star a)}{%
      \FormulaRefAuto{SemigroupAssociativityAxiom}[ax]{%
        1,
        \FormulaRefAuto{SemigroupPowerClosure}[thm]{1,2,3,5},
        \FormulaRefAuto{SemigroupPowerClosure}[thm]{%
          1,2,7,\FormulaRefAuto{PeanoSuccNonzero}[ax]{4}},2}}
    \proofstepwidestarbreak[1,2,4]{%
      a^{\Succ(k)+1}=a^{\Succ(k)}\star a}{%
      \FormulaRefAuto{SemigroupPowerSuccessor}[thm]{%
        1,2,7,\FormulaRefAuto{PeanoSuccNonzero}[ax]{4}}}
    \proofstepwidestarbreak[1,2,3,4,5]{%
      a^m\star(a^{\Succ(k)}\star a)
      =a^m\star a^{\Succ(\Succ(k))}}{%
      \rIE{\FormulaRefAuto{PeanoAddOneSucc}[thm]{7},
        \FormulaRefAuto{a=b\vdash b=a}[thm]{15}}}
    \proofstepwidestarbreak[1,2,3,4,5,6]{%
      a^{m+\Succ(\Succ(k))}=a^m\star a^{\Succ(\Succ(k))}}{%
      \rIE{11,\FormulaRefAuto{a = b,\, b = c \vdash a = c}[thm]{%
        12,\FormulaRefAuto{a = b,\, b = c \vdash a = c}[thm]{%
          13,\FormulaRefAuto{a = b,\, b = c \vdash a = c}[thm]{14,16}}}}}
  \closeproofpartwideindR

  \proofpartwide{Induktionsschluss und positiver Exponent}
    \proofstepwidestar[1]{\SemigroupAlg{A}{\star}}{\rA}
    \proofstepwidestar[2]{a\in A}{\rA}
    \proofstepwidestar[3]{m\in\mathbb N}{\rA}
    \proofstepwidestar[4]{n\in\mathbb N}{\rA}
    \proofstepwidestar[5]{m\neq0}{\rA}
    \proofstepwidestar[6]{n\neq0}{\rA}
    \proofstepwidestarbreak[1,2,3,5]{%
      \forall k\in\mathbb N\,
      \bigl(a^{m+\Succ(k)}=a^m\star a^{\Succ(k)}\bigr)}{%
      \FormulaRefAuto{PeanoInductionRule}[ax]{%
        \FormulaRefAuto{SemigroupPowerAdditionBase}[thm]{1,2,3,5},
        \FormulaRefAuto{SemigroupPowerAdditionStep}[thm]{1,2,3,5}}}
    \proofstepwidestar[4,6]{n=\Succ(\Pred(n))}{%
      \FormulaRefAuto{PeanoPredNonzeroValue}[thm]{4,6}}
    \proofstepwidestar[4,6]{\Pred(n)\in\mathbb N}{%
      \FormulaRefAuto{PeanoPredNatural}[thm]{4,6}}
    \proofstepwidestar[1,2,3,4,5]{%
      a^{m+\Succ(\Pred(n))}=a^m\star a^{\Succ(\Pred(n))}}{%
      \rRE{9,\rUE{7}}}
    \proofstepwidestar[1,2,3,4,5,6]{a^{m+n}=a^m\star a^n}{\rIE{8,10}}
  \closeproofpartwide
\end{tabproofsplitwide}

\FormulaThmDeltaK[Potenzen desselben Elements kommutieren]{%
  \begin{aligned}[t]
    &\SemigroupAlg{A}{\star}\dsep a\in A\dsep
      m,n\in\mathbb N\dsep m\neq0\dsep n\neq0\vdash{}\\[-2pt]
    &a^m\star a^n=a^n\star a^m
  \end{aligned}
}{SemigroupPowersCommute}{%
  \DeltaRow{Mengen}{A\dsep a\dsep m\dsep n}
  \DeltaRow{Funktionen}{\star}
}
\begin{tabproofwide}
  \proofstepwidestar[1]{\SemigroupAlg{A}{\star}}{\rA}
  \proofstepwidestar[2]{a\in A}{\rA}
  \proofstepwidestar[3]{m\in\mathbb N}{\rA}
  \proofstepwidestar[4]{n\in\mathbb N}{\rA}
  \proofstepwidestar[5]{m\neq0}{\rA}
  \proofstepwidestar[6]{n\neq0}{\rA}
  \proofstepwide[1,2,3,4,5,6]{a^m\star a^n}{=}{a^{m+n}}{%
    \FormulaRefAuto{a=b\vdash b=a}[thm]{%
      \FormulaRefAuto{SemigroupPowerAdditionLaw}[thm]{1,2,3,4,5,6}}}
  \proofstepwide[3,4]{}{=}{a^{n+m}}{%
    \rIE{\FormulaRefAuto{PeanoAddCommutative}[thm]{3,4}}}
  \proofstepwide[1,2,3,4,5,6]{}{=}{a^n\star a^m}{%
    \FormulaRefAuto{SemigroupPowerAdditionLaw}[thm]{1,2,4,3,6,5}}
  \proofstepwide[1,2,3,4,5,6]{a^m\star a^n}{=}{a^n\star a^m}{\rChain{7,8,9}}
\end{tabproofwide}

\section{Idempotente Elemente}

Ein Element ist idempotent, wenn seine zweite Potenz bereits wieder das
Element selbst ist. F\"ur eine feste Halbgruppe sammeln wir alle solchen
Elemente in einer ausgezeichneten Teilmenge des Tr\"agers.

\FormulaDefDeltaK[Idempotentenmenge einer Halbgruppe]{%
  \SemigroupAlg{A}{\star}\vdash
  E_{A,\star}\coloneqq
  \bigl\{e\in A\bigm|e\star e=e\bigr\}%
}{SemigroupIdempotentSetDef}{%
  \DeltaRow{Mengen}{A\dsep e}
  \DeltaRow{Funktionen}{\star}
  \DeltaRow{\textbf{Neues Symbol}}{}
  \DeltaPrem{Idempotentenmenge}{E_{A,\star}}
}

\FormulaThmDeltaK[Elementkriterium f\"ur Idempotente]{%
  \SemigroupAlg{A}{\star}\vdash
  e\in E_{A,\star}\leftrightarrow
  \bigl(e\in A\land e\star e=e\bigr)%
}{SemigroupIdempotentMembership}{%
  \DeltaRow{Mengen}{A\dsep e}
  \DeltaRow{Funktionen}{\star}
}
\begin{tabproofwide}
  \proofstepwidestar[1]{\SemigroupAlg{A}{\star}}{\rA}
  \proofstepwidestar[1]{%
    E_{A,\star}=\{u\in A\mid u\star u=u\}}{%
    \FormulaRefAuto{SemigroupIdempotentSetDef}[def]{1}}
  \proofstepwidestar[1]{%
    e\in E_{A,\star}\leftrightarrow
    \bigl(e\in A\land e\star e=e\bigr)}{%
    \rIE{2,\FormulaRefAuto{x\in\{u\in A\mid P(u)\}\eqvdash
      x\in A\land P(x)}[thm]}}
\end{tabproofwide}

\FormulaThmDeltaK[Ein idempotentes Element besitzt nur eine positive Potenz]{%
  \begin{aligned}[t]
    &\SemigroupAlg{A}{\star}\dsep e\in E_{A,\star}\dsep
      n\in\mathbb N\dsep n\neq0\vdash{}\\[-2pt]
    &e^n=e
  \end{aligned}
}{SemigroupIdempotentPower}{%
  \DeltaRow{Mengen}{A\dsep e\dsep n\dsep k}
  \DeltaRow{Funktionen}{\star}
}
\begin{tabproofsplitwide}
  \proofpartwideindR[Induktionsanfang]{%
    \SemigroupAlg{A}{\star}\dsep e\in E_{A,\star}\vdash e^{\Succ(0)}=e
  }{%
    \SemigroupAlg{A}{\star}\dsep e\in E_{A,\star}\vdash e^{\Succ(0)}=e
  }[SemigroupIdempotentPowerBase]
    \proofstepwidestar[1]{\SemigroupAlg{A}{\star}}{\rA}
    \proofstepwidestar[2]{e\in E_{A,\star}}{\rA}
    \proofstepwidestar[1,2]{e\in A}{%
      \rAEa{\FormulaRefAuto{SemigroupIdempotentMembership}[thm]{1,2}}}
    \proofstepwidestar[1,2]{e^1=e}{%
      \FormulaRefAuto{SemigroupPowerOne}[thm]{1,3}}
    \proofstepwidestar[]{1=\Succ(0)}{\FormulaRefAuto{PeanoOne}[def]}
    \proofstepwidestar[1,2]{e^{\Succ(0)}=e}{\rIE{5,4}}
  \closeproofpartwideindR

  \proofpartwideindR[Induktionsschritt]{%
    \begin{aligned}[t]
      &\SemigroupAlg{A}{\star}\dsep e\in E_{A,\star}\dsep
        k\in\mathbb N\dsep e^{\Succ(k)}=e\vdash{}\\[-2pt]
      &e^{\Succ(\Succ(k))}=e
    \end{aligned}
  }{%
    \begin{aligned}[t]
      &\SemigroupAlg{A}{\star}\dsep e\in E_{A,\star}\dsep
        k\in\mathbb N\dsep e^{\Succ(k)}=e\vdash{}\\[-2pt]
      &e^{\Succ(\Succ(k))}=e
    \end{aligned}
  }[SemigroupIdempotentPowerStep]
    \proofstepwidestar[1]{\SemigroupAlg{A}{\star}}{\rA}
    \proofstepwidestar[2]{e\in E_{A,\star}}{\rA}
    \proofstepwidestar[3]{k\in\mathbb N}{\rA}
    \proofstepwidestar[4]{e^{\Succ(k)}=e}{\rA}
    \proofstepwidestar[1,2]{e\in A\land e\star e=e}{%
      \FormulaRefAuto{SemigroupIdempotentMembership}[thm]{1,2}}
    \proofstepwidestar[1,2]{e\in A}{\rAEa{5}}
    \proofstepwidestar[1,2]{e\star e=e}{\rAEb{5}}
    \proofstepwidestar[3]{\Succ(k)\in\mathbb N}{%
      \FormulaRefAuto{PeanoSuccClosure}[ax]{3}}
    \proofstepwidestar[3]{\Succ(k)\neq0}{%
      \FormulaRefAuto{PeanoSuccNonzero}[ax]{3}}
    \proofstepwidestar[1,2,3]{%
      e^{\Succ(k)+1}=e^{\Succ(k)}\star e}{%
      \FormulaRefAuto{SemigroupPowerSuccessor}[thm]{1,6,8,9}}
    \proofstepwidestar[1,2,3,4]{e^{\Succ(k)}\star e=e\star e}{\rIE{4}}
    \proofstepwidestar[3]{\Succ(k)+1=\Succ(\Succ(k))}{%
      \FormulaRefAuto{PeanoAddOneSucc}[thm]{8}}
    \proofstepwidestar[1,2,3,4]{e^{\Succ(\Succ(k))}=e}{%
      \rIE{12,\FormulaRefAuto{a = b,\, b = c \vdash a = c}[thm]{%
        10,\FormulaRefAuto{a = b,\, b = c \vdash a = c}[thm]{11,7}}}}
  \closeproofpartwideindR

  \proofpartwide{Induktionsschluss}
    \proofstepwidestar[1]{\SemigroupAlg{A}{\star}}{\rA}
    \proofstepwidestar[2]{e\in E_{A,\star}}{\rA}
    \proofstepwidestar[3]{n\in\mathbb N}{\rA}
    \proofstepwidestar[4]{n\neq0}{\rA}
    \proofstepwidestar[1,2]{%
      \forall k\in\mathbb N\,\bigl(e^{\Succ(k)}=e\bigr)}{%
      \FormulaRefAuto{PeanoInductionRule}[ax]{%
        \FormulaRefAuto{SemigroupIdempotentPowerBase}[thm]{1,2},
        \FormulaRefAuto{SemigroupIdempotentPowerStep}[thm]{1,2}}}
    \proofstepwidestar[3,4]{n=\Succ(\Pred(n))}{%
      \FormulaRefAuto{PeanoPredNonzeroValue}[thm]{3,4}}
    \proofstepwidestar[3,4]{\Pred(n)\in\mathbb N}{%
      \FormulaRefAuto{PeanoPredNatural}[thm]{3,4}}
    \proofstepwidestar[1,2,3]{e^{\Succ(\Pred(n))}=e}{\rRE{7,\rUE{5}}}
    \proofstepwidestar[1,2,3,4]{e^n=e}{\rIE{6,8}}
  \closeproofpartwide
\end{tabproofsplitwide}

\section{Unterhalbgruppen}

Eine Unterhalbgruppe erbt die Assoziativit\"at von der umgebenden
Halbgruppe. Deshalb gen\"ugen Nichtleerheit und Abgeschlossenheit unter der
eingeschr\"ankten Operation.

Eine \emph{Unterhalbgruppe} ist dabei noch keine Gruppe: Gemeint ist nur
eine nichtleere, produktabgeschlossene Teilmenge. Echte Untergruppen mit
neutralem Element und Inversen werden erst nach Einführung des
Gruppenbegriffs behandelt.

Für ihren Träger verwenden wir die in der Standardliteratur übliche
Bezeichnung (U); das Prädikat schreiben wir als
\(\mathsf{UHGr}(U;A,\star)\).

\FormulaDefDeltaK[Unterhalbgruppe]{%
  \begin{aligned}[t]
    &\SemigroupAlg{A}{\star}\vdash{}\\[-2pt]
    &\mathsf{UHGr}(U;A,\star)\coloneqq
      U\subseteq A\land U\neq\varnothing\land
      \forall x\in U\,\forall y\in U\,(x\star y\in U)
  \end{aligned}
}{SemigroupSubsemigroupDef}{%
  \DeltaRow{Mengen}{A\dsep U\dsep x\dsep y}
  \DeltaRow{Funktionen}{\star}
  \DeltaRow{\textbf{Neues Symbol}}{}
  \DeltaPrem{Unterhalbgruppe}{\mathsf{UHGr}(U;A,\star)}
}

\FormulaThmDeltaK[Tr\"agermenge einer Unterhalbgruppe]{%
  \SemigroupAlg{A}{\star}\dsep\mathsf{UHGr}(U;A,\star)
  \vdash U\subseteq A%
}{SemigroupSubsemigroupSubset}{%
  \DeltaRow{Mengen}{A\dsep U}
  \DeltaRow{Funktionen}{\star}
}
\begin{tabproofwide}
  \proofstepwidestar[1]{\SemigroupAlg{A}{\star}}{\rA}
  \proofstepwidestar[2]{\mathsf{UHGr}(U;A,\star)}{\rA}
  \proofstepwidestar[1,2]{%
    U\subseteq A\land U\neq\varnothing\land
    \forall x\in U\,\forall y\in U\,(x\star y\in U)}{%
    \FormulaRefAuto{SemigroupSubsemigroupDef}[def]{1,2}}
  \proofstepwidestar[1,2]{U\subseteq A}{\rAEa{3}}
\end{tabproofwide}

\FormulaThmDeltaK[Nichtleerheit einer Unterhalbgruppe]{%
  \SemigroupAlg{A}{\star}\dsep\mathsf{UHGr}(U;A,\star)
  \vdash U\neq\varnothing%
}{SemigroupSubsemigroupNonempty}{%
  \DeltaRow{Mengen}{A\dsep U}
  \DeltaRow{Funktionen}{\star}
}
\begin{tabproofwide}
  \proofstepwidestar[1]{\SemigroupAlg{A}{\star}}{\rA}
  \proofstepwidestar[2]{\mathsf{UHGr}(U;A,\star)}{\rA}
  \proofstepwidestar[1,2]{%
    U\subseteq A\land U\neq\varnothing\land
    \forall x\in U\,\forall y\in U\,(x\star y\in U)}{%
    \FormulaRefAuto{SemigroupSubsemigroupDef}[def]{1,2}}
  \proofstepwidestar[1,2]{%
    U\neq\varnothing\land
    \forall x\in U\,\forall y\in U\,(x\star y\in U)}{\rAEb{3}}
  \proofstepwidestar[1,2]{U\neq\varnothing}{\rAEa{4}}
\end{tabproofwide}

\FormulaThmDeltaK[Produktabschluss einer Unterhalbgruppe]{%
  \begin{aligned}[t]
    &\SemigroupAlg{A}{\star}\dsep\mathsf{UHGr}(U;A,\star)\dsep
      x,y\in U\vdash{}\\[-2pt]
    &x\star y\in U
  \end{aligned}
}{SemigroupSubsemigroupClosure}{%
  \DeltaRow{Mengen}{A\dsep U\dsep x\dsep y}
  \DeltaRow{Funktionen}{\star}
}
\begin{tabproofwide}
  \proofstepwidestar[1]{\SemigroupAlg{A}{\star}}{\rA}
  \proofstepwidestar[2]{\mathsf{UHGr}(U;A,\star)}{\rA}
  \proofstepwidestar[3]{x\in U}{\rA}
  \proofstepwidestar[4]{y\in U}{\rA}
  \proofstepwidestar[1,2]{%
    U\subseteq A\land U\neq\varnothing\land
    \forall u\in U\,\forall v\in U\,(u\star v\in U)}{%
    \FormulaRefAuto{SemigroupSubsemigroupDef}[def]{1,2}}
  \proofstepwidestar[1,2]{%
    \forall u\in U\,\forall v\in U\,(u\star v\in U)}{\rAEb{\rAEb{5}}}
  \proofstepwidestar[1,2,3,4]{x\star y\in U}{\rRE{4,\rUE{\rRE{3,\rUE{6}}}}}
\end{tabproofwide}

\FormulaThmDeltaK[Eine Unterhalbgruppe ist eine Halbgruppe]{%
  \SemigroupAlg{A}{\star}\dsep\mathsf{UHGr}(U;A,\star)
  \vdash\SemigroupAlg{U}{\star}%
}{SubsemigroupFormsSemigroup}{%
  \DeltaRow{Mengen}{A\dsep U\dsep x\dsep y\dsep z}
  \DeltaRow{Funktionen}{\star}
}
\begin{tabproofsplitwide}
  \proofpartwideindR[Abgeschlossenheit]{%
    \SemigroupAlg{A}{\star}\dsep\mathsf{UHGr}(U;A,\star)\dsep
    x,y\in U\vdash x\star y\in U
  }{%
    \SemigroupAlg{A}{\star}\dsep\mathsf{UHGr}(U;A,\star)\dsep
    x,y\in U\vdash x\star y\in U
  }[SubsemigroupClosurePart]
    \proofstepwidestar[1]{\SemigroupAlg{A}{\star}}{\rA}
    \proofstepwidestar[2]{\mathsf{UHGr}(U;A,\star)}{\rA}
    \proofstepwidestar[3]{x\in U}{\rA}
    \proofstepwidestar[4]{y\in U}{\rA}
    \proofstepwidestar[1,2,3,4]{x\star y\in U}{%
      \FormulaRefAuto{SemigroupSubsemigroupClosure}[thm]{1,2,3,4}}
  \closeproofpartwideindR

  \proofpartwideindR[Vererbte Assoziativit\"at]{%
    \begin{aligned}[t]
      &\SemigroupAlg{A}{\star}\dsep\mathsf{UHGr}(U;A,\star)\dsep
        x,y,z\in U\vdash{}\\[-2pt]
      &(x\star y)\star z=x\star(y\star z)
    \end{aligned}
  }{%
    \begin{aligned}[t]
      &\SemigroupAlg{A}{\star}\dsep\mathsf{UHGr}(U;A,\star)\dsep
        x,y,z\in U\vdash{}\\[-2pt]
      &(x\star y)\star z=x\star(y\star z)
    \end{aligned}
  }[SubsemigroupAssociativityPart]
    \proofstepwidestar[1]{\SemigroupAlg{A}{\star}}{\rA}
    \proofstepwidestar[2]{\mathsf{UHGr}(U;A,\star)}{\rA}
    \proofstepwidestar[3]{x\in U}{\rA}
    \proofstepwidestar[4]{y\in U}{\rA}
    \proofstepwidestar[5]{z\in U}{\rA}
    \proofstepwidestar[1,2]{U\subseteq A}{%
      \FormulaRefAuto{SemigroupSubsemigroupSubset}[thm]{1,2}}
    \proofstepwidestar[1,2,3]{x\in A}{%
      \FormulaRefAuto{A\subseteq B,\,x\in A\vdash x\in B}[thm]{6,3}}
    \proofstepwidestar[1,2,4]{y\in A}{%
      \FormulaRefAuto{A\subseteq B,\,x\in A\vdash x\in B}[thm]{6,4}}
    \proofstepwidestar[1,2,5]{z\in A}{%
      \FormulaRefAuto{A\subseteq B,\,x\in A\vdash x\in B}[thm]{6,5}}
    \proofstepwidestar[1,2,3,4,5]{%
      (x\star y)\star z=x\star(y\star z)}{%
      \FormulaRefAuto{SemigroupAssociativityAxiom}[ax]{1,7,8,9}}
  \closeproofpartwideindR

  \proofpartwide{Schluss}
    \proofstepwidestar[1]{\SemigroupAlg{A}{\star}}{\rA}
    \proofstepwidestar[2]{\mathsf{UHGr}(U;A,\star)}{\rA}
    \proofstepwidestar[1,2]{\SemigroupAlg{U}{\star}}{%
      \FormulaRefAuto{SemigroupDef}[def]{%
        \FormulaRefAuto{SubsemigroupClosurePart}[thm]{1,2},
        \FormulaRefAuto{SubsemigroupAssociativityPart}[thm]{1,2}}}
  \closeproofpartwide
\end{tabproofsplitwide}

\FormulaThmDeltaK[Der Tr\"ager ist eine Unterhalbgruppe]{%
  \SemigroupAlg{A}{\star}\dsep A\neq\varnothing
  \vdash\mathsf{UHGr}(A;A,\star)%
}{SemigroupCarrierIsSubsemigroup}{%
  \DeltaRow{Mengen}{A\dsep x\dsep y}
  \DeltaRow{Funktionen}{\star}
}
\begin{tabproofwide}
  \proofstepwidestar[1]{\SemigroupAlg{A}{\star}}{\rA}
  \proofstepwidestar[2]{A\neq\varnothing}{\rA}
  \proofstepwidestar[]{A\subseteq A}{%
    \FormulaRefAuto{A\subseteq A}[thm]}
  \proofstepwidestar[4]{x\in A}{\rA}
  \proofstepwidestar[5]{y\in A}{\rA}
  \proofstepwidestar[1,4,5]{x\star y\in A}{%
    \FormulaRefAuto{SemigroupClosureAxiom}[ax]{1,4,5}}
  \proofstepwidestar[1]{%
    \forall x\in A\,\forall y\in A\,(x\star y\in A)}{%
    \rUI{\rRI{4,\rUI{\rRI{5,6}}}}}
  \proofstepwidestar[1,2]{%
    A\subseteq A\land A\neq\varnothing\land
    \forall x\in A\,\forall y\in A\,(x\star y\in A)}{%
    \rAI{3,\rAI{2,7}}}
  \proofstepwidestar[1,2]{\mathsf{UHGr}(A;A,\star)}{%
    \FormulaRefAuto{SemigroupSubsemigroupDef}[def]{1,8}}
\end{tabproofwide}

\subsection{Erzeugte Unterhalbgruppen}

\FormulaDefDeltaK[Von einer nichtleeren Menge erzeugte Unterhalbgruppe]{%
  \begin{aligned}[t]
    &\SemigroupAlg{A}{\star}\dsep X\subseteq A\dsep
      X\neq\varnothing\vdash{}\\[-2pt]
    &\langle X\rangle_{A,\star}\coloneqq
      \bigcap_{\mathsf{UHGr}(U;A,\star)\land X\subseteq U}U
  \end{aligned}
}{SemigroupGeneratedSubsemigroupDef}{%
  \DeltaRow{Mengen}{A\dsep U\dsep X}
  \DeltaRow{Funktionen}{\star}
  \DeltaRow{\textbf{Neues Symbol}}{}
  \DeltaPrem{Erzeugte Unterhalbgruppe}{\langle X\rangle_{A,\star}}
}

\FormulaThmDeltaK[Der Tr\"ager enth\"alt die Erzeuger]{%
  \begin{aligned}[t]
    &\SemigroupAlg{A}{\star}\dsep X\subseteq A\dsep
      X\neq\varnothing\vdash{}\\[-2pt]
    &X\subseteq\langle X\rangle_{A,\star}
  \end{aligned}
}{GeneratedSubsemigroupContainsGenerators}{%
  \DeltaRow{Mengen}{A\dsep U\dsep X\dsep x\dsep u}
  \DeltaRow{Funktionen}{\star}
}
\begin{tabproofwide}
  \proofstepwidestar[1]{\SemigroupAlg{A}{\star}}{\rA}
  \proofstepwidestar[2]{X\subseteq A}{\rA}
  \proofstepwidestar[3]{X\neq\varnothing}{\rA}
  \proofstepwidestar[3]{\exists u\,(u\in X)}{%
    \FormulaRefAuto{A\neq\varnothing\eqvdash\exists x\,(x\in A)}[thm]{3}}
  \proofstepwidestar[5]{u\in X}{\rA}
  \proofstepwidestar[2,5]{u\in A}{%
    \FormulaRefAuto{A\subseteq B,\,x\in A\vdash x\in B}[thm]{2,5}}
  \proofstepwidestar[2,5]{A\neq\varnothing}{%
    \FormulaRefAuto{a\in A\vdash A\neq\varnothing}[thm]{6}}
  \proofstepwidestar[2,3]{A\neq\varnothing}{\rEE{4,5,7}}
  \proofstepwidestar[1,2,3]{\mathsf{UHGr}(A;A,\star)}{%
    \FormulaRefAuto{SemigroupCarrierIsSubsemigroup}[thm]{1,8}}
  \proofstepwidestar[1,2,3]{%
    \mathsf{UHGr}(A;A,\star)\land X\subseteq A}{\rAI{9,2}}
  \proofstepwidestar[11]{x\in X}{\rA}
  \proofstepwidestar[12]{%
    \mathsf{UHGr}(U;A,\star)\land X\subseteq U}{\rA}
  \proofstepwidestar[12]{X\subseteq U}{\rAEb{12}}
  \proofstepwidestar[11,12]{x\in U}{%
    \FormulaRefAuto{A\subseteq B,\,x\in A\vdash x\in B}[thm]{13,11}}
  \proofstepwidestar[11]{%
    \forall U\,\Bigl(
      \mathsf{UHGr}(U;A,\star)\land X\subseteq U
      \rightarrow x\in U\Bigr)}{\rUI{\rRI{12,14}}}
  \proofstepwidestar[1,2,3,11]{%
    x\in\bigcap_{\mathsf{UHGr}(U;A,\star)\land X\subseteq U}U}{%
    \FormulaRefAuto{P(A)\vdash x\in\bigcap_{P(B)}B
      \leftrightarrow\forall C\,(P(C)\rightarrow x\in C)}[thm]{10,15}}
  \proofstepwidestar[1,2,3]{%
    \langle X\rangle_{A,\star}=
    \bigcap_{\mathsf{UHGr}(U;A,\star)\land X\subseteq U}U}{%
    \FormulaRefAuto{SemigroupGeneratedSubsemigroupDef}[def]{1,2,3}}
  \proofstepwidestar[1,2,3,11]{x\in\langle X\rangle_{A,\star}}{%
    \rIE{\FormulaRefAuto{a=b\vdash b=a}[thm]{17},16}}
  \proofstepwidestar[1,2,3]{X\subseteq\langle X\rangle_{A,\star}}{%
    \FormulaRefAuto{A\subseteq B\coloneq
      \forall x\,(x\in A\rightarrow x\in B)}[def]{\rUI{\rRI{11,18}}}}
\end{tabproofwide}

\FormulaThmDeltaK[Minimalit\"at der erzeugten Unterhalbgruppe]{%
  \begin{aligned}[t]
    &\SemigroupAlg{A}{\star}\dsep X\subseteq A\dsep X\neq\varnothing\dsep
      \mathsf{UHGr}(U;A,\star)\dsep X\subseteq U\vdash{}\\[-2pt]
    &\langle X\rangle_{A,\star}\subseteq U
  \end{aligned}
}{GeneratedSubsemigroupMinimal}{%
  \DeltaRow{Mengen}{A\dsep U\dsep V\dsep X}
  \DeltaRow{Funktionen}{\star}
}
\begin{tabproofwide}
  \proofstepwidestar[1]{\SemigroupAlg{A}{\star}}{\rA}
  \proofstepwidestar[2]{X\subseteq A}{\rA}
  \proofstepwidestar[3]{X\neq\varnothing}{\rA}
  \proofstepwidestar[4]{\mathsf{UHGr}(U;A,\star)}{\rA}
  \proofstepwidestar[5]{X\subseteq U}{\rA}
  \proofstepwidestar[4,5]{%
    \mathsf{UHGr}(U;A,\star)\land X\subseteq U}{\rAI{4,5}}
  \proofstepwidestar[4,5]{%
    \bigcap_{\mathsf{UHGr}(V;A,\star)\land X\subseteq V}V
    \subseteq U}{%
    \FormulaRefAuto{P(C)\vdash\bigcap_{P(A)}A\subseteq C}[thm]{6}}
  \proofstepwidestar[1,2,3]{%
    \langle X\rangle_{A,\star}=
    \bigcap_{\mathsf{UHGr}(V;A,\star)\land X\subseteq V}V}{%
    \FormulaRefAuto{SemigroupGeneratedSubsemigroupDef}[def]{1,2,3}}
  \proofstepwidestar[1,2,3,4,5]{\langle X\rangle_{A,\star}\subseteq U}{%
    \rIE{8,7}}
\end{tabproofwide}

\FormulaThmDeltaK[Produktabschluss der erzeugten Unterhalbgruppe]{%
  \begin{aligned}[t]
    &\SemigroupAlg{A}{\star}\dsep X\subseteq A\dsep X\neq\varnothing\dsep
      x,y\in\langle X\rangle_{A,\star}\vdash{}\\[-2pt]
    &x\star y\in\langle X\rangle_{A,\star}
  \end{aligned}
}{GeneratedSubsemigroupClosure}{%
  \DeltaRow{Mengen}{A\dsep U\dsep X\dsep x\dsep y\dsep u}
  \DeltaRow{Funktionen}{\star}
}
\begin{tabproofwide}
  \proofstepwidestar[1]{\SemigroupAlg{A}{\star}}{\rA}
  \proofstepwidestar[2]{X\subseteq A}{\rA}
  \proofstepwidestar[3]{X\neq\varnothing}{\rA}
  \proofstepwidestar[4]{x\in\langle X\rangle_{A,\star}}{\rA}
  \proofstepwidestar[5]{y\in\langle X\rangle_{A,\star}}{\rA}
  \proofstepwidestar[3]{\exists u\,(u\in X)}{%
    \FormulaRefAuto{A\neq\varnothing\eqvdash\exists u\,(u\in A)}[thm]{3}}
  \proofstepwidestar[7]{u\in X}{\rA}
  \proofstepwidestar[2,7]{u\in A}{%
    \FormulaRefAuto{A\subseteq B,\,x\in A\vdash x\in B}[thm]{2,7}}
  \proofstepwidestar[2,7]{A\neq\varnothing}{%
    \FormulaRefAuto{a\in A\vdash A\neq\varnothing}[thm]{8}}
  \proofstepwidestar[2,3]{A\neq\varnothing}{\rEE{6,7,9}}
  \proofstepwidestar[1,2,3]{\mathsf{UHGr}(A;A,\star)}{%
    \FormulaRefAuto{SemigroupCarrierIsSubsemigroup}[thm]{1,10}}
  \proofstepwidestar[1,2,3]{%
    \mathsf{UHGr}(A;A,\star)\land X\subseteq A}{\rAI{11,2}}
  \proofstepwidestar[1,2,3]{%
    \langle X\rangle_{A,\star}=
    \bigcap_{\mathsf{UHGr}(U;A,\star)\land X\subseteq U}U}{%
    \FormulaRefAuto{SemigroupGeneratedSubsemigroupDef}[def]{1,2,3}}
  \proofstepwidestar[1,2,3,4]{%
    x\in\bigcap_{\mathsf{UHGr}(U;A,\star)\land X\subseteq U}U}{\rIE{13,4}}
  \proofstepwidestar[1,2,3,5]{%
    y\in\bigcap_{\mathsf{UHGr}(U;A,\star)\land X\subseteq U}U}{\rIE{13,5}}
  \proofstepwidestar[1,2,3,4]{%
    \forall U\,\Bigl(
      \mathsf{UHGr}(U;A,\star)\land X\subseteq U
      \rightarrow x\in U\Bigr)}{%
    \FormulaRefAuto{P(A)\vdash x\in\bigcap_{P(B)}B
      \leftrightarrow\forall C\,(P(C)\rightarrow x\in C)}[thm]{12,14}}
  \proofstepwidestar[1,2,3,5]{%
    \forall U\,\Bigl(
      \mathsf{UHGr}(U;A,\star)\land X\subseteq U
      \rightarrow y\in U\Bigr)}{%
    \FormulaRefAuto{P(A)\vdash x\in\bigcap_{P(B)}B
      \leftrightarrow\forall C\,(P(C)\rightarrow x\in C)}[thm]{12,15}}
  \proofstepwidestar[18]{%
    \mathsf{UHGr}(U;A,\star)\land X\subseteq U}{\rA}
  \proofstepwidestar[18]{\mathsf{UHGr}(U;A,\star)}{\rAEa{18}}
  \proofstepwidestar[1,2,3,4,18]{x\in U}{\rRE{18,\rUE{16}}}
  \proofstepwidestar[1,2,3,5,18]{y\in U}{\rRE{18,\rUE{17}}}
  \proofstepwidestar[1,2,3,4,5,18]{x\star y\in U}{%
    \FormulaRefAuto{SemigroupSubsemigroupClosure}[thm]{1,19,20,21}}
  \proofstepwidestar[1,2,3,4,5]{%
    \forall U\,\Bigl(
      \mathsf{UHGr}(U;A,\star)\land X\subseteq U
      \rightarrow x\star y\in U\Bigr)}{\rUI{\rRI{18,22}}}
  \proofstepwidestar[1,2,3,4,5]{%
    x\star y\in
    \bigcap_{\mathsf{UHGr}(U;A,\star)\land X\subseteq U}U}{%
    \FormulaRefAuto{P(A)\vdash x\in\bigcap_{P(B)}B
      \leftrightarrow\forall C\,(P(C)\rightarrow x\in C)}[thm]{12,23}}
  \proofstepwidestar[1,2,3,4,5]{%
    x\star y\in\langle X\rangle_{A,\star}}{%
    \rIE{\FormulaRefAuto{a=b\vdash b=a}[thm]{13},24}}
\end{tabproofwide}

\FormulaThmDeltaK[Das Erzeugnis ist eine Unterhalbgruppe]{%
  \begin{aligned}[t]
    &\SemigroupAlg{A}{\star}\dsep X\subseteq A\dsep
      X\neq\varnothing\vdash{}\\[-2pt]
    &\mathsf{UHGr}(\langle X\rangle_{A,\star};A,\star)
  \end{aligned}
}{GeneratedSubsemigroupIsSubsemigroup}{%
  \DeltaRow{Mengen}{A\dsep X\dsep x\dsep y\dsep u}
  \DeltaRow{Funktionen}{\star}
}
\begin{tabproofsplitwide}
  \proofpartwideindR[Träger und Nichtleerheit des Erzeugnisses]{%
    \begin{aligned}[t]
      &\SemigroupAlg{A}{\star}\dsep X\subseteq A\dsep
        X\neq\varnothing\vdash{}\\[-2pt]
      &\langle X\rangle_{A,\star}\subseteq A\land
        \langle X\rangle_{A,\star}\neq\varnothing
    \end{aligned}
  }{%
    \begin{aligned}[t]
      &\SemigroupAlg{A}{\star}\dsep X\subseteq A\dsep
        X\neq\varnothing\vdash{}\\[-2pt]
      &\langle X\rangle_{A,\star}\subseteq A\land
        \langle X\rangle_{A,\star}\neq\varnothing
    \end{aligned}%
  }[GeneratedSubsemigroupCarrierAndNonempty]
    \proofstepwidestar[1]{\SemigroupAlg{A}{\star}}{\rA}
    \proofstepwidestar[2]{X\subseteq A}{\rA}
    \proofstepwidestar[3]{X\neq\varnothing}{\rA}
    \proofstepwidestar[3]{\exists u\,(u\in X)}{%
      \FormulaRefAuto{A\neq\varnothing\eqvdash
        \exists x\,(x\in A)}[thm]{3}}
    \proofstepwidestar[5]{u\in X}{\rA}
    \proofstepwidestar[2,5]{u\in A}{%
      \FormulaRefAuto{A\subseteq B,\,x\in A\vdash x\in B}[thm]{2,5}}
    \proofstepwidestar[2,5]{A\neq\varnothing}{%
      \FormulaRefAuto{a\in A\vdash A\neq\varnothing}[thm]{6}}
    \proofstepwidestar[2,3]{A\neq\varnothing}{\rEE{4,5,7}}
    \proofstepwidestar[1,2,3]{\mathsf{UHGr}(A;A,\star)}{%
      \FormulaRefAuto{SemigroupCarrierIsSubsemigroup}[thm]{1,8}}
    \proofstepwidestarbreak[1,2,3]{%
      \langle X\rangle_{A,\star}\subseteq A}{%
      \FormulaRefAuto{GeneratedSubsemigroupMinimal}[thm]{1,2,3,9,2}}
    \proofstepwidestar[1,2,3]{%
      X\subseteq\langle X\rangle_{A,\star}}{%
      \FormulaRefAuto{GeneratedSubsemigroupContainsGenerators}[thm]{1,2,3}}
    \proofstepwidestarbreak[2,3]{%
      \exists u\,(u\in\langle X\rangle_{A,\star})}{%
      \FormulaRefAuto{A\subseteq B\dsep A\neq\varnothing
        \vdash\exists a\in B\,a\in A}[thm]{11,3}}
    \proofstepwidestar[1,2,3]{%
      \langle X\rangle_{A,\star}\neq\varnothing}{%
      \FormulaRefAuto{A\neq\varnothing\eqvdash
        \exists x\,(x\in A)}[thm]{12}}
    \proofstepwidestarbreak[1,2,3]{%
      \langle X\rangle_{A,\star}\subseteq A\land
      \langle X\rangle_{A,\star}\neq\varnothing}{\rAI{10,13}}
  \closeproofpartwideindR

  \proofpartwideindR[Produktabschluss des Erzeugnisses]{%
    \begin{aligned}[t]
      &\SemigroupAlg{A}{\star}\dsep X\subseteq A\dsep
        X\neq\varnothing\vdash{}\\[-2pt]
      &\forall x\in\langle X\rangle_{A,\star}\,
        \forall y\in\langle X\rangle_{A,\star}\,
        (x\star y\in\langle X\rangle_{A,\star})
    \end{aligned}
  }{%
    \begin{aligned}[t]
      &\SemigroupAlg{A}{\star}\dsep X\subseteq A\dsep
        X\neq\varnothing\vdash{}\\[-2pt]
      &\forall x\in\langle X\rangle_{A,\star}\,
        \forall y\in\langle X\rangle_{A,\star}\,
        (x\star y\in\langle X\rangle_{A,\star})
    \end{aligned}%
  }[GeneratedSubsemigroupProductClosurePart]
    \proofstepwidestar[1]{\SemigroupAlg{A}{\star}}{\rA}
    \proofstepwidestar[2]{X\subseteq A}{\rA}
    \proofstepwidestar[3]{X\neq\varnothing}{\rA}
    \proofstepwidestar[4]{x\in\langle X\rangle_{A,\star}}{\rA}
    \proofstepwidestar[5]{y\in\langle X\rangle_{A,\star}}{\rA}
    \proofstepwidestarbreak[1,2,3,4,5]{%
      x\star y\in\langle X\rangle_{A,\star}}{%
      \FormulaRefAuto{GeneratedSubsemigroupClosure}[thm]{1,2,3,4,5}}
    \proofstepwidestarbreak[1,2,3]{%
      \forall x\in\langle X\rangle_{A,\star}\,
      \forall y\in\langle X\rangle_{A,\star}\,
      (x\star y\in\langle X\rangle_{A,\star})}{%
      \rUI{\rRI{4,\rUI{\rRI{5,6}}}}}
  \closeproofpartwideindR

  \proofpartwideindR[Schluss]{%
    \begin{aligned}[t]
      &\SemigroupAlg{A}{\star}\dsep X\subseteq A\dsep
        X\neq\varnothing\vdash{}\\[-2pt]
      &\mathsf{UHGr}(\langle X\rangle_{A,\star};A,\star)
    \end{aligned}
  }{%
    \begin{aligned}[t]
      &\SemigroupAlg{A}{\star}\dsep X\subseteq A\dsep
        X\neq\varnothing\vdash{}\\[-2pt]
      &\mathsf{UHGr}(\langle X\rangle_{A,\star};A,\star)
    \end{aligned}%
  }[GeneratedSubsemigroupConclusion]
    \proofstepwidestar[1]{\SemigroupAlg{A}{\star}}{\rA}
    \proofstepwidestar[2]{X\subseteq A}{\rA}
    \proofstepwidestar[3]{X\neq\varnothing}{\rA}
    \proofstepwidestarbreak[1,2,3]{%
      \begin{aligned}[t]
        &\langle X\rangle_{A,\star}\subseteq A\land{}\\[-2pt]
        &\langle X\rangle_{A,\star}\neq\varnothing
      \end{aligned}}{%
      \FormulaRefAuto{GeneratedSubsemigroupCarrierAndNonempty}[thm]{1,2,3}}
    \proofstepwidestar[1,2,3]{%
      \langle X\rangle_{A,\star}\subseteq A}{\rAEa{4}}
    \proofstepwidestar[1,2,3]{%
      \langle X\rangle_{A,\star}\neq\varnothing}{\rAEb{4}}
    \proofstepwidestarbreak[1,2,3]{%
      \begin{aligned}[t]
        &\forall x\in\langle X\rangle_{A,\star}\,\\[-2pt]
        &\quad\forall y\in\langle X\rangle_{A,\star}\,
          (x\star y\in\langle X\rangle_{A,\star})
      \end{aligned}}{%
      \FormulaRefAuto{GeneratedSubsemigroupProductClosurePart}[thm]{1,2,3}}
    \proofstepwidestarbreak[1,2,3]{%
      \begin{aligned}[t]
        &\langle X\rangle_{A,\star}\subseteq A\land
          \langle X\rangle_{A,\star}\neq\varnothing\land{}\\[-2pt]
        &\forall x\in\langle X\rangle_{A,\star}\,\\[-2pt]
        &\quad\forall y\in\langle X\rangle_{A,\star}\,
          (x\star y\in\langle X\rangle_{A,\star})
      \end{aligned}}{%
      \rAI{5,\rAI{6,7}}}
    \proofstepwidestarbreak[1,2,3]{%
      \mathsf{UHGr}(\langle X\rangle_{A,\star};A,\star)}{%
      \FormulaRefAuto{SemigroupSubsemigroupDef}[def]{1,8}}
  \closeproofpartwideindR
\end{tabproofsplitwide}

\subsection{Monogene Unterhalbgruppen}

\FormulaDefDeltaK[Von einem Element erzeugte Unterhalbgruppe]{%
  \SemigroupAlg{A}{\star}\dsep a\in A\vdash
  \langle a\rangle_{A,\star}\coloneqq
  \langle\{a\}\rangle_{A,\star}%
}{SemigroupMonogenicSubsemigroupDef}{%
  \DeltaRow{Mengen}{A\dsep a}
  \DeltaRow{Funktionen}{\star}
  \DeltaRow{\textbf{Neues Symbol}}{}
  \DeltaPrem{Monogene Unterhalbgruppe}{\langle a\rangle_{A,\star}}
}

\FormulaThmDeltaK[Das Monogen ist eine Unterhalbgruppe]{%
  \SemigroupAlg{A}{\star}\dsep a\in A\vdash
  \mathsf{UHGr}(\langle a\rangle_{A,\star};A,\star)%
}{MonogenicSubsemigroupIsSubsemigroup}{%
  \DeltaRow{Mengen}{A\dsep a}
  \DeltaRow{Funktionen}{\star}
}
\begin{tabproofwide}
  \proofstepwidestar[1]{\SemigroupAlg{A}{\star}}{\rA}
  \proofstepwidestar[2]{a\in A}{\rA}
  \proofstepwidestar[2]{\{a\}\subseteq A}{%
    \FormulaRefAuto{a\in A\vdash\{a\}\subseteq A}[thm]{2}}
  \proofstepwidestar[]{a\in\{a\}}{\FormulaRefAuto{a\in\{a\}}[thm]}
  \proofstepwidestar[]{\{a\}\neq\varnothing}{%
    \FormulaRefAuto{a\in A\vdash A\neq\varnothing}[thm]{4}}
  \proofstepwidestar[1,2]{%
    \mathsf{UHGr}(\langle\{a\}\rangle_{A,\star};A,\star)}{%
    \FormulaRefAuto{GeneratedSubsemigroupIsSubsemigroup}[thm]{1,3,5}}
  \proofstepwidestar[1,2]{%
    \langle a\rangle_{A,\star}=\langle\{a\}\rangle_{A,\star}}{%
    \FormulaRefAuto{SemigroupMonogenicSubsemigroupDef}[def]{1,2}}
  \proofstepwidestar[1,2]{%
    \mathsf{UHGr}(\langle a\rangle_{A,\star};A,\star)}{\rIE{7,6}}
\end{tabproofwide}

\FormulaDefDeltaK[Menge der positiven Potenzen]{%
  \SemigroupAlg{A}{\star}\dsep a\in A\vdash
  P_{A,\star}(a)\coloneqq
  \bigl\{x\in A\bigm|
    \exists n\in\mathbb N\,(n\neq0\land x=a^n)\bigr\}%
}{SemigroupPositivePowerSetDef}{%
  \DeltaRow{Mengen}{A\dsep a\dsep x\dsep n}
  \DeltaRow{Funktionen}{\star}
  \DeltaRow{\textbf{Neues Symbol}}{}
  \DeltaPrem{Positive Potenzmenge}{P_{A,\star}(a)}
}

\FormulaThmDeltaK[Elementkriterium der positiven Potenzmenge]{%
  \begin{aligned}[t]
    &\SemigroupAlg{A}{\star}\dsep a\in A\vdash{}\\[-2pt]
    &x\in P_{A,\star}(a)\leftrightarrow
      \Bigl(x\in A\land
        \exists n\in\mathbb N\,(n\neq0\land x=a^n)\Bigr)
  \end{aligned}
}{SemigroupPositivePowerSetMembership}{%
  \DeltaRow{Mengen}{A\dsep a\dsep x\dsep n}
  \DeltaRow{Funktionen}{\star}
}
\begin{tabproofwide}
  \proofstepwidestar[1]{\SemigroupAlg{A}{\star}}{\rA}
  \proofstepwidestar[2]{a\in A}{\rA}
  \proofstepwidestar[1,2]{%
    P_{A,\star}(a)=\{u\in A\mid
      \exists n\in\mathbb N\,(n\neq0\land u=a^n)\}}{%
    \FormulaRefAuto{SemigroupPositivePowerSetDef}[def]{1,2}}
  \proofstepwidestar[1,2]{%
    x\in P_{A,\star}(a)\leftrightarrow
    \Bigl(x\in A\land
      \exists n\in\mathbb N\,(n\neq0\land x=a^n)\Bigr)}{%
    \rIE{3,\FormulaRefAuto{x\in\{u\in A\mid P(u)\}\eqvdash
      x\in A\land P(x)}[thm]}}
\end{tabproofwide}

\FormulaThmDeltaK[Die positiven Potenzen bilden eine Unterhalbgruppe]{%
  \SemigroupAlg{A}{\star}\dsep a\in A\vdash
  \mathsf{UHGr}(P_{A,\star}(a);A,\star)%
}{SemigroupPositivePowersFormSubsemigroup}{%
  \DeltaRow{Mengen}{A\dsep a\dsep x\dsep y\dsep m\dsep n}
  \DeltaRow{Funktionen}{\star}
}
\begin{tabproofsplitwide}
  \proofpartwideindR[Teilmenge und Nichtleerheit]{%
    \begin{aligned}[t]
      &\SemigroupAlg{A}{\star}\dsep a\in A\vdash{}\\[-2pt]
      &P_{A,\star}(a)\subseteq A\land
        P_{A,\star}(a)\neq\varnothing
    \end{aligned}
  }{%
    \begin{aligned}[t]
      &\SemigroupAlg{A}{\star}\dsep a\in A\vdash{}\\[-2pt]
      &P_{A,\star}(a)\subseteq A\land
        P_{A,\star}(a)\neq\varnothing
    \end{aligned}%
  }[SemigroupPositivePowerSetCarrierAndNonempty]
    \proofstepwidestar[1]{\SemigroupAlg{A}{\star}}{\rA}
    \proofstepwidestar[2]{a\in A}{\rA}
    \proofstepwidestar[]{P_{A,\star}(a)\subseteq A}{%
      \FormulaRefAuto{\{x\in A\mid P(x)\}\subseteq A}[thm]}
    \proofstepwidestar[]{1\in\mathbb N}{%
      \FormulaRefAuto{PeanoOneInN}[thm]}
    \proofstepwidestar[]{1\neq0}{%
      \rIE{\FormulaRefAuto{PeanoOne}[def],
        \FormulaRefAuto{PeanoSuccNonzero}[ax]{%
          \FormulaRefAuto{PeanoZero}[ax]}}}
    \proofstepwidestar[1,2]{a^1=a}{%
      \FormulaRefAuto{SemigroupPowerOne}[thm]{1,2}}
    \proofstepwidestarbreak[1,2]{%
      a\in A\land
      \exists n\in\mathbb N\,(n\neq0\land a=a^n)}{%
      \rAI{2,\rEI{4,
        \rAI{5,\FormulaRefAuto{a=b\vdash b=a}[thm]{6}}}}}
    \proofstepwidestar[1,2]{a\in P_{A,\star}(a)}{%
      \FormulaRefAuto{SemigroupPositivePowerSetMembership}[thm]{1,2,7}}
    \proofstepwidestar[1,2]{P_{A,\star}(a)\neq\varnothing}{%
      \FormulaRefAuto{a\in A\vdash A\neq\varnothing}[thm]{8}}
    \proofstepwidestarbreak[1,2]{%
      P_{A,\star}(a)\subseteq A\land
      P_{A,\star}(a)\neq\varnothing}{\rAI{3,9}}
  \closeproofpartwideindR

  \proofpartwideindR[Produktabschluss der positiven Potenzen]{%
    \begin{aligned}[t]
      &\SemigroupAlg{A}{\star}\dsep a\in A\vdash{}\\[-2pt]
      &\forall x\in P_{A,\star}(a)\,
        \forall y\in P_{A,\star}(a)\,
        (x\star y\in P_{A,\star}(a))
    \end{aligned}
  }{%
    \begin{aligned}[t]
      &\SemigroupAlg{A}{\star}\dsep a\in A\vdash{}\\[-2pt]
      &\forall x\in P_{A,\star}(a)\,
        \forall y\in P_{A,\star}(a)\,
        (x\star y\in P_{A,\star}(a))
    \end{aligned}%
  }[SemigroupPositivePowerSetProductClosure]
    \proofstepwidestar[1]{\SemigroupAlg{A}{\star}}{\rA}
    \proofstepwidestar[2]{a\in A}{\rA}
    \proofstepwidestar[3]{x\in P_{A,\star}(a)}{\rA}
    \proofstepwidestar[4]{y\in P_{A,\star}(a)}{\rA}
    \proofstepwidestarbreak[1,2,3]{%
      \exists m\in\mathbb N\,(m\neq0\land x=a^m)}{%
      \rAEb{\FormulaRefAuto{SemigroupPositivePowerSetMembership}[thm]{%
        1,2,3}}}
    \proofstepwidestarbreak[1,2,4]{%
      \exists n\in\mathbb N\,(n\neq0\land y=a^n)}{%
      \rAEb{\FormulaRefAuto{SemigroupPositivePowerSetMembership}[thm]{%
        1,2,4}}}
    \proofstepwidestar[7]{%
      m\in\mathbb N\land(m\neq0\land x=a^m)}{\rA}
    \proofstepwidestar[8]{%
      n\in\mathbb N\land(n\neq0\land y=a^n)}{\rA}
    \proofstepwidestar[7]{m\in\mathbb N}{\rAEa{7}}
    \proofstepwidestar[7]{m\neq0}{\rAEa{\rAEb{7}}}
    \proofstepwidestar[7]{x=a^m}{\rAEb{\rAEb{7}}}
    \proofstepwidestar[8]{n\in\mathbb N}{\rAEa{8}}
    \proofstepwidestar[8]{n\neq0}{\rAEa{\rAEb{8}}}
    \proofstepwidestar[8]{y=a^n}{\rAEb{\rAEb{8}}}
    \proofstepwidestar[7,8]{m+n\in\mathbb N}{%
      \FormulaRefAuto{PeanoAddClosure}[thm]{9,12}}
    \proofstepwidestar[7]{\Pred(m)\in\mathbb N}{%
      \FormulaRefAuto{PeanoPredNatural}[thm]{9,10}}
    \proofstepwidestar[7]{m=\Succ(\Pred(m))}{%
      \FormulaRefAuto{PeanoPredNonzeroValue}[thm]{9,10}}
    \proofstepwidestarbreak[7,8]{%
      \Succ(\Pred(m))+n=\Succ(\Pred(m)+n)}{%
      \FormulaRefAuto{PeanoAddSuccLeft}[thm]{16,12}}
    \proofstepwidestar[7,8]{\Pred(m)+n\in\mathbb N}{%
      \FormulaRefAuto{PeanoAddClosure}[thm]{16,12}}
    \proofstepwidestar[7,8]{\Succ(\Pred(m)+n)\neq0}{%
      \FormulaRefAuto{PeanoSuccNonzero}[ax]{19}}
    \proofstepwidestar[7,8]{%
      m+n=\Succ(\Pred(m)+n)}{\rIE{17,18}}
    \proofstepwidestar[7,8]{m+n\neq0}{\rIE{21,20}}
    \proofstepwidestarbreak[1,2,7,8]{%
      a^{m+n}=a^m\star a^n}{%
      \FormulaRefAuto{SemigroupPowerAdditionLaw}[thm]{%
        1,2,9,12,10,13}}
    \proofstepwidestar[1,2,7,8]{%
      x\star y=a^{m+n}}{\rIE{11,14,23}}
    \proofstepwidestar[1,2,7,8]{x\star y\in A}{%
      \FormulaRefAuto{SemigroupPowerClosure}[thm]{1,2,15,22}}
    \proofstepwidestarbreak[1,2,7,8]{%
      x\star y\in A\land
      \exists k\in\mathbb N\,(k\neq0\land x\star y=a^k)}{%
      \rAI{25,\rEI{15,\rAI{22,24}}}}
    \proofstepwidestarbreak[1,2,7,8]{%
      x\star y\in P_{A,\star}(a)}{%
      \FormulaRefAuto{SemigroupPositivePowerSetMembership}[thm]{1,2,26}}
    \proofstepwidestar[1,2,3,4]{%
      x\star y\in P_{A,\star}(a)}{%
      \rEE{5,7,\rEE{6,8,27}}}
    \proofstepwidestarbreak[1,2]{%
      \forall x\in P_{A,\star}(a)\,
      \forall y\in P_{A,\star}(a)\,
      (x\star y\in P_{A,\star}(a))}{%
      \rUI{\rRI{3,\rUI{\rRI{4,28}}}}}
  \closeproofpartwideindR

  \proofpartwideindR[Schluss]{%
    \begin{aligned}[t]
      &\SemigroupAlg{A}{\star}\dsep a\in A\vdash{}\\[-2pt]
      &\mathsf{UHGr}(P_{A,\star}(a);A,\star)
    \end{aligned}
  }{%
    \begin{aligned}[t]
      &\SemigroupAlg{A}{\star}\dsep a\in A\vdash{}\\[-2pt]
      &\mathsf{UHGr}(P_{A,\star}(a);A,\star)
    \end{aligned}%
  }[SemigroupPositivePowerSetSubsemigroupConclusion]
    \proofstepwidestar[1]{\SemigroupAlg{A}{\star}}{\rA}
    \proofstepwidestar[2]{a\in A}{\rA}
    \proofstepwidestarbreak[1,2]{%
      \begin{aligned}[t]
        &P_{A,\star}(a)\subseteq A\land{}\\[-2pt]
        &P_{A,\star}(a)\neq\varnothing
      \end{aligned}}{%
      \FormulaRefAuto{SemigroupPositivePowerSetCarrierAndNonempty}[thm]{1,2}}
    \proofstepwidestar[1,2]{P_{A,\star}(a)\subseteq A}{\rAEa{3}}
    \proofstepwidestar[1,2]{P_{A,\star}(a)\neq\varnothing}{\rAEb{3}}
    \proofstepwidestarbreak[1,2]{%
      \begin{aligned}[t]
        &\forall x\in P_{A,\star}(a)\,\\[-2pt]
        &\quad\forall y\in P_{A,\star}(a)\,
          (x\star y\in P_{A,\star}(a))
      \end{aligned}}{%
      \FormulaRefAuto{SemigroupPositivePowerSetProductClosure}[thm]{1,2}}
    \proofstepwidestarbreak[1,2]{%
      \begin{aligned}[t]
        &P_{A,\star}(a)\subseteq A\land
          P_{A,\star}(a)\neq\varnothing\land{}\\[-2pt]
        &\forall x\in P_{A,\star}(a)\,\\[-2pt]
        &\quad\forall y\in P_{A,\star}(a)\,
          (x\star y\in P_{A,\star}(a))
      \end{aligned}}{%
      \rAI{4,\rAI{5,6}}}
    \proofstepwidestarbreak[1,2]{%
      \mathsf{UHGr}(P_{A,\star}(a);A,\star)}{%
      \FormulaRefAuto{SemigroupSubsemigroupDef}[def]{1,7}}
  \closeproofpartwideindR
\end{tabproofsplitwide}

\FormulaThmDeltaK[Beschreibung der monogenen Unterhalbgruppe durch Potenzen]{%
  \SemigroupAlg{A}{\star}\dsep a\in A\vdash
  \langle a\rangle_{A,\star}=P_{A,\star}(a)%
}{SemigroupMonogenicPowerDescription}{%
  \DeltaRow{Mengen}{A\dsep a\dsep x\dsep n\dsep k}
  \DeltaRow{Funktionen}{\star}
}
\begin{tabproofsplitwide}
  \proofpartwideindR[Positive Potenzen liegen im Monogen]{%
    \begin{aligned}[t]
      &\SemigroupAlg{A}{\star}\dsep a\in A\dsep k\in\mathbb N\vdash{}\\[-2pt]
      &a^{\Succ(k)}\in\langle a\rangle_{A,\star}
    \end{aligned}
  }{%
    \SemigroupAlg{A}{\star}\dsep a\in A\dsep k\in\mathbb N
    \vdash a^{\Succ(k)}\in\langle a\rangle_{A,\star}
  }[SemigroupPowerInMonogenicSubsemigroup]
    \proofstepwidestar[1]{\SemigroupAlg{A}{\star}}{\rA}
    \proofstepwidestar[2]{a\in A}{\rA}
    \proofstepwidestar[3]{k\in\mathbb N}{\rA}
    \proofstepwidestar[2]{\{a\}\subseteq A}{%
      \FormulaRefAuto{a\in A\vdash\{a\}\subseteq A}[thm]{2}}
    \proofstepwidestar[]{\{a\}\neq\varnothing}{%
      \FormulaRefAuto{a\in A\vdash A\neq\varnothing}[thm]{%
        \FormulaRefAuto{a\in\{a\}}[thm]}}
    \proofstepwidestar[1,2]{a\in\langle\{a\}\rangle_{A,\star}}{%
      \FormulaRefAuto{GeneratedSubsemigroupContainsGenerators}[thm]{%
        1,4,5,\FormulaRefAuto{a\in\{a\}}[thm]}}
    \proofstepwidestar[1,2]{a\in\langle a\rangle_{A,\star}}{%
      \rIE{\FormulaRefAuto{a=b\vdash b=a}[thm]{%
        \FormulaRefAuto{SemigroupMonogenicSubsemigroupDef}[def]{1,2}},6}}
    \proofstepwidestar[1,2]{a^{\Succ(0)}=a}{%
      \rIE{\FormulaRefAuto{PeanoOne}[def],
        \FormulaRefAuto{SemigroupPowerOne}[thm]{1,2}}}
    \proofstepwidestar[1,2]{a^{\Succ(0)}\in\langle a\rangle_{A,\star}}{%
      \rIE{\FormulaRefAuto{a=b\vdash b=a}[thm]{8},7}}
    \proofstepwidestar[1,2]{%
      \mathsf{UHGr}(\langle a\rangle_{A,\star};A,\star)}{%
      \FormulaRefAuto{MonogenicSubsemigroupIsSubsemigroup}[thm]{1,2}}
    \proofstepwidestar[1,2,3]{a^{\Succ(k)}\in A}{%
      \FormulaRefAuto{SemigroupPowerClosure}[thm]{%
        1,2,\FormulaRefAuto{PeanoSuccClosure}[ax]{3},
        \FormulaRefAuto{PeanoSuccNonzero}[ax]{3}}}
    \proofstepwidestar[12]{a^{\Succ(k)}\in\langle a\rangle_{A,\star}}{\rA}
    \proofstepwidestar[1,2,3,12]{%
      a^{\Succ(k)+1}=a^{\Succ(k)}\star a}{%
      \FormulaRefAuto{SemigroupPowerSuccessor}[thm]{%
        1,2,\FormulaRefAuto{PeanoSuccClosure}[ax]{3},
        \FormulaRefAuto{PeanoSuccNonzero}[ax]{3}}}
    \proofstepwidestar[1,2,3,12]{%
      a^{\Succ(k)}\star a\in\langle a\rangle_{A,\star}}{%
      \FormulaRefAuto{SemigroupSubsemigroupClosure}[thm]{1,10,12,7}}
    \proofstepwidestar[1,2,3,12]{%
      a^{\Succ(\Succ(k))}\in\langle a\rangle_{A,\star}}{%
      \rIE{\FormulaRefAuto{PeanoAddOneSucc}[thm]{%
          \FormulaRefAuto{PeanoSuccClosure}[ax]{3}},\rIE{13,14}}}
    \proofstepwidestar[1,2]{%
      \forall k\in\mathbb N\,
      \bigl(a^{\Succ(k)}\in\langle a\rangle_{A,\star}\bigr)}{%
      \FormulaRefAuto{PeanoInductionRule}[ax]{9,12,15}}
    \proofstepwidestar[1,2,3]{%
      a^{\Succ(k)}\in\langle a\rangle_{A,\star}}{\rRE{3,\rUE{16}}}
  \closeproofpartwideindR

  \proofpartwideindR[Monogen liegt in der Menge positiver Potenzen]{%
    \begin{aligned}[t]
      &\SemigroupAlg{A}{\star}\dsep a\in A\vdash{}\\[-2pt]
      &\langle a\rangle_{A,\star}\subseteq P_{A,\star}(a)
    \end{aligned}
  }{%
    \SemigroupAlg{A}{\star}\dsep a\in A\vdash
    \langle a\rangle_{A,\star}\subseteq P_{A,\star}(a)%
  }[SemigroupMonogenicSubsemigroupInPositivePowers]
    \proofstepwidestar[1]{\SemigroupAlg{A}{\star}}{\rA}
    \proofstepwidestar[2]{a\in A}{\rA}
    \proofstepwidestarbreak[1,2]{\mathsf{UHGr}(P_{A,\star}(a);A,\star)}{%
      \FormulaRefAuto{SemigroupPositivePowersFormSubsemigroup}[thm]{1,2}}
    \proofstepwidestar[]{1\in\mathbb N}{\FormulaRefAuto{PeanoOneInN}[thm]}
    \proofstepwidestar[]{1\neq0}{%
      \rIE{\FormulaRefAuto{PeanoOne}[def],
        \FormulaRefAuto{PeanoSuccNonzero}[ax]{\FormulaRefAuto{PeanoZero}[ax]}}}
    \proofstepwidestar[1,2]{a^1=a}{%
      \FormulaRefAuto{SemigroupPowerOne}[thm]{1,2}}
    \proofstepwidestarbreak[1,2]{%
      \exists n\in\mathbb N\,(n\neq0\land a=a^n)}{%
      \rEI{4,\rAI{5,\FormulaRefAuto{a=b\vdash b=a}[thm]{6}}}}
    \proofstepwidestar[1,2]{%
      a\in A\land\exists n\in\mathbb N\,(n\neq0\land a=a^n)}{%
      \rAI{2,7}}
    \proofstepwidestarbreak[1,2]{a\in P_{A,\star}(a)}{%
      \FormulaRefAuto{SemigroupPositivePowerSetMembership}[thm]{1,2,8}}
    \proofstepwidestarbreak[1,2]{\{a\}\subseteq P_{A,\star}(a)}{%
      \FormulaRefAuto{a\in A\vdash\{a\}\subseteq A}[thm]{9}}
    \proofstepwidestar[2]{\{a\}\subseteq A}{%
      \FormulaRefAuto{a\in A\vdash\{a\}\subseteq A}[thm]{2}}
    \proofstepwidestar[]{\{a\}\neq\varnothing}{%
      \FormulaRefAuto{a\in A\vdash A\neq\varnothing}[thm]{%
        \FormulaRefAuto{a\in\{a\}}[thm]}}
    \proofstepwidestarbreak[1,2]{%
      \langle\{a\}\rangle_{A,\star}\subseteq P_{A,\star}(a)}{%
      \FormulaRefAuto{GeneratedSubsemigroupMinimal}[thm]{1,11,12,3,10}}
    \proofstepwidestarbreak[1,2]{%
      \langle a\rangle_{A,\star}\subseteq P_{A,\star}(a)}{%
      \rIE{\FormulaRefAuto{SemigroupMonogenicSubsemigroupDef}[def]{1,2},13}}
  \closeproofpartwideindR

  \proofpartwideindR[Menge positiver Potenzen liegt im Monogen]{%
    \begin{aligned}[t]
      &\SemigroupAlg{A}{\star}\dsep a\in A\vdash{}\\[-2pt]
      &P_{A,\star}(a)\subseteq\langle a\rangle_{A,\star}
    \end{aligned}
  }{%
    \SemigroupAlg{A}{\star}\dsep a\in A\vdash
    P_{A,\star}(a)\subseteq\langle a\rangle_{A,\star}%
  }[SemigroupPositivePowersInMonogenicSubsemigroup]
    \proofstepwidestar[1]{\SemigroupAlg{A}{\star}}{\rA}
    \proofstepwidestar[2]{a\in A}{\rA}
    \proofstepwidestar[3]{x\in P_{A,\star}(a)}{\rA}
    \proofstepwidestarbreak[1,2,3]{%
      \exists n\in\mathbb N\,(n\neq0\land x=a^n)}{%
      \rAEb{\FormulaRefAuto{SemigroupPositivePowerSetMembership}[thm]{1,2,3}}}
    \proofstepwidestar[5]{n\in\mathbb N\land(n\neq0\land x=a^n)}{\rA}
    \proofstepwidestar[5]{n\in\mathbb N}{\rAEa{5}}
    \proofstepwidestar[5]{n\neq0}{\rAEa{\rAEb{5}}}
    \proofstepwidestar[5]{x=a^n}{\rAEb{\rAEb{5}}}
    \proofstepwidestar[5]{n=\Succ(\Pred(n))}{%
      \FormulaRefAuto{PeanoPredNonzeroValue}[thm]{6,7}}
    \proofstepwidestar[5]{\Pred(n)\in\mathbb N}{%
      \FormulaRefAuto{PeanoPredNatural}[thm]{6,7}}
    \proofstepwidestarbreak[1,2,5]{%
      a^{\Succ(\Pred(n))}\in\langle a\rangle_{A,\star}}{%
      \FormulaRefAuto{SemigroupPowerInMonogenicSubsemigroup}[thm]{1,2,10}}
    \proofstepwidestar[1,2,5]{a^n\in\langle a\rangle_{A,\star}}{\rIE{9,11}}
    \proofstepwidestar[1,2,5]{x\in\langle a\rangle_{A,\star}}{\rIE{8,12}}
    \proofstepwidestar[1,2,3]{x\in\langle a\rangle_{A,\star}}{\rEE{4,5,13}}
    \proofstepwidestarbreak[1,2]{%
      P_{A,\star}(a)\subseteq\langle a\rangle_{A,\star}}{%
      \FormulaRefAuto{A\subseteq B\coloneq
        \forall x\,(x\in A\rightarrow x\in B)}[def]{\rUI{\rRI{3,14}}}}
  \closeproofpartwideindR

  \proofpartwideindR[Gleichheit aus beiden Inklusionen]{%
    \begin{aligned}[t]
      &\SemigroupAlg{A}{\star}\dsep a\in A\vdash{}\\[-2pt]
      &\langle a\rangle_{A,\star}=P_{A,\star}(a)
    \end{aligned}
  }{%
    \SemigroupAlg{A}{\star}\dsep a\in A\vdash
    \langle a\rangle_{A,\star}=P_{A,\star}(a)%
  }[SemigroupMonogenicPowerDescriptionConclusion]
    \proofstepwidestar[1]{\SemigroupAlg{A}{\star}}{\rA}
    \proofstepwidestar[2]{a\in A}{\rA}
    \proofstepwidestarbreak[1,2]{%
      \langle a\rangle_{A,\star}\subseteq P_{A,\star}(a)}{%
      \FormulaRefAuto{SemigroupMonogenicSubsemigroupInPositivePowers}[thm]{1,2}}
    \proofstepwidestarbreak[1,2]{%
      P_{A,\star}(a)\subseteq\langle a\rangle_{A,\star}}{%
      \FormulaRefAuto{SemigroupPositivePowersInMonogenicSubsemigroup}[thm]{1,2}}
    \proofstepwidestarbreak[1,2]{%
      \langle a\rangle_{A,\star}=P_{A,\star}(a)}{%
      \FormulaRefAuto{A\subseteq B,\,B\subseteq A\vdash A=B}[thm]{3,4}}
  \closeproofpartwideindR
\end{tabproofsplitwide}

\section{Ideale}

Im Unterschied zu Unterhalbgruppen werden Ideale durch Multiplikation mit
beliebigen Elementen der umgebenden Halbgruppe absorbiert. Wir verlangen
ausdr\"ucklich Nichtleerheit; andernfalls w\"are die leere Menge stets das
kleinste Ideal und die sp\"atere Minimalidealtheorie w\"urde entarten.

\FormulaDefDeltaK[Linksideal]{%
  \begin{aligned}[t]
    &\SemigroupAlg{A}{\star}\vdash{}\\[-2pt]
    &\mathsf{LId}(I;A,\star)\coloneqq
      I\subseteq A\land I\neq\varnothing\land
      \forall a\in A\,\forall x\in I\,(a\star x\in I)
  \end{aligned}
}{SemigroupLeftIdealDef}{%
  \DeltaRow{Mengen}{A\dsep I\dsep a\dsep x}
  \DeltaRow{Funktionen}{\star}
  \DeltaRow{\textbf{Neues Symbol}}{}
  \DeltaPrem{Linksideal}{\mathsf{LId}(I;A,\star)}
}

\FormulaDefDeltaK[Rechtsideal]{%
  \begin{aligned}[t]
    &\SemigroupAlg{A}{\star}\vdash{}\\[-2pt]
    &\mathsf{RId}(I;A,\star)\coloneqq
      I\subseteq A\land I\neq\varnothing\land
      \forall x\in I\,\forall a\in A\,(x\star a\in I)
  \end{aligned}
}{SemigroupRightIdealDef}{%
  \DeltaRow{Mengen}{A\dsep I\dsep a\dsep x}
  \DeltaRow{Funktionen}{\star}
  \DeltaRow{\textbf{Neues Symbol}}{}
  \DeltaPrem{Rechtsideal}{\mathsf{RId}(I;A,\star)}
}

\FormulaDefDeltaK[Zweiseitiges Ideal]{%
  \SemigroupAlg{A}{\star}\vdash
  \mathsf{Id}(I;A,\star)\coloneqq
  \mathsf{LId}(I;A,\star)\land\mathsf{RId}(I;A,\star)%
}{SemigroupIdealDef}{%
  \DeltaRow{Mengen}{A\dsep I}
  \DeltaRow{Funktionen}{\star}
  \DeltaRow{\textbf{Neues Symbol}}{}
  \DeltaPrem{Zweiseitiges Ideal}{\mathsf{Id}(I;A,\star)}
}

\FormulaThmDeltaK[Tr\"agermenge eines Ideals]{%
  \SemigroupAlg{A}{\star}\dsep\mathsf{Id}(I;A,\star)
  \vdash I\subseteq A%
}{SemigroupIdealSubset}{%
  \DeltaRow{Mengen}{A\dsep I}
  \DeltaRow{Funktionen}{\star}
}
\begin{tabproofwide}
  \proofstepwidestar[1]{\SemigroupAlg{A}{\star}}{\rA}
  \proofstepwidestar[2]{\mathsf{Id}(I;A,\star)}{\rA}
  \proofstepwidestar[1,2]{%
    \mathsf{LId}(I;A,\star)\land\mathsf{RId}(I;A,\star)}{%
    \FormulaRefAuto{SemigroupIdealDef}[def]{1,2}}
  \proofstepwidestar[1,2]{\mathsf{LId}(I;A,\star)}{\rAEa{3}}
  \proofstepwidestar[1,2]{%
    I\subseteq A\land I\neq\varnothing\land
    \forall a\in A\,\forall x\in I\,(a\star x\in I)}{%
    \FormulaRefAuto{SemigroupLeftIdealDef}[def]{1,4}}
  \proofstepwidestar[1,2]{I\subseteq A}{\rAEa{5}}
\end{tabproofwide}

\FormulaThmDeltaK[Nichtleerheit eines Ideals]{%
  \SemigroupAlg{A}{\star}\dsep\mathsf{Id}(I;A,\star)
  \vdash I\neq\varnothing%
}{SemigroupIdealNonempty}{%
  \DeltaRow{Mengen}{A\dsep I}
  \DeltaRow{Funktionen}{\star}
}
\begin{tabproofwide}
  \proofstepwidestar[1]{\SemigroupAlg{A}{\star}}{\rA}
  \proofstepwidestar[2]{\mathsf{Id}(I;A,\star)}{\rA}
  \proofstepwidestar[1,2]{\mathsf{LId}(I;A,\star)}{%
    \rAEa{\FormulaRefAuto{SemigroupIdealDef}[def]{1,2}}}
  \proofstepwidestar[1,2]{%
    I\subseteq A\land I\neq\varnothing\land
    \forall a\in A\,\forall x\in I\,(a\star x\in I)}{%
    \FormulaRefAuto{SemigroupLeftIdealDef}[def]{1,3}}
  \proofstepwidestar[1,2]{I\neq\varnothing}{\rAEa{\rAEb{4}}}
\end{tabproofwide}

\FormulaThmDeltaK[Linksabsorption eines Ideals]{%
  \SemigroupAlg{A}{\star}\dsep\mathsf{Id}(I;A,\star)\dsep
  a\in A\dsep x\in I\vdash a\star x\in I%
}{SemigroupIdealLeftAbsorption}{%
  \DeltaRow{Mengen}{A\dsep I\dsep a\dsep x}
  \DeltaRow{Funktionen}{\star}
}
\begin{tabproofwide}
  \proofstepwidestar[1]{\SemigroupAlg{A}{\star}}{\rA}
  \proofstepwidestar[2]{\mathsf{Id}(I;A,\star)}{\rA}
  \proofstepwidestar[3]{a\in A}{\rA}
  \proofstepwidestar[4]{x\in I}{\rA}
  \proofstepwidestar[1,2]{\mathsf{LId}(I;A,\star)}{%
    \rAEa{\FormulaRefAuto{SemigroupIdealDef}[def]{1,2}}}
  \proofstepwidestar[1,2]{%
    \forall u\in A\,\forall v\in I\,(u\star v\in I)}{%
    \rAEb{\rAEb{\FormulaRefAuto{SemigroupLeftIdealDef}[def]{1,5}}}}
  \proofstepwidestar[1,2,3,4]{a\star x\in I}{\rRE{4,\rUE{\rRE{3,\rUE{6}}}}}
\end{tabproofwide}

\FormulaThmDeltaK[Rechtsabsorption eines Ideals]{%
  \SemigroupAlg{A}{\star}\dsep\mathsf{Id}(I;A,\star)\dsep
  x\in I\dsep a\in A\vdash x\star a\in I%
}{SemigroupIdealRightAbsorption}{%
  \DeltaRow{Mengen}{A\dsep I\dsep a\dsep x}
  \DeltaRow{Funktionen}{\star}
}
\begin{tabproofwide}
  \proofstepwidestar[1]{\SemigroupAlg{A}{\star}}{\rA}
  \proofstepwidestar[2]{\mathsf{Id}(I;A,\star)}{\rA}
  \proofstepwidestar[3]{x\in I}{\rA}
  \proofstepwidestar[4]{a\in A}{\rA}
  \proofstepwidestar[1,2]{\mathsf{RId}(I;A,\star)}{%
    \rAEb{\FormulaRefAuto{SemigroupIdealDef}[def]{1,2}}}
  \proofstepwidestar[1,2]{%
    \forall u\in I\,\forall v\in A\,(u\star v\in I)}{%
    \rAEb{\rAEb{\FormulaRefAuto{SemigroupRightIdealDef}[def]{1,5}}}}
  \proofstepwidestar[1,2,3,4]{x\star a\in I}{\rRE{4,\rUE{\rRE{3,\rUE{6}}}}}
\end{tabproofwide}

\FormulaThmDeltaK[Der nichtleere Tr\"ager ist ein Ideal]{%
  \SemigroupAlg{A}{\star}\dsep A\neq\varnothing
  \vdash\mathsf{Id}(A;A,\star)%
}{SemigroupCarrierIsIdeal}{%
  \DeltaRow{Mengen}{A\dsep a\dsep x}
  \DeltaRow{Funktionen}{\star}
}
\begin{tabproofwide}
  \proofstepwidestar[1]{\SemigroupAlg{A}{\star}}{\rA}
  \proofstepwidestar[2]{A\neq\varnothing}{\rA}
  \proofstepwidestar[]{A\subseteq A}{\FormulaRefAuto{A\subseteq A}[thm]}
  \proofstepwidestar[4]{a\in A}{\rA}
  \proofstepwidestar[5]{x\in A}{\rA}
  \proofstepwidestar[1,4,5]{a\star x\in A}{%
    \FormulaRefAuto{SemigroupClosureAxiom}[ax]{1,4,5}}
  \proofstepwidestar[1]{\forall a\in A\,\forall x\in A\,(a\star x\in A)}{%
    \rUI{\rRI{4,\rUI{\rRI{5,6}}}}}
  \proofstepwidestar[1,2]{\mathsf{LId}(A;A,\star)}{%
    \FormulaRefAuto{SemigroupLeftIdealDef}[def]{1,\rAI{3,\rAI{2,7}}}}
  \proofstepwidestar[1]{\forall x\in A\,\forall a\in A\,(x\star a\in A)}{%
    \rUI{\rRI{5,\rUI{\rRI{4,6}}}}}
  \proofstepwidestar[1,2]{\mathsf{RId}(A;A,\star)}{%
    \FormulaRefAuto{SemigroupRightIdealDef}[def]{1,\rAI{3,\rAI{2,9}}}}
  \proofstepwidestar[1,2]{\mathsf{Id}(A;A,\star)}{%
    \FormulaRefAuto{SemigroupIdealDef}[def]{1,\rAI{8,10}}}
\end{tabproofwide}

\FormulaThmDeltaK[Schnitt zweier Ideale]{%
  \begin{aligned}[t]
    &\SemigroupAlg{A}{\star}\dsep\mathsf{Id}(I;A,\star)\dsep
      \mathsf{Id}(J;A,\star)\vdash{}\\[-2pt]
    &\mathsf{Id}(I\cap J;A,\star)
  \end{aligned}
}{SemigroupIdealIntersection}{%
  \DeltaRow{Mengen}{A\dsep I\dsep J\dsep a\dsep x\dsep i\dsep j}
  \DeltaRow{Funktionen}{\star}
}
Der Schnitt ist nichtleer, weil das Produkt eines Elements von (I) mit
einem Element von (J) zugleich in beiden Idealen liegt.  Die linke und
die rechte Absorption werden danach komponentenweise vererbt.

\begin{tabproofsplitwide}
  \proofpartwideindR[Nichtleerheit]{%
    \SemigroupAlg{A}{\star}\dsep\mathsf{Id}(I;A,\star)\dsep
    \mathsf{Id}(J;A,\star)\vdash I\cap J\neq\varnothing
  }{%
    \SemigroupAlg{A}{\star}\dsep\mathsf{Id}(I;A,\star)\dsep
    \mathsf{Id}(J;A,\star)\vdash I\cap J\neq\varnothing
  }[SemigroupIdealIntersectionNonempty]
    \proofstepwidestar[1]{\SemigroupAlg{A}{\star}}{\rA}
    \proofstepwidestar[2]{\mathsf{Id}(I;A,\star)}{\rA}
    \proofstepwidestar[3]{\mathsf{Id}(J;A,\star)}{\rA}
    \proofstepwidestar[1,2]{I\subseteq A}{%
      \FormulaRefAuto{SemigroupIdealSubset}[thm]{1,2}}
    \proofstepwidestar[1,3]{J\subseteq A}{%
      \FormulaRefAuto{SemigroupIdealSubset}[thm]{1,3}}
    \proofstepwidestar[1,2]{I\neq\varnothing}{%
      \FormulaRefAuto{SemigroupIdealNonempty}[thm]{1,2}}
    \proofstepwidestar[1,3]{J\neq\varnothing}{%
      \FormulaRefAuto{SemigroupIdealNonempty}[thm]{1,3}}
    \proofstepwidestar[1,2]{\exists i\,(i\in I)}{%
      \FormulaRefAuto{A\neq\varnothing\eqvdash\exists x\,(x\in A)}[thm]{6}}
    \proofstepwidestar[1,3]{\exists j\,(j\in J)}{%
      \FormulaRefAuto{A\neq\varnothing\eqvdash\exists x\,(x\in A)}[thm]{7}}
    \proofstepwidestar[10]{i\in I}{\rA}
    \proofstepwidestar[11]{j\in J}{\rA}
    \proofstepwidestar[1,2,10]{i\in A}{%
      \FormulaRefAuto{A\subseteq B,\,x\in A\vdash x\in B}[thm]{4,10}}
    \proofstepwidestar[1,3,11]{j\in A}{%
      \FormulaRefAuto{A\subseteq B,\,x\in A\vdash x\in B}[thm]{5,11}}
    \proofstepwidestar[1,3,10,11]{i\star j\in J}{%
      \FormulaRefAuto{SemigroupIdealLeftAbsorption}[thm]{1,3,12,11}}
    \proofstepwidestar[1,2,10,11]{i\star j\in I}{%
      \FormulaRefAuto{SemigroupIdealRightAbsorption}[thm]{1,2,10,13}}
    \proofstepwidestar[1,2,3,10,11]{i\star j\in I\cap J}{%
      \FormulaRefAuto{x\in A\cap B\eqvdash x\in A\land x\in B}[thm]{%
        \rAI{15,14}}}
    \proofstepwidestar[1,2,3,10,11]{I\cap J\neq\varnothing}{%
      \FormulaRefAuto{a\in A\vdash A\neq\varnothing}[thm]{16}}
    \proofstepwidestar[1,2,3,10]{I\cap J\neq\varnothing}{%
      \rEE{9,11,17}}
    \proofstepwidestar[1,2,3]{I\cap J\neq\varnothing}{%
      \rEE{8,10,18}}
  \closeproofpartwideindR

  \proofpartwideindR[Linksidealeigenschaft]{%
    \SemigroupAlg{A}{\star}\dsep\mathsf{Id}(I;A,\star)\dsep
    \mathsf{Id}(J;A,\star)\vdash\mathsf{LId}(I\cap J;A,\star)
  }{%
    \SemigroupAlg{A}{\star}\dsep\mathsf{Id}(I;A,\star)\dsep
    \mathsf{Id}(J;A,\star)\vdash\mathsf{LId}(I\cap J;A,\star)
  }[SemigroupIdealIntersectionLeftIdeal]
    \proofstepwidestar[1]{\SemigroupAlg{A}{\star}}{\rA}
    \proofstepwidestar[2]{\mathsf{Id}(I;A,\star)}{\rA}
    \proofstepwidestar[3]{\mathsf{Id}(J;A,\star)}{\rA}
    \proofstepwidestar[1,2]{I\subseteq A}{%
      \FormulaRefAuto{SemigroupIdealSubset}[thm]{1,2}}
    \proofstepwidestar[]{I\cap J\subseteq I}{%
      \FormulaRefAuto{A\cap B\subseteq A}[thm]}
    \proofstepwidestar[1,2]{I\cap J\subseteq A}{%
      \FormulaRefAuto{A\subseteq B\dsep B\subseteq C\vdash
        A\subseteq C}[thm]{5,4}}
    \proofstepwidestar[1,2,3]{I\cap J\neq\varnothing}{%
      \FormulaRefAuto{SemigroupIdealIntersectionNonempty}[thm]{1,2,3}}
    \proofstepwidestar[8]{a\in A}{\rA}
    \proofstepwidestar[9]{x\in I\cap J}{\rA}
    \proofstepwidestar[9]{x\in I}{%
      \FormulaRefAuto{x\in A\cap B\vdash x\in A}[thm]{9}}
    \proofstepwidestar[9]{x\in J}{%
      \FormulaRefAuto{x\in A\cap B\vdash x\in B}[thm]{9}}
    \proofstepwidestar[1,2,8,9]{a\star x\in I}{%
      \FormulaRefAuto{SemigroupIdealLeftAbsorption}[thm]{1,2,8,10}}
    \proofstepwidestar[1,3,8,9]{a\star x\in J}{%
      \FormulaRefAuto{SemigroupIdealLeftAbsorption}[thm]{1,3,8,11}}
    \proofstepwidestar[1,2,3,8,9]{a\star x\in I\cap J}{%
      \FormulaRefAuto{x\in A\cap B\eqvdash x\in A\land x\in B}[thm]{%
        \rAI{12,13}}}
    \proofstepwidestar[1,2,3]{%
      \forall a\in A\,\forall x\in I\cap J\,(a\star x\in I\cap J)}{%
      \rUI{\rRI{8,\rUI{\rRI{9,14}}}}}
    \proofstepwidestar[1,2,3]{\mathsf{LId}(I\cap J;A,\star)}{%
      \FormulaRefAuto{SemigroupLeftIdealDef}[def]{1,\rAI{6,\rAI{7,15}}}}
  \closeproofpartwideindR

  \proofpartwideindR[Rechtsidealeigenschaft]{%
    \SemigroupAlg{A}{\star}\dsep\mathsf{Id}(I;A,\star)\dsep
    \mathsf{Id}(J;A,\star)\vdash\mathsf{RId}(I\cap J;A,\star)
  }{%
    \SemigroupAlg{A}{\star}\dsep\mathsf{Id}(I;A,\star)\dsep
    \mathsf{Id}(J;A,\star)\vdash\mathsf{RId}(I\cap J;A,\star)
  }[SemigroupIdealIntersectionRightIdeal]
    \proofstepwidestar[1]{\SemigroupAlg{A}{\star}}{\rA}
    \proofstepwidestar[2]{\mathsf{Id}(I;A,\star)}{\rA}
    \proofstepwidestar[3]{\mathsf{Id}(J;A,\star)}{\rA}
    \proofstepwidestar[1,2]{I\subseteq A}{%
      \FormulaRefAuto{SemigroupIdealSubset}[thm]{1,2}}
    \proofstepwidestar[]{I\cap J\subseteq I}{%
      \FormulaRefAuto{A\cap B\subseteq A}[thm]}
    \proofstepwidestar[1,2]{I\cap J\subseteq A}{%
      \FormulaRefAuto{A\subseteq B\dsep B\subseteq C\vdash
        A\subseteq C}[thm]{5,4}}
    \proofstepwidestar[1,2,3]{I\cap J\neq\varnothing}{%
      \FormulaRefAuto{SemigroupIdealIntersectionNonempty}[thm]{1,2,3}}
    \proofstepwidestar[8]{x\in I\cap J}{\rA}
    \proofstepwidestar[9]{a\in A}{\rA}
    \proofstepwidestar[8]{x\in I}{%
      \FormulaRefAuto{x\in A\cap B\vdash x\in A}[thm]{8}}
    \proofstepwidestar[8]{x\in J}{%
      \FormulaRefAuto{x\in A\cap B\vdash x\in B}[thm]{8}}
    \proofstepwidestar[1,2,8,9]{x\star a\in I}{%
      \FormulaRefAuto{SemigroupIdealRightAbsorption}[thm]{1,2,10,9}}
    \proofstepwidestar[1,3,8,9]{x\star a\in J}{%
      \FormulaRefAuto{SemigroupIdealRightAbsorption}[thm]{1,3,11,9}}
    \proofstepwidestar[1,2,3,8,9]{x\star a\in I\cap J}{%
      \FormulaRefAuto{x\in A\cap B\eqvdash x\in A\land x\in B}[thm]{%
        \rAI{12,13}}}
    \proofstepwidestar[1,2,3]{%
      \forall x\in I\cap J\,\forall a\in A\,(x\star a\in I\cap J)}{%
      \rUI{\rRI{8,\rUI{\rRI{9,14}}}}}
    \proofstepwidestar[1,2,3]{\mathsf{RId}(I\cap J;A,\star)}{%
      \FormulaRefAuto{SemigroupRightIdealDef}[def]{1,\rAI{6,\rAI{7,15}}}}
  \closeproofpartwideindR

  \proofpartwide{Schluss}
    \proofstepwidestar[1]{\SemigroupAlg{A}{\star}}{\rA}
    \proofstepwidestar[2]{\mathsf{Id}(I;A,\star)}{\rA}
    \proofstepwidestar[3]{\mathsf{Id}(J;A,\star)}{\rA}
    \proofstepwidestar[1,2,3]{\mathsf{LId}(I\cap J;A,\star)}{%
      \FormulaRefAuto{SemigroupIdealIntersectionLeftIdeal}[thm]{1,2,3}}
    \proofstepwidestar[1,2,3]{\mathsf{RId}(I\cap J;A,\star)}{%
      \FormulaRefAuto{SemigroupIdealIntersectionRightIdeal}[thm]{1,2,3}}
    \proofstepwidestar[1,2,3]{\mathsf{Id}(I\cap J;A,\star)}{%
      \FormulaRefAuto{SemigroupIdealDef}[def]{1,\rAI{4,5}}}
  \closeproofpartwide
\end{tabproofsplitwide}

\subsection{Hauptideale}

Statt die symbolische Einheitsadjunktion \(A^1\) informell zu benutzen,
definieren wir Hauptideale unmittelbar durch ihre Elemente. Damit sind auch
die F\"alle ohne linken oder rechten Faktor ausdr\"ucklich erfasst.

\FormulaDefDeltaK[Linkshauptideal]{%
  \SemigroupAlg{A}{\star}\dsep a\in A\vdash
  L_{A,\star}(a)\coloneqq
  \bigl\{b\in A\bigm|b=a\lor
    \exists x\in A\,(b=x\star a)\bigr\}%
}{SemigroupPrincipalLeftIdealDef}{%
  \DeltaRow{Mengen}{A\dsep a\dsep b\dsep x}
  \DeltaRow{Funktionen}{\star}
  \DeltaRow{\textbf{Neues Symbol}}{}
  \DeltaPrem{Linkshauptideal}{L_{A,\star}(a)}
}

\FormulaDefDeltaK[Rechtshauptideal]{%
  \SemigroupAlg{A}{\star}\dsep a\in A\vdash
  R_{A,\star}(a)\coloneqq
  \bigl\{b\in A\bigm|b=a\lor
    \exists y\in A\,(b=a\star y)\bigr\}%
}{SemigroupPrincipalRightIdealDef}{%
  \DeltaRow{Mengen}{A\dsep a\dsep b\dsep y}
  \DeltaRow{Funktionen}{\star}
  \DeltaRow{\textbf{Neues Symbol}}{}
  \DeltaPrem{Rechtshauptideal}{R_{A,\star}(a)}
}

\FormulaDefDeltaK[Zweiseitiges Hauptideal]{%
  \begin{aligned}[t]
    &\SemigroupAlg{A}{\star}\dsep a\in A\vdash{}\\[-2pt]
    &J_{A,\star}(a)\coloneqq
      \Bigl\{b\in A\Bigm|
        b=a\lor\exists x\in A\,(b=x\star a)\\[-2pt]
    &\hspace{93pt}{}\lor\exists y\in A\,(b=a\star y)
        \lor\exists x,y\in A\,(b=(x\star a)\star y)
      \Bigr\}
  \end{aligned}
}{SemigroupPrincipalTwoSidedIdealDef}{%
  \DeltaRow{Mengen}{A\dsep a\dsep b\dsep x\dsep y}
  \DeltaRow{Funktionen}{\star}
  \DeltaRow{\textbf{Neues Symbol}}{}
  \DeltaPrem{Zweiseitiges Hauptideal}{J_{A,\star}(a)}
}

\FormulaThmDeltaK[Elementkriterium des Linkshauptideals]{%
  \begin{aligned}[t]
    &\SemigroupAlg{A}{\star}\dsep a\in A\vdash{}\\[-2pt]
    &b\in L_{A,\star}(a)\leftrightarrow
      \bigl(b\in A\land(b=a\lor\exists x\in A\,(b=x\star a))\bigr)
  \end{aligned}
}{SemigroupPrincipalLeftIdealMembership}{%
  \DeltaRow{Mengen}{A\dsep a\dsep b\dsep x}
  \DeltaRow{Funktionen}{\star}
}
\begin{tabproofwide}
  \proofstepwidestar[1]{\SemigroupAlg{A}{\star}}{\rA}
  \proofstepwidestar[2]{a\in A}{\rA}
  \proofstepwidestar[1,2]{%
    L_{A,\star}(a)=\{u\in A\mid
      u=a\lor\exists x\in A\,(u=x\star a)\}}{%
    \FormulaRefAuto{SemigroupPrincipalLeftIdealDef}[def]{1,2}}
  \proofstepwidestar[1,2]{%
    b\in L_{A,\star}(a)\leftrightarrow
    \bigl(b\in A\land(b=a\lor\exists x\in A\,(b=x\star a))\bigr)}{%
    \rIE{3,\FormulaRefAuto{x\in\{u\in A\mid P(u)\}\eqvdash
      x\in A\land P(x)}[thm]}}
\end{tabproofwide}

\FormulaThmDeltaK[Elementkriterium des Rechtshauptideals]{%
  \begin{aligned}[t]
    &\SemigroupAlg{A}{\star}\dsep a\in A\vdash{}\\[-2pt]
    &b\in R_{A,\star}(a)\leftrightarrow
      \bigl(b\in A\land(b=a\lor\exists y\in A\,(b=a\star y))\bigr)
  \end{aligned}
}{SemigroupPrincipalRightIdealMembership}{%
  \DeltaRow{Mengen}{A\dsep a\dsep b\dsep y}
  \DeltaRow{Funktionen}{\star}
}
\begin{tabproofwide}
  \proofstepwidestar[1]{\SemigroupAlg{A}{\star}}{\rA}
  \proofstepwidestar[2]{a\in A}{\rA}
  \proofstepwidestar[1,2]{%
    R_{A,\star}(a)=\{u\in A\mid
      u=a\lor\exists y\in A\,(u=a\star y)\}}{%
    \FormulaRefAuto{SemigroupPrincipalRightIdealDef}[def]{1,2}}
  \proofstepwidestar[1,2]{%
    b\in R_{A,\star}(a)\leftrightarrow
    \bigl(b\in A\land(b=a\lor\exists y\in A\,(b=a\star y))\bigr)}{%
    \rIE{3,\FormulaRefAuto{x\in\{u\in A\mid P(u)\}\eqvdash
      x\in A\land P(x)}[thm]}}
\end{tabproofwide}

\FormulaThmDeltaK[Elementkriterium des zweiseitigen Hauptideals]{%
  \begin{aligned}[t]
    &\SemigroupAlg{A}{\star}\dsep a\in A\vdash{}\\[-2pt]
    &b\in J_{A,\star}(a)\leftrightarrow
      \Bigl(b\in A\land\bigl(
        b=a\lor\exists x\in A\,(b=x\star a)\\[-2pt]
    &\hspace{112pt}{}\lor\exists y\in A\,(b=a\star y)
        \lor\exists x,y\in A\,(b=(x\star a)\star y)
      \bigr)\Bigr)
  \end{aligned}
}{SemigroupPrincipalTwoSidedIdealMembership}{%
  \DeltaRow{Mengen}{A\dsep a\dsep b\dsep x\dsep y}
  \DeltaRow{Funktionen}{\star}
}
\begin{tabproofwide}
  \proofstepwidestar[1]{\SemigroupAlg{A}{\star}}{\rA}
  \proofstepwidestar[2]{a\in A}{\rA}
  \proofstepwidestarbreak[1,2]{%
    \begin{aligned}[t]
      J_{A,\star}(a)=\Bigl\{u\in A\Bigm|{}&
        u=a\lor\exists x\in A\,(u=x\star a)\\[-2pt]
        &{}\lor\exists y\in A\,(u=a\star y)\\[-2pt]
        &{}\lor\exists x,y\in A\,(u=(x\star a)\star y)\Bigr\}
    \end{aligned}}{%
    \FormulaRefAuto{SemigroupPrincipalTwoSidedIdealDef}[def]{1,2}}
  \proofstepwidestarbreak[1,2]{%
    \begin{aligned}[t]
      b\in J_{A,\star}(a)\leftrightarrow
      \Bigl(b\in A\land\bigl({}&
        b=a\lor\exists x\in A\,(b=x\star a)\\[-2pt]
        &{}\lor\exists y\in A\,(b=a\star y)\\[-2pt]
        &{}\lor\exists x,y\in A\,(b=(x\star a)\star y)\bigr)\Bigr)
    \end{aligned}}{%
    \rIE{3,\FormulaRefAuto{x\in\{u\in A\mid P(u)\}\eqvdash
      x\in A\land P(x)}[thm]}}
\end{tabproofwide}

\FormulaThmDeltaK[Der Erzeuger liegt in seinen Hauptidealen]{%
  \begin{aligned}[t]
    &\SemigroupAlg{A}{\star}\dsep a\in A\vdash{}\\[-2pt]
    &a\in L_{A,\star}(a)\land
      a\in R_{A,\star}(a)\land
      a\in J_{A,\star}(a)
  \end{aligned}
}{SemigroupPrincipalIdealsContainGenerator}{%
  \DeltaRow{Mengen}{A\dsep a}
  \DeltaRow{Funktionen}{\star}
}
\begin{tabproofwide}
  \proofstepwidestar[1]{\SemigroupAlg{A}{\star}}{\rA}
  \proofstepwidestar[2]{a\in A}{\rA}
  \proofstepwidestar[1,2]{a\in L_{A,\star}(a)}{%
    \FormulaRefAuto{SemigroupPrincipalLeftIdealMembership}[thm]{%
      1,2,\rAI{2,\rOIa{\rII}}}}
  \proofstepwidestar[1,2]{a\in R_{A,\star}(a)}{%
    \FormulaRefAuto{SemigroupPrincipalRightIdealMembership}[thm]{%
      1,2,\rAI{2,\rOIa{\rII}}}}
  \proofstepwidestar[1,2]{a\in J_{A,\star}(a)}{%
    \FormulaRefAuto{SemigroupPrincipalTwoSidedIdealMembership}[thm]{%
      1,2,\rAI{2,\rOIa{\rII}}}}
  \proofstepwidestar[1,2]{%
    a\in L_{A,\star}(a)\land a\in R_{A,\star}(a)\land
    a\in J_{A,\star}(a)}{\rAI{3,\rAI{4,5}}}
\end{tabproofwide}

\FormulaThmDeltaK[Das Linkshauptideal ist ein Linksideal]{%
  \SemigroupAlg{A}{\star}\dsep a\in A\vdash
  \mathsf{LId}(L_{A,\star}(a);A,\star)%
}{SemigroupPrincipalLeftIdealIsLeftIdeal}{%
  \DeltaRow{Mengen}{A\dsep a\dsep b\dsep s\dsep x}
  \DeltaRow{Funktionen}{\star}
}
\begin{tabproofwide}
  \proofstepwidestar[1]{\SemigroupAlg{A}{\star}}{\rA}
  \proofstepwidestar[2]{a\in A}{\rA}
  \proofstepwidestar[]{L_{A,\star}(a)\subseteq A}{%
    \FormulaRefAuto{\{x\in A\mid P(x)\}\subseteq A}[thm]}
  \proofstepwidestar[1,2]{a\in L_{A,\star}(a)}{%
    \rAEa{\FormulaRefAuto{SemigroupPrincipalIdealsContainGenerator}[thm]{1,2}}}
  \proofstepwidestar[1,2]{L_{A,\star}(a)\neq\varnothing}{%
    \FormulaRefAuto{a\in A\vdash A\neq\varnothing}[thm]{4}}
  \proofstepwidestar[6]{s\in A}{\rA}
  \proofstepwidestar[7]{b\in L_{A,\star}(a)}{\rA}
  \proofstepwidestar[1,2,7]{b\in A}{%
    \rAEa{\FormulaRefAuto{SemigroupPrincipalLeftIdealMembership}[thm]{1,2,7}}}
  \proofstepwidestar[1,2,6,7]{s\star b\in A}{%
    \FormulaRefAuto{SemigroupClosureAxiom}[ax]{1,6,8}}
  \proofstepwidestar[1,2,7]{b=a\lor\exists x\in A\,(b=x\star a)}{%
    \rAEb{\FormulaRefAuto{SemigroupPrincipalLeftIdealMembership}[thm]{1,2,7}}}
  \proofstepwidestar[11]{b=a}{\rA}
  \proofstepwidestar[6,11]{s\star b=s\star a}{\rIE{11}}
  \proofstepwidestar[6,11]{\exists u\in A\,(s\star b=u\star a)}{%
    \rEI{6,12}}
  \proofstepwidestar[1,2,6,7,11]{s\star b\in L_{A,\star}(a)}{%
    \FormulaRefAuto{SemigroupPrincipalLeftIdealMembership}[thm]{%
      1,2,\rAI{9,\rOIb{13}}}}
  \proofstepwidestar[15]{\exists x\in A\,(b=x\star a)}{\rA}
  \proofstepwidestar[16]{x\in A\land b=x\star a}{\rA}
  \proofstepwidestar[16]{x\in A}{\rAEa{16}}
  \proofstepwidestar[16]{b=x\star a}{\rAEb{16}}
  \proofstepwidestar[1,6,16]{s\star x\in A}{%
    \FormulaRefAuto{SemigroupClosureAxiom}[ax]{1,6,17}}
  \proofstepwidestar[1,2,6,16]{s\star b=(s\star x)\star a}{%
    \rIE{18,\FormulaRefAuto{SemigroupAssociativityAxiom}[ax]{1,6,17,2}}}
  \proofstepwidestar[1,2,6,16]{%
    \exists u\in A\,(s\star b=u\star a)}{\rEI{19,20}}
  \proofstepwidestar[1,2,6,7,16]{s\star b\in L_{A,\star}(a)}{%
    \FormulaRefAuto{SemigroupPrincipalLeftIdealMembership}[thm]{%
      1,2,\rAI{9,\rOIb{21}}}}
  \proofstepwidestar[1,2,6,7,15]{s\star b\in L_{A,\star}(a)}{%
    \rEE{15,16,22}}
  \proofstepwidestar[1,2,6,7]{s\star b\in L_{A,\star}(a)}{%
    \rOE{10,11,14,15,23}}
  \proofstepwidestar[1,2]{%
    \forall s\in A\,\forall b\in L_{A,\star}(a)\,
      (s\star b\in L_{A,\star}(a))}{%
    \rUI{\rRI{6,\rUI{\rRI{7,24}}}}}
  \proofstepwidestar[1,2]{\mathsf{LId}(L_{A,\star}(a);A,\star)}{%
    \FormulaRefAuto{SemigroupLeftIdealDef}[def]{1,\rAI{3,\rAI{5,25}}}}
\end{tabproofwide}

\FormulaThmDeltaK[Das Rechtshauptideal ist ein Rechtsideal]{%
  \SemigroupAlg{A}{\star}\dsep a\in A\vdash
  \mathsf{RId}(R_{A,\star}(a);A,\star)%
}{SemigroupPrincipalRightIdealIsRightIdeal}{%
  \DeltaRow{Mengen}{A\dsep a\dsep b\dsep s\dsep y}
  \DeltaRow{Funktionen}{\star}
}
\begin{tabproofwide}
  \proofstepwidestar[1]{\SemigroupAlg{A}{\star}}{\rA}
  \proofstepwidestar[2]{a\in A}{\rA}
  \proofstepwidestar[]{R_{A,\star}(a)\subseteq A}{%
    \FormulaRefAuto{\{x\in A\mid P(x)\}\subseteq A}[thm]}
  \proofstepwidestar[1,2]{a\in R_{A,\star}(a)}{%
    \rAEa{\rAEb{\FormulaRefAuto{SemigroupPrincipalIdealsContainGenerator}[thm]{1,2}}}}
  \proofstepwidestar[1,2]{R_{A,\star}(a)\neq\varnothing}{%
    \FormulaRefAuto{a\in A\vdash A\neq\varnothing}[thm]{4}}
  \proofstepwidestar[6]{b\in R_{A,\star}(a)}{\rA}
  \proofstepwidestar[7]{s\in A}{\rA}
  \proofstepwidestar[1,2,6]{b\in A}{%
    \rAEa{\FormulaRefAuto{SemigroupPrincipalRightIdealMembership}[thm]{1,2,6}}}
  \proofstepwidestar[1,2,6,7]{b\star s\in A}{%
    \FormulaRefAuto{SemigroupClosureAxiom}[ax]{1,8,7}}
  \proofstepwidestar[1,2,6]{b=a\lor\exists y\in A\,(b=a\star y)}{%
    \rAEb{\FormulaRefAuto{SemigroupPrincipalRightIdealMembership}[thm]{1,2,6}}}
  \proofstepwidestar[11]{b=a}{\rA}
  \proofstepwidestar[7,11]{b\star s=a\star s}{\rIE{11}}
  \proofstepwidestar[7,11]{\exists u\in A\,(b\star s=a\star u)}{\rEI{7,12}}
  \proofstepwidestar[1,2,6,7,11]{b\star s\in R_{A,\star}(a)}{%
    \FormulaRefAuto{SemigroupPrincipalRightIdealMembership}[thm]{%
      1,2,\rAI{9,\rOIb{13}}}}
  \proofstepwidestar[15]{\exists y\in A\,(b=a\star y)}{\rA}
  \proofstepwidestar[16]{y\in A\land b=a\star y}{\rA}
  \proofstepwidestar[16]{y\in A}{\rAEa{16}}
  \proofstepwidestar[16]{b=a\star y}{\rAEb{16}}
  \proofstepwidestar[1,7,16]{y\star s\in A}{%
    \FormulaRefAuto{SemigroupClosureAxiom}[ax]{1,17,7}}
  \proofstepwidestar[1,2,7,16]{b\star s=a\star(y\star s)}{%
    \rIE{18,\FormulaRefAuto{SemigroupAssociativityAxiom}[ax]{1,2,17,7}}}
  \proofstepwidestar[1,2,7,16]{\exists u\in A\,(b\star s=a\star u)}{\rEI{19,20}}
  \proofstepwidestar[1,2,6,7,16]{b\star s\in R_{A,\star}(a)}{%
    \FormulaRefAuto{SemigroupPrincipalRightIdealMembership}[thm]{%
      1,2,\rAI{9,\rOIb{21}}}}
  \proofstepwidestar[1,2,6,7,15]{b\star s\in R_{A,\star}(a)}{\rEE{15,16,22}}
  \proofstepwidestar[1,2,6,7]{b\star s\in R_{A,\star}(a)}{%
    \rOE{10,11,14,15,23}}
  \proofstepwidestar[1,2]{%
    \forall b\in R_{A,\star}(a)\,\forall s\in A\,
      (b\star s\in R_{A,\star}(a))}{%
    \rUI{\rRI{6,\rUI{\rRI{7,24}}}}}
  \proofstepwidestar[1,2]{\mathsf{RId}(R_{A,\star}(a);A,\star)}{%
    \FormulaRefAuto{SemigroupRightIdealDef}[def]{1,\rAI{3,\rAI{5,25}}}}
\end{tabproofwide}

\FormulaThmDeltaK[Minimalität des Linkshauptideals]{%
  \begin{aligned}[t]
    &\SemigroupAlg{A}{\star}\dsep a\in A\dsep
      \mathsf{LId}(I;A,\star)\dsep a\in I\vdash{}\\[-2pt]
    &L_{A,\star}(a)\subseteq I
  \end{aligned}
}{SemigroupPrincipalLeftIdealMinimal}{%
  \DeltaRow{Mengen}{A\dsep I\dsep a\dsep b\dsep x}
  \DeltaRow{Funktionen}{\star}
}
\begin{tabproofwide}
  \proofstepwidestar[1]{\SemigroupAlg{A}{\star}}{\rA}
  \proofstepwidestar[2]{a\in A}{\rA}
  \proofstepwidestar[3]{\mathsf{LId}(I;A,\star)}{\rA}
  \proofstepwidestar[4]{a\in I}{\rA}
  \proofstepwidestar[5]{b\in L_{A,\star}(a)}{\rA}
  \proofstepwidestar[1,2,5]{b=a\lor\exists x\in A\,(b=x\star a)}{%
    \rAEb{\FormulaRefAuto{SemigroupPrincipalLeftIdealMembership}[thm]{1,2,5}}}
  \proofstepwidestar[7]{b=a}{\rA}
  \proofstepwidestar[4,7]{b\in I}{\rIE{\FormulaRefAuto{a=b\vdash b=a}[thm]{7},4}}
  \proofstepwidestar[9]{\exists x\in A\,(b=x\star a)}{\rA}
  \proofstepwidestar[10]{x\in A\land b=x\star a}{\rA}
  \proofstepwidestar[10]{x\in A}{\rAEa{10}}
  \proofstepwidestar[10]{b=x\star a}{\rAEb{10}}
  \proofstepwidestar[1,3]{%
    \forall u\in A\,\forall v\in I\,(u\star v\in I)}{%
    \rAEb{\rAEb{\FormulaRefAuto{SemigroupLeftIdealDef}[def]{1,3}}}}
  \proofstepwidestar[1,3,4,10]{x\star a\in I}{\rRE{4,\rUE{\rRE{11,\rUE{13}}}}}
  \proofstepwidestar[1,3,4,10]{b\in I}{\rIE{12,14}}
  \proofstepwidestar[1,3,4,9]{b\in I}{\rEE{9,10,15}}
  \proofstepwidestar[1,2,3,4,5]{b\in I}{\rOE{6,7,8,9,16}}
  \proofstepwidestar[1,2,3,4]{L_{A,\star}(a)\subseteq I}{%
    \FormulaRefAuto{A\subseteq B\coloneq
      \forall x\,(x\in A\rightarrow x\in B)}[def]{\rUI{\rRI{5,17}}}}
\end{tabproofwide}

\FormulaThmDeltaK[Minimalität des Rechtshauptideals]{%
  \begin{aligned}[t]
    &\SemigroupAlg{A}{\star}\dsep a\in A\dsep
      \mathsf{RId}(I;A,\star)\dsep a\in I\vdash{}\\[-2pt]
    &R_{A,\star}(a)\subseteq I
  \end{aligned}
}{SemigroupPrincipalRightIdealMinimal}{%
  \DeltaRow{Mengen}{A\dsep I\dsep a\dsep b\dsep y}
  \DeltaRow{Funktionen}{\star}
}
\begin{tabproofwide}
  \proofstepwidestar[1]{\SemigroupAlg{A}{\star}}{\rA}
  \proofstepwidestar[2]{a\in A}{\rA}
  \proofstepwidestar[3]{\mathsf{RId}(I;A,\star)}{\rA}
  \proofstepwidestar[4]{a\in I}{\rA}
  \proofstepwidestar[5]{b\in R_{A,\star}(a)}{\rA}
  \proofstepwidestar[1,2,5]{b=a\lor\exists y\in A\,(b=a\star y)}{%
    \rAEb{\FormulaRefAuto{SemigroupPrincipalRightIdealMembership}[thm]{1,2,5}}}
  \proofstepwidestar[7]{b=a}{\rA}
  \proofstepwidestar[4,7]{b\in I}{\rIE{\FormulaRefAuto{a=b\vdash b=a}[thm]{7},4}}
  \proofstepwidestar[9]{\exists y\in A\,(b=a\star y)}{\rA}
  \proofstepwidestar[10]{y\in A\land b=a\star y}{\rA}
  \proofstepwidestar[10]{y\in A}{\rAEa{10}}
  \proofstepwidestar[10]{b=a\star y}{\rAEb{10}}
  \proofstepwidestar[1,3]{%
    \forall u\in I\,\forall v\in A\,(u\star v\in I)}{%
    \rAEb{\rAEb{\FormulaRefAuto{SemigroupRightIdealDef}[def]{1,3}}}}
  \proofstepwidestar[1,3,4,10]{a\star y\in I}{\rRE{11,\rUE{\rRE{4,\rUE{13}}}}}
  \proofstepwidestar[1,3,4,10]{b\in I}{\rIE{12,14}}
  \proofstepwidestar[1,3,4,9]{b\in I}{\rEE{9,10,15}}
  \proofstepwidestar[1,2,3,4,5]{b\in I}{\rOE{6,7,8,9,16}}
  \proofstepwidestar[1,2,3,4]{R_{A,\star}(a)\subseteq I}{%
    \FormulaRefAuto{A\subseteq B\coloneq
      \forall x\,(x\in A\rightarrow x\in B)}[def]{\rUI{\rRI{5,17}}}}
\end{tabproofwide}

\FormulaThmDeltaK[Das zweiseitige Hauptideal ist ein Ideal]{%
  \SemigroupAlg{A}{\star}\dsep a\in A\vdash
  \mathsf{Id}(J_{A,\star}(a);A,\star)%
}{SemigroupPrincipalTwoSidedIdealIsIdeal}{%
  \DeltaRow{Mengen}{A\dsep a\dsep b\dsep s\dsep x\dsep y}
  \DeltaRow{Funktionen}{\star}
}
Nach dem Elementkriterium besitzt jedes Element des Hauptideals eine von
vier Formen: Es ist der Erzeuger selbst, hat nur einen linken oder rechten
Faktor oder hat beide Faktoren.  Multiplikation von links beziehungsweise
rechts erh\"alt jede dieser Formen; die beiden Absorptionseigenschaften
werden deshalb getrennt bewiesen und erst im Schluss zusammengef\"uhrt.

\begin{tabproofsplitwide}
  \proofpartwideindR[Linksidealeigenschaft]{%
    \SemigroupAlg{A}{\star}\dsep a\in A\vdash
    \mathsf{LId}(J_{A,\star}(a);A,\star)
  }{%
    \SemigroupAlg{A}{\star}\dsep a\in A\vdash
    \mathsf{LId}(J_{A,\star}(a);A,\star)
  }[SemigroupPrincipalTwoSidedIdealLeftPart]
    \proofstepwidestar[1]{\SemigroupAlg{A}{\star}}{\rA}
    \proofstepwidestar[2]{a\in A}{\rA}
    \proofstepwidestar[]{J_{A,\star}(a)\subseteq A}{%
      \FormulaRefAuto{\{x\in A\mid P(x)\}\subseteq A}[thm]}
    \proofstepwidestar[1,2]{a\in J_{A,\star}(a)}{%
      \rAEb{\rAEb{\FormulaRefAuto{SemigroupPrincipalIdealsContainGenerator}[thm]{1,2}}}}
    \proofstepwidestar[1,2]{J_{A,\star}(a)\neq\varnothing}{%
      \FormulaRefAuto{a\in A\vdash A\neq\varnothing}[thm]{4}}
    \proofstepwidestar[6]{s\in A}{\rA}
    \proofstepwidestar[7]{b\in J_{A,\star}(a)}{\rA}
    \proofstepwidestar[1,2,7]{b\in A}{%
      \rAEa{\FormulaRefAuto{SemigroupPrincipalTwoSidedIdealMembership}[thm]{1,2,7}}}
    \proofstepwidestar[1,2,6,7]{s\star b\in A}{%
      \FormulaRefAuto{SemigroupClosureAxiom}[ax]{1,6,8}}
    \proofstepwidestarbreak[1,2,7]{%
      \begin{aligned}[t]
        b=a\lor{}&\exists x\in A\,(b=x\star a)\\[-2pt]
        &{}\lor\exists y\in A\,(b=a\star y)\\[-2pt]
        &{}\lor\exists x,y\in A\,(b=(x\star a)\star y)
      \end{aligned}}{%
      \rAEb{\FormulaRefAuto{SemigroupPrincipalTwoSidedIdealMembership}[thm]{1,2,7}}}
    \proofstepwidestar[11]{b=a}{\rA}
    \proofstepwidestar[6,11]{\exists u\in A\,(s\star b=u\star a)}{%
      \rEI{6,\rIE{11}}}
    \proofstepwidestar[1,2,6,7,11]{s\star b\in J_{A,\star}(a)}{%
      \FormulaRefAuto{SemigroupPrincipalTwoSidedIdealMembership}[thm]{%
        1,2,\rAI{9,\rOIb{12}}}}
    \proofstepwidestar[14]{\exists x\in A\,(b=x\star a)}{\rA}
    \proofstepwidestar[15]{x\in A\land b=x\star a}{\rA}
    \proofstepwidestar[15]{x\in A}{\rAEa{15}}
    \proofstepwidestar[15]{b=x\star a}{\rAEb{15}}
    \proofstepwidestar[1,6,15]{s\star x\in A}{%
      \FormulaRefAuto{SemigroupClosureAxiom}[ax]{1,6,16}}
    \proofstepwidestar[1,2,6,15]{s\star b=(s\star x)\star a}{%
      \rIE{17,\FormulaRefAuto{SemigroupAssociativityAxiom}[ax]{1,6,16,2}}}
    \proofstepwidestar[1,2,6,15]{\exists u\in A\,(s\star b=u\star a)}{\rEI{18,19}}
    \proofstepwidestar[1,2,6,7,15]{s\star b\in J_{A,\star}(a)}{%
      \FormulaRefAuto{SemigroupPrincipalTwoSidedIdealMembership}[thm]{%
        1,2,\rAI{9,\rOIb{20}}}}
    \proofstepwidestar[22]{\exists y\in A\,(b=a\star y)}{\rA}
    \proofstepwidestar[23]{y\in A\land b=a\star y}{\rA}
    \proofstepwidestar[23]{y\in A}{\rAEa{23}}
    \proofstepwidestar[23]{b=a\star y}{\rAEb{23}}
    \proofstepwidestar[1,2,6,23]{s\star b=(s\star a)\star y}{%
      \rIE{25,\FormulaRefAuto{SemigroupAssociativityAxiom}[ax]{1,6,2,24}}}
    \proofstepwidestar[1,2,6,23]{%
      \exists u,v\in A\,(s\star b=(u\star a)\star v)}{\rEI{6,\rEI{24,26}}}
    \proofstepwidestarbreak[1,2,6,7,23]{s\star b\in J_{A,\star}(a)}{%
      \FormulaRefAuto{SemigroupPrincipalTwoSidedIdealMembership}[thm]{%
        1,2,\rAI{9,\rOIb{\rOIb{\rOIb{27}}}}}}
    \proofstepwidestar[29]{\exists x,y\in A\,(b=(x\star a)\star y)}{\rA}
    \proofstepwidestar[30]{x\in A\land(y\in A\land b=(x\star a)\star y)}{\rA}
    \proofstepwidestar[30]{x\in A}{\rAEa{30}}
    \proofstepwidestar[30]{y\in A}{\rAEa{\rAEb{30}}}
    \proofstepwidestar[30]{b=(x\star a)\star y}{\rAEb{\rAEb{30}}}
    \proofstepwidestar[1,6,30]{s\star x\in A}{%
      \FormulaRefAuto{SemigroupClosureAxiom}[ax]{1,6,31}}
    \proofstepwidestarbreak[1,2,6,30]{%
      s\star b=((s\star x)\star a)\star y}{%
      \rIE{33,\FormulaRefAuto{a = b,\, b = c \vdash a = c}[thm]{%
        \FormulaRefAuto{SemigroupAssociativityAxiom}[ax]{%
          1,6,\FormulaRefAuto{SemigroupClosureAxiom}[ax]{1,31,2},32},
        \rIE{\FormulaRefAuto{SemigroupAssociativityAxiom}[ax]{1,6,31,2}}}}}
    \proofstepwidestar[1,2,6,30]{%
      \exists u,v\in A\,(s\star b=(u\star a)\star v)}{\rEI{34,\rEI{32,35}}}
    \proofstepwidestarbreak[1,2,6,7,30]{s\star b\in J_{A,\star}(a)}{%
      \FormulaRefAuto{SemigroupPrincipalTwoSidedIdealMembership}[thm]{%
        1,2,\rAI{9,\rOIb{\rOIb{\rOIb{36}}}}}}
    \proofstepwidestar[1,2,6,7,29]{s\star b\in J_{A,\star}(a)}{\rEE{29,30,37}}
    \proofstepwidestar[1,2,6,7,22]{s\star b\in J_{A,\star}(a)}{\rEE{22,23,28}}
    \proofstepwidestar[1,2,6,7,14]{s\star b\in J_{A,\star}(a)}{\rEE{14,15,21}}
    \proofstepwidestar[1,2,6,7]{s\star b\in J_{A,\star}(a)}{%
      \rOE{10,11,13,14,39,22,38,29,37}}
    \proofstepwidestar[1,2]{%
      \forall s\in A\,\forall b\in J_{A,\star}(a)\,
      (s\star b\in J_{A,\star}(a))}{%
      \rUI{\rRI{6,\rUI{\rRI{7,40}}}}}
    \proofstepwidestar[1,2]{\mathsf{LId}(J_{A,\star}(a);A,\star)}{%
      \FormulaRefAuto{SemigroupLeftIdealDef}[def]{1,\rAI{3,\rAI{5,41}}}}
  \closeproofpartwideindR

  \proofpartwideindR[Rechtsidealeigenschaft]{%
    \SemigroupAlg{A}{\star}\dsep a\in A\vdash
    \mathsf{RId}(J_{A,\star}(a);A,\star)
  }{%
    \SemigroupAlg{A}{\star}\dsep a\in A\vdash
    \mathsf{RId}(J_{A,\star}(a);A,\star)
  }[SemigroupPrincipalTwoSidedIdealRightPart]
    \proofstepwidestar[1]{\SemigroupAlg{A}{\star}}{\rA}
    \proofstepwidestar[2]{a\in A}{\rA}
    \proofstepwidestar[]{J_{A,\star}(a)\subseteq A}{%
      \FormulaRefAuto{\{x\in A\mid P(x)\}\subseteq A}[thm]}
    \proofstepwidestar[1,2]{a\in J_{A,\star}(a)}{%
      \rAEb{\rAEb{\FormulaRefAuto{SemigroupPrincipalIdealsContainGenerator}[thm]{1,2}}}}
    \proofstepwidestar[1,2]{J_{A,\star}(a)\neq\varnothing}{%
      \FormulaRefAuto{a\in A\vdash A\neq\varnothing}[thm]{4}}
    \proofstepwidestar[6]{b\in J_{A,\star}(a)}{\rA}
    \proofstepwidestar[7]{s\in A}{\rA}
    \proofstepwidestar[1,2,6]{b\in A}{%
      \rAEa{\FormulaRefAuto{SemigroupPrincipalTwoSidedIdealMembership}[thm]{1,2,6}}}
    \proofstepwidestar[1,2,6,7]{b\star s\in A}{%
      \FormulaRefAuto{SemigroupClosureAxiom}[ax]{1,8,7}}
    \proofstepwidestarbreak[1,2,6]{%
      \begin{aligned}[t]
        b=a\lor{}&\exists x\in A\,(b=x\star a)\\[-2pt]
        &{}\lor\exists y\in A\,(b=a\star y)\\[-2pt]
        &{}\lor\exists x,y\in A\,(b=(x\star a)\star y)
      \end{aligned}}{%
      \rAEb{\FormulaRefAuto{SemigroupPrincipalTwoSidedIdealMembership}[thm]{1,2,6}}}
    \proofstepwidestar[11]{b=a}{\rA}
    \proofstepwidestar[7,11]{\exists v\in A\,(b\star s=a\star v)}{\rEI{7,\rIE{11}}}
    \proofstepwidestarbreak[1,2,6,7,11]{b\star s\in J_{A,\star}(a)}{%
      \FormulaRefAuto{SemigroupPrincipalTwoSidedIdealMembership}[thm]{%
        1,2,\rAI{9,\rOIb{\rOIb{12}}}}}
    \proofstepwidestar[14]{\exists x\in A\,(b=x\star a)}{\rA}
    \proofstepwidestar[15]{x\in A\land b=x\star a}{\rA}
    \proofstepwidestar[15]{x\in A}{\rAEa{15}}
    \proofstepwidestar[15]{b=x\star a}{\rAEb{15}}
    \proofstepwidestar[1,2,7,15]{b\star s=(x\star a)\star s}{\rIE{17}}
    \proofstepwidestar[1,2,7,15]{%
      \exists u,v\in A\,(b\star s=(u\star a)\star v)}{\rEI{16,\rEI{7,18}}}
    \proofstepwidestarbreak[1,2,6,7,15]{b\star s\in J_{A,\star}(a)}{%
      \FormulaRefAuto{SemigroupPrincipalTwoSidedIdealMembership}[thm]{%
        1,2,\rAI{9,\rOIb{\rOIb{\rOIb{19}}}}}}
    \proofstepwidestar[21]{\exists y\in A\,(b=a\star y)}{\rA}
    \proofstepwidestar[22]{y\in A\land b=a\star y}{\rA}
    \proofstepwidestar[22]{y\in A}{\rAEa{22}}
    \proofstepwidestar[22]{b=a\star y}{\rAEb{22}}
    \proofstepwidestar[1,7,22]{y\star s\in A}{%
      \FormulaRefAuto{SemigroupClosureAxiom}[ax]{1,23,7}}
    \proofstepwidestar[1,2,7,22]{b\star s=a\star(y\star s)}{%
      \rIE{24,\FormulaRefAuto{SemigroupAssociativityAxiom}[ax]{1,2,23,7}}}
    \proofstepwidestar[1,2,7,22]{\exists v\in A\,(b\star s=a\star v)}{\rEI{25,26}}
    \proofstepwidestarbreak[1,2,6,7,22]{b\star s\in J_{A,\star}(a)}{%
      \FormulaRefAuto{SemigroupPrincipalTwoSidedIdealMembership}[thm]{%
        1,2,\rAI{9,\rOIb{\rOIb{27}}}}}
    \proofstepwidestar[29]{\exists x,y\in A\,(b=(x\star a)\star y)}{\rA}
    \proofstepwidestar[30]{x\in A\land(y\in A\land b=(x\star a)\star y)}{\rA}
    \proofstepwidestar[30]{x\in A}{\rAEa{30}}
    \proofstepwidestar[30]{y\in A}{\rAEa{\rAEb{30}}}
    \proofstepwidestar[30]{b=(x\star a)\star y}{\rAEb{\rAEb{30}}}
    \proofstepwidestar[1,7,30]{y\star s\in A}{%
      \FormulaRefAuto{SemigroupClosureAxiom}[ax]{1,32,7}}
    \proofstepwidestarbreak[1,2,7,30]{b\star s=(x\star a)\star(y\star s)}{%
      \rIE{33,\FormulaRefAuto{SemigroupAssociativityAxiom}[ax]{%
        1,\FormulaRefAuto{SemigroupClosureAxiom}[ax]{1,31,2},32,7}}}
    \proofstepwidestar[1,2,7,30]{%
      \exists u,v\in A\,(b\star s=(u\star a)\star v)}{\rEI{31,\rEI{34,35}}}
    \proofstepwidestarbreak[1,2,6,7,30]{b\star s\in J_{A,\star}(a)}{%
      \FormulaRefAuto{SemigroupPrincipalTwoSidedIdealMembership}[thm]{%
        1,2,\rAI{9,\rOIb{\rOIb{\rOIb{36}}}}}}
    \proofstepwidestar[1,2,6,7,29]{b\star s\in J_{A,\star}(a)}{\rEE{29,30,37}}
    \proofstepwidestar[1,2,6,7,21]{b\star s\in J_{A,\star}(a)}{\rEE{21,22,28}}
    \proofstepwidestar[1,2,6,7,14]{b\star s\in J_{A,\star}(a)}{\rEE{14,15,20}}
    \proofstepwidestar[1,2,6,7]{b\star s\in J_{A,\star}(a)}{%
      \rOE{10,11,13,14,39,21,38,29,37}}
    \proofstepwidestar[1,2]{%
      \forall b\in J_{A,\star}(a)\,\forall s\in A\,
      (b\star s\in J_{A,\star}(a))}{%
      \rUI{\rRI{6,\rUI{\rRI{7,40}}}}}
    \proofstepwidestar[1,2]{\mathsf{RId}(J_{A,\star}(a);A,\star)}{%
      \FormulaRefAuto{SemigroupRightIdealDef}[def]{1,\rAI{3,\rAI{5,41}}}}
  \closeproofpartwideindR

  \proofpartwide{Schluss}
    \proofstepwidestar[1]{\SemigroupAlg{A}{\star}}{\rA}
    \proofstepwidestar[2]{a\in A}{\rA}
    \proofstepwidestar[1,2]{\mathsf{LId}(J_{A,\star}(a);A,\star)}{%
      \FormulaRefAuto{SemigroupPrincipalTwoSidedIdealLeftPart}[thm]{1,2}}
    \proofstepwidestar[1,2]{\mathsf{RId}(J_{A,\star}(a);A,\star)}{%
      \FormulaRefAuto{SemigroupPrincipalTwoSidedIdealRightPart}[thm]{1,2}}
    \proofstepwidestar[1,2]{\mathsf{Id}(J_{A,\star}(a);A,\star)}{%
      \FormulaRefAuto{SemigroupIdealDef}[def]{1,\rAI{3,4}}}
  \closeproofpartwide
\end{tabproofsplitwide}

\FormulaThmDeltaK[Minimalit\"at des zweiseitigen Hauptideals]{%
  \begin{aligned}[t]
    &\SemigroupAlg{A}{\star}\dsep a\in A\dsep
      \mathsf{Id}(I;A,\star)\dsep a\in I\vdash{}\\[-2pt]
    &J_{A,\star}(a)\subseteq I
  \end{aligned}
}{SemigroupPrincipalTwoSidedIdealMinimal}{%
  \DeltaRow{Mengen}{A\dsep I\dsep a\dsep b\dsep x\dsep y}
  \DeltaRow{Funktionen}{\star}
}
Das Elementkriterium liefert vier und nur vier Formen.  In jedem Fall folgt
die Zugeh\"origkeit zu (I) unmittelbar aus (a\in I) und der linken
beziehungsweise rechten Idealeigenschaft.

\begin{tabproofsplitwide}
  \proofpartwideindR[Erzeugerfall]{%
    \SemigroupAlg{A}{\star}\dsep a\in A\dsep\mathsf{Id}(I;A,\star)\dsep
    a\in I\dsep b=a\vdash b\in I
  }{%
    \SemigroupAlg{A}{\star}\dsep a\in A\dsep\mathsf{Id}(I;A,\star)\dsep
    a\in I\dsep b=a\vdash b\in I
  }[SemigroupPrincipalTwoSidedIdealMinimalGeneratorCase]
    \proofstepwidestar[1]{\SemigroupAlg{A}{\star}}{\rA}
    \proofstepwidestar[2]{a\in A}{\rA}
    \proofstepwidestar[3]{\mathsf{Id}(I;A,\star)}{\rA}
    \proofstepwidestar[4]{a\in I}{\rA}
    \proofstepwidestar[5]{b=a}{\rA}
    \proofstepwidestar[4,5]{b\in I}{%
      \rIE{\FormulaRefAuto{a=b\vdash b=a}[thm]{5},4}}
  \closeproofpartwideindR

  \proofpartwideindR[Fall eines linken Faktors]{%
    \begin{aligned}[t]
      &\SemigroupAlg{A}{\star}\dsep a\in A\dsep
        \mathsf{Id}(I;A,\star)\dsep a\in I\dsep{}\\[-2pt]
      &\exists x\in A\,(b=x\star a)\vdash b\in I
    \end{aligned}
  }{%
    \SemigroupAlg{A}{\star}\dsep a\in A\dsep
    \mathsf{Id}(I;A,\star)\dsep a\in I\dsep
    \exists x\in A\,(b=x\star a)\vdash b\in I
  }[SemigroupPrincipalTwoSidedIdealMinimalLeftFactorCase]
    \proofstepwidestar[1]{\SemigroupAlg{A}{\star}}{\rA}
    \proofstepwidestar[2]{a\in A}{\rA}
    \proofstepwidestar[3]{\mathsf{Id}(I;A,\star)}{\rA}
    \proofstepwidestar[4]{a\in I}{\rA}
    \proofstepwidestar[5]{\exists x\in A\,(b=x\star a)}{\rA}
    \proofstepwidestar[6]{x\in A\land b=x\star a}{\rA}
    \proofstepwidestar[6]{x\in A}{\rAEa{6}}
    \proofstepwidestar[6]{b=x\star a}{\rAEb{6}}
    \proofstepwidestar[1,3,4,6]{x\star a\in I}{%
      \FormulaRefAuto{SemigroupIdealLeftAbsorption}[thm]{1,3,7,4}}
    \proofstepwidestar[1,3,4,6]{b\in I}{\rIE{8,9}}
    \proofstepwidestar[1,3,4,5]{b\in I}{\rEE{5,6,10}}
  \closeproofpartwideindR

  \proofpartwideindR[Fall eines rechten Faktors]{%
    \begin{aligned}[t]
      &\SemigroupAlg{A}{\star}\dsep a\in A\dsep
        \mathsf{Id}(I;A,\star)\dsep a\in I\dsep{}\\[-2pt]
      &\exists y\in A\,(b=a\star y)\vdash b\in I
    \end{aligned}
  }{%
    \SemigroupAlg{A}{\star}\dsep a\in A\dsep
    \mathsf{Id}(I;A,\star)\dsep a\in I\dsep
    \exists y\in A\,(b=a\star y)\vdash b\in I
  }[SemigroupPrincipalTwoSidedIdealMinimalRightFactorCase]
    \proofstepwidestar[1]{\SemigroupAlg{A}{\star}}{\rA}
    \proofstepwidestar[2]{a\in A}{\rA}
    \proofstepwidestar[3]{\mathsf{Id}(I;A,\star)}{\rA}
    \proofstepwidestar[4]{a\in I}{\rA}
    \proofstepwidestar[5]{\exists y\in A\,(b=a\star y)}{\rA}
    \proofstepwidestar[6]{y\in A\land b=a\star y}{\rA}
    \proofstepwidestar[6]{y\in A}{\rAEa{6}}
    \proofstepwidestar[6]{b=a\star y}{\rAEb{6}}
    \proofstepwidestar[1,3,4,6]{a\star y\in I}{%
      \FormulaRefAuto{SemigroupIdealRightAbsorption}[thm]{1,3,4,7}}
    \proofstepwidestar[1,3,4,6]{b\in I}{\rIE{8,9}}
    \proofstepwidestar[1,3,4,5]{b\in I}{\rEE{5,6,10}}
  \closeproofpartwideindR

  \proofpartwideindR[Fall zweier Faktoren]{%
    \begin{aligned}[t]
      &\SemigroupAlg{A}{\star}\dsep a\in A\dsep
        \mathsf{Id}(I;A,\star)\dsep a\in I\dsep{}\\[-2pt]
      &\exists x,y\in A\,(b=(x\star a)\star y)\vdash b\in I
    \end{aligned}
  }{%
    \begin{aligned}[t]
      &\SemigroupAlg{A}{\star}\dsep a\in A\dsep
        \mathsf{Id}(I;A,\star)\dsep a\in I\dsep{}\\[-2pt]
      &\exists x,y\in A\,(b=(x\star a)\star y)\vdash b\in I
    \end{aligned}
  }[SemigroupPrincipalTwoSidedIdealMinimalTwoFactorCase]
    \proofstepwidestar[1]{\SemigroupAlg{A}{\star}}{\rA}
    \proofstepwidestar[2]{a\in A}{\rA}
    \proofstepwidestar[3]{\mathsf{Id}(I;A,\star)}{\rA}
    \proofstepwidestar[4]{a\in I}{\rA}
    \proofstepwidestar[5]{\exists x,y\in A\,(b=(x\star a)\star y)}{\rA}
    \proofstepwidestar[6]{%
      x\in A\land\bigl(y\in A\land b=(x\star a)\star y\bigr)}{\rA}
    \proofstepwidestar[6]{x\in A}{\rAEa{6}}
    \proofstepwidestar[6]{y\in A}{\rAEa{\rAEb{6}}}
    \proofstepwidestar[6]{b=(x\star a)\star y}{\rAEb{\rAEb{6}}}
    \proofstepwidestar[1,3,4,6]{x\star a\in I}{%
      \FormulaRefAuto{SemigroupIdealLeftAbsorption}[thm]{1,3,7,4}}
    \proofstepwidestar[1,3,4,6]{(x\star a)\star y\in I}{%
      \FormulaRefAuto{SemigroupIdealRightAbsorption}[thm]{1,3,10,8}}
    \proofstepwidestar[1,3,4,6]{b\in I}{\rIE{9,11}}
    \proofstepwidestar[1,3,4,5]{b\in I}{\rEE{5,6,12}}
  \closeproofpartwideindR

  \proofpartwide{Fallunterscheidung und Schluss}
    \proofstepwidestar[1]{\SemigroupAlg{A}{\star}}{\rA}
    \proofstepwidestar[2]{a\in A}{\rA}
    \proofstepwidestar[3]{\mathsf{Id}(I;A,\star)}{\rA}
    \proofstepwidestar[4]{a\in I}{\rA}
    \proofstepwidestar[5]{b\in J_{A,\star}(a)}{\rA}
    \proofstepwidestarbreak[1,2,5]{%
      \begin{aligned}[t]
        b=a\lor{}&\exists x\in A\,(b=x\star a)\\[-2pt]
        &{}\lor\exists y\in A\,(b=a\star y)\\[-2pt]
        &{}\lor\exists x,y\in A\,(b=(x\star a)\star y)
      \end{aligned}}{%
      \rAEb{\FormulaRefAuto{SemigroupPrincipalTwoSidedIdealMembership}[thm]{1,2,5}}}
    \proofstepwidestar[7]{b=a}{\rA}
    \proofstepwidestar[1,2,3,4,7]{b\in I}{%
      \FormulaRefAuto{SemigroupPrincipalTwoSidedIdealMinimalGeneratorCase}[thm]{%
        1,2,3,4,7}}
    \proofstepwidestar[9]{\exists x\in A\,(b=x\star a)}{\rA}
    \proofstepwidestar[1,2,3,4,9]{b\in I}{%
      \FormulaRefAuto{SemigroupPrincipalTwoSidedIdealMinimalLeftFactorCase}[thm]{%
        1,2,3,4,9}}
    \proofstepwidestar[11]{\exists y\in A\,(b=a\star y)}{\rA}
    \proofstepwidestar[1,2,3,4,11]{b\in I}{%
      \FormulaRefAuto{SemigroupPrincipalTwoSidedIdealMinimalRightFactorCase}[thm]{%
        1,2,3,4,11}}
    \proofstepwidestar[13]{\exists x,y\in A\,(b=(x\star a)\star y)}{\rA}
    \proofstepwidestar[1,2,3,4,13]{b\in I}{%
      \FormulaRefAuto{SemigroupPrincipalTwoSidedIdealMinimalTwoFactorCase}[thm]{%
        1,2,3,4,13}}
    \proofstepwidestar[1,2,3,4,5]{b\in I}{%
      \rOE{6,7,8,9,10,11,12,13,14}}
    \proofstepwidestar[1,2,3,4]{J_{A,\star}(a)\subseteq I}{%
      \FormulaRefAuto{A\subseteq B\coloneq
        \forall x\,(x\in A\rightarrow x\in B)}[def]{\rUI{\rRI{5,15}}}}
  \closeproofpartwide
\end{tabproofsplitwide}

\section{Teilbarkeit in Halbgruppen}

Ohne neutrales Element wäre die Bedingung \(b=a\star c\) nicht notwendig
reflexiv. Deshalb nehmen die folgenden Definitionen den Fall \(a=b\)
ausdrücklich auf. Die zuvor definierten Hauptideale erfassen genau diese
fehlenden Faktoren; eine unbestimmte Schreibweise wie \(A^1\) wird daher
nicht benötigt.

\FormulaDefDeltaK[Linksteilbarkeit in einer Halbgruppe]{%
  \SemigroupLeftDivides{A,\star}{a}{b}%
}{SemigroupLeftDividesDef}{%
  \DeltaRow{Mengen}{A\dsep a\dsep b\dsep c}
  \DeltaRow{Funktionen}{\star}
  \DeltaPrem{Halbgruppe}{\SemigroupAlg{A}{\star}}
  \DeltaRow{Elemente}{a,b\in A}
  \DeltaRow{Linksteilbarkeit}{%
    b=a\lor\exists c\in A\,(b=a\star c)}
  \DeltaRow{\textbf{Neues Symbol}}{}
  \DeltaPrem{Linksteilbarkeit}{%
    \SemigroupLeftDivides{A,\star}{a}{b}}
}

\FormulaDefDeltaK[Rechtsteilbarkeit in einer Halbgruppe]{%
  \SemigroupRightDivides{A,\star}{a}{b}%
}{SemigroupRightDividesDef}{%
  \DeltaRow{Mengen}{A\dsep a\dsep b\dsep c}
  \DeltaRow{Funktionen}{\star}
  \DeltaPrem{Halbgruppe}{\SemigroupAlg{A}{\star}}
  \DeltaRow{Elemente}{a,b\in A}
  \DeltaRow{Rechtsteilbarkeit}{%
    b=a\lor\exists c\in A\,(b=c\star a)}
  \DeltaRow{\textbf{Neues Symbol}}{}
  \DeltaPrem{Rechtsteilbarkeit}{%
    \SemigroupRightDivides{A,\star}{a}{b}}
}

\FormulaDefDeltaK[Zweiseitige Teilbarkeit in einer Halbgruppe]{%
  \SemigroupDivides{A,\star}{a}{b}%
}{SemigroupTwoSidedDividesDef}{%
  \DeltaRow{Mengen}{A\dsep a\dsep b\dsep x\dsep y}
  \DeltaRow{Funktionen}{\star}
  \DeltaPrem{Halbgruppe}{\SemigroupAlg{A}{\star}}
  \DeltaRow{Elemente}{a,b\in A}
  \DeltaRow{Zweiseitige Teilbarkeit}{%
    \begin{aligned}[t]
      &b=a\lor\exists x\in A\,(b=x\star a)\\[-2pt]
      &{}\lor\exists y\in A\,(b=a\star y)\\[-2pt]
      &{}\lor\exists x,y\in A\,
        \bigl(b=(x\star a)\star y\bigr)
    \end{aligned}}
  \DeltaRow{\textbf{Neues Symbol}}{}
  \DeltaPrem{Zweiseitige Teilbarkeit}{%
    \SemigroupDivides{A,\star}{a}{b}}
}

\FormulaThmDeltaK[Linksteilbarkeit als Mitgliedschaft im Rechtshauptideal]{%
  \begin{aligned}[t]
    &\SemigroupAlg{A}{\star}\dsep a,b\in A\vdash{}\\[-2pt]
    &\SemigroupLeftDivides{A,\star}{a}{b}
      \leftrightarrow b\in R_{A,\star}(a)
  \end{aligned}
}{SemigroupLeftDividesPrincipalRightIdeal}{%
  \DeltaRow{Mengen}{A\dsep a\dsep b\dsep c}
  \DeltaRow{Funktionen}{\star}
}
\begin{tabproofwide}
  \proofstepwidestar[1]{\SemigroupAlg{A}{\star}}{\rA}
  \proofstepwidestar[2]{a\in A}{\rA}
  \proofstepwidestar[3]{b\in A}{\rA}
  \proofstepwidestar[4]{\SemigroupLeftDivides{A,\star}{a}{b}}{\rA}
  \proofstepwidestar[1,2,3,4]{b=a\lor\exists c\in A\,(b=a\star c)}{%
    \FormulaRefAuto{SemigroupLeftDividesDef}[def]{1,2,3,4}}
  \proofstepwidestar[1,2,3,4]{b\in R_{A,\star}(a)}{%
    \FormulaRefAuto{SemigroupPrincipalRightIdealMembership}[thm]{%
      1,2,\rAI{3,5}}}
  \proofstepwidestar[1,2,3]{%
    \SemigroupLeftDivides{A,\star}{a}{b}\rightarrow
    b\in R_{A,\star}(a)}{\rRI{4,6}}
  \proofstepwidestar[8]{b\in R_{A,\star}(a)}{\rA}
  \proofstepwidestar[1,2,8]{b=a\lor\exists c\in A\,(b=a\star c)}{%
    \rAEb{\FormulaRefAuto{SemigroupPrincipalRightIdealMembership}[thm]{1,2,8}}}
  \proofstepwidestar[1,2,3,8]{\SemigroupLeftDivides{A,\star}{a}{b}}{%
    \FormulaRefAuto{SemigroupLeftDividesDef}[def]{1,2,3,9}}
  \proofstepwidestar[1,2,3]{%
    b\in R_{A,\star}(a)\rightarrow
    \SemigroupLeftDivides{A,\star}{a}{b}}{\rRI{8,10}}
  \proofstepwidestar[1,2,3]{%
    \SemigroupLeftDivides{A,\star}{a}{b}
    \leftrightarrow b\in R_{A,\star}(a)}{\rLRI{7,11}}
\end{tabproofwide}

\FormulaThmDeltaK[Rechtsteilbarkeit als Mitgliedschaft im Linkshauptideal]{%
  \begin{aligned}[t]
    &\SemigroupAlg{A}{\star}\dsep a,b\in A\vdash{}\\[-2pt]
    &\SemigroupRightDivides{A,\star}{a}{b}
      \leftrightarrow b\in L_{A,\star}(a)
  \end{aligned}
}{SemigroupRightDividesPrincipalLeftIdeal}{%
  \DeltaRow{Mengen}{A\dsep a\dsep b\dsep c}
  \DeltaRow{Funktionen}{\star}
}
\begin{tabproofwide}
  \proofstepwidestar[1]{\SemigroupAlg{A}{\star}}{\rA}
  \proofstepwidestar[2]{a\in A}{\rA}
  \proofstepwidestar[3]{b\in A}{\rA}
  \proofstepwidestar[4]{\SemigroupRightDivides{A,\star}{a}{b}}{\rA}
  \proofstepwidestar[1,2,3,4]{b=a\lor\exists c\in A\,(b=c\star a)}{%
    \FormulaRefAuto{SemigroupRightDividesDef}[def]{1,2,3,4}}
  \proofstepwidestar[1,2,3,4]{b\in L_{A,\star}(a)}{%
    \FormulaRefAuto{SemigroupPrincipalLeftIdealMembership}[thm]{%
      1,2,\rAI{3,5}}}
  \proofstepwidestar[1,2,3]{%
    \SemigroupRightDivides{A,\star}{a}{b}\rightarrow
    b\in L_{A,\star}(a)}{\rRI{4,6}}
  \proofstepwidestar[8]{b\in L_{A,\star}(a)}{\rA}
  \proofstepwidestar[1,2,8]{b=a\lor\exists c\in A\,(b=c\star a)}{%
    \rAEb{\FormulaRefAuto{SemigroupPrincipalLeftIdealMembership}[thm]{1,2,8}}}
  \proofstepwidestar[1,2,3,8]{\SemigroupRightDivides{A,\star}{a}{b}}{%
    \FormulaRefAuto{SemigroupRightDividesDef}[def]{1,2,3,9}}
  \proofstepwidestar[1,2,3]{%
    b\in L_{A,\star}(a)\rightarrow
    \SemigroupRightDivides{A,\star}{a}{b}}{\rRI{8,10}}
  \proofstepwidestar[1,2,3]{%
    \SemigroupRightDivides{A,\star}{a}{b}
    \leftrightarrow b\in L_{A,\star}(a)}{\rLRI{7,11}}
\end{tabproofwide}

\FormulaThmDeltaK[Zweiseitige Teilbarkeit als Hauptidealmitgliedschaft]{%
  \begin{aligned}[t]
    &\SemigroupAlg{A}{\star}\dsep a,b\in A\vdash{}\\[-2pt]
    &\SemigroupDivides{A,\star}{a}{b}
      \leftrightarrow b\in J_{A,\star}(a)
  \end{aligned}
}{SemigroupTwoSidedDividesPrincipalIdeal}{%
  \DeltaRow{Mengen}{A\dsep a\dsep b\dsep x\dsep y}
  \DeltaRow{Funktionen}{\star}
}
\begin{tabproofwide}
  \proofstepwidestar[1]{\SemigroupAlg{A}{\star}}{\rA}
  \proofstepwidestar[2]{a\in A}{\rA}
  \proofstepwidestar[3]{b\in A}{\rA}
  \proofstepwidestar[4]{\SemigroupDivides{A,\star}{a}{b}}{\rA}
  \proofstepwidestarbreak[1,2,3,4]{%
    \begin{aligned}[t]
      b=a\lor{}&\exists x\in A\,(b=x\star a)\\[-2pt]
      &{}\lor\exists y\in A\,(b=a\star y)\\[-2pt]
      &{}\lor\exists x,y\in A\,(b=(x\star a)\star y)
    \end{aligned}}{%
    \FormulaRefAuto{SemigroupTwoSidedDividesDef}[def]{1,2,3,4}}
  \proofstepwidestar[1,2,3,4]{b\in J_{A,\star}(a)}{%
    \FormulaRefAuto{SemigroupPrincipalTwoSidedIdealMembership}[thm]{%
      1,2,\rAI{3,5}}}
  \proofstepwidestar[1,2,3]{%
    \SemigroupDivides{A,\star}{a}{b}\rightarrow
    b\in J_{A,\star}(a)}{\rRI{4,6}}
  \proofstepwidestar[8]{b\in J_{A,\star}(a)}{\rA}
  \proofstepwidestarbreak[1,2,8]{%
    \begin{aligned}[t]
      b=a\lor{}&\exists x\in A\,(b=x\star a)\\[-2pt]
      &{}\lor\exists y\in A\,(b=a\star y)\\[-2pt]
      &{}\lor\exists x,y\in A\,(b=(x\star a)\star y)
    \end{aligned}}{%
    \rAEb{\FormulaRefAuto{SemigroupPrincipalTwoSidedIdealMembership}[thm]{1,2,8}}}
  \proofstepwidestar[1,2,3,8]{\SemigroupDivides{A,\star}{a}{b}}{%
    \FormulaRefAuto{SemigroupTwoSidedDividesDef}[def]{1,2,3,9}}
  \proofstepwidestar[1,2,3]{%
    b\in J_{A,\star}(a)\rightarrow
    \SemigroupDivides{A,\star}{a}{b}}{\rRI{8,10}}
  \proofstepwidestar[1,2,3]{%
    \SemigroupDivides{A,\star}{a}{b}
    \leftrightarrow b\in J_{A,\star}(a)}{\rLRI{7,11}}
\end{tabproofwide}

\FormulaThmDeltaK[Reflexivität der drei Teilbarkeiten]{%
  \begin{aligned}[t]
    &\SemigroupAlg{A}{\star}\dsep a\in A\vdash{}\\[-2pt]
    &\SemigroupLeftDivides{A,\star}{a}{a}\land
      \SemigroupRightDivides{A,\star}{a}{a}\land
      \SemigroupDivides{A,\star}{a}{a}
  \end{aligned}
}{SemigroupDivisibilityReflexive}{%
  \DeltaRow{Mengen}{A\dsep a}
  \DeltaRow{Funktionen}{\star}
}
\begin{tabproofwide}
  \proofstepwidestar[1]{\SemigroupAlg{A}{\star}}{\rA}
  \proofstepwidestar[2]{a\in A}{\rA}
  \proofstepwidestar[1,2]{\SemigroupLeftDivides{A,\star}{a}{a}}{%
    \FormulaRefAuto{SemigroupLeftDividesDef}[def]{1,2,2,\rOIa{\rII}}}
  \proofstepwidestar[1,2]{\SemigroupRightDivides{A,\star}{a}{a}}{%
    \FormulaRefAuto{SemigroupRightDividesDef}[def]{1,2,2,\rOIa{\rII}}}
  \proofstepwidestar[1,2]{\SemigroupDivides{A,\star}{a}{a}}{%
    \FormulaRefAuto{SemigroupTwoSidedDividesDef}[def]{1,2,2,\rOIa{\rII}}}
  \proofstepwidestar[1,2]{%
    \SemigroupLeftDivides{A,\star}{a}{a}\land
    \SemigroupRightDivides{A,\star}{a}{a}\land
    \SemigroupDivides{A,\star}{a}{a}}{\rAI{3,\rAI{4,5}}}
\end{tabproofwide}

\FormulaThmDeltaK[Transitivität der Linksteilbarkeit]{%
  \begin{aligned}[t]
    &\SemigroupAlg{A}{\star}\dsep a,b,c\in A\dsep
      \SemigroupLeftDivides{A,\star}{a}{b}\dsep
      \SemigroupLeftDivides{A,\star}{b}{c}\vdash{}\\[-2pt]
    &\SemigroupLeftDivides{A,\star}{a}{c}
  \end{aligned}
}{SemigroupLeftDividesTransitive}{%
  \DeltaRow{Mengen}{A\dsep a\dsep b\dsep c}
  \DeltaRow{Funktionen}{\star}
}
\begin{tabproofwide}
  \proofstepwidestar[1]{\SemigroupAlg{A}{\star}}{\rA}
  \proofstepwidestar[2]{a\in A}{\rA}
  \proofstepwidestar[3]{b\in A}{\rA}
  \proofstepwidestar[4]{c\in A}{\rA}
  \proofstepwidestar[5]{\SemigroupLeftDivides{A,\star}{a}{b}}{\rA}
  \proofstepwidestar[6]{\SemigroupLeftDivides{A,\star}{b}{c}}{\rA}
  \proofstepwidestar[1,2,3,5]{b\in R_{A,\star}(a)}{%
    \FormulaRefAuto{SemigroupLeftDividesPrincipalRightIdeal}[thm]{1,2,3,5}}
  \proofstepwidestar[1,3,4,6]{c\in R_{A,\star}(b)}{%
    \FormulaRefAuto{SemigroupLeftDividesPrincipalRightIdeal}[thm]{1,3,4,6}}
  \proofstepwidestar[1,2]{\mathsf{RId}(R_{A,\star}(a);A,\star)}{%
    \FormulaRefAuto{SemigroupPrincipalRightIdealIsRightIdeal}[thm]{1,2}}
  \proofstepwidestar[1,2,3,5]{R_{A,\star}(b)\subseteq R_{A,\star}(a)}{%
    \FormulaRefAuto{SemigroupPrincipalRightIdealMinimal}[thm]{1,3,9,7}}
  \proofstepwidestar[1,2,3,4,5,6]{c\in R_{A,\star}(a)}{%
    \FormulaRefAuto{A\subseteq B,\,x\in A\vdash x\in B}[thm]{10,8}}
  \proofstepwidestar[1,2,3,4,5,6]{\SemigroupLeftDivides{A,\star}{a}{c}}{%
    \FormulaRefAuto{SemigroupLeftDividesPrincipalRightIdeal}[thm]{1,2,4,11}}
\end{tabproofwide}

\FormulaThmDeltaK[Transitivität der Rechtsteilbarkeit]{%
  \begin{aligned}[t]
    &\SemigroupAlg{A}{\star}\dsep a,b,c\in A\dsep
      \SemigroupRightDivides{A,\star}{a}{b}\dsep
      \SemigroupRightDivides{A,\star}{b}{c}\vdash{}\\[-2pt]
    &\SemigroupRightDivides{A,\star}{a}{c}
  \end{aligned}
}{SemigroupRightDividesTransitive}{%
  \DeltaRow{Mengen}{A\dsep a\dsep b\dsep c}
  \DeltaRow{Funktionen}{\star}
}
\begin{tabproofwide}
  \proofstepwidestar[1]{\SemigroupAlg{A}{\star}}{\rA}
  \proofstepwidestar[2]{a\in A}{\rA}
  \proofstepwidestar[3]{b\in A}{\rA}
  \proofstepwidestar[4]{c\in A}{\rA}
  \proofstepwidestar[5]{\SemigroupRightDivides{A,\star}{a}{b}}{\rA}
  \proofstepwidestar[6]{\SemigroupRightDivides{A,\star}{b}{c}}{\rA}
  \proofstepwidestar[1,2,3,5]{b\in L_{A,\star}(a)}{%
    \FormulaRefAuto{SemigroupRightDividesPrincipalLeftIdeal}[thm]{1,2,3,5}}
  \proofstepwidestar[1,3,4,6]{c\in L_{A,\star}(b)}{%
    \FormulaRefAuto{SemigroupRightDividesPrincipalLeftIdeal}[thm]{1,3,4,6}}
  \proofstepwidestar[1,2]{\mathsf{LId}(L_{A,\star}(a);A,\star)}{%
    \FormulaRefAuto{SemigroupPrincipalLeftIdealIsLeftIdeal}[thm]{1,2}}
  \proofstepwidestar[1,2,3,5]{L_{A,\star}(b)\subseteq L_{A,\star}(a)}{%
    \FormulaRefAuto{SemigroupPrincipalLeftIdealMinimal}[thm]{1,3,9,7}}
  \proofstepwidestar[1,2,3,4,5,6]{c\in L_{A,\star}(a)}{%
    \FormulaRefAuto{A\subseteq B,\,x\in A\vdash x\in B}[thm]{10,8}}
  \proofstepwidestar[1,2,3,4,5,6]{\SemigroupRightDivides{A,\star}{a}{c}}{%
    \FormulaRefAuto{SemigroupRightDividesPrincipalLeftIdeal}[thm]{1,2,4,11}}
\end{tabproofwide}

\FormulaThmDeltaK[Transitivität der zweiseitigen Teilbarkeit]{%
  \begin{aligned}[t]
    &\SemigroupAlg{A}{\star}\dsep a,b,c\in A\dsep
      \SemigroupDivides{A,\star}{a}{b}\dsep
      \SemigroupDivides{A,\star}{b}{c}\vdash{}\\[-2pt]
    &\SemigroupDivides{A,\star}{a}{c}
  \end{aligned}
}{SemigroupTwoSidedDividesTransitive}{%
  \DeltaRow{Mengen}{A\dsep a\dsep b\dsep c}
  \DeltaRow{Funktionen}{\star}
}
\begin{tabproofwide}
  \proofstepwidestar[1]{\SemigroupAlg{A}{\star}}{\rA}
  \proofstepwidestar[2]{a\in A}{\rA}
  \proofstepwidestar[3]{b\in A}{\rA}
  \proofstepwidestar[4]{c\in A}{\rA}
  \proofstepwidestar[5]{\SemigroupDivides{A,\star}{a}{b}}{\rA}
  \proofstepwidestar[6]{\SemigroupDivides{A,\star}{b}{c}}{\rA}
  \proofstepwidestar[1,2,3,5]{b\in J_{A,\star}(a)}{%
    \FormulaRefAuto{SemigroupTwoSidedDividesPrincipalIdeal}[thm]{1,2,3,5}}
  \proofstepwidestar[1,3,4,6]{c\in J_{A,\star}(b)}{%
    \FormulaRefAuto{SemigroupTwoSidedDividesPrincipalIdeal}[thm]{1,3,4,6}}
  \proofstepwidestar[1,2]{\mathsf{Id}(J_{A,\star}(a);A,\star)}{%
    \FormulaRefAuto{SemigroupPrincipalTwoSidedIdealIsIdeal}[thm]{1,2}}
  \proofstepwidestar[1,2,3,5]{J_{A,\star}(b)\subseteq J_{A,\star}(a)}{%
    \FormulaRefAuto{SemigroupPrincipalTwoSidedIdealMinimal}[thm]{1,3,9,7}}
  \proofstepwidestar[1,2,3,4,5,6]{c\in J_{A,\star}(a)}{%
    \FormulaRefAuto{A\subseteq B,\,x\in A\vdash x\in B}[thm]{10,8}}
  \proofstepwidestar[1,2,3,4,5,6]{\SemigroupDivides{A,\star}{a}{c}}{%
    \FormulaRefAuto{SemigroupTwoSidedDividesPrincipalIdeal}[thm]{1,2,4,11}}
\end{tabproofwide}

Linksteilbarkeit, Rechtsteilbarkeit und zweiseitige Teilbarkeit sind
reflexiv und transitiv.  In einer kommutativen Halbgruppe stimmen die drei
Relationen überein.  Für Monoide ist die ausdrückliche Alternative \(b=a\)
überflüssig, weil das neutrale Element als fehlender Faktor dient.

\section{Green-Relationen}

Die Green-Relationen vergleichen Elemente danach, welche Hauptideale sie
erzeugen. Die Indizes (A,\star) bleiben sichtbar, damit die Relation stets
eindeutig zu ihrer umgebenden Halbgruppe gehört.

\FormulaDefDeltaK[Greensche \(\mathcal L\)-Relation]{%
  \SemigroupAlg{A}{\star}\dsep a,b\in A\vdash
  a\mathrel{\mathcal L_{A,\star}}b\coloneqq
  L_{A,\star}(a)=L_{A,\star}(b)%
}{SemigroupGreenLDef}{%
  \DeltaRow{Mengen}{A\dsep a\dsep b}
  \DeltaRow{Funktionen}{\star}
  \DeltaRow{Relationsoperatoren}{\mathcal L_{A,\star}}
}

\FormulaDefDeltaK[Greensche \(\mathcal R\)-Relation]{%
  \SemigroupAlg{A}{\star}\dsep a,b\in A\vdash
  a\mathrel{\mathcal R_{A,\star}}b\coloneqq
  R_{A,\star}(a)=R_{A,\star}(b)%
}{SemigroupGreenRDef}{%
  \DeltaRow{Mengen}{A\dsep a\dsep b}
  \DeltaRow{Funktionen}{\star}
  \DeltaRow{Relationsoperatoren}{\mathcal R_{A,\star}}
}

\FormulaDefDeltaK[Greensche \(\mathcal J\)-Relation]{%
  \SemigroupAlg{A}{\star}\dsep a,b\in A\vdash
  a\mathrel{\mathcal J_{A,\star}}b\coloneqq
  J_{A,\star}(a)=J_{A,\star}(b)%
}{SemigroupGreenJDef}{%
  \DeltaRow{Mengen}{A\dsep a\dsep b}
  \DeltaRow{Funktionen}{\star}
  \DeltaRow{Relationsoperatoren}{\mathcal J_{A,\star}}
}

\FormulaDefDeltaK[Greensche \(\mathcal H\)-Relation]{%
  \begin{aligned}[t]
    &\SemigroupAlg{A}{\star}\dsep a,b\in A\vdash{}\\[-2pt]
    &a\mathrel{\mathcal H_{A,\star}}b\coloneqq
      \bigl(a\mathrel{\mathcal L_{A,\star}}b\land
            a\mathrel{\mathcal R_{A,\star}}b\bigr)
  \end{aligned}
}{SemigroupGreenHDef}{%
  \DeltaRow{Mengen}{A\dsep a\dsep b}
  \DeltaRow{Funktionen}{\star}
  \DeltaRow{Relationsoperatoren}{%
    \mathcal L_{A,\star}\dsep\mathcal R_{A,\star}\dsep
    \mathcal H_{A,\star}}
}

\FormulaThmDeltaK[\(\mathcal L\) als gegenseitige Rechtsteilbarkeit]{%
  \begin{aligned}[t]
    &\SemigroupAlg{A}{\star}\dsep a,b\in A\vdash{}\\[-2pt]
    &a\mathrel{\mathcal L_{A,\star}}b\leftrightarrow
      \Bigl(\SemigroupRightDivides{A,\star}{a}{b}\land
            \SemigroupRightDivides{A,\star}{b}{a}\Bigr)
  \end{aligned}
}{SemigroupGreenLMutualRightDivisibility}{%
  \DeltaRow{Mengen}{A\dsep a\dsep b}
  \DeltaRow{Funktionen}{\star}
  \DeltaRow{Relationsoperatoren}{\mathcal L_{A,\star}}
}
\begin{tabproofwide}
  \proofstepwidestar[1]{\SemigroupAlg{A}{\star}}{\rA}
  \proofstepwidestar[2]{a\in A}{\rA}
  \proofstepwidestar[3]{b\in A}{\rA}
  \proofstepwidestar[4]{a\mathrel{\mathcal L_{A,\star}}b}{\rA}
  \proofstepwidestar[1,2,3,4]{L_{A,\star}(a)=L_{A,\star}(b)}{%
    \FormulaRefAuto{SemigroupGreenLDef}[def]{1,2,3,4}}
  \proofstepwidestar[1,3]{b\in L_{A,\star}(b)}{%
    \rAEa{\FormulaRefAuto{SemigroupPrincipalIdealsContainGenerator}[thm]{1,3}}}
  \proofstepwidestar[1,2,3,4]{b\in L_{A,\star}(a)}{%
    \rIE{\FormulaRefAuto{a=b\vdash b=a}[thm]{5},6}}
  \proofstepwidestar[1,2,3,4]{\SemigroupRightDivides{A,\star}{a}{b}}{%
    \FormulaRefAuto{SemigroupRightDividesPrincipalLeftIdeal}[thm]{1,2,3,7}}
  \proofstepwidestar[1,2]{a\in L_{A,\star}(a)}{%
    \rAEa{\FormulaRefAuto{SemigroupPrincipalIdealsContainGenerator}[thm]{1,2}}}
  \proofstepwidestar[1,2,3,4]{a\in L_{A,\star}(b)}{\rIE{5,9}}
  \proofstepwidestar[1,2,3,4]{\SemigroupRightDivides{A,\star}{b}{a}}{%
    \FormulaRefAuto{SemigroupRightDividesPrincipalLeftIdeal}[thm]{1,3,2,10}}
  \proofstepwidestar[1,2,3,4]{%
    \SemigroupRightDivides{A,\star}{a}{b}\land
    \SemigroupRightDivides{A,\star}{b}{a}}{\rAI{8,11}}
  \proofstepwidestar[1,2,3]{%
    a\mathrel{\mathcal L_{A,\star}}b\rightarrow
    \bigl(\SemigroupRightDivides{A,\star}{a}{b}\land
          \SemigroupRightDivides{A,\star}{b}{a}\bigr)}{\rRI{4,12}}
  \proofstepwidestar[14]{%
    \SemigroupRightDivides{A,\star}{a}{b}\land
    \SemigroupRightDivides{A,\star}{b}{a}}{\rA}
  \proofstepwidestar[14]{\SemigroupRightDivides{A,\star}{a}{b}}{\rAEa{14}}
  \proofstepwidestar[14]{\SemigroupRightDivides{A,\star}{b}{a}}{\rAEb{14}}
  \proofstepwidestar[1,2,3,14]{b\in L_{A,\star}(a)}{%
    \FormulaRefAuto{SemigroupRightDividesPrincipalLeftIdeal}[thm]{1,2,3,15}}
  \proofstepwidestar[1,2,3,14]{a\in L_{A,\star}(b)}{%
    \FormulaRefAuto{SemigroupRightDividesPrincipalLeftIdeal}[thm]{1,3,2,16}}
  \proofstepwidestar[1,2]{\mathsf{LId}(L_{A,\star}(a);A,\star)}{%
    \FormulaRefAuto{SemigroupPrincipalLeftIdealIsLeftIdeal}[thm]{1,2}}
  \proofstepwidestar[1,3]{\mathsf{LId}(L_{A,\star}(b);A,\star)}{%
    \FormulaRefAuto{SemigroupPrincipalLeftIdealIsLeftIdeal}[thm]{1,3}}
  \proofstepwidestar[1,2,3,14]{L_{A,\star}(a)\subseteq L_{A,\star}(b)}{%
    \FormulaRefAuto{SemigroupPrincipalLeftIdealMinimal}[thm]{1,2,20,18}}
  \proofstepwidestar[1,2,3,14]{L_{A,\star}(b)\subseteq L_{A,\star}(a)}{%
    \FormulaRefAuto{SemigroupPrincipalLeftIdealMinimal}[thm]{1,3,19,17}}
  \proofstepwidestar[1,2,3,14]{L_{A,\star}(a)=L_{A,\star}(b)}{%
    \FormulaRefAuto{A\subseteq B,\,B\subseteq A\vdash A=B}[thm]{21,22}}
  \proofstepwidestar[1,2,3,14]{a\mathrel{\mathcal L_{A,\star}}b}{%
    \FormulaRefAuto{SemigroupGreenLDef}[def]{1,2,3,23}}
  \proofstepwidestar[1,2,3]{%
    \bigl(\SemigroupRightDivides{A,\star}{a}{b}\land
          \SemigroupRightDivides{A,\star}{b}{a}\bigr)
    \rightarrow a\mathrel{\mathcal L_{A,\star}}b}{\rRI{14,24}}
  \proofstepwidestar[1,2,3]{%
    a\mathrel{\mathcal L_{A,\star}}b\leftrightarrow
    \bigl(\SemigroupRightDivides{A,\star}{a}{b}\land
          \SemigroupRightDivides{A,\star}{b}{a}\bigr)}{\rLRI{13,25}}
\end{tabproofwide}

\FormulaThmDeltaK[\(\mathcal R\) als gegenseitige Linksteilbarkeit]{%
  \begin{aligned}[t]
    &\SemigroupAlg{A}{\star}\dsep a,b\in A\vdash{}\\[-2pt]
    &a\mathrel{\mathcal R_{A,\star}}b\leftrightarrow
      \Bigl(\SemigroupLeftDivides{A,\star}{a}{b}\land
            \SemigroupLeftDivides{A,\star}{b}{a}\Bigr)
  \end{aligned}
}{SemigroupGreenRMutualLeftDivisibility}{%
  \DeltaRow{Mengen}{A\dsep a\dsep b}
  \DeltaRow{Funktionen}{\star}
  \DeltaRow{Relationsoperatoren}{\mathcal R_{A,\star}}
}
\begin{tabproofwide}
  \proofstepwidestar[1]{\SemigroupAlg{A}{\star}}{\rA}
  \proofstepwidestar[2]{a\in A}{\rA}
  \proofstepwidestar[3]{b\in A}{\rA}
  \proofstepwidestar[4]{a\mathrel{\mathcal R_{A,\star}}b}{\rA}
  \proofstepwidestar[1,2,3,4]{R_{A,\star}(a)=R_{A,\star}(b)}{%
    \FormulaRefAuto{SemigroupGreenRDef}[def]{1,2,3,4}}
  \proofstepwidestar[1,3]{b\in R_{A,\star}(b)}{%
    \rAEa{\rAEb{\FormulaRefAuto{SemigroupPrincipalIdealsContainGenerator}[thm]{1,3}}}}
  \proofstepwidestar[1,2,3,4]{b\in R_{A,\star}(a)}{%
    \rIE{\FormulaRefAuto{a=b\vdash b=a}[thm]{5},6}}
  \proofstepwidestar[1,2,3,4]{\SemigroupLeftDivides{A,\star}{a}{b}}{%
    \FormulaRefAuto{SemigroupLeftDividesPrincipalRightIdeal}[thm]{1,2,3,7}}
  \proofstepwidestar[1,2]{a\in R_{A,\star}(a)}{%
    \rAEa{\rAEb{\FormulaRefAuto{SemigroupPrincipalIdealsContainGenerator}[thm]{1,2}}}}
  \proofstepwidestar[1,2,3,4]{a\in R_{A,\star}(b)}{\rIE{5,9}}
  \proofstepwidestar[1,2,3,4]{\SemigroupLeftDivides{A,\star}{b}{a}}{%
    \FormulaRefAuto{SemigroupLeftDividesPrincipalRightIdeal}[thm]{1,3,2,10}}
  \proofstepwidestar[1,2,3,4]{%
    \SemigroupLeftDivides{A,\star}{a}{b}\land
    \SemigroupLeftDivides{A,\star}{b}{a}}{\rAI{8,11}}
  \proofstepwidestar[1,2,3]{%
    a\mathrel{\mathcal R_{A,\star}}b\rightarrow
    \bigl(\SemigroupLeftDivides{A,\star}{a}{b}\land
          \SemigroupLeftDivides{A,\star}{b}{a}\bigr)}{\rRI{4,12}}
  \proofstepwidestar[14]{%
    \SemigroupLeftDivides{A,\star}{a}{b}\land
    \SemigroupLeftDivides{A,\star}{b}{a}}{\rA}
  \proofstepwidestar[14]{\SemigroupLeftDivides{A,\star}{a}{b}}{\rAEa{14}}
  \proofstepwidestar[14]{\SemigroupLeftDivides{A,\star}{b}{a}}{\rAEb{14}}
  \proofstepwidestar[1,2,3,14]{b\in R_{A,\star}(a)}{%
    \FormulaRefAuto{SemigroupLeftDividesPrincipalRightIdeal}[thm]{1,2,3,15}}
  \proofstepwidestar[1,2,3,14]{a\in R_{A,\star}(b)}{%
    \FormulaRefAuto{SemigroupLeftDividesPrincipalRightIdeal}[thm]{1,3,2,16}}
  \proofstepwidestar[1,2]{\mathsf{RId}(R_{A,\star}(a);A,\star)}{%
    \FormulaRefAuto{SemigroupPrincipalRightIdealIsRightIdeal}[thm]{1,2}}
  \proofstepwidestar[1,3]{\mathsf{RId}(R_{A,\star}(b);A,\star)}{%
    \FormulaRefAuto{SemigroupPrincipalRightIdealIsRightIdeal}[thm]{1,3}}
  \proofstepwidestar[1,2,3,14]{R_{A,\star}(a)\subseteq R_{A,\star}(b)}{%
    \FormulaRefAuto{SemigroupPrincipalRightIdealMinimal}[thm]{1,2,20,18}}
  \proofstepwidestar[1,2,3,14]{R_{A,\star}(b)\subseteq R_{A,\star}(a)}{%
    \FormulaRefAuto{SemigroupPrincipalRightIdealMinimal}[thm]{1,3,19,17}}
  \proofstepwidestar[1,2,3,14]{R_{A,\star}(a)=R_{A,\star}(b)}{%
    \FormulaRefAuto{A\subseteq B,\,B\subseteq A\vdash A=B}[thm]{21,22}}
  \proofstepwidestar[1,2,3,14]{a\mathrel{\mathcal R_{A,\star}}b}{%
    \FormulaRefAuto{SemigroupGreenRDef}[def]{1,2,3,23}}
  \proofstepwidestar[1,2,3]{%
    \bigl(\SemigroupLeftDivides{A,\star}{a}{b}\land
          \SemigroupLeftDivides{A,\star}{b}{a}\bigr)
    \rightarrow a\mathrel{\mathcal R_{A,\star}}b}{\rRI{14,24}}
  \proofstepwidestar[1,2,3]{%
    a\mathrel{\mathcal R_{A,\star}}b\leftrightarrow
    \bigl(\SemigroupLeftDivides{A,\star}{a}{b}\land
          \SemigroupLeftDivides{A,\star}{b}{a}\bigr)}{\rLRI{13,25}}
\end{tabproofwide}

\FormulaThmDeltaK[\(\mathcal J\) als gegenseitige zweiseitige Teilbarkeit]{%
  \begin{aligned}[t]
    &\SemigroupAlg{A}{\star}\dsep a,b\in A\vdash{}\\[-2pt]
    &a\mathrel{\mathcal J_{A,\star}}b\leftrightarrow
      \Bigl(\SemigroupDivides{A,\star}{a}{b}\land
            \SemigroupDivides{A,\star}{b}{a}\Bigr)
  \end{aligned}
}{SemigroupGreenJMutualTwoSidedDivisibility}{%
  \DeltaRow{Mengen}{A\dsep a\dsep b}
  \DeltaRow{Funktionen}{\star}
  \DeltaRow{Relationsoperatoren}{\mathcal J_{A,\star}}
}
\begin{tabproofwide}
  \proofstepwidestar[1]{\SemigroupAlg{A}{\star}}{\rA}
  \proofstepwidestar[2]{a\in A}{\rA}
  \proofstepwidestar[3]{b\in A}{\rA}
  \proofstepwidestar[4]{a\mathrel{\mathcal J_{A,\star}}b}{\rA}
  \proofstepwidestar[1,2,3,4]{J_{A,\star}(a)=J_{A,\star}(b)}{%
    \FormulaRefAuto{SemigroupGreenJDef}[def]{1,2,3,4}}
  \proofstepwidestar[1,3]{b\in J_{A,\star}(b)}{%
    \rAEb{\rAEb{\FormulaRefAuto{SemigroupPrincipalIdealsContainGenerator}[thm]{1,3}}}}
  \proofstepwidestar[1,2,3,4]{b\in J_{A,\star}(a)}{%
    \rIE{\FormulaRefAuto{a=b\vdash b=a}[thm]{5},6}}
  \proofstepwidestar[1,2,3,4]{\SemigroupDivides{A,\star}{a}{b}}{%
    \FormulaRefAuto{SemigroupTwoSidedDividesPrincipalIdeal}[thm]{1,2,3,7}}
  \proofstepwidestar[1,2]{a\in J_{A,\star}(a)}{%
    \rAEb{\rAEb{\FormulaRefAuto{SemigroupPrincipalIdealsContainGenerator}[thm]{1,2}}}}
  \proofstepwidestar[1,2,3,4]{a\in J_{A,\star}(b)}{\rIE{5,9}}
  \proofstepwidestar[1,2,3,4]{\SemigroupDivides{A,\star}{b}{a}}{%
    \FormulaRefAuto{SemigroupTwoSidedDividesPrincipalIdeal}[thm]{1,3,2,10}}
  \proofstepwidestar[1,2,3,4]{%
    \SemigroupDivides{A,\star}{a}{b}\land
    \SemigroupDivides{A,\star}{b}{a}}{\rAI{8,11}}
  \proofstepwidestar[1,2,3]{%
    a\mathrel{\mathcal J_{A,\star}}b\rightarrow
    \bigl(\SemigroupDivides{A,\star}{a}{b}\land
          \SemigroupDivides{A,\star}{b}{a}\bigr)}{\rRI{4,12}}
  \proofstepwidestar[14]{%
    \SemigroupDivides{A,\star}{a}{b}\land
    \SemigroupDivides{A,\star}{b}{a}}{\rA}
  \proofstepwidestar[14]{\SemigroupDivides{A,\star}{a}{b}}{\rAEa{14}}
  \proofstepwidestar[14]{\SemigroupDivides{A,\star}{b}{a}}{\rAEb{14}}
  \proofstepwidestar[1,2,3,14]{b\in J_{A,\star}(a)}{%
    \FormulaRefAuto{SemigroupTwoSidedDividesPrincipalIdeal}[thm]{1,2,3,15}}
  \proofstepwidestar[1,2,3,14]{a\in J_{A,\star}(b)}{%
    \FormulaRefAuto{SemigroupTwoSidedDividesPrincipalIdeal}[thm]{1,3,2,16}}
  \proofstepwidestar[1,2]{\mathsf{Id}(J_{A,\star}(a);A,\star)}{%
    \FormulaRefAuto{SemigroupPrincipalTwoSidedIdealIsIdeal}[thm]{1,2}}
  \proofstepwidestar[1,3]{\mathsf{Id}(J_{A,\star}(b);A,\star)}{%
    \FormulaRefAuto{SemigroupPrincipalTwoSidedIdealIsIdeal}[thm]{1,3}}
  \proofstepwidestar[1,2,3,14]{J_{A,\star}(a)\subseteq J_{A,\star}(b)}{%
    \FormulaRefAuto{SemigroupPrincipalTwoSidedIdealMinimal}[thm]{1,2,20,18}}
  \proofstepwidestar[1,2,3,14]{J_{A,\star}(b)\subseteq J_{A,\star}(a)}{%
    \FormulaRefAuto{SemigroupPrincipalTwoSidedIdealMinimal}[thm]{1,3,19,17}}
  \proofstepwidestar[1,2,3,14]{J_{A,\star}(a)=J_{A,\star}(b)}{%
    \FormulaRefAuto{A\subseteq B,\,B\subseteq A\vdash A=B}[thm]{21,22}}
  \proofstepwidestar[1,2,3,14]{a\mathrel{\mathcal J_{A,\star}}b}{%
    \FormulaRefAuto{SemigroupGreenJDef}[def]{1,2,3,23}}
  \proofstepwidestar[1,2,3]{%
    \bigl(\SemigroupDivides{A,\star}{a}{b}\land
          \SemigroupDivides{A,\star}{b}{a}\bigr)
    \rightarrow a\mathrel{\mathcal J_{A,\star}}b}{\rRI{14,24}}
  \proofstepwidestar[1,2,3]{%
    a\mathrel{\mathcal J_{A,\star}}b\leftrightarrow
    \bigl(\SemigroupDivides{A,\star}{a}{b}\land
          \SemigroupDivides{A,\star}{b}{a}\bigr)}{\rLRI{13,25}}
\end{tabproofwide}

\FormulaThmDeltaK[\(\mathcal H\) als gegenseitige einseitige Teilbarkeit]{%
  \begin{aligned}[t]
    &\SemigroupAlg{A}{\star}\dsep a,b\in A\vdash{}\\[-2pt]
    &a\mathrel{\mathcal H_{A,\star}}b\leftrightarrow
      \Bigl(
        \SemigroupRightDivides{A,\star}{a}{b}\land
        \SemigroupRightDivides{A,\star}{b}{a}\land
        \SemigroupLeftDivides{A,\star}{a}{b}\land
        \SemigroupLeftDivides{A,\star}{b}{a}\Bigr)
  \end{aligned}
}{SemigroupGreenHMutualOneSidedDivisibility}{%
  \DeltaRow{Mengen}{A\dsep a\dsep b}
  \DeltaRow{Funktionen}{\star}
  \DeltaRow{Relationsoperatoren}{\mathcal H_{A,\star}}
}
Die Definition von \(\mathcal H\) zerf\"allt in eine \(\mathcal L\)- und
eine \(\mathcal R\)-Bedingung.  Die beiden Richtungen bestehen daher nur
darin, diese Bedingungen in die zugeh\"origen Teilbarkeitsaussagen zu
\"ubersetzen beziehungsweise wieder zusammenzusetzen.

\begin{tabproofsplitwide}
  \proofpartwideindR[Hinrichtung]{%
    \begin{aligned}[t]
      &\SemigroupAlg{A}{\star}\dsep a,b\in A\dsep
        a\mathrel{\mathcal H_{A,\star}}b\vdash{}\\[-2pt]
      &\SemigroupRightDivides{A,\star}{a}{b}\land
        \SemigroupRightDivides{A,\star}{b}{a}\land
        \SemigroupLeftDivides{A,\star}{a}{b}\land
        \SemigroupLeftDivides{A,\star}{b}{a}
    \end{aligned}
  }{%
    \begin{aligned}[t]
      &\SemigroupAlg{A}{\star}\dsep a,b\in A\dsep
        a\mathrel{\mathcal H_{A,\star}}b\vdash{}\\[-2pt]
      &\SemigroupRightDivides{A,\star}{a}{b}\land
        \SemigroupRightDivides{A,\star}{b}{a}\land
        \SemigroupLeftDivides{A,\star}{a}{b}\land
        \SemigroupLeftDivides{A,\star}{b}{a}
    \end{aligned}
  }[SemigroupGreenHMutualOneSidedDivisibilityForward]
    \proofstepwidestar[1]{\SemigroupAlg{A}{\star}}{\rA}
    \proofstepwidestar[2]{a\in A}{\rA}
    \proofstepwidestar[3]{b\in A}{\rA}
    \proofstepwidestar[4]{a\mathrel{\mathcal H_{A,\star}}b}{\rA}
    \proofstepwidestar[1,2,3,4]{%
      a\mathrel{\mathcal L_{A,\star}}b\land
      a\mathrel{\mathcal R_{A,\star}}b}{%
      \FormulaRefAuto{SemigroupGreenHDef}[def]{1,2,3,4}}
    \proofstepwidestar[1,2,3,4]{%
      \SemigroupRightDivides{A,\star}{a}{b}\land
      \SemigroupRightDivides{A,\star}{b}{a}}{%
      \FormulaRefAuto{SemigroupGreenLMutualRightDivisibility}[thm]{%
        1,2,3,\rAEa{5}}}
    \proofstepwidestar[1,2,3,4]{%
      \SemigroupLeftDivides{A,\star}{a}{b}\land
      \SemigroupLeftDivides{A,\star}{b}{a}}{%
      \FormulaRefAuto{SemigroupGreenRMutualLeftDivisibility}[thm]{%
        1,2,3,\rAEb{5}}}
    \proofstepwidestarbreak[1,2,3,4]{%
      \begin{aligned}[t]
        &\SemigroupRightDivides{A,\star}{a}{b}\land
          \SemigroupRightDivides{A,\star}{b}{a}\land{}\\[-2pt]
        &\SemigroupLeftDivides{A,\star}{a}{b}\land
          \SemigroupLeftDivides{A,\star}{b}{a}
      \end{aligned}}{%
      \rAI{\rAEa{6},\rAI{\rAEb{6},\rAI{\rAEa{7},\rAEb{7}}}}}
  \closeproofpartwideindR

  \proofpartwideindR[R\"uckrichtung]{%
    \begin{aligned}[t]
      &\SemigroupAlg{A}{\star}\dsep a,b\in A\dsep{}\\[-2pt]
      &\SemigroupRightDivides{A,\star}{a}{b}\land
        \SemigroupRightDivides{A,\star}{b}{a}\land
        \SemigroupLeftDivides{A,\star}{a}{b}\land
        \SemigroupLeftDivides{A,\star}{b}{a}
        \vdash a\mathrel{\mathcal H_{A,\star}}b
    \end{aligned}
  }{%
    \begin{aligned}[t]
      &\SemigroupAlg{A}{\star}\dsep a,b\in A\dsep{}\\[-2pt]
      &\SemigroupRightDivides{A,\star}{a}{b}\land
        \SemigroupRightDivides{A,\star}{b}{a}\land
        \SemigroupLeftDivides{A,\star}{a}{b}\land
        \SemigroupLeftDivides{A,\star}{b}{a}
        \vdash a\mathrel{\mathcal H_{A,\star}}b
    \end{aligned}
  }[SemigroupGreenHMutualOneSidedDivisibilityBackward]
    \proofstepwidestar[1]{\SemigroupAlg{A}{\star}}{\rA}
    \proofstepwidestar[2]{a\in A}{\rA}
    \proofstepwidestar[3]{b\in A}{\rA}
    \proofstepwidestar[4]{%
      \begin{aligned}[t]
        &\SemigroupRightDivides{A,\star}{a}{b}\land
          \SemigroupRightDivides{A,\star}{b}{a}\land{}\\[-2pt]
        &\SemigroupLeftDivides{A,\star}{a}{b}\land
          \SemigroupLeftDivides{A,\star}{b}{a}
      \end{aligned}}{\rA}
    \proofstepwidestar[4]{%
      \SemigroupRightDivides{A,\star}{a}{b}\land
      \SemigroupRightDivides{A,\star}{b}{a}}{%
      \rAI{\rAEa{4},\rAEa{\rAEb{4}}}}
    \proofstepwidestar[4]{%
      \SemigroupLeftDivides{A,\star}{a}{b}\land
      \SemigroupLeftDivides{A,\star}{b}{a}}{%
      \rAI{\rAEa{\rAEb{\rAEb{4}}},\rAEb{\rAEb{\rAEb{4}}}}}
    \proofstepwidestar[1,2,3,4]{a\mathrel{\mathcal L_{A,\star}}b}{%
      \FormulaRefAuto{SemigroupGreenLMutualRightDivisibility}[thm]{1,2,3,5}}
    \proofstepwidestar[1,2,3,4]{a\mathrel{\mathcal R_{A,\star}}b}{%
      \FormulaRefAuto{SemigroupGreenRMutualLeftDivisibility}[thm]{1,2,3,6}}
    \proofstepwidestar[1,2,3,4]{a\mathrel{\mathcal H_{A,\star}}b}{%
      \FormulaRefAuto{SemigroupGreenHDef}[def]{1,2,3,\rAI{7,8}}}
  \closeproofpartwideindR

  \proofpartwide{Schluss}
    \proofstepwidestar[1]{\SemigroupAlg{A}{\star}}{\rA}
    \proofstepwidestar[2]{a\in A}{\rA}
    \proofstepwidestar[3]{b\in A}{\rA}
    \proofstepwidestar[4]{a\mathrel{\mathcal H_{A,\star}}b}{\rA}
    \proofstepwidestarbreak[1,2,3,4]{%
      \begin{aligned}[t]
        &\SemigroupRightDivides{A,\star}{a}{b}\land
          \SemigroupRightDivides{A,\star}{b}{a}\land{}\\[-2pt]
        &\SemigroupLeftDivides{A,\star}{a}{b}\land
          \SemigroupLeftDivides{A,\star}{b}{a}
      \end{aligned}}{%
      \FormulaRefAuto{SemigroupGreenHMutualOneSidedDivisibilityForward}[thm]{%
        1,2,3,4}}
    \proofstepwidestar[1,2,3]{%
      \begin{aligned}[t]
        a\mathrel{\mathcal H_{A,\star}}b\rightarrow\bigl(&
          \SemigroupRightDivides{A,\star}{a}{b}\land
          \SemigroupRightDivides{A,\star}{b}{a}\land{}\\[-2pt]
        &\SemigroupLeftDivides{A,\star}{a}{b}\land
          \SemigroupLeftDivides{A,\star}{b}{a}\bigr)
      \end{aligned}}{\rRI{4,5}}
    \proofstepwidestar[7]{%
      \begin{aligned}[t]
        &\SemigroupRightDivides{A,\star}{a}{b}\land
          \SemigroupRightDivides{A,\star}{b}{a}\land{}\\[-2pt]
        &\SemigroupLeftDivides{A,\star}{a}{b}\land
          \SemigroupLeftDivides{A,\star}{b}{a}
      \end{aligned}}{\rA}
    \proofstepwidestar[1,2,3,7]{a\mathrel{\mathcal H_{A,\star}}b}{%
      \FormulaRefAuto{SemigroupGreenHMutualOneSidedDivisibilityBackward}[thm]{%
        1,2,3,7}}
    \proofstepwidestar[1,2,3]{%
      \begin{aligned}[t]
        \bigl(&\SemigroupRightDivides{A,\star}{a}{b}\land
          \SemigroupRightDivides{A,\star}{b}{a}\land{}\\[-2pt]
        &\SemigroupLeftDivides{A,\star}{a}{b}\land
          \SemigroupLeftDivides{A,\star}{b}{a}\bigr)
          \rightarrow a\mathrel{\mathcal H_{A,\star}}b
      \end{aligned}}{\rRI{7,8}}
    \proofstepwidestarbreak[1,2,3]{%
      \begin{aligned}[t]
        a\mathrel{\mathcal H_{A,\star}}b\leftrightarrow\bigl(&
          \SemigroupRightDivides{A,\star}{a}{b}\land
          \SemigroupRightDivides{A,\star}{b}{a}\land{}\\[-2pt]
        &\SemigroupLeftDivides{A,\star}{a}{b}\land
          \SemigroupLeftDivides{A,\star}{b}{a}\bigr)
      \end{aligned}}{\rLRI{6,9}}
  \closeproofpartwide
\end{tabproofsplitwide}

\subsection{Äquivalenzeigenschaften}

\FormulaThmDeltaK[Reflexivität von \(\mathcal L\)]{%
  \SemigroupAlg{A}{\star}\dsep a\in A\vdash
  a\mathrel{\mathcal L_{A,\star}}a%
}{SemigroupGreenLReflexive}{%
  \DeltaRow{Mengen}{A\dsep a}
  \DeltaRow{Funktionen}{\star}
  \DeltaRow{Relationsoperatoren}{\mathcal L_{A,\star}}
}
\begin{tabproofwide}
  \proofstepwidestar[1]{\SemigroupAlg{A}{\star}}{\rA}
  \proofstepwidestar[2]{a\in A}{\rA}
  \proofstepwidestar[]{L_{A,\star}(a)=L_{A,\star}(a)}{\rII}
  \proofstepwidestar[1,2]{a\mathrel{\mathcal L_{A,\star}}a}{%
    \FormulaRefAuto{SemigroupGreenLDef}[def]{1,2,2,3}}
\end{tabproofwide}

\FormulaThmDeltaK[Symmetrie von \(\mathcal L\)]{%
  \begin{aligned}[t]
    &\SemigroupAlg{A}{\star}\dsep a,b\in A\dsep
      a\mathrel{\mathcal L_{A,\star}}b\vdash{}\\[-2pt]
    &b\mathrel{\mathcal L_{A,\star}}a
  \end{aligned}
}{SemigroupGreenLSymmetric}{%
  \DeltaRow{Mengen}{A\dsep a\dsep b}
  \DeltaRow{Funktionen}{\star}
  \DeltaRow{Relationsoperatoren}{\mathcal L_{A,\star}}
}
\begin{tabproofwide}
  \proofstepwidestar[1]{\SemigroupAlg{A}{\star}}{\rA}
  \proofstepwidestar[2]{a\in A}{\rA}
  \proofstepwidestar[3]{b\in A}{\rA}
  \proofstepwidestar[4]{a\mathrel{\mathcal L_{A,\star}}b}{\rA}
  \proofstepwidestar[1,2,3,4]{L_{A,\star}(a)=L_{A,\star}(b)}{%
    \FormulaRefAuto{SemigroupGreenLDef}[def]{1,2,3,4}}
  \proofstepwidestar[1,2,3,4]{L_{A,\star}(b)=L_{A,\star}(a)}{%
    \FormulaRefAuto{a=b\vdash b=a}[thm]{5}}
  \proofstepwidestar[1,2,3,4]{b\mathrel{\mathcal L_{A,\star}}a}{%
    \FormulaRefAuto{SemigroupGreenLDef}[def]{1,3,2,6}}
\end{tabproofwide}

\FormulaThmDeltaK[Transitivität von \(\mathcal L\)]{%
  \begin{aligned}[t]
    &\SemigroupAlg{A}{\star}\dsep a,b,c\in A\dsep
      a\mathrel{\mathcal L_{A,\star}}b\dsep
      b\mathrel{\mathcal L_{A,\star}}c\vdash{}\\[-2pt]
    &a\mathrel{\mathcal L_{A,\star}}c
  \end{aligned}
}{SemigroupGreenLTransitive}{%
  \DeltaRow{Mengen}{A\dsep a\dsep b\dsep c}
  \DeltaRow{Funktionen}{\star}
  \DeltaRow{Relationsoperatoren}{\mathcal L_{A,\star}}
}
\begin{tabproofwide}
  \proofstepwidestar[1]{\SemigroupAlg{A}{\star}}{\rA}
  \proofstepwidestar[2]{a\in A}{\rA}
  \proofstepwidestar[3]{b\in A}{\rA}
  \proofstepwidestar[4]{c\in A}{\rA}
  \proofstepwidestar[5]{a\mathrel{\mathcal L_{A,\star}}b}{\rA}
  \proofstepwidestar[6]{b\mathrel{\mathcal L_{A,\star}}c}{\rA}
  \proofstepwidestar[1,2,3,5]{L_{A,\star}(a)=L_{A,\star}(b)}{%
    \FormulaRefAuto{SemigroupGreenLDef}[def]{1,2,3,5}}
  \proofstepwidestar[1,3,4,6]{L_{A,\star}(b)=L_{A,\star}(c)}{%
    \FormulaRefAuto{SemigroupGreenLDef}[def]{1,3,4,6}}
  \proofstepwidestar[1,2,3,4,5,6]{L_{A,\star}(a)=L_{A,\star}(c)}{%
    \FormulaRefAuto{a=b\dsep b=c\vdash a=c}[thm]{7,8}}
  \proofstepwidestar[1,2,3,4,5,6]{a\mathrel{\mathcal L_{A,\star}}c}{%
    \FormulaRefAuto{SemigroupGreenLDef}[def]{1,2,4,9}}
\end{tabproofwide}

\FormulaThmDeltaK[\(\mathcal L\) ist eine Äquivalenzrelation]{%
  \SemigroupAlg{A}{\star}\vdash\EqRel{A,\mathcal L_{A,\star}}%
}{SemigroupGreenLIsEquivalence}{%
  \DeltaRow{Mengen}{A\dsep a\dsep b\dsep c}
  \DeltaRow{Funktionen}{\star}
  \DeltaRow{Relationsoperatoren}{\mathcal L_{A,\star}}
}
\begin{tabproofwide}
  \proofstepwidestar[1]{\SemigroupAlg{A}{\star}}{\rA}
  \proofstepwidestar[1]{%
    \forall a\in A\,
      \bigl(a\mathrel{\mathcal L_{A,\star}}a\bigr)}{%
    \FormulaRefAuto{SemigroupGreenLReflexive}[thm]{1}}
  \proofstepwidestar[1]{%
    \forall a,b\in A\,
      \bigl(a\mathrel{\mathcal L_{A,\star}}b\rightarrow
        b\mathrel{\mathcal L_{A,\star}}a\bigr)}{%
    \FormulaRefAuto{SemigroupGreenLSymmetric}[thm]{1}}
  \proofstepwidestar[1]{%
    \forall a,b,c\in A\,
      \Bigl(\bigl(a\mathrel{\mathcal L_{A,\star}}b\land
        b\mathrel{\mathcal L_{A,\star}}c\bigr)
        \rightarrow a\mathrel{\mathcal L_{A,\star}}c\Bigr)}{%
    \FormulaRefAuto{SemigroupGreenLTransitive}[thm]{1}}
  \proofstepwidestar[1]{\EqRel{A,\mathcal L_{A,\star}}}{%
    \FormulaRefAuto{EqRelDef}[def]{2,3,4}}
\end{tabproofwide}

\FormulaThmDeltaK[Reflexivität von \(\mathcal R\)]{%
  \SemigroupAlg{A}{\star}\dsep a\in A\vdash
  a\mathrel{\mathcal R_{A,\star}}a%
}{SemigroupGreenRReflexive}{%
  \DeltaRow{Mengen}{A\dsep a}
  \DeltaRow{Funktionen}{\star}
  \DeltaRow{Relationsoperatoren}{\mathcal R_{A,\star}}
}
\begin{tabproofwide}
  \proofstepwidestar[1]{\SemigroupAlg{A}{\star}}{\rA}
  \proofstepwidestar[2]{a\in A}{\rA}
  \proofstepwidestar[]{R_{A,\star}(a)=R_{A,\star}(a)}{\rII}
  \proofstepwidestar[1,2]{a\mathrel{\mathcal R_{A,\star}}a}{%
    \FormulaRefAuto{SemigroupGreenRDef}[def]{1,2,2,3}}
\end{tabproofwide}

\FormulaThmDeltaK[Symmetrie von \(\mathcal R\)]{%
  \SemigroupAlg{A}{\star}\dsep a,b\in A\dsep
  a\mathrel{\mathcal R_{A,\star}}b\vdash
  b\mathrel{\mathcal R_{A,\star}}a%
}{SemigroupGreenRSymmetric}{%
  \DeltaRow{Mengen}{A\dsep a\dsep b}
  \DeltaRow{Funktionen}{\star}
  \DeltaRow{Relationsoperatoren}{\mathcal R_{A,\star}}
}
\begin{tabproofwide}
  \proofstepwidestar[1]{\SemigroupAlg{A}{\star}}{\rA}
  \proofstepwidestar[2]{a\in A}{\rA}
  \proofstepwidestar[3]{b\in A}{\rA}
  \proofstepwidestar[4]{a\mathrel{\mathcal R_{A,\star}}b}{\rA}
  \proofstepwidestar[1,2,3,4]{R_{A,\star}(a)=R_{A,\star}(b)}{%
    \FormulaRefAuto{SemigroupGreenRDef}[def]{1,2,3,4}}
  \proofstepwidestar[1,2,3,4]{R_{A,\star}(b)=R_{A,\star}(a)}{%
    \FormulaRefAuto{a=b\vdash b=a}[thm]{5}}
  \proofstepwidestar[1,2,3,4]{b\mathrel{\mathcal R_{A,\star}}a}{%
    \FormulaRefAuto{SemigroupGreenRDef}[def]{1,3,2,6}}
\end{tabproofwide}

\FormulaThmDeltaK[Transitivität von \(\mathcal R\)]{%
  \begin{aligned}[t]
    &\SemigroupAlg{A}{\star}\dsep a,b,c\in A\dsep
      a\mathrel{\mathcal R_{A,\star}}b\dsep
      b\mathrel{\mathcal R_{A,\star}}c\vdash{}\\[-2pt]
    &a\mathrel{\mathcal R_{A,\star}}c
  \end{aligned}
}{SemigroupGreenRTransitive}{%
  \DeltaRow{Mengen}{A\dsep a\dsep b\dsep c}
  \DeltaRow{Funktionen}{\star}
  \DeltaRow{Relationsoperatoren}{\mathcal R_{A,\star}}
}
\begin{tabproofwide}
  \proofstepwidestar[1]{\SemigroupAlg{A}{\star}}{\rA}
  \proofstepwidestar[2]{a\in A}{\rA}
  \proofstepwidestar[3]{b\in A}{\rA}
  \proofstepwidestar[4]{c\in A}{\rA}
  \proofstepwidestar[5]{a\mathrel{\mathcal R_{A,\star}}b}{\rA}
  \proofstepwidestar[6]{b\mathrel{\mathcal R_{A,\star}}c}{\rA}
  \proofstepwidestar[1,2,3,5]{R_{A,\star}(a)=R_{A,\star}(b)}{%
    \FormulaRefAuto{SemigroupGreenRDef}[def]{1,2,3,5}}
  \proofstepwidestar[1,3,4,6]{R_{A,\star}(b)=R_{A,\star}(c)}{%
    \FormulaRefAuto{SemigroupGreenRDef}[def]{1,3,4,6}}
  \proofstepwidestar[1,2,3,4,5,6]{R_{A,\star}(a)=R_{A,\star}(c)}{%
    \FormulaRefAuto{a=b\dsep b=c\vdash a=c}[thm]{7,8}}
  \proofstepwidestar[1,2,3,4,5,6]{a\mathrel{\mathcal R_{A,\star}}c}{%
    \FormulaRefAuto{SemigroupGreenRDef}[def]{1,2,4,9}}
\end{tabproofwide}

\FormulaThmDeltaK[\(\mathcal R\) ist eine Äquivalenzrelation]{%
  \SemigroupAlg{A}{\star}\vdash\EqRel{A,\mathcal R_{A,\star}}%
}{SemigroupGreenRIsEquivalence}{%
  \DeltaRow{Mengen}{A\dsep a\dsep b\dsep c}
  \DeltaRow{Funktionen}{\star}
  \DeltaRow{Relationsoperatoren}{\mathcal R_{A,\star}}
}
\begin{tabproofwide}
  \proofstepwidestar[1]{\SemigroupAlg{A}{\star}}{\rA}
  \proofstepwidestar[1]{%
    \forall a\in A\,
      \bigl(a\mathrel{\mathcal R_{A,\star}}a\bigr)}{%
    \FormulaRefAuto{SemigroupGreenRReflexive}[thm]{1}}
  \proofstepwidestar[1]{%
    \forall a,b\in A\,
      \bigl(a\mathrel{\mathcal R_{A,\star}}b\rightarrow
        b\mathrel{\mathcal R_{A,\star}}a\bigr)}{%
    \FormulaRefAuto{SemigroupGreenRSymmetric}[thm]{1}}
  \proofstepwidestar[1]{%
    \forall a,b,c\in A\,
      \Bigl(\bigl(a\mathrel{\mathcal R_{A,\star}}b\land
        b\mathrel{\mathcal R_{A,\star}}c\bigr)
        \rightarrow a\mathrel{\mathcal R_{A,\star}}c\Bigr)}{%
    \FormulaRefAuto{SemigroupGreenRTransitive}[thm]{1}}
  \proofstepwidestar[1]{\EqRel{A,\mathcal R_{A,\star}}}{%
    \FormulaRefAuto{EqRelDef}[def]{2,3,4}}
\end{tabproofwide}

\FormulaThmDeltaK[Reflexivität von \(\mathcal J\)]{%
  \SemigroupAlg{A}{\star}\dsep a\in A\vdash
  a\mathrel{\mathcal J_{A,\star}}a%
}{SemigroupGreenJReflexive}{%
  \DeltaRow{Mengen}{A\dsep a}
  \DeltaRow{Funktionen}{\star}
  \DeltaRow{Relationsoperatoren}{\mathcal J_{A,\star}}
}
\begin{tabproofwide}
  \proofstepwidestar[1]{\SemigroupAlg{A}{\star}}{\rA}
  \proofstepwidestar[2]{a\in A}{\rA}
  \proofstepwidestar[]{J_{A,\star}(a)=J_{A,\star}(a)}{\rII}
  \proofstepwidestar[1,2]{a\mathrel{\mathcal J_{A,\star}}a}{%
    \FormulaRefAuto{SemigroupGreenJDef}[def]{1,2,2,3}}
\end{tabproofwide}

\FormulaThmDeltaK[Symmetrie von \(\mathcal J\)]{%
  \SemigroupAlg{A}{\star}\dsep a,b\in A\dsep
  a\mathrel{\mathcal J_{A,\star}}b\vdash
  b\mathrel{\mathcal J_{A,\star}}a%
}{SemigroupGreenJSymmetric}{%
  \DeltaRow{Mengen}{A\dsep a\dsep b}
  \DeltaRow{Funktionen}{\star}
  \DeltaRow{Relationsoperatoren}{\mathcal J_{A,\star}}
}
\begin{tabproofwide}
  \proofstepwidestar[1]{\SemigroupAlg{A}{\star}}{\rA}
  \proofstepwidestar[2]{a\in A}{\rA}
  \proofstepwidestar[3]{b\in A}{\rA}
  \proofstepwidestar[4]{a\mathrel{\mathcal J_{A,\star}}b}{\rA}
  \proofstepwidestar[1,2,3,4]{J_{A,\star}(a)=J_{A,\star}(b)}{%
    \FormulaRefAuto{SemigroupGreenJDef}[def]{1,2,3,4}}
  \proofstepwidestar[1,2,3,4]{J_{A,\star}(b)=J_{A,\star}(a)}{%
    \FormulaRefAuto{a=b\vdash b=a}[thm]{5}}
  \proofstepwidestar[1,2,3,4]{b\mathrel{\mathcal J_{A,\star}}a}{%
    \FormulaRefAuto{SemigroupGreenJDef}[def]{1,3,2,6}}
\end{tabproofwide}

\FormulaThmDeltaK[Transitivität von \(\mathcal J\)]{%
  \begin{aligned}[t]
    &\SemigroupAlg{A}{\star}\dsep a,b,c\in A\dsep
      a\mathrel{\mathcal J_{A,\star}}b\dsep
      b\mathrel{\mathcal J_{A,\star}}c\vdash{}\\[-2pt]
    &a\mathrel{\mathcal J_{A,\star}}c
  \end{aligned}
}{SemigroupGreenJTransitive}{%
  \DeltaRow{Mengen}{A\dsep a\dsep b\dsep c}
  \DeltaRow{Funktionen}{\star}
  \DeltaRow{Relationsoperatoren}{\mathcal J_{A,\star}}
}
\begin{tabproofwide}
  \proofstepwidestar[1]{\SemigroupAlg{A}{\star}}{\rA}
  \proofstepwidestar[2]{a\in A}{\rA}
  \proofstepwidestar[3]{b\in A}{\rA}
  \proofstepwidestar[4]{c\in A}{\rA}
  \proofstepwidestar[5]{a\mathrel{\mathcal J_{A,\star}}b}{\rA}
  \proofstepwidestar[6]{b\mathrel{\mathcal J_{A,\star}}c}{\rA}
  \proofstepwidestar[1,2,3,5]{J_{A,\star}(a)=J_{A,\star}(b)}{%
    \FormulaRefAuto{SemigroupGreenJDef}[def]{1,2,3,5}}
  \proofstepwidestar[1,3,4,6]{J_{A,\star}(b)=J_{A,\star}(c)}{%
    \FormulaRefAuto{SemigroupGreenJDef}[def]{1,3,4,6}}
  \proofstepwidestar[1,2,3,4,5,6]{J_{A,\star}(a)=J_{A,\star}(c)}{%
    \FormulaRefAuto{a=b\dsep b=c\vdash a=c}[thm]{7,8}}
  \proofstepwidestar[1,2,3,4,5,6]{a\mathrel{\mathcal J_{A,\star}}c}{%
    \FormulaRefAuto{SemigroupGreenJDef}[def]{1,2,4,9}}
\end{tabproofwide}

\FormulaThmDeltaK[\(\mathcal J\) ist eine Äquivalenzrelation]{%
  \SemigroupAlg{A}{\star}\vdash\EqRel{A,\mathcal J_{A,\star}}%
}{SemigroupGreenJIsEquivalence}{%
  \DeltaRow{Mengen}{A\dsep a\dsep b\dsep c}
  \DeltaRow{Funktionen}{\star}
  \DeltaRow{Relationsoperatoren}{\mathcal J_{A,\star}}
}
\begin{tabproofwide}
  \proofstepwidestar[1]{\SemigroupAlg{A}{\star}}{\rA}
  \proofstepwidestar[1]{%
    \forall a\in A\,
      \bigl(a\mathrel{\mathcal J_{A,\star}}a\bigr)}{%
    \FormulaRefAuto{SemigroupGreenJReflexive}[thm]{1}}
  \proofstepwidestar[1]{%
    \forall a,b\in A\,
      \bigl(a\mathrel{\mathcal J_{A,\star}}b\rightarrow
        b\mathrel{\mathcal J_{A,\star}}a\bigr)}{%
    \FormulaRefAuto{SemigroupGreenJSymmetric}[thm]{1}}
  \proofstepwidestar[1]{%
    \forall a,b,c\in A\,
      \Bigl(\bigl(a\mathrel{\mathcal J_{A,\star}}b\land
        b\mathrel{\mathcal J_{A,\star}}c\bigr)
        \rightarrow a\mathrel{\mathcal J_{A,\star}}c\Bigr)}{%
    \FormulaRefAuto{SemigroupGreenJTransitive}[thm]{1}}
  \proofstepwidestar[1]{\EqRel{A,\mathcal J_{A,\star}}}{%
    \FormulaRefAuto{EqRelDef}[def]{2,3,4}}
\end{tabproofwide}

\FormulaThmDeltaK[Reflexivität von \(\mathcal H\)]{%
  \SemigroupAlg{A}{\star}\dsep a\in A\vdash
  a\mathrel{\mathcal H_{A,\star}}a%
}{SemigroupGreenHReflexive}{%
  \DeltaRow{Mengen}{A\dsep a}
  \DeltaRow{Funktionen}{\star}
  \DeltaRow{Relationsoperatoren}{\mathcal H_{A,\star}}
}
\begin{tabproofwide}
  \proofstepwidestar[1]{\SemigroupAlg{A}{\star}}{\rA}
  \proofstepwidestar[2]{a\in A}{\rA}
  \proofstepwidestar[1,2]{a\mathrel{\mathcal L_{A,\star}}a}{%
    \FormulaRefAuto{SemigroupGreenLReflexive}[thm]{1,2}}
  \proofstepwidestar[1,2]{a\mathrel{\mathcal R_{A,\star}}a}{%
    \FormulaRefAuto{SemigroupGreenRReflexive}[thm]{1,2}}
  \proofstepwidestar[1,2]{a\mathrel{\mathcal H_{A,\star}}a}{%
    \FormulaRefAuto{SemigroupGreenHDef}[def]{1,2,2,\rAI{3,4}}}
\end{tabproofwide}

\FormulaThmDeltaK[Symmetrie von \(\mathcal H\)]{%
  \SemigroupAlg{A}{\star}\dsep a,b\in A\dsep
  a\mathrel{\mathcal H_{A,\star}}b\vdash
  b\mathrel{\mathcal H_{A,\star}}a%
}{SemigroupGreenHSymmetric}{%
  \DeltaRow{Mengen}{A\dsep a\dsep b}
  \DeltaRow{Funktionen}{\star}
  \DeltaRow{Relationsoperatoren}{\mathcal H_{A,\star}}
}
\begin{tabproofwide}
  \proofstepwidestar[1]{\SemigroupAlg{A}{\star}}{\rA}
  \proofstepwidestar[2]{a\in A}{\rA}
  \proofstepwidestar[3]{b\in A}{\rA}
  \proofstepwidestar[4]{a\mathrel{\mathcal H_{A,\star}}b}{\rA}
  \proofstepwidestar[1,2,3,4]{%
    a\mathrel{\mathcal L_{A,\star}}b\land
    a\mathrel{\mathcal R_{A,\star}}b}{%
    \FormulaRefAuto{SemigroupGreenHDef}[def]{1,2,3,4}}
  \proofstepwidestar[1,2,3,4]{b\mathrel{\mathcal L_{A,\star}}a}{%
    \FormulaRefAuto{SemigroupGreenLSymmetric}[thm]{1,2,3,\rAEa{5}}}
  \proofstepwidestar[1,2,3,4]{b\mathrel{\mathcal R_{A,\star}}a}{%
    \FormulaRefAuto{SemigroupGreenRSymmetric}[thm]{1,2,3,\rAEb{5}}}
  \proofstepwidestar[1,2,3,4]{b\mathrel{\mathcal H_{A,\star}}a}{%
    \FormulaRefAuto{SemigroupGreenHDef}[def]{1,3,2,\rAI{6,7}}}
\end{tabproofwide}

\FormulaThmDeltaK[Transitivität von \(\mathcal H\)]{%
  \begin{aligned}[t]
    &\SemigroupAlg{A}{\star}\dsep a,b,c\in A\dsep
      a\mathrel{\mathcal H_{A,\star}}b\dsep
      b\mathrel{\mathcal H_{A,\star}}c\vdash{}\\[-2pt]
    &a\mathrel{\mathcal H_{A,\star}}c
  \end{aligned}
}{SemigroupGreenHTransitive}{%
  \DeltaRow{Mengen}{A\dsep a\dsep b\dsep c}
  \DeltaRow{Funktionen}{\star}
  \DeltaRow{Relationsoperatoren}{\mathcal H_{A,\star}}
}
\begin{tabproofwide}
  \proofstepwidestar[1]{\SemigroupAlg{A}{\star}}{\rA}
  \proofstepwidestar[2]{a\in A}{\rA}
  \proofstepwidestar[3]{b\in A}{\rA}
  \proofstepwidestar[4]{c\in A}{\rA}
  \proofstepwidestar[5]{a\mathrel{\mathcal H_{A,\star}}b}{\rA}
  \proofstepwidestar[6]{b\mathrel{\mathcal H_{A,\star}}c}{\rA}
  \proofstepwidestar[1,2,3,5]{%
    a\mathrel{\mathcal L_{A,\star}}b\land
    a\mathrel{\mathcal R_{A,\star}}b}{%
    \FormulaRefAuto{SemigroupGreenHDef}[def]{1,2,3,5}}
  \proofstepwidestar[1,3,4,6]{%
    b\mathrel{\mathcal L_{A,\star}}c\land
    b\mathrel{\mathcal R_{A,\star}}c}{%
    \FormulaRefAuto{SemigroupGreenHDef}[def]{1,3,4,6}}
  \proofstepwidestar[1,2,3,4,5,6]{a\mathrel{\mathcal L_{A,\star}}c}{%
    \FormulaRefAuto{SemigroupGreenLTransitive}[thm]{1,2,3,4,\rAEa{7},\rAEa{8}}}
  \proofstepwidestar[1,2,3,4,5,6]{a\mathrel{\mathcal R_{A,\star}}c}{%
    \FormulaRefAuto{SemigroupGreenRTransitive}[thm]{1,2,3,4,\rAEb{7},\rAEb{8}}}
  \proofstepwidestar[1,2,3,4,5,6]{a\mathrel{\mathcal H_{A,\star}}c}{%
    \FormulaRefAuto{SemigroupGreenHDef}[def]{1,2,4,\rAI{9,10}}}
\end{tabproofwide}

\FormulaThmDeltaK[\(\mathcal H\) ist eine Äquivalenzrelation]{%
  \SemigroupAlg{A}{\star}\vdash\EqRel{A,\mathcal H_{A,\star}}%
}{SemigroupGreenHIsEquivalence}{%
  \DeltaRow{Mengen}{A\dsep a\dsep b\dsep c}
  \DeltaRow{Funktionen}{\star}
  \DeltaRow{Relationsoperatoren}{\mathcal H_{A,\star}}
}
\begin{tabproofwide}
  \proofstepwidestar[1]{\SemigroupAlg{A}{\star}}{\rA}
  \proofstepwidestar[1]{%
    \forall a\in A\,
      \bigl(a\mathrel{\mathcal H_{A,\star}}a\bigr)}{%
    \FormulaRefAuto{SemigroupGreenHReflexive}[thm]{1}}
  \proofstepwidestar[1]{%
    \forall a,b\in A\,
      \bigl(a\mathrel{\mathcal H_{A,\star}}b\rightarrow
        b\mathrel{\mathcal H_{A,\star}}a\bigr)}{%
    \FormulaRefAuto{SemigroupGreenHSymmetric}[thm]{1}}
  \proofstepwidestar[1]{%
    \forall a,b,c\in A\,
      \Bigl(\bigl(a\mathrel{\mathcal H_{A,\star}}b\land
        b\mathrel{\mathcal H_{A,\star}}c\bigr)
        \rightarrow a\mathrel{\mathcal H_{A,\star}}c\Bigr)}{%
    \FormulaRefAuto{SemigroupGreenHTransitive}[thm]{1}}
  \proofstepwidestar[1]{\EqRel{A,\mathcal H_{A,\star}}}{%
    \FormulaRefAuto{EqRelDef}[def]{2,3,4}}
\end{tabproofwide}

\section{Morphismen von Halbgruppen}

\FormulaDefDeltaK[Begriff des Halbgruppenhomomorphismus]{%
  \SemigroupHom{f}{A,\star}{B,\diamond}%
}{SemigroupHomDef}{%
  \DeltaRow{Mengen}{A\dsep B\dsep x\dsep y}
  \DeltaRow{Funktionen}{f\dsep\star\dsep\diamond}
  \DeltaRow{\textbf{Axiome}}{}
  \DeltaPrem{Quellhalbgruppe}{\SemigroupAlg{A}{\star}}
  \DeltaPrem{Zielhalbgruppe}{\SemigroupAlg{B}{\diamond}}
  \DeltaPrem{Abbildung}{f\colon A\to B}
  \DeltaRow{Operationserhaltung}{%
    x,y\in A\vdash
    f(x\star y)=f(x)\diamond f(y)}
  \DeltaRow{\textbf{Neues Symbol}}{}
  \DeltaPrem{Halbgruppenhomomorphismus}{%
    \SemigroupHom{f}{A,\star}{B,\diamond}}
}

\FormulaAxiomDeltaK[Quellhalbgruppe]{%
  \SemigroupHom{f}{A,\star}{B,\diamond}
  \vdash\SemigroupAlg{A}{\star}%
}{SemigroupHomSourceAxiom}{%
  \DeltaRow{Mengen}{A\dsep B}
  \DeltaRow{Funktionen}{f\dsep\star\dsep\diamond}
}

\FormulaAxiomDeltaK[Zielhalbgruppe]{%
  \SemigroupHom{f}{A,\star}{B,\diamond}
  \vdash\SemigroupAlg{B}{\diamond}%
}{SemigroupHomTargetAxiom}{%
  \DeltaRow{Mengen}{A\dsep B}
  \DeltaRow{Funktionen}{f\dsep\star\dsep\diamond}
}

\FormulaAxiomDeltaK[Operationserhaltung eines Halbgruppenhomomorphismus]{%
  \SemigroupHom{f}{A,\star}{B,\diamond}\dsep x,y\in A
  \vdash f(x\star y)=f(x)\diamond f(y)%
}{SemigroupHomOperationAxiom}{%
  \DeltaRow{Mengen}{A\dsep B\dsep x\dsep y}
  \DeltaRow{Funktionen}{f\dsep\star\dsep\diamond}
}

\FormulaAxiomDeltaK[Trägerabbildung eines Halbgruppenhomomorphismus]{%
  \SemigroupHom{f}{A,\star}{B,\diamond}
  \vdash f\colon A\to B%
}{SemigroupHomFunctionAxiom}{%
  \DeltaRow{Mengen}{A\dsep B}
  \DeltaRow{Funktionen}{f\dsep\star\dsep\diamond}
}

\FormulaDefDeltaK[Begriff des Halbgruppenisomorphismus]{%
  \SemigroupIso{f}{A,\star}{B,\diamond}%
}{SemigroupIsoDef}{%
  \DeltaRow{Mengen}{A\dsep B}
  \DeltaRow{Funktionen}{f\dsep\star\dsep\diamond}
  \DeltaRow{\textbf{Axiome}}{}
  \DeltaPrem{Halbgruppenhomomorphismus}{%
    \SemigroupHom{f}{A,\star}{B,\diamond}}
  \DeltaPrem{Bijektivität}{f\colon A\bij B}
  \DeltaRow{\textbf{Neues Symbol}}{}
  \DeltaPrem{Halbgruppenisomorphismus}{%
    \SemigroupIso{f}{A,\star}{B,\diamond}}
}

\FormulaAxiomDeltaK[Homomorphismuszugriff]{%
  \SemigroupIso{f}{A,\star}{B,\diamond}
  \vdash\SemigroupHom{f}{A,\star}{B,\diamond}%
}{SemigroupIsoHomAxiom}{%
  \DeltaRow{Mengen}{A\dsep B}
  \DeltaRow{Funktionen}{f\dsep\star\dsep\diamond}
}

\FormulaAxiomDeltaK[Bijektivitätszugriff]{%
  \SemigroupIso{f}{A,\star}{B,\diamond}
  \vdash f\colon A\bij B%
}{SemigroupIsoBijectiveAxiom}{%
  \DeltaRow{Mengen}{A\dsep B}
  \DeltaRow{Funktionen}{f\dsep\star\dsep\diamond}
}

\subsection{Grundlegende Eigenschaften}

\FormulaThmDeltaK[Komposition von Halbgruppenhomomorphismen]{%
  \begin{aligned}[t]
    &\SemigroupHom{f}{A,\star}{B,\diamond}\dsep
      \SemigroupHom{g}{B,\diamond}{C,\bullet}\\
    &\vdash\SemigroupHom{g\circ f}{A,\star}{C,\bullet}
  \end{aligned}
}{SemigroupHomComposition}{%
  \DeltaRow{Mengen}{A\dsep B\dsep C\dsep x\dsep y}
  \DeltaRow{Funktionen}{f\dsep g\dsep\star\dsep\diamond\dsep\bullet}
}
\begin{tabproofsplitwide}
  \proofpartwideindR[Komposition der Trägerabbildungen]{%
    \begin{aligned}[t]
      &\SemigroupHom{f}{A,\star}{B,\diamond}\dsep
        \SemigroupHom{g}{B,\diamond}{C,\bullet}\vdash{}\\[-2pt]
      &g\circ f\colon A\to C
    \end{aligned}
  }{%
    \SemigroupHom{f}{A,\star}{B,\diamond}\dsep
    \SemigroupHom{g}{B,\diamond}{C,\bullet}
    \vdash g\circ f\colon A\to C%
  }[SemigroupHomCompositionFunction]
    \proofstepwidestar[1]{\SemigroupHom{f}{A,\star}{B,\diamond}}{\rA}
    \proofstepwidestar[2]{\SemigroupHom{g}{B,\diamond}{C,\bullet}}{\rA}
    \proofstepwidestar[1]{f\colon A\to B}{%
      \FormulaRefAuto{SemigroupHomFunctionAxiom}[ax]{1}}
    \proofstepwidestar[2]{g\colon B\to C}{%
      \FormulaRefAuto{SemigroupHomFunctionAxiom}[ax]{2}}
    \proofstepwidestar[1,2]{g\circ f\colon A\to C}{%
      \FormulaRefAuto{G\circ F\colon A\to C}[thm]{3,4}}
  \closeproofpartwideindR

  \proofpartwideindR[Einsetzen der ersten Homomorphie]{%
    \begin{aligned}[t]
      &\SemigroupHom{f}{A,\star}{B,\diamond}\dsep
        \SemigroupHom{g}{B,\diamond}{C,\bullet}\dsep x,y\in A\vdash{}\\[-2pt]
      &(g\circ f)(x\star y)=g\bigl(f(x)\diamond f(y)\bigr)
    \end{aligned}
  }{%
    \SemigroupHom{f}{A,\star}{B,\diamond}\dsep
    \SemigroupHom{g}{B,\diamond}{C,\bullet}\dsep x,y\in A
    \vdash (g\circ f)(x\star y)=g\bigl(f(x)\diamond f(y)\bigr)%
  }[SemigroupHomCompositionFirstOperation]
    \proofstepwidestar[1]{\SemigroupHom{f}{A,\star}{B,\diamond}}{\rA}
    \proofstepwidestar[2]{\SemigroupHom{g}{B,\diamond}{C,\bullet}}{\rA}
    \proofstepwidestar[3]{x\in A}{\rA}
    \proofstepwidestar[4]{y\in A}{\rA}
    \proofstepwidestar[1]{f\colon A\to B}{%
      \FormulaRefAuto{SemigroupHomFunctionAxiom}[ax]{1}}
    \proofstepwidestar[2]{g\colon B\to C}{%
      \FormulaRefAuto{SemigroupHomFunctionAxiom}[ax]{2}}
    \proofstepwidestar[1]{\SemigroupAlg{A}{\star}}{%
      \FormulaRefAuto{SemigroupHomSourceAxiom}[ax]{1}}
    \proofstepwidestar[1,3,4]{x\star y\in A}{%
      \FormulaRefAuto{SemigroupClosureAxiom}[ax]{7,3,4}}
    \proofstepwidestar[1,3,4]{f(x\star y)=f(x)\diamond f(y)}{%
      \FormulaRefAuto{SemigroupHomOperationAxiom}[ax]{1,3,4}}
    \proofstepwide[1,2,3,4]{(g\circ f)(x\star y)}{=}{g(f(x\star y))}{%
      \FormulaRefAuto{%
        x\in A \vdash (F\circ G)(x)=F(G(x))%
      }[thm]{5,6,8}}
    \proofstepwide[1,2,3,4]{}{=}{g\bigl(f(x)\diamond f(y)\bigr)}{%
      \rIE{9}}
    \proofstepwidestar[1,2,3,4]{%
      (g\circ f)(x\star y)=g\bigl(f(x)\diamond f(y)\bigr)}{%
      \rChain{10,11}}
  \closeproofpartwideindR

  \proofpartwideindR[Einsetzen der zweiten Homomorphie]{%
    \begin{aligned}[t]
      &\SemigroupHom{f}{A,\star}{B,\diamond}\dsep
        \SemigroupHom{g}{B,\diamond}{C,\bullet}\dsep x,y\in A\vdash{}\\[-2pt]
      &g\bigl(f(x)\diamond f(y)\bigr)
        =(g\circ f)(x)\bullet(g\circ f)(y)
    \end{aligned}
  }{%
    \SemigroupHom{f}{A,\star}{B,\diamond}\dsep
    \SemigroupHom{g}{B,\diamond}{C,\bullet}\dsep x,y\in A
    \vdash g\bigl(f(x)\diamond f(y)\bigr)
      =(g\circ f)(x)\bullet(g\circ f)(y)%
  }[SemigroupHomCompositionSecondOperation]
    \proofstepwidestar[1]{\SemigroupHom{f}{A,\star}{B,\diamond}}{\rA}
    \proofstepwidestar[2]{\SemigroupHom{g}{B,\diamond}{C,\bullet}}{\rA}
    \proofstepwidestar[3]{x\in A}{\rA}
    \proofstepwidestar[4]{y\in A}{\rA}
    \proofstepwidestar[1]{f\colon A\to B}{%
      \FormulaRefAuto{SemigroupHomFunctionAxiom}[ax]{1}}
    \proofstepwidestar[2]{g\colon B\to C}{%
      \FormulaRefAuto{SemigroupHomFunctionAxiom}[ax]{2}}
    \proofstepwidestar[1,3]{f(x)\in B}{%
      \FormulaRefAuto{%
        F\colon A\to B\dsep x\in A\vdash F(x)\in B%
      }[thm]{5,3}}
    \proofstepwidestar[1,4]{f(y)\in B}{%
      \FormulaRefAuto{%
        F\colon A\to B\dsep x\in A\vdash F(x)\in B%
      }[thm]{5,4}}
    \proofstepwidestar[1,2,3]{g(f(x))=(g\circ f)(x)}{%
      \FormulaRefAuto{%
        x\in A \vdash F(G(x))=(F\circ G)(x)%
      }[thm]{5,6,3}}
    \proofstepwidestar[1,2,4]{g(f(y))=(g\circ f)(y)}{%
      \FormulaRefAuto{%
        x\in A \vdash F(G(x))=(F\circ G)(x)%
      }[thm]{5,6,4}}
    \proofstepwide[1,2,3,4]{g\bigl(f(x)\diamond f(y)\bigr)}{=}{%
      g(f(x))\bullet g(f(y))}{%
      \FormulaRefAuto{SemigroupHomOperationAxiom}[ax]{2,7,8}}
    \proofstepwide[1,2,3,4]{}{=}{(g\circ f)(x)\bullet g(f(y))}{\rIE{9}}
    \proofstepwide[1,2,3,4]{}{=}{%
      (g\circ f)(x)\bullet(g\circ f)(y)}{\rIE{10}}
    \proofstepwidestar[1,2,3,4]{%
      g\bigl(f(x)\diamond f(y)\bigr)
      =(g\circ f)(x)\bullet(g\circ f)(y)}{%
      \rChain{11,13}}
  \closeproofpartwideindR

  \proofpartwideindR[Zusammenführung]{%
    \begin{aligned}[t]
      &\SemigroupHom{f}{A,\star}{B,\diamond}\dsep
        \SemigroupHom{g}{B,\diamond}{C,\bullet}\vdash{}\\[-2pt]
      &\SemigroupHom{g\circ f}{A,\star}{C,\bullet}
    \end{aligned}
  }{%
    \SemigroupHom{f}{A,\star}{B,\diamond}\dsep
    \SemigroupHom{g}{B,\diamond}{C,\bullet}
    \vdash \SemigroupHom{g\circ f}{A,\star}{C,\bullet}%
  }[SemigroupHomCompositionConclusion]
    \proofstepwidestar[1]{\SemigroupHom{f}{A,\star}{B,\diamond}}{\rA}
    \proofstepwidestar[2]{\SemigroupHom{g}{B,\diamond}{C,\bullet}}{\rA}
    \proofstepwidestar[1]{\SemigroupAlg{A}{\star}}{%
      \FormulaRefAuto{SemigroupHomSourceAxiom}[ax]{1}}
    \proofstepwidestar[2]{\SemigroupAlg{C}{\bullet}}{%
      \FormulaRefAuto{SemigroupHomTargetAxiom}[ax]{2}}
    \proofstepwidestar[1,2]{g\circ f\colon A\to C}{%
      \FormulaRefAuto{SemigroupHomCompositionFunction}[thm]{1,2}}
    \proofstepwidestar[6]{x\in A}{\rA}
    \proofstepwidestar[7]{y\in A}{\rA}
    \proofstepwidestar[1,2,6,7]{%
      (g\circ f)(x\star y)=g\bigl(f(x)\diamond f(y)\bigr)}{%
      \FormulaRefAuto{SemigroupHomCompositionFirstOperation}[thm]{1,2,6,7}}
    \proofstepwidestar[1,2,6,7]{%
      g\bigl(f(x)\diamond f(y)\bigr)
      =(g\circ f)(x)\bullet(g\circ f)(y)}{%
      \FormulaRefAuto{SemigroupHomCompositionSecondOperation}[thm]{1,2,6,7}}
    \proofstepwidestar[1,2,6,7]{%
      (g\circ f)(x\star y)=(g\circ f)(x)\bullet(g\circ f)(y)}{%
      \rChain{8,9}}
    \proofstepwidestar[1,2]{%
      \SemigroupHom{g\circ f}{A,\star}{C,\bullet}}{%
      \FormulaRefAuto{SemigroupHomDef}[def]{3,4,5,10}}
  \closeproofpartwideindR
\end{tabproofsplitwide}

\FormulaThmDeltaK[Die Identität ist ein Halbgruppenisomorphismus]{%
  \SemigroupAlg{A}{\star}
  \vdash\SemigroupIso{\Id_A}{A,\star}{A,\star}%
}{SemigroupIdentityIsomorphism}{%
  \DeltaRow{Mengen}{A\dsep x\dsep y}
  \DeltaRow{Funktionen}{\star}
}
\begin{tabproofwide}
  \proofstepwidestar[1]{\SemigroupAlg{A}{\star}}{\rA}
  \proofstepwidestar[]{\Id_A\colon A\to A}{%
    \FormulaRefAuto{IdA}[thm]}
  \proofstepwidestar[2]{x\in A}{\rA}
  \proofstepwidestar[3]{y\in A}{\rA}
  \proofstepwidestar[1,2,3]{x\star y\in A}{%
    \FormulaRefAuto{SemigroupClosureAxiom}[ax]{1,2,3}}
  \proofstepwidestar[2]{\Id_A(x)=x}{%
    \FormulaRefAuto{x\in A\vdash \Id_A(x)=x}{2}}
  \proofstepwidestar[2]{x=\Id_A(x)}{%
    \FormulaRefAuto{a=b\vdash b=a}{6}}
  \proofstepwidestar[3]{\Id_A(y)=y}{%
    \FormulaRefAuto{x\in A\vdash \Id_A(x)=x}{3}}
  \proofstepwidestar[3]{y=\Id_A(y)}{%
    \FormulaRefAuto{a=b\vdash b=a}{8}}
  \proofstepwide[1,2,3]{\Id_A(x\star y)}{=}{x\star y}{%
    \FormulaRefAuto{x\in A\vdash \Id_A(x)=x}{5}}
  \proofstepwide[1,2,3]{}{=}{\Id_A(x)\star y}{\rIE{7}}
  \proofstepwide[1,2,3]{}{=}{\Id_A(x)\star\Id_A(y)}{\rIE{9}}
  \proofstepwidestar[1,2,3]{%
    \Id_A(x\star y)=\Id_A(x)\star\Id_A(y)}{\rChain{10,12}}
  \proofstepwidestar[1]{%
    \SemigroupHom{\Id_A}{A,\star}{A,\star}}{%
    \FormulaRefAuto{SemigroupHomDef}[def]{1,1,2,13}}
  \proofstepwidestar[]{\Id_A\colon A\bij A}{%
    \FormulaRefAuto{\Id_A\colon A\bij A}[thm]}
  \proofstepwidestar[1]{%
    \SemigroupIso{\Id_A}{A,\star}{A,\star}}{%
    \FormulaRefAuto{SemigroupIsoDef}[def]{14,15}}
\end{tabproofwide}

\FormulaThmDeltaK[Komposition von Halbgruppenisomorphismen]{%
  \begin{aligned}[t]
    &\SemigroupIso{f}{A,\star}{B,\diamond}\dsep
      \SemigroupIso{g}{B,\diamond}{C,\bullet}\\
    &\vdash\SemigroupIso{g\circ f}{A,\star}{C,\bullet}
  \end{aligned}
}{SemigroupIsoComposition}{%
  \DeltaRow{Mengen}{A\dsep B\dsep C}
  \DeltaRow{Funktionen}{f\dsep g\dsep\star\dsep\diamond\dsep\bullet}
}
\begin{tabproof}
  \proofstep{1}{\SemigroupIso{f}{A,\star}{B,\diamond}}{\rA}
  \proofstep{2}{\SemigroupIso{g}{B,\diamond}{C,\bullet}}{\rA}
  \proofstep{1}{\SemigroupHom{f}{A,\star}{B,\diamond}}{%
    \FormulaRefAuto{SemigroupIsoHomAxiom}[ax]{1}}
  \proofstep{2}{\SemigroupHom{g}{B,\diamond}{C,\bullet}}{%
    \FormulaRefAuto{SemigroupIsoHomAxiom}[ax]{2}}
  \proofstep{1,2}{\SemigroupHom{g\circ f}{A,\star}{C,\bullet}}{%
    \FormulaRefAuto{SemigroupHomComposition}[thm]{3,4}}
  \proofstep{1}{f\colon A\bij B}{%
    \FormulaRefAuto{SemigroupIsoBijectiveAxiom}[ax]{1}}
  \proofstep{2}{g\colon B\bij C}{%
    \FormulaRefAuto{SemigroupIsoBijectiveAxiom}[ax]{2}}
  \proofstep{1,2}{g\circ f\colon A\bij C}{%
    \FormulaRefAuto{F\colon A\bij B \dsep G\colon B\bij C\vdash
      G\circ F\colon A\bij C}[thm]{6,7}}
  \proofstep{1,2}{\SemigroupIso{g\circ f}{A,\star}{C,\bullet}}{%
    \FormulaRefAuto{SemigroupIsoDef}[def]{5,8}}
\end{tabproof}

\section{Koordinatenisomorphismen von Halbgruppen}

Die folgende Konstruktion trennt die beiden Koordinaten eines Produkts:
Die erste Koordinate eines Produkts stammt vom linken, die zweite vom
rechten Faktor.  F\"ur ein festes Element \(e\) werden die daf\"ur ben\"otigten
Koordinaten durch Rechts- beziehungsweise Linksmultiplikation mit \(e\)
gewonnen.  Der Ausdruck \emph{\(e\)-Koordinatenisomorphismus} ist die in
diesem Werk verwendete Bezeichnung f\"ur diese spezielle Darstellung.

\subsection{Axiomatische Definition}

\FormulaDefDeltaK[Koordinatentr\"ager bez\"uglich eines
Halbgruppenelements]{%
  \begin{aligned}[t]
    &\SemigroupAlg{A}{\star}\dsep e\in A\vdash{}\\[-2pt]
    &\SemigroupCoordinateCarrier{A}{\star}{e}\coloneqq
      \bigl\{p\in A\bigm|\exists x\in A\,(p=x\star e)\bigr\}
      \times
      \bigl\{q\in A\bigm|\exists y\in A\,(q=e\star y)\bigr\}
  \end{aligned}%
}{SemigroupCoordinateCarrierDef}{%
  \DeltaRow{Mengen}{A\dsep e\dsep p\dsep q\dsep x\dsep y}
  \DeltaRow{Funktionen}{\star}
  \DeltaRow{\textbf{Neues Symbol}}{}
  \DeltaPrem{Koordinatentr\"ager}{%
    \SemigroupCoordinateCarrier{A}{\star}{e}}
}

\FormulaDefDeltaK[Koordinatenprodukt]{%
  \begin{aligned}[t]
    &\SemigroupAlg{A}{\star}\dsep e\in A\vdash{}\\[-2pt]
    &\SemigroupCoordinateProductOp{A}{\star}{e}\colon
      \SemigroupCoordinateCarrier{A}{\star}{e}
      \times\SemigroupCoordinateCarrier{A}{\star}{e}
      \longrightarrow\SemigroupCoordinateCarrier{A}{\star}{e},\\[-2pt]
    &(p,q),(r,s)\in\SemigroupCoordinateCarrier{A}{\star}{e}\vdash
      (p,q)\SemigroupCoordinateProductOp{A}{\star}{e}(r,s)
      \coloneqq(p,s)
  \end{aligned}%
}{SemigroupCoordinateProductDef}{%
  \DeltaRow{Mengen}{A\dsep e\dsep p\dsep q\dsep r\dsep s}
  \DeltaRow{Funktionen}{\star\dsep
    \SemigroupCoordinateProductOp{A}{\star}{e}}
  \DeltaRow{\textbf{Neues Symbol}}{}
  \DeltaPrem{Koordinatenprodukt}{%
    \SemigroupCoordinateProductOp{A}{\star}{e}}
}

\FormulaDefDeltaK[\(e\)-Koordinatenabbildung]{%
  \begin{aligned}[t]
    &\SemigroupAlg{A}{\star}\dsep e\in A\vdash{}\\[-2pt]
    &\SemigroupCoordinateMap{A}{\star}{e}\colon
      A\longrightarrow\SemigroupCoordinateCarrier{A}{\star}{e},\\[-2pt]
    &\forall x\in A\quad
      \SemigroupCoordinateMap{A}{\star}{e}(x)
      \coloneqq(x\star e,e\star x)
  \end{aligned}%
}{SemigroupCoordinateMapDef}{%
  \DeltaRow{Mengen}{A\dsep e\dsep x}
  \DeltaRow{Funktionen}{\star\dsep
    \SemigroupCoordinateMap{A}{\star}{e}}
  \DeltaRow{\textbf{Neues Symbol}}{}
  \DeltaPrem{\(e\)-Koordinatenabbildung}{%
    \SemigroupCoordinateMap{A}{\star}{e}}
}

\FormulaDefDeltaK[Multiplikationsabbildung der \(e\)-Koordinaten]{%
  \begin{aligned}[t]
    &\SemigroupAlg{A}{\star}\dsep e\in A\vdash{}\\[-2pt]
    &\SemigroupCoordinateMultiplicationMap{A}{\star}{e}\colon
      \SemigroupCoordinateCarrier{A}{\star}{e}\longrightarrow A,\\[-2pt]
    &(p,q)\in\SemigroupCoordinateCarrier{A}{\star}{e}\vdash
      \SemigroupCoordinateMultiplicationMap{A}{\star}{e}(p,q)
      \coloneqq p\star q
  \end{aligned}%
}{SemigroupCoordinateMultiplicationMapDef}{%
  \DeltaRow{Mengen}{A\dsep e\dsep p\dsep q}
  \DeltaRow{Funktionen}{\star\dsep
    \SemigroupCoordinateMultiplicationMap{A}{\star}{e}}
  \DeltaRow{\textbf{Neues Symbol}}{}
  \DeltaPrem{Multiplikationsabbildung}{%
    \SemigroupCoordinateMultiplicationMap{A}{\star}{e}}
}

Der neue Strukturbegriff wird nicht als lange Konjunktion notiert.  Seine
vier Bestandteile werden vielmehr als getrennte Axiome eingef\"uhrt und sind
anschlie\ss end einzeln zug\"anglich.

\FormulaDefDeltaK[Begriff des \(e\)-Koordinatenisomorphismus]{%
  \SemigroupCoordinateIso{A}{\star}{e}%
}{SemigroupCoordinateIsoDef}{%
  \DeltaRow{Mengen}{A\dsep e\dsep x\dsep y}
  \DeltaRow{Funktionen}{\star\dsep
    \SemigroupCoordinateMap{A}{\star}{e}\dsep
    \SemigroupCoordinateProductOp{A}{\star}{e}}
  \DeltaRow{\textbf{Axiome}}{}
  \DeltaPrem{Quellhalbgruppe}{\SemigroupAlg{A}{\star}}
  \DeltaPrem{Basiselement}{e\in A}
  \DeltaPrem{Bijektive Koordinatenabbildung}{%
    \SemigroupCoordinateMap{A}{\star}{e}\colon
    A\bij\SemigroupCoordinateCarrier{A}{\star}{e}}
  \DeltaRow{Operationserhaltung}{%
    \begin{gathered}
      x,y\in A\vdash{}\\[-2pt]
      \SemigroupCoordinateMap{A}{\star}{e}(x\star y)
      =\SemigroupCoordinateMap{A}{\star}{e}(x)
        \SemigroupCoordinateProductOp{A}{\star}{e}
        \SemigroupCoordinateMap{A}{\star}{e}(y)
    \end{gathered}}
  \DeltaRow{\textbf{Neues Symbol}}{}
  \DeltaPrem{\(e\)-Koordinatenisomorphismus}{%
    \SemigroupCoordinateIso{A}{\star}{e}}
}

\FormulaAxiomDeltaK[Zugriff auf die Quellhalbgruppe eines
Koordinatenisomorphismus]{%
  \SemigroupCoordinateIso{A}{\star}{e}
  \vdash\SemigroupAlg{A}{\star}%
}{SemigroupCoordinateIsoSourceAxiom}{%
  \DeltaRow{Mengen}{A\dsep e}
  \DeltaRow{Funktionen}{\star}
}

\FormulaAxiomDeltaK[Zugriff auf das Basiselement eines
Koordinatenisomorphismus]{%
  \SemigroupCoordinateIso{A}{\star}{e}\vdash e\in A%
}{SemigroupCoordinateIsoBaseElementAxiom}{%
  \DeltaRow{Mengen}{A\dsep e}
  \DeltaRow{Funktionen}{\star}
}

\FormulaAxiomDeltaK[Zugriff auf die Bijektivit\"at der
Koordinatenabbildung]{%
  \SemigroupCoordinateIso{A}{\star}{e}\vdash
  \SemigroupCoordinateMap{A}{\star}{e}\colon
  A\bij\SemigroupCoordinateCarrier{A}{\star}{e}%
}{SemigroupCoordinateIsoBijectiveAxiom}{%
  \DeltaRow{Mengen}{A\dsep e}
  \DeltaRow{Funktionen}{\star\dsep
    \SemigroupCoordinateMap{A}{\star}{e}}
}

\FormulaAxiomDeltaK[Zugriff auf die Operationserhaltung der
Koordinatenabbildung]{%
  \begin{aligned}[t]
    &\SemigroupCoordinateIso{A}{\star}{e}\dsep x,y\in A\vdash{}\\[-2pt]
    &\SemigroupCoordinateMap{A}{\star}{e}(x\star y)
      =\SemigroupCoordinateMap{A}{\star}{e}(x)
       \SemigroupCoordinateProductOp{A}{\star}{e}
       \SemigroupCoordinateMap{A}{\star}{e}(y)
  \end{aligned}%
}{SemigroupCoordinateIsoOperationAxiom}{%
  \DeltaRow{Mengen}{A\dsep e\dsep x\dsep y}
  \DeltaRow{Funktionen}{\star\dsep
    \SemigroupCoordinateMap{A}{\star}{e}\dsep
    \SemigroupCoordinateProductOp{A}{\star}{e}}
}

\subsection{Grundlegende Eigenschaften}

\FormulaThmDeltaK[Elementkriterium des Koordinatentr\"agers]{%
  \begin{aligned}[t]
    &\SemigroupAlg{A}{\star}\dsep e\in A\vdash{}\\[-2pt]
    &(p,q)\in\SemigroupCoordinateCarrier{A}{\star}{e}
      \leftrightarrow\\[-2pt]
    &\quad\Bigl(p\in A\land\exists x\in A\,(p=x\star e)\Bigr)
      \land
      \Bigl(q\in A\land\exists y\in A\,(q=e\star y)\Bigr)
  \end{aligned}%
}{SemigroupCoordinateCarrierMembership}{%
  \DeltaRow{Mengen}{A\dsep e\dsep p\dsep q\dsep x\dsep y}
  \DeltaRow{Funktionen}{\star}
}
\begin{proof}
Setze in \FormulaRefAuto{SemigroupCoordinateCarrierDef}[def] die
Definition des Koordinatentr\"agers ein.  Das Elementkriterium des
kartesischen Produkts und zweimal das Elementkriterium einer
ausgesonderten Teilmenge liefern genau die beiden Konjunkte der rechten
Seite.
\end{proof}

\FormulaThmDeltaK[Der Koordinatentr\"ager ist nichtleer]{%
  \SemigroupAlg{A}{\star}\dsep e\in A\vdash
  \SemigroupCoordinateCarrier{A}{\star}{e}\neq\varnothing%
}{SemigroupCoordinateCarrierNonempty}{%
  \DeltaRow{Mengen}{A\dsep e}
  \DeltaRow{Funktionen}{\star}
}
\begin{proof}
Nach der Abgeschlossenheit liegt \(e\star e\) in \(A\).  In beiden
Faktoren des Koordinatentr\"agers dient \(e\) als Zeuge, also liegt
\((e\star e,e\star e)\) nach
\FormulaRefAuto{SemigroupCoordinateCarrierMembership}[thm] im
Koordinatentr\"ager.  Dieser ist daher nichtleer.
\end{proof}

\FormulaThmDeltaK[Auswertung des Koordinatenprodukts]{%
  \begin{aligned}[t]
    &\SemigroupAlg{A}{\star}\dsep e\in A\dsep
      (p,q),(r,s)\in\SemigroupCoordinateCarrier{A}{\star}{e}\vdash{}\\[-2pt]
    &(p,q)\SemigroupCoordinateProductOp{A}{\star}{e}(r,s)=(p,s)
  \end{aligned}%
}{SemigroupCoordinateProductEvaluation}{%
  \DeltaRow{Mengen}{A\dsep e\dsep p\dsep q\dsep r\dsep s}
  \DeltaRow{Funktionen}{\star\dsep
    \SemigroupCoordinateProductOp{A}{\star}{e}}
}
\begin{proof}
Dies ist die Auswertungsregel aus
\FormulaRefAuto{SemigroupCoordinateProductDef}[def].
\end{proof}

\FormulaThmDeltaK[Das Koordinatenprodukt bildet eine Halbgruppe]{%
  \SemigroupAlg{A}{\star}\dsep e\in A\vdash
  \SemigroupAlg{\SemigroupCoordinateCarrier{A}{\star}{e}}
    {\SemigroupCoordinateProductOp{A}{\star}{e}}%
}{SemigroupCoordinateProductSemigroup}{%
  \DeltaRow{Mengen}{A\dsep e\dsep p\dsep q\dsep r\dsep s\dsep
    t\dsep u}
  \DeltaRow{Funktionen}{\star\dsep
    \SemigroupCoordinateProductOp{A}{\star}{e}}
}
\begin{proof}
Die Abgeschlossenheit ist Teil der Typisierung in
\FormulaRefAuto{SemigroupCoordinateProductDef}[def].  Sind
\((p,q),(r,s),(t,u)\) Elemente des Koordinatentr\"agers, so liefert
\FormulaRefAuto{SemigroupCoordinateProductEvaluation}[thm] zweimal
\[
 \bigl((p,q)\SemigroupCoordinateProductOp{A}{\star}{e}(r,s)\bigr)
      \SemigroupCoordinateProductOp{A}{\star}{e}(t,u)
 =(p,u)
 =(p,q)\SemigroupCoordinateProductOp{A}{\star}{e}
      \bigl((r,s)\SemigroupCoordinateProductOp{A}{\star}{e}(t,u)\bigr).
\]
Damit ist das Koordinatenprodukt assoziativ; zusammen mit der
Abgeschlossenheit folgt die Behauptung aus
\FormulaRefAuto{SemigroupDef}[def].
\end{proof}

\FormulaThmDeltaK[Idempotenz des Koordinatenprodukts]{%
  \begin{aligned}[t]
    &\SemigroupAlg{A}{\star}\dsep e\in A\dsep
      u\in\SemigroupCoordinateCarrier{A}{\star}{e}\vdash{}\\[-2pt]
    &u\SemigroupCoordinateProductOp{A}{\star}{e}u=u
  \end{aligned}%
}{SemigroupCoordinateProductIdempotence}{%
  \DeltaRow{Mengen}{A\dsep e\dsep u\dsep p\dsep q}
  \DeltaRow{Funktionen}{\star\dsep
    \SemigroupCoordinateProductOp{A}{\star}{e}}
}
\begin{proof}
Nach dem Elementkriterium des kartesischen Produkts hat \(u\) die Form
\((p,q)\).  Die Auswertungsregel ergibt
\[
 (p,q)\SemigroupCoordinateProductOp{A}{\star}{e}(p,q)=(p,q).
\]
\end{proof}

\FormulaThmDeltaK[Dreitermgesetz des Koordinatenprodukts]{%
  \begin{aligned}[t]
    &\SemigroupAlg{A}{\star}\dsep e\in A\dsep
      u,v,w\in\SemigroupCoordinateCarrier{A}{\star}{e}\vdash{}\\[-2pt]
    &(u\SemigroupCoordinateProductOp{A}{\star}{e}v)
       \SemigroupCoordinateProductOp{A}{\star}{e}w
      =u\SemigroupCoordinateProductOp{A}{\star}{e}w
  \end{aligned}%
}{SemigroupCoordinateProductThreeTermLaw}{%
  \DeltaRow{Mengen}{A\dsep e\dsep u\dsep v\dsep w\dsep
    p\dsep q\dsep r\dsep s\dsep t\dsep z}
  \DeltaRow{Funktionen}{\star\dsep
    \SemigroupCoordinateProductOp{A}{\star}{e}}
}
\begin{proof}
Schreibe \(u=(p,q)\), \(v=(r,s)\) und \(w=(t,z)\).  Zweimalige
Anwendung von
\FormulaRefAuto{SemigroupCoordinateProductEvaluation}[thm] ergibt auf
beiden Seiten \((p,z)\).
\end{proof}

\FormulaThmDeltaK[Rechtecksidentit\"at des Koordinatenprodukts]{%
  \begin{aligned}[t]
    &\SemigroupAlg{A}{\star}\dsep e\in A\dsep
      u,v\in\SemigroupCoordinateCarrier{A}{\star}{e}\vdash{}\\[-2pt]
    &(u\SemigroupCoordinateProductOp{A}{\star}{e}v)
       \SemigroupCoordinateProductOp{A}{\star}{e}u=u
  \end{aligned}%
}{SemigroupCoordinateProductRectangularIdentity}{%
  \DeltaRow{Mengen}{A\dsep e\dsep u\dsep v}
  \DeltaRow{Funktionen}{\star\dsep
    \SemigroupCoordinateProductOp{A}{\star}{e}}
}
\begin{tabproofwide}
  \proofstepwidestar[1]{\SemigroupAlg{A}{\star}}{\rA}
  \proofstepwidestar[2]{e\in A}{\rA}
  \proofstepwidestar[3]{u,v\in
    \SemigroupCoordinateCarrier{A}{\star}{e}}{\rA}
  \proofstepwidestar[1,2,3]{%
    (u\SemigroupCoordinateProductOp{A}{\star}{e}v)
      \SemigroupCoordinateProductOp{A}{\star}{e}u
    =u\SemigroupCoordinateProductOp{A}{\star}{e}u}{%
    \FormulaRefAuto{SemigroupCoordinateProductThreeTermLaw}[thm]{1,2,3}}
  \proofstepwidestar[1,2,3]{%
    u\SemigroupCoordinateProductOp{A}{\star}{e}u=u}{%
    \FormulaRefAuto{SemigroupCoordinateProductIdempotence}[thm]{1,2,3}}
  \proofstepwidestar[1,2,3]{%
    (u\SemigroupCoordinateProductOp{A}{\star}{e}v)
      \SemigroupCoordinateProductOp{A}{\star}{e}u=u}{\rChain{4,5}}
\end{tabproofwide}

\FormulaThmDeltaK[Die \(e\)-Koordinatenabbildung ist eine Abbildung]{%
  \SemigroupAlg{A}{\star}\dsep e\in A\vdash
  \SemigroupCoordinateMap{A}{\star}{e}\colon
  A\to\SemigroupCoordinateCarrier{A}{\star}{e}%
}{SemigroupCoordinateMapFunction}{%
  \DeltaRow{Mengen}{A\dsep e}
  \DeltaRow{Funktionen}{\star\dsep
    \SemigroupCoordinateMap{A}{\star}{e}}
}
\begin{proof}
F\"ur \(x\in A\) liegen \(x\star e\) und \(e\star x\) nach der
Abgeschlossenheit in \(A\).  In den beiden Faktoren des
Koordinatentr\"agers dient jeweils \(x\) als Zeuge.  Somit liegt
\((x\star e,e\star x)\) im Koordinatentr\"ager.  Die Behauptung folgt aus
\FormulaRefAuto{SemigroupCoordinateMapDef}[def].
\end{proof}

\FormulaThmDeltaK[Auswertung der \(e\)-Koordinatenabbildung]{%
  \SemigroupAlg{A}{\star}\dsep e,x\in A\vdash
  \SemigroupCoordinateMap{A}{\star}{e}(x)=(x\star e,e\star x)%
}{SemigroupCoordinateMapEvaluation}{%
  \DeltaRow{Mengen}{A\dsep e\dsep x}
  \DeltaRow{Funktionen}{\star\dsep
    \SemigroupCoordinateMap{A}{\star}{e}}
}
\begin{proof}
Dies ist die Auswertungsregel aus
\FormulaRefAuto{SemigroupCoordinateMapDef}[def].
\end{proof}

\FormulaThmDeltaK[Die Multiplikationsabbildung ist eine Abbildung]{%
  \SemigroupAlg{A}{\star}\dsep e\in A\vdash
  \SemigroupCoordinateMultiplicationMap{A}{\star}{e}\colon
  \SemigroupCoordinateCarrier{A}{\star}{e}\to A%
}{SemigroupCoordinateMultiplicationMapFunction}{%
  \DeltaRow{Mengen}{A\dsep e\dsep p\dsep q}
  \DeltaRow{Funktionen}{\star\dsep
    \SemigroupCoordinateMultiplicationMap{A}{\star}{e}}
}
\begin{proof}
Ist \((p,q)\) im Koordinatentr\"ager, so sind \(p,q\in A\).  Daher liegt
\(p\star q\) nach der Abgeschlossenheit in \(A\).  Dies ist genau die
Wohldefiniertheit der Abbildung aus
\FormulaRefAuto{SemigroupCoordinateMultiplicationMapDef}[def].
\end{proof}

\FormulaThmDeltaK[Auswertung der Multiplikationsabbildung]{%
  \begin{aligned}[t]
    &\SemigroupAlg{A}{\star}\dsep e\in A\dsep
      (p,q)\in\SemigroupCoordinateCarrier{A}{\star}{e}\vdash{}\\[-2pt]
    &\SemigroupCoordinateMultiplicationMap{A}{\star}{e}(p,q)=p\star q
  \end{aligned}%
}{SemigroupCoordinateMultiplicationMapEvaluation}{%
  \DeltaRow{Mengen}{A\dsep e\dsep p\dsep q}
  \DeltaRow{Funktionen}{\star\dsep
    \SemigroupCoordinateMultiplicationMap{A}{\star}{e}}
}
\begin{proof}
Dies ist die Auswertungsregel aus
\FormulaRefAuto{SemigroupCoordinateMultiplicationMapDef}[def].
\end{proof}

\FormulaThmDeltaK[Der Koordinatenisomorphismus ist ein
Halbgruppenisomorphismus]{%
  \SemigroupCoordinateIso{A}{\star}{e}\vdash
  \SemigroupIso{\SemigroupCoordinateMap{A}{\star}{e}}
    {A,\star}
    {\SemigroupCoordinateCarrier{A}{\star}{e},
     \SemigroupCoordinateProductOp{A}{\star}{e}}%
}{SemigroupCoordinateIsoIsSemigroupIso}{%
  \DeltaRow{Mengen}{A\dsep e\dsep x\dsep y}
  \DeltaRow{Funktionen}{\star\dsep
    \SemigroupCoordinateMap{A}{\star}{e}\dsep
    \SemigroupCoordinateProductOp{A}{\star}{e}}
}
\begin{tabproofwide}
  \proofstepwidestar[1]{\SemigroupCoordinateIso{A}{\star}{e}}{\rA}
  \proofstepwidestar[1]{\SemigroupAlg{A}{\star}}{%
    \FormulaRefAuto{SemigroupCoordinateIsoSourceAxiom}[ax]{1}}
  \proofstepwidestar[1]{e\in A}{%
    \FormulaRefAuto{SemigroupCoordinateIsoBaseElementAxiom}[ax]{1}}
  \proofstepwidestar[1]{%
    \SemigroupAlg{\SemigroupCoordinateCarrier{A}{\star}{e}}
      {\SemigroupCoordinateProductOp{A}{\star}{e}}}{%
    \FormulaRefAuto{SemigroupCoordinateProductSemigroup}[thm]{2,3}}
  \proofstepwidestar[1]{%
    \SemigroupCoordinateMap{A}{\star}{e}\colon
    A\bij\SemigroupCoordinateCarrier{A}{\star}{e}}{%
    \FormulaRefAuto{SemigroupCoordinateIsoBijectiveAxiom}[ax]{1}}
  \proofstepwidestar[1]{%
    \SemigroupCoordinateMap{A}{\star}{e}\colon
    A\to\SemigroupCoordinateCarrier{A}{\star}{e}}{%
    \FormulaRefAuto{SemigroupCoordinateMapFunction}[thm]{2,3}}
  \proofstepwidestarbreak[1]{%
    \forall x\in A\,\forall y\in A\,
    \Bigl(
      \SemigroupCoordinateMap{A}{\star}{e}(x\star y)
      =\SemigroupCoordinateMap{A}{\star}{e}(x)
       \SemigroupCoordinateProductOp{A}{\star}{e}
       \SemigroupCoordinateMap{A}{\star}{e}(y)
    \Bigr)}{%
    \text{Allgemeinheit von }x,y\text{ und }%
    \FormulaRefAuto{SemigroupCoordinateIsoOperationAxiom}[ax]{1}}
  \proofstepwidestar[1]{%
    \SemigroupHom{\SemigroupCoordinateMap{A}{\star}{e}}
      {A,\star}
      {\SemigroupCoordinateCarrier{A}{\star}{e},
       \SemigroupCoordinateProductOp{A}{\star}{e}}}{%
    \FormulaRefAuto{SemigroupHomDef}[def]{2,4,6,7}}
  \proofstepwidestar[1]{%
    \SemigroupIso{\SemigroupCoordinateMap{A}{\star}{e}}
      {A,\star}
      {\SemigroupCoordinateCarrier{A}{\star}{e},
       \SemigroupCoordinateProductOp{A}{\star}{e}}}{%
    \FormulaRefAuto{SemigroupIsoDef}[def]{8,5}}
\end{tabproofwide}

\FormulaThmDeltaK[Idempotenz- und Dreitermkriterium f\"ur den
Koordinatenisomorphismus]{%
  \begin{aligned}[t]
    &\SemigroupAlg{A}{\star}\dsep e\in A\dsep
      \forall x\in A\,(x\star x=x)\dsep{}\\[-2pt]
    &\forall x,y,z\in A\,\bigl((x\star y)\star z=x\star z\bigr)
      \vdash\SemigroupCoordinateIso{A}{\star}{e}
  \end{aligned}%
}{SemigroupCoordinateIsoCriterion}{%
  \DeltaRow{Mengen}{A\dsep e\dsep p\dsep q\dsep r\dsep s\dsep
    x\dsep y\dsep z}
  \DeltaRow{Funktionen}{\star\dsep
    \SemigroupCoordinateMap{A}{\star}{e}\dsep
    \SemigroupCoordinateMultiplicationMap{A}{\star}{e}\dsep
    \SemigroupCoordinateProductOp{A}{\star}{e}}
}
\begin{proof}
Schreibe kurz

\[
 K=\SemigroupCoordinateCarrier{A}{\star}{e},\qquad
 \chi=\SemigroupCoordinateMap{A}{\star}{e},\qquad
 \mu=\SemigroupCoordinateMultiplicationMap{A}{\star}{e}.
\]

Nach
\FormulaRefAuto{SemigroupCoordinateMapFunction}[thm] und
\FormulaRefAuto{SemigroupCoordinateMultiplicationMapFunction}[thm]
sind \(\chi\colon A\to K\) und \(\mu\colon K\to A\) Abbildungen.  F\"ur
\(x\in A\) liefern Assoziativit\"at, Dreitermgesetz und Idempotenz die
vollst\"andig geklammerte Kette
\[
\begin{aligned}
 \mu(\chi(x))
 &=(x\star e)\star(e\star x)\\
 &=\bigl((x\star e)\star e\bigr)\star x\\
 &=(x\star e)\star x\\
 &=x\star x=x.
\end{aligned}
\]
Ist \((p,q)\in K\), so gibt es \(a,b\in A\) mit
\(p=a\star e\) und \(q=e\star b\).  Nach dem Dreitermgesetz,
der Assoziativit\"at und der Idempotenz gilt
\[
 p\star e=p,\qquad e\star q=q,\qquad
 (p\star q)\star e=p\star e,
 \qquad e\star(p\star q)=(e\star p)\star q=e\star q.
\]
Mit den Auswertungsregeln von \(\chi\) und \(\mu\) folgt daher
\[
 \chi(\mu(p,q))
 =\bigl((p\star q)\star e,e\star(p\star q)\bigr)
 =(p,q).
\]
Die Allgemeinheit der eingesetzten Elemente und
\FormulaRefAuto{FunctionExtensionality}[thm] liefern aus den beiden
Gleichungen
\[
 \mu\circ\chi=\Id_A,
 \qquad
 \chi\circ\mu=\Id_K.
\]
Der Satz
\FormulaRefAuto{F\circ G = \Id_A\dsep G\circ F =\Id_B
  \vdash G\colon A\bij B}[thm]
\emph{(zwei wechselseitige Inversengleichungen liefern eine
Bijektion)} ergibt daher \(\chi\colon A\bij K\).

F\"ur \(x,y\in A\) ergibt sich schlie\ss lich
\[
\begin{aligned}
 \chi(x\star y)
 &=\bigl((x\star y)\star e,e\star(x\star y)\bigr)\\
 &=\bigl(x\star e,(e\star x)\star y\bigr)\\
 &=\bigl(x\star e,e\star y\bigr)
  =\chi(x)\SemigroupCoordinateProductOp{A}{\star}{e}\chi(y).
\end{aligned}
\]
In der zweiten Zeile wurden das Dreitermgesetz und die Assoziativit\"at,
in der dritten Zeile erneut das Dreitermgesetz verwendet.  Die
Quellhalbgruppe, das Basiselement, die soeben bewiesene Bijektivit\"at und
die zuletzt gezeigte Operationserhaltung sind genau die vier Axiome aus
\FormulaRefAuto{SemigroupCoordinateIsoDef}[def].
\end{proof}

\FormulaThmDeltaK[Idempotenz aus einem Koordinatenisomorphismus]{%
  \SemigroupCoordinateIso{A}{\star}{e}\vdash
  \forall x\in A\,(x\star x=x)%
}{SemigroupCoordinateIsoIdempotence}{%
  \DeltaRow{Mengen}{A\dsep e\dsep x}
  \DeltaRow{Funktionen}{\star\dsep
    \SemigroupCoordinateMap{A}{\star}{e}\dsep
    \SemigroupCoordinateProductOp{A}{\star}{e}}
}
\begin{proof}
Der abgeleitete Halbgruppenisomorphismus aus
\FormulaRefAuto{SemigroupCoordinateIsoIsSemigroupIso}[thm] ist injektiv
und erh\"alt die Operation.  F\"ur \(x\in A\) ist sein Bild ein Element
des Koordinatentr\"agers.  Nach
\FormulaRefAuto{SemigroupCoordinateProductIdempotence}[thm] gilt daher
\[
 \begin{aligned}
 &\SemigroupCoordinateMap{A}{\star}{e}(x\star x)\\
 &\quad=
 \SemigroupCoordinateMap{A}{\star}{e}(x)
 \SemigroupCoordinateProductOp{A}{\star}{e}
 \SemigroupCoordinateMap{A}{\star}{e}(x)
 =\SemigroupCoordinateMap{A}{\star}{e}(x).
 \end{aligned}
\]
Die Injektivit\"at der Koordinatenabbildung liefert \(x\star x=x\);
die Allgemeinheit
von \(x\) ergibt die Behauptung.
\end{proof}

\FormulaThmDeltaK[Dreitermgesetz aus einem Koordinatenisomorphismus]{%
  \begin{aligned}[t]
    &\SemigroupCoordinateIso{A}{\star}{e}\vdash{}\\[-2pt]
    &\forall x,y,z\in A\,
      \bigl((x\star y)\star z=x\star z\bigr)
  \end{aligned}%
}{SemigroupCoordinateIsoThreeTermLaw}{%
  \DeltaRow{Mengen}{A\dsep e\dsep x\dsep y\dsep z}
  \DeltaRow{Funktionen}{\star\dsep
    \SemigroupCoordinateMap{A}{\star}{e}\dsep
    \SemigroupCoordinateProductOp{A}{\star}{e}}
}
\begin{proof}
Wie im vorigen Satz ist die Koordinatenabbildung ein injektiver
Homomorphismus.  F\"ur \(x,y,z\in A\) liefert das Dreitermgesetz des
Koordinatenprodukts
\[
\begin{aligned}
 &\SemigroupCoordinateMap{A}{\star}{e}
   \bigl((x\star y)\star z\bigr)\\
 &\quad=\Bigl(
   \SemigroupCoordinateMap{A}{\star}{e}(x)
   \SemigroupCoordinateProductOp{A}{\star}{e}
   \SemigroupCoordinateMap{A}{\star}{e}(y)
   \Bigr)
   \SemigroupCoordinateProductOp{A}{\star}{e}
   \SemigroupCoordinateMap{A}{\star}{e}(z)\\
 &\quad=\SemigroupCoordinateMap{A}{\star}{e}(x)
   \SemigroupCoordinateProductOp{A}{\star}{e}
   \SemigroupCoordinateMap{A}{\star}{e}(z)
   =\SemigroupCoordinateMap{A}{\star}{e}(x\star z).
\end{aligned}
\]
Injektivit\"at und Allgemeinheit der drei Elemente ergeben die Aussage.
\end{proof}

\FormulaThmDeltaK[Rechtecksidentit\"at aus einem
Koordinatenisomorphismus]{%
  \begin{aligned}[t]
    &\SemigroupCoordinateIso{A}{\star}{e}\vdash{}\\[-2pt]
    &\forall x,y\in A\,\bigl((x\star y)\star x=x\bigr)
  \end{aligned}%
}{SemigroupCoordinateIsoRectangularIdentity}{%
  \DeltaRow{Mengen}{A\dsep e\dsep x\dsep y}
  \DeltaRow{Funktionen}{\star}
}
\begin{proof}
Nach \FormulaRefAuto{SemigroupCoordinateIsoThreeTermLaw}[thm] und
\FormulaRefAuto{SemigroupCoordinateIsoIdempotence}[thm] gilt f\"ur
\(x,y\in A\)
\[
 (x\star y)\star x=x\star x=x.
\]
Die Allgemeinheit von \(x,y\) liefert die Behauptung.
\end{proof}

\section{Morphismen von Potenzhalbgruppen}

\FormulaThmDeltaK[Die Einermengenkopie ist zur Ausgangshalbgruppe
isomorph]{%
  \SemigroupAlg{A}{\star}\vdash
  \SemigroupIso{\eta_A}
    {A,\star}{\operatorname{Sing}(A),\star_{\mathcal P}}%
}{SemigroupSingletonCopyIsomorphism}{%
  \DeltaRow{Mengen}{A\dsep a\dsep b\dsep X\dsep z}
  \DeltaRow{Funktionen}{\star\dsep\star_{\mathcal P}\dsep\eta_A}
}
\begin{proof}
Seien \(a,b\in A\). Ihre Einermengen liegen in \(\PowNE A\). Ferner ist
\(a\star b\in A\) nach \FormulaRefAuto{SemigroupClosureAxiom}[ax]. Nach
\FormulaRefAuto{PowerSemigroupProductMembership}[thm] gilt für jedes \(z\)
\[
  z\in\{a\}\star_{\mathcal P}\{b\}
  \quad\Longleftrightarrow\quad
  z=a\star b
  \quad\Longleftrightarrow\quad
  z\in\{a\star b\}.
\]
Die Extensionalität liefert somit
\[
  \{a\}\star_{\mathcal P}\{b\}=\{a\star b\}.       \tag{1}
\]
Da \(a\star b\in A\), liegt die rechte Seite von (1) in
\(\operatorname{Sing}(A)\). Damit ist
\(\operatorname{Sing}(A)\) unter \(\star_{\mathcal P}\) abgeschlossen.
Jedes \(X\in\operatorname{Sing}(A)\) ist eine nichtleere Teilmenge von
\(A\), also gilt \(\operatorname{Sing}(A)\subseteq\PowNE A\). Die
Assoziativität aus \FormulaRefAuto{NonemptyPowerSemigroup}[thm] schränkt sich
daher auf \(\operatorname{Sing}(A)\) ein. Nach
\FormulaRefAuto{SemigroupDef}[def] ist folglich
\(\SemigroupAlg{\operatorname{Sing}(A)}{\star_{\mathcal P}}\) erfüllt.

Nach \FormulaRefAuto{SemigroupSingletonCopyDef}[def] gilt
\(\eta_A(a)=\{a\}\). Gleichung (1) ist deshalb genau
\[
  \eta_A(a\star b)
   =\eta_A(a)\star_{\mathcal P}\eta_A(b).
\]
Zusammen mit der Typisierung von \(\eta_A\), der vorausgesetzten
Quellhalbgruppe und der eben bewiesenen Zielhalbgruppe ergibt
\FormulaRefAuto{SemigroupHomDef}[def], dass \(\eta_A\) ein
Halbgruppenhomomorphismus ist. Aus \(\eta_A(a)=\eta_A(b)\), also
\(\{a\}=\{b\}\), folgt \(a=b\); und jedes Element von
\(\operatorname{Sing}(A)\) hat nach seiner Definition die Form \(\eta_A(a)\)
für ein \(a\in A\). Somit ist \(\eta_A\) bijektiv. Mit
\FormulaRefAuto{SemigroupIsoDef}[def] folgt die behauptete Isomorphie.
\end{proof}

\FormulaThmDeltaK[Ein Homomorphismus erhält das Mengenprodukt]{%
  \begin{aligned}[t]
    &\SemigroupHom{f}{A,\star}{B,\diamond}\dsep
      X,Y\in\PowNE{A}\vdash{}\\[-2pt]
    &f[X\star_{\mathcal P}Y]
      =f[X]\diamond_{\mathcal P}f[Y]
  \end{aligned}%
}{SemigroupHomPreservesPowerProduct}{%
  \DeltaRow{Mengen}{A\dsep B\dsep X\dsep Y\dsep
    u\dsep v\dsep w\dsep x\dsep y\dsep z}
  \DeltaRow{Funktionen}{f\dsep\star\dsep\diamond\dsep
    \star_{\mathcal P}\dsep\diamond_{\mathcal P}}
}
\begin{tabproofsplitwide}
  \proofpartwideindR[Bilder fester Faktoren]{%
    \begin{aligned}[t]
      &\SemigroupHom{f}{A,\star}{B,\diamond}\dsep
        X,Y\in\PowNE{A}\dsep x\in X\dsep y\in Y\vdash{}\\[-2pt]
      &f(x)\diamond f(y)\in f[X]\diamond_{\mathcal P}f[Y]
    \end{aligned}
  }{%
    \begin{aligned}[t]
      &\SemigroupHom{f}{A,\star}{B,\diamond}\dsep
        X,Y\in\PowNE{A}\dsep x\in X\dsep y\in Y\vdash{}\\[-2pt]
      &f(x)\diamond f(y)\in f[X]\diamond_{\mathcal P}f[Y]
    \end{aligned}
  }[PowerHomImagesOfFactors]
    \proofstepwidestar[1]{\SemigroupHom{f}{A,\star}{B,\diamond}}{\rA}
    \proofstepwidestar[2]{X\in\PowNE{A}}{\rA}
    \proofstepwidestar[3]{Y\in\PowNE{A}}{\rA}
    \proofstepwidestar[4]{x\in X}{\rA}
    \proofstepwidestar[5]{y\in Y}{\rA}
    \proofstepwidestar[1]{\SemigroupAlg{B}{\diamond}}{%
      \FormulaRefAuto{SemigroupHomTargetAxiom}[ax]{1}}
    \proofstepwidestar[1]{f\colon A\to B}{%
      \FormulaRefAuto{SemigroupHomFunctionAxiom}[ax]{1}}
    \proofstepwidestar[1,2]{f[X]\in\PowNE{B}}{%
      \FormulaRefAuto{NonemptyDirectImage}[thm]{7,2}}
    \proofstepwidestar[1,3]{f[Y]\in\PowNE{B}}{%
      \FormulaRefAuto{NonemptyDirectImage}[thm]{7,3}}
    \proofstepwidestar[2]{X\subseteq A}{%
      \rAEa{\FormulaRefAuto{NonemptyPowersetMembership}[thm]{2}}}
    \proofstepwidestar[3]{Y\subseteq A}{%
      \rAEa{\FormulaRefAuto{NonemptyPowersetMembership}[thm]{3}}}
    \proofstepwidestar[1,2,4]{f(x)\in f[X]}{%
      \FormulaRefAuto{C\subseteq A\dsep F\colon A\to B\dsep
        x\in C\vdash F(x)\in F[C]}[thm]{10,7,4}}
    \proofstepwidestar[1,3,5]{f(y)\in f[Y]}{%
      \FormulaRefAuto{C\subseteq A\dsep F\colon A\to B\dsep
        x\in C\vdash F(x)\in F[C]}[thm]{11,7,5}}
    \proofstepwidestar[1,2,3,4,5]{%
      f(x)\diamond f(y)\in f[X]\diamond_{\mathcal P}f[Y]}{%
      \FormulaRefAuto{PowerSemigroupProductContains}[thm]{6,8,9,12,13}}
  \closeproofpartwideindR

  \proofpartwideindR[Vorwärtsrichtung mit festen Zeugen]{%
    \begin{aligned}[t]
      &\SemigroupHom{f}{A,\star}{B,\diamond}\dsep X,Y\in\PowNE{A}
        \dsep x\in X\dsep y\in Y\dsep{}\\[-2pt]
      &w=x\star y\dsep z=f(w)
        \vdash z\in f[X]\diamond_{\mathcal P}f[Y]
    \end{aligned}
  }{%
    \begin{aligned}[t]
      &\SemigroupHom{f}{A,\star}{B,\diamond}\dsep X,Y\in\PowNE{A}
        \dsep x\in X\dsep y\in Y\dsep{}\\[-2pt]
      &w=x\star y\dsep z=f(w)
        \vdash z\in f[X]\diamond_{\mathcal P}f[Y]
    \end{aligned}
  }[PowerHomForwardFixedWitnesses]
    \proofstepwidestar[1]{\SemigroupHom{f}{A,\star}{B,\diamond}}{\rA}
    \proofstepwidestar[2]{X\in\PowNE{A}}{\rA}
    \proofstepwidestar[3]{Y\in\PowNE{A}}{\rA}
    \proofstepwidestar[4]{x\in X}{\rA}
    \proofstepwidestar[5]{y\in Y}{\rA}
    \proofstepwidestar[6]{w=x\star y}{\rA}
    \proofstepwidestar[7]{z=f(w)}{\rA}
    \proofstepwidestar[2,4]{x\in A}{%
      \FormulaRefAuto{PowerSemigroupCarrierElement}[thm]{2,4}}
    \proofstepwidestar[3,5]{y\in A}{%
      \FormulaRefAuto{PowerSemigroupCarrierElement}[thm]{3,5}}
    \proofstepwidestar[6,7]{z=f(x\star y)}{\rIE{6,7}}
    \proofstepwidestar[1,2,3,4,5]{%
      f(x\star y)=f(x)\diamond f(y)}{%
      \FormulaRefAuto{SemigroupHomOperationAxiom}[ax]{1,8,9}}
    \proofstepwidestar[1,2,3,4,5,6,7]{z=f(x)\diamond f(y)}{%
      \rChain{10,11}}
    \proofstepwidestar[1,2,3,4,5]{%
      f(x)\diamond f(y)\in f[X]\diamond_{\mathcal P}f[Y]}{%
      \FormulaRefAuto{PowerHomImagesOfFactors}[thm]{1,2,3,4,5}}
    \proofstepwidestar[1,2,3,4,5,6,7]{%
      z\in f[X]\diamond_{\mathcal P}f[Y]}{\rIE{12,13}}
  \closeproofpartwideindR

  \proofpartwideindR[Vorwärtsrichtung mit einem Produktzeugen]{%
    \begin{aligned}[t]
      &\SemigroupHom{f}{A,\star}{B,\diamond}\dsep X,Y\in\PowNE{A}
        \dsep w\in X\star_{\mathcal P}Y\dsep z=f(w)\vdash{}\\[-2pt]
      &z\in f[X]\diamond_{\mathcal P}f[Y]
    \end{aligned}
  }{%
    \begin{aligned}[t]
      &\SemigroupHom{f}{A,\star}{B,\diamond}\dsep X,Y\in\PowNE{A}
        \dsep w\in X\star_{\mathcal P}Y\dsep z=f(w)\vdash{}\\[-2pt]
      &z\in f[X]\diamond_{\mathcal P}f[Y]
    \end{aligned}
  }[PowerHomForwardProductWitness]
    \proofstepwidestar[1]{\SemigroupHom{f}{A,\star}{B,\diamond}}{\rA}
    \proofstepwidestar[2]{X\in\PowNE{A}}{\rA}
    \proofstepwidestar[3]{Y\in\PowNE{A}}{\rA}
    \proofstepwidestar[4]{w\in X\star_{\mathcal P}Y}{\rA}
    \proofstepwidestar[5]{z=f(w)}{\rA}
    \proofstepwidestar[1]{\SemigroupAlg{A}{\star}}{%
      \FormulaRefAuto{SemigroupHomSourceAxiom}[ax]{1}}
    \proofstepwidestar[1,2,3,4]{%
      \exists x\in X\,\exists y\in Y\,(w=x\star y)}{%
      \FormulaRefAuto{PowerSemigroupProductWitnesses}[thm]{6,2,3,4}}
    \proofstepwidestar[8]{x\in X\land\exists y\in Y\,(w=x\star y)}{\rA}
    \proofstepwidestar[8]{x\in X}{\rAEa{8}}
    \proofstepwidestar[8]{\exists y\in Y\,(w=x\star y)}{\rAEb{8}}
    \proofstepwidestar[11]{y\in Y\land w=x\star y}{\rA}
    \proofstepwidestar[11]{y\in Y}{\rAEa{11}}
    \proofstepwidestar[11]{w=x\star y}{\rAEb{11}}
    \proofstepwidestar[1,2,3,5,8,11]{%
      z\in f[X]\diamond_{\mathcal P}f[Y]}{%
      \FormulaRefAuto{PowerHomForwardFixedWitnesses}[thm]{%
        1,2,3,9,12,13,5}}
    \proofstepwidestar[1,2,3,5,8]{%
      z\in f[X]\diamond_{\mathcal P}f[Y]}{\rEE{10,11,14}}
    \proofstepwidestar[1,2,3,4,5]{%
      z\in f[X]\diamond_{\mathcal P}f[Y]}{\rEE{7,8,15}}
  \closeproofpartwideindR

  \proofpartwideindR[Vorwärtsrichtung für Bildelemente]{%
    \begin{aligned}[t]
      &\SemigroupHom{f}{A,\star}{B,\diamond}\dsep X,Y\in\PowNE{A}
        \dsep z\in f[X\star_{\mathcal P}Y]\vdash{}\\[-2pt]
      &z\in f[X]\diamond_{\mathcal P}f[Y]
    \end{aligned}
  }{%
    \begin{aligned}[t]
      &\SemigroupHom{f}{A,\star}{B,\diamond}\dsep X,Y\in\PowNE{A}
        \dsep z\in f[X\star_{\mathcal P}Y]\vdash{}\\[-2pt]
      &z\in f[X]\diamond_{\mathcal P}f[Y]
    \end{aligned}
  }[PowerHomForwardElement]
    \proofstepwidestar[1]{\SemigroupHom{f}{A,\star}{B,\diamond}}{\rA}
    \proofstepwidestar[2]{X\in\PowNE{A}}{\rA}
    \proofstepwidestar[3]{Y\in\PowNE{A}}{\rA}
    \proofstepwidestar[4]{z\in f[X\star_{\mathcal P}Y]}{\rA}
    \proofstepwidestar[1]{\SemigroupAlg{A}{\star}}{%
      \FormulaRefAuto{SemigroupHomSourceAxiom}[ax]{1}}
    \proofstepwidestar[1]{f\colon A\to B}{%
      \FormulaRefAuto{SemigroupHomFunctionAxiom}[ax]{1}}
    \proofstepwidestar[1,2,3]{X\star_{\mathcal P}Y\subseteq A}{%
      \FormulaRefAuto{PowerSemigroupProductSubset}[thm]{5,2,3}}
    \proofstepwidestar[1,2,3,4]{%
      \exists w\in X\star_{\mathcal P}Y\,(z=f(w))}{%
      \FormulaRefAuto{C\subseteq A\dsep F\colon A\to B\dsep
        y\in F[C]\vdash\exists x\in C\,y=F(x)}[thm]{7,6,4}}
    \proofstepwidestar[9]{w\in X\star_{\mathcal P}Y\land z=f(w)}{\rA}
    \proofstepwidestar[9]{w\in X\star_{\mathcal P}Y}{\rAEa{9}}
    \proofstepwidestar[9]{z=f(w)}{\rAEb{9}}
    \proofstepwidestar[1,2,3,9]{z\in f[X]\diamond_{\mathcal P}f[Y]}{%
      \FormulaRefAuto{PowerHomForwardProductWitness}[thm]{1,2,3,10,11}}
    \proofstepwidestar[1,2,3,4]{z\in f[X]\diamond_{\mathcal P}f[Y]}{%
      \rEE{8,9,12}}
  \closeproofpartwideindR

  \proofpartwideindR[Erste Teilmengenrichtung]{%
    \begin{aligned}[t]
      &\SemigroupHom{f}{A,\star}{B,\diamond}\dsep X,Y\in\PowNE{A}
        \vdash{}\\[-2pt]
      &f[X\star_{\mathcal P}Y]
        \subseteq f[X]\diamond_{\mathcal P}f[Y]
    \end{aligned}
  }{%
    \begin{aligned}[t]
      &\SemigroupHom{f}{A,\star}{B,\diamond}\dsep X,Y\in\PowNE{A}
        \vdash{}\\[-2pt]
      &f[X\star_{\mathcal P}Y]
        \subseteq f[X]\diamond_{\mathcal P}f[Y]
    \end{aligned}
  }[PowerHomForwardSubset]
    \proofstepwidestar[1]{\SemigroupHom{f}{A,\star}{B,\diamond}}{\rA}
    \proofstepwidestar[2]{X\in\PowNE{A}}{\rA}
    \proofstepwidestar[3]{Y\in\PowNE{A}}{\rA}
    \proofstepwidestar[4]{z\in f[X\star_{\mathcal P}Y]}{\rA}
    \proofstepwidestar[1,2,3,4]{z\in f[X]\diamond_{\mathcal P}f[Y]}{%
      \FormulaRefAuto{PowerHomForwardElement}[thm]{1,2,3,4}}
    \proofstepwidestar[1,2,3]{%
      \forall z\bigl(z\in f[X\star_{\mathcal P}Y]\rightarrow
      z\in f[X]\diamond_{\mathcal P}f[Y]\bigr)}{\rUI{\rRI{4,5}}}
    \proofstepwidestar[1,2,3]{%
      f[X\star_{\mathcal P}Y]\subseteq
      f[X]\diamond_{\mathcal P}f[Y]}{%
      \FormulaRefAuto{A\subseteq B\coloneq
        \forall x\,(x\in A\rightarrow x\in B)}[def]{6}}
  \closeproofpartwideindR

  \proofpartwideindR[Rückwärtsrichtung der Gleichheit]{%
    \begin{aligned}[t]
      &\SemigroupHom{f}{A,\star}{B,\diamond}\dsep x,y\in A\dsep
        u=f(x)\dsep v=f(y)\dsep z=u\diamond v\vdash{}\\[-2pt]
      &z=f(x\star y)
    \end{aligned}
  }{%
    \begin{aligned}[t]
      &\SemigroupHom{f}{A,\star}{B,\diamond}\dsep x,y\in A\dsep
        u=f(x)\dsep v=f(y)\dsep z=u\diamond v\vdash{}\\[-2pt]
      &z=f(x\star y)
    \end{aligned}
  }[PowerHomBackwardEquality]
    \proofstepwidestar[1]{\SemigroupHom{f}{A,\star}{B,\diamond}}{\rA}
    \proofstepwidestar[2]{x\in A}{\rA}
    \proofstepwidestar[3]{y\in A}{\rA}
    \proofstepwidestar[4]{u=f(x)}{\rA}
    \proofstepwidestar[5]{v=f(y)}{\rA}
    \proofstepwidestar[6]{z=u\diamond v}{\rA}
    \proofstepwidestar[1,2,3]{f(x\star y)=f(x)\diamond f(y)}{%
      \FormulaRefAuto{SemigroupHomOperationAxiom}[ax]{1,2,3}}
    \proofstepwidestar[1,2,3]{f(x)\diamond f(y)=f(x\star y)}{%
      \FormulaRefAuto{a=b\vdash b=a}[thm]{7}}
    \proofstepwidestar[4,6]{z=f(x)\diamond v}{\rIE{4,6}}
    \proofstepwidestar[4,5,6]{z=f(x)\diamond f(y)}{\rIE{5,9}}
    \proofstepwidestar[1,2,3,4,5,6]{z=f(x\star y)}{\rChain{10,8}}
  \closeproofpartwideindR

  \proofpartwideindR[Rückwärtsrichtung mit festen Zeugen]{%
    \begin{aligned}[t]
      &\SemigroupHom{f}{A,\star}{B,\diamond}\dsep X,Y\in\PowNE{A}
        \dsep x\in X\dsep y\in Y\dsep{}\\[-2pt]
      &u=f(x)\dsep v=f(y)\dsep z=u\diamond v
        \vdash z\in f[X\star_{\mathcal P}Y]
    \end{aligned}
  }{%
    \begin{aligned}[t]
      &\SemigroupHom{f}{A,\star}{B,\diamond}\dsep X,Y\in\PowNE{A}
        \dsep x\in X\dsep y\in Y\dsep{}\\[-2pt]
      &u=f(x)\dsep v=f(y)\dsep z=u\diamond v
        \vdash z\in f[X\star_{\mathcal P}Y]
    \end{aligned}
  }[PowerHomBackwardFixedWitnesses]
    \proofstepwidestar[1]{\SemigroupHom{f}{A,\star}{B,\diamond}}{\rA}
    \proofstepwidestar[2]{X\in\PowNE{A}}{\rA}
    \proofstepwidestar[3]{Y\in\PowNE{A}}{\rA}
    \proofstepwidestar[4]{x\in X}{\rA}
    \proofstepwidestar[5]{y\in Y}{\rA}
    \proofstepwidestar[6]{u=f(x)}{\rA}
    \proofstepwidestar[7]{v=f(y)}{\rA}
    \proofstepwidestar[8]{z=u\diamond v}{\rA}
    \proofstepwidestar[1]{\SemigroupAlg{A}{\star}}{%
      \FormulaRefAuto{SemigroupHomSourceAxiom}[ax]{1}}
    \proofstepwidestar[1]{f\colon A\to B}{%
      \FormulaRefAuto{SemigroupHomFunctionAxiom}[ax]{1}}
    \proofstepwidestar[2,4]{x\in A}{%
      \FormulaRefAuto{PowerSemigroupCarrierElement}[thm]{2,4}}
    \proofstepwidestar[3,5]{y\in A}{%
      \FormulaRefAuto{PowerSemigroupCarrierElement}[thm]{3,5}}
    \proofstepwidestar[1,2,3,4,5]{x\star y\in X\star_{\mathcal P}Y}{%
      \FormulaRefAuto{PowerSemigroupProductContains}[thm]{9,2,3,4,5}}
    \proofstepwidestar[1,2,3]{X\star_{\mathcal P}Y\subseteq A}{%
      \FormulaRefAuto{PowerSemigroupProductSubset}[thm]{9,2,3}}
    \proofstepwidestar[1,2,3,4,5]{f(x\star y)\in f[X\star_{\mathcal P}Y]}{%
      \FormulaRefAuto{C\subseteq A\dsep F\colon A\to B\dsep
        x\in C\vdash F(x)\in F[C]}[thm]{14,10,13}}
    \proofstepwidestar[1,2,3,4,5,6,7,8]{z=f(x\star y)}{%
      \FormulaRefAuto{PowerHomBackwardEquality}[thm]{1,11,12,6,7,8}}
    \proofstepwidestar[1,2,3,4,5,6,7,8]{z\in f[X\star_{\mathcal P}Y]}{%
      \rIE{16,15}}
  \closeproofpartwideindR

  \proofpartwideindR[Rückwärtsrichtung mit rechtem Bildzeugen]{%
    \begin{aligned}[t]
      &\SemigroupHom{f}{A,\star}{B,\diamond}\dsep X,Y\in\PowNE{A}
        \dsep x\in X\dsep u=f(x)\dsep v\in f[Y]\dsep{}\\[-2pt]
      &z=u\diamond v\vdash z\in f[X\star_{\mathcal P}Y]
    \end{aligned}
  }{%
    \begin{aligned}[t]
      &\SemigroupHom{f}{A,\star}{B,\diamond}\dsep X,Y\in\PowNE{A}
        \dsep x\in X\dsep u=f(x)\dsep v\in f[Y]\dsep{}\\[-2pt]
      &z=u\diamond v\vdash z\in f[X\star_{\mathcal P}Y]
    \end{aligned}
  }[PowerHomBackwardRightWitness]
    \proofstepwidestar[1]{\SemigroupHom{f}{A,\star}{B,\diamond}}{\rA}
    \proofstepwidestar[2]{X\in\PowNE{A}}{\rA}
    \proofstepwidestar[3]{Y\in\PowNE{A}}{\rA}
    \proofstepwidestar[4]{x\in X}{\rA}
    \proofstepwidestar[5]{u=f(x)}{\rA}
    \proofstepwidestar[6]{v\in f[Y]}{\rA}
    \proofstepwidestar[7]{z=u\diamond v}{\rA}
    \proofstepwidestar[1]{f\colon A\to B}{%
      \FormulaRefAuto{SemigroupHomFunctionAxiom}[ax]{1}}
    \proofstepwidestar[3]{Y\subseteq A}{%
      \rAEa{\FormulaRefAuto{NonemptyPowersetMembership}[thm]{3}}}
    \proofstepwidestar[1,3,6]{\exists y\in Y\,(v=f(y))}{%
      \FormulaRefAuto{C\subseteq A\dsep F\colon A\to B\dsep
        y\in F[C]\vdash\exists x\in C\,y=F(x)}[thm]{9,8,6}}
    \proofstepwidestar[11]{y\in Y\land v=f(y)}{\rA}
    \proofstepwidestar[11]{y\in Y}{\rAEa{11}}
    \proofstepwidestar[11]{v=f(y)}{\rAEb{11}}
    \proofstepwidestar[1,2,3,4,5,7,11]{z\in f[X\star_{\mathcal P}Y]}{%
      \FormulaRefAuto{PowerHomBackwardFixedWitnesses}[thm]{%
        1,2,3,4,12,5,13,7}}
    \proofstepwidestar[1,2,3,4,5,6,7]{z\in f[X\star_{\mathcal P}Y]}{%
      \rEE{10,11,14}}
  \closeproofpartwideindR

  \proofpartwideindR[Rückwärtsrichtung mit zwei Bildzeugen]{%
    \begin{aligned}[t]
      &\SemigroupHom{f}{A,\star}{B,\diamond}\dsep X,Y\in\PowNE{A}
        \dsep u\in f[X]\dsep v\in f[Y]\dsep z=u\diamond v\vdash{}\\[-2pt]
      &z\in f[X\star_{\mathcal P}Y]
    \end{aligned}
  }{%
    \begin{aligned}[t]
      &\SemigroupHom{f}{A,\star}{B,\diamond}\dsep X,Y\in\PowNE{A}
        \dsep u\in f[X]\dsep v\in f[Y]\dsep z=u\diamond v\vdash{}\\[-2pt]
      &z\in f[X\star_{\mathcal P}Y]
    \end{aligned}
  }[PowerHomBackwardImageWitnesses]
    \proofstepwidestar[1]{\SemigroupHom{f}{A,\star}{B,\diamond}}{\rA}
    \proofstepwidestar[2]{X\in\PowNE{A}}{\rA}
    \proofstepwidestar[3]{Y\in\PowNE{A}}{\rA}
    \proofstepwidestar[4]{u\in f[X]}{\rA}
    \proofstepwidestar[5]{v\in f[Y]}{\rA}
    \proofstepwidestar[6]{z=u\diamond v}{\rA}
    \proofstepwidestar[1]{f\colon A\to B}{%
      \FormulaRefAuto{SemigroupHomFunctionAxiom}[ax]{1}}
    \proofstepwidestar[2]{X\subseteq A}{%
      \rAEa{\FormulaRefAuto{NonemptyPowersetMembership}[thm]{2}}}
    \proofstepwidestar[1,2,4]{\exists x\in X\,(u=f(x))}{%
      \FormulaRefAuto{C\subseteq A\dsep F\colon A\to B\dsep
        y\in F[C]\vdash\exists x\in C\,y=F(x)}[thm]{8,7,4}}
    \proofstepwidestar[10]{x\in X\land u=f(x)}{\rA}
    \proofstepwidestar[10]{x\in X}{\rAEa{10}}
    \proofstepwidestar[10]{u=f(x)}{\rAEb{10}}
    \proofstepwidestar[1,2,3,5,6,10]{z\in f[X\star_{\mathcal P}Y]}{%
      \FormulaRefAuto{PowerHomBackwardRightWitness}[thm]{%
        1,2,3,11,12,5,6}}
    \proofstepwidestar[1,2,3,4,5,6]{z\in f[X\star_{\mathcal P}Y]}{%
      \rEE{9,10,13}}
  \closeproofpartwideindR

  \proofpartwideindR[Auflösung des zweiten Produktzeugen]{%
    \begin{aligned}[t]
      &\SemigroupHom{f}{A,\star}{B,\diamond}\dsep X,Y\in\PowNE{A}
        \dsep u\in f[X]\dsep{}\\[-2pt]
      &\exists v\in f[Y]\,(z=u\diamond v)
        \vdash z\in f[X\star_{\mathcal P}Y]
    \end{aligned}
  }{%
    \begin{aligned}[t]
      &\SemigroupHom{f}{A,\star}{B,\diamond}\dsep X,Y\in\PowNE{A}
        \dsep u\in f[X]\dsep{}\\[-2pt]
      &\exists v\in f[Y]\,(z=u\diamond v)
        \vdash z\in f[X\star_{\mathcal P}Y]
    \end{aligned}
  }[PowerHomBackwardSecondProductWitness]
    \proofstepwidestar[1]{\SemigroupHom{f}{A,\star}{B,\diamond}}{\rA}
    \proofstepwidestar[2]{X\in\PowNE{A}}{\rA}
    \proofstepwidestar[3]{Y\in\PowNE{A}}{\rA}
    \proofstepwidestar[4]{u\in f[X]}{\rA}
    \proofstepwidestar[5]{\exists v\in f[Y]\,(z=u\diamond v)}{\rA}
    \proofstepwidestar[6]{v\in f[Y]\land z=u\diamond v}{\rA}
    \proofstepwidestar[6]{v\in f[Y]}{\rAEa{6}}
    \proofstepwidestar[6]{z=u\diamond v}{\rAEb{6}}
    \proofstepwidestar[1,2,3,4,6]{z\in f[X\star_{\mathcal P}Y]}{%
      \FormulaRefAuto{PowerHomBackwardImageWitnesses}[thm]{1,2,3,4,7,8}}
    \proofstepwidestar[1,2,3,4,5]{z\in f[X\star_{\mathcal P}Y]}{%
      \rEE{5,6,9}}
  \closeproofpartwideindR

  \proofpartwideindR[Rückwärtsrichtung für Produktelemente]{%
    \begin{aligned}[t]
      &\SemigroupHom{f}{A,\star}{B,\diamond}\dsep X,Y\in\PowNE{A}
        \dsep z\in f[X]\diamond_{\mathcal P}f[Y]\vdash{}\\[-2pt]
      &z\in f[X\star_{\mathcal P}Y]
    \end{aligned}
  }{%
    \begin{aligned}[t]
      &\SemigroupHom{f}{A,\star}{B,\diamond}\dsep X,Y\in\PowNE{A}
        \dsep z\in f[X]\diamond_{\mathcal P}f[Y]\vdash{}\\[-2pt]
      &z\in f[X\star_{\mathcal P}Y]
    \end{aligned}
  }[PowerHomBackwardElement]
    \proofstepwidestar[1]{\SemigroupHom{f}{A,\star}{B,\diamond}}{\rA}
    \proofstepwidestar[2]{X\in\PowNE{A}}{\rA}
    \proofstepwidestar[3]{Y\in\PowNE{A}}{\rA}
    \proofstepwidestar[4]{z\in f[X]\diamond_{\mathcal P}f[Y]}{\rA}
    \proofstepwidestar[1]{\SemigroupAlg{B}{\diamond}}{%
      \FormulaRefAuto{SemigroupHomTargetAxiom}[ax]{1}}
    \proofstepwidestar[1]{f\colon A\to B}{%
      \FormulaRefAuto{SemigroupHomFunctionAxiom}[ax]{1}}
    \proofstepwidestar[1,2]{f[X]\in\PowNE{B}}{%
      \FormulaRefAuto{NonemptyDirectImage}[thm]{6,2}}
    \proofstepwidestar[1,3]{f[Y]\in\PowNE{B}}{%
      \FormulaRefAuto{NonemptyDirectImage}[thm]{6,3}}
    \proofstepwidestar[1,2,3,4]{%
      \exists u\in f[X]\,\exists v\in f[Y]\,(z=u\diamond v)}{%
      \FormulaRefAuto{PowerSemigroupProductWitnesses}[thm]{5,7,8,4}}
    \proofstepwidestar[10]{%
      u\in f[X]\land\exists v\in f[Y]\,(z=u\diamond v)}{\rA}
    \proofstepwidestar[10]{u\in f[X]}{\rAEa{10}}
    \proofstepwidestar[10]{\exists v\in f[Y]\,(z=u\diamond v)}{\rAEb{10}}
    \proofstepwidestar[1,2,3,10]{z\in f[X\star_{\mathcal P}Y]}{%
      \FormulaRefAuto{PowerHomBackwardSecondProductWitness}[thm]{%
        1,2,3,11,12}}
    \proofstepwidestar[1,2,3,4]{z\in f[X\star_{\mathcal P}Y]}{%
      \rEE{9,10,13}}
  \closeproofpartwideindR

  \proofpartwideindR[Zweite Teilmengenrichtung]{%
    \begin{aligned}[t]
      &\SemigroupHom{f}{A,\star}{B,\diamond}\dsep X,Y\in\PowNE{A}
        \vdash{}\\[-2pt]
      &f[X]\diamond_{\mathcal P}f[Y]
        \subseteq f[X\star_{\mathcal P}Y]
    \end{aligned}
  }{%
    \begin{aligned}[t]
      &\SemigroupHom{f}{A,\star}{B,\diamond}\dsep X,Y\in\PowNE{A}
        \vdash{}\\[-2pt]
      &f[X]\diamond_{\mathcal P}f[Y]
        \subseteq f[X\star_{\mathcal P}Y]
    \end{aligned}
  }[PowerHomBackwardSubset]
    \proofstepwidestar[1]{\SemigroupHom{f}{A,\star}{B,\diamond}}{\rA}
    \proofstepwidestar[2]{X\in\PowNE{A}}{\rA}
    \proofstepwidestar[3]{Y\in\PowNE{A}}{\rA}
    \proofstepwidestar[4]{z\in f[X]\diamond_{\mathcal P}f[Y]}{\rA}
    \proofstepwidestar[1,2,3,4]{z\in f[X\star_{\mathcal P}Y]}{%
      \FormulaRefAuto{PowerHomBackwardElement}[thm]{1,2,3,4}}
    \proofstepwidestar[1,2,3]{%
      \forall z\bigl(z\in f[X]\diamond_{\mathcal P}f[Y]\rightarrow
      z\in f[X\star_{\mathcal P}Y]\bigr)}{\rUI{\rRI{4,5}}}
    \proofstepwidestar[1,2,3]{%
      f[X]\diamond_{\mathcal P}f[Y]\subseteq
      f[X\star_{\mathcal P}Y]}{%
      \FormulaRefAuto{A\subseteq B\coloneq
        \forall x\,(x\in A\rightarrow x\in B)}[def]{6}}
  \closeproofpartwideindR

  \proofpartwide{Schluss}
    \proofstepwidestar[1]{\SemigroupHom{f}{A,\star}{B,\diamond}}{\rA}
    \proofstepwidestar[2]{X\in\PowNE{A}}{\rA}
    \proofstepwidestar[3]{Y\in\PowNE{A}}{\rA}
    \proofstepwidestar[1,2,3]{%
      f[X\star_{\mathcal P}Y]\subseteq
      f[X]\diamond_{\mathcal P}f[Y]}{%
      \FormulaRefAuto{PowerHomForwardSubset}[thm]{1,2,3}}
    \proofstepwidestar[1,2,3]{%
      f[X]\diamond_{\mathcal P}f[Y]\subseteq
      f[X\star_{\mathcal P}Y]}{%
      \FormulaRefAuto{PowerHomBackwardSubset}[thm]{1,2,3}}
    \proofstepwidestar[1,2,3]{%
      f[X\star_{\mathcal P}Y]=f[X]\diamond_{\mathcal P}f[Y]}{%
      \FormulaRefAuto{A\subseteq B, B\subseteq A\vdash A=B}[thm]{4,5}}
  \closeproofpartwide
\end{tabproofsplitwide}

\FormulaThmDeltaK[Operationserhaltung der induzierten Potenzabbildung]{%
  \begin{aligned}[t]
    &\SemigroupHom{f}{A,\star}{B,\diamond}\dsep
      X,Y\in\PowNE{A}\vdash{}\\[-2pt]
    &\PowMapNE{f}(X\star_{\mathcal P}Y)
      =\PowMapNE{f}(X)\diamond_{\mathcal P}\PowMapNE{f}(Y)
  \end{aligned}%
}{PowerSemigroupInducedMapOperation}{%
  \DeltaRow{Mengen}{A\dsep B\dsep X\dsep Y}
  \DeltaRow{Funktionen}{f\dsep\star\dsep\diamond\dsep
    \star_{\mathcal P}\dsep\diamond_{\mathcal P}\dsep
    \PowMap{f}\dsep\PowMapNE{f}}
}
\begin{tabproofwide}
  \proofstepwidestar[1]{\SemigroupHom{f}{A,\star}{B,\diamond}}{\rA}
  \proofstepwidestar[2]{X\in\PowNE{A}}{\rA}
  \proofstepwidestar[3]{Y\in\PowNE{A}}{\rA}
  \proofstepwidestar[1]{\SemigroupAlg{A}{\star}}{%
    \FormulaRefAuto{SemigroupHomSourceAxiom}[ax]{1}}
  \proofstepwidestar[1]{f\colon A\to B}{%
    \FormulaRefAuto{SemigroupHomFunctionAxiom}[ax]{1}}
  \proofstepwidestar[1,2,3]{X\star_{\mathcal P}Y\in\PowNE{A}}{%
    \FormulaRefAuto{PowerSemigroupProductClosure}[thm]{4,2,3}}
  \proofstepwidestar[1,2]{%
    \PowMapNE{f}(X)=f[X]}{%
    \FormulaRefAuto{NonemptyPowerMapValue}[thm]{5,2}}
  \proofstepwidestar[1,2]{%
    f[X]=\PowMapNE{f}(X)}{%
    \FormulaRefAuto{a=b\vdash b=a}[thm]{7}}
  \proofstepwidestar[1,3]{%
    \PowMapNE{f}(Y)=f[Y]}{%
    \FormulaRefAuto{NonemptyPowerMapValue}[thm]{5,3}}
  \proofstepwidestar[1,3]{%
    f[Y]=\PowMapNE{f}(Y)}{%
    \FormulaRefAuto{a=b\vdash b=a}[thm]{9}}
  \proofstepwide[1,2,3]{%
    \PowMapNE{f}(X\star_{\mathcal P}Y)}{=}{f[X\star_{\mathcal P}Y]}{%
    \FormulaRefAuto{NonemptyPowerMapValue}[thm]{5,6}}
  \proofstepwide[1,2,3]{}{=}{f[X]\diamond_{\mathcal P}f[Y]}{%
    \FormulaRefAuto{SemigroupHomPreservesPowerProduct}[thm]{1,2,3}}
  \proofstepwide[1,2,3]{}{=}{%
    \PowMapNE{f}(X)\diamond_{\mathcal P}f[Y]}{\rLRS{8}}
  \proofstepwide[1,2,3]{}{=}{%
    \PowMapNE{f}(X)\diamond_{\mathcal P}\PowMapNE{f}(Y)}{\rLRS{10}}
  \proofstepwide[1,2,3]{%
    \PowMapNE{f}(X\star_{\mathcal P}Y)}{=}{%
    \PowMapNE{f}(X)\diamond_{\mathcal P}\PowMapNE{f}(Y)}{\rChain{11,14}}
\end{tabproofwide}

\FormulaThmDeltaK[Isomorphismen induzieren Potenzhalbgruppenisomorphismen]{%
  \begin{aligned}[t]
    &\SemigroupIso{f}{A,\star}{B,\diamond}\vdash{}\\[-2pt]
    &\SemigroupIso{\PowMapNE{f}}
      {\PowNE{A},\star_{\mathcal P}}
      {\PowNE{B},\diamond_{\mathcal P}}
  \end{aligned}%
}{SemigroupIsoInducesPowerSemigroupIso}{%
  \DeltaRow{Mengen}{A\dsep B\dsep X\dsep Y}
  \DeltaRow{Funktionen}{f\dsep\star\dsep\diamond\dsep
    \star_{\mathcal P}\dsep\diamond_{\mathcal P}\dsep
    \PowMap{f}\dsep\PowMapNE{f}}
}
\begin{tabproofwide}
  \proofstepwidestar[1]{\SemigroupIso{f}{A,\star}{B,\diamond}}{\rA}
  \proofstepwidestar[1]{\SemigroupHom{f}{A,\star}{B,\diamond}}{%
    \FormulaRefAuto{SemigroupIsoHomAxiom}[ax]{1}}
  \proofstepwidestar[1]{f\colon A\bij B}{%
    \FormulaRefAuto{SemigroupIsoBijectiveAxiom}[ax]{1}}
  \proofstepwidestar[1]{\SemigroupAlg{A}{\star}}{%
    \FormulaRefAuto{SemigroupHomSourceAxiom}[ax]{2}}
  \proofstepwidestar[1]{\SemigroupAlg{B}{\diamond}}{%
    \FormulaRefAuto{SemigroupHomTargetAxiom}[ax]{2}}
  \proofstepwidestar[1]{\SemigroupAlg{\PowNE{A}}{\star_{\mathcal P}}}{%
    \FormulaRefAuto{NonemptyPowerSemigroup}[thm]{4}}
  \proofstepwidestar[1]{\SemigroupAlg{\PowNE{B}}{\diamond_{\mathcal P}}}{%
    \FormulaRefAuto{NonemptyPowerSemigroup}[thm]{5}}
  \proofstepwidestar[1]{%
    \PowMapNE{f}\colon\PowNE{A}\bij\PowNE{B}}{%
    \FormulaRefAuto{NonemptyPowerMapBijective}[thm]{3}}
  \proofstepwidestar[9]{X\in\PowNE{A}}{\rA}
  \proofstepwidestar[10]{Y\in\PowNE{A}}{\rA}
  \proofstepwidestar[1,9,10]{%
    \PowMapNE{f}(X\star_{\mathcal P}Y)
      =\PowMapNE{f}(X)\diamond_{\mathcal P}\PowMapNE{f}(Y)}{%
    \FormulaRefAuto{PowerSemigroupInducedMapOperation}[thm]{2,9,10}}
  \proofstepwidestar[1]{%
    \PowMapNE{f}\colon\PowNE{A}\to\PowNE{B}}{%
    \FormulaRefAuto{Bijektive Funktion}[def]{8}}
  \proofstepwidestar[1]{%
    \SemigroupHom{\PowMapNE{f}}
      {\PowNE{A},\star_{\mathcal P}}
      {\PowNE{B},\diamond_{\mathcal P}}}{%
    \FormulaRefAuto{SemigroupHomDef}[def]{6,7,12,11}}
  \proofstepwidestar[1]{%
    \SemigroupIso{\PowMapNE{f}}
      {\PowNE{A},\star_{\mathcal P}}
      {\PowNE{B},\diamond_{\mathcal P}}}{%
    \FormulaRefAuto{SemigroupIsoDef}[def]{13,8}}
\end{tabproofwide}

\section{Das Isomorphieproblem für Potenzhalbgruppen}

Zwei Halbgruppen heißen \emph{global isomorph}, wenn ihre
Potenzhalbgruppen isomorph sind.  Die folgende Richtung gilt ohne jede
Endlichkeitsvoraussetzung.

\FormulaThmDeltaK[Die isomorphe Richtung des Potenzhalbgruppenproblems]{%
  \begin{aligned}[t]
    &\exists f\,\SemigroupIso{f}{A,\star}{B,\diamond}\vdash{}\\[-2pt]
    &\exists F\,\SemigroupIso{F}
      {\PowNE{A},\star_{\mathcal P}}
      {\PowNE{B},\diamond_{\mathcal P}}
  \end{aligned}%
}{SemigroupIsoImpliesPowerSemigroupIso}{%
  \DeltaRow{Mengen}{A\dsep B}
  \DeltaRow{Funktionen}{F\dsep f\dsep\star\dsep\diamond\dsep
    \star_{\mathcal P}\dsep\diamond_{\mathcal P}\dsep
    \PowMap{f}\dsep\PowMapNE{f}}
}
\begin{tabproofwide}
  \proofstepwidestar[1]{\exists f\,\SemigroupIso{f}{A,\star}{B,\diamond}}{\rA}
  \proofstepwidestar[2]{\SemigroupIso{f}{A,\star}{B,\diamond}}{\rA}
  \proofstepwidestar[2]{%
    \SemigroupIso{\PowMapNE{f}}
      {\PowNE{A},\star_{\mathcal P}}
      {\PowNE{B},\diamond_{\mathcal P}}}{%
    \FormulaRefAuto{SemigroupIsoInducesPowerSemigroupIso}[thm]{2}}
  \proofstepwidestar[2]{%
    \exists F\,\SemigroupIso{F}
      {\PowNE{A},\star_{\mathcal P}}
      {\PowNE{B},\diamond_{\mathcal P}}}{\rEI{3}}
  \proofstepwidestar[1]{%
    \exists F\,\SemigroupIso{F}
      {\PowNE{A},\star_{\mathcal P}}
      {\PowNE{B},\diamond_{\mathcal P}}}{\rEE{1,2,4}}
\end{tabproofwide}

\subsection{Inflationen und translationelle Ununterscheidbarkeit}

Eine Erweiterung kann neue Elemente enthalten, ohne neue Multiplikationswerte
zu erzeugen.  Der folgende Begriff fasst diese Situation zusammen.  Die
Abbildung \(r\) zieht jedes neue Element auf ein altes Element zurück; die
Multiplikation der Erweiterung hängt nur von diesen Rückzugswerten ab.

\FormulaDefDeltaK[Inflation einer Halbgruppe]{%
  \mathsf{Infl}_{r}\bigl((A,\star),(B,\diamond)\bigr)%
}{SemigroupInflationDef}{%
  \DeltaRow{Mengen}{A\dsep B\dsep x\dsep y\dsep a}
  \DeltaRow{Funktionen}{r\dsep\star\dsep\diamond}
  \DeltaRow{\textbf{Axiome}}{}
  \DeltaPrem{Quellhalbgruppe}{\SemigroupAlg{A}{\star}}
  \DeltaPrem{Zielhalbgruppe}{\SemigroupAlg{B}{\diamond}}
  \DeltaPrem{Erweiterung}{A\subseteq B}
  \DeltaPrem{Rückzug}{r\colon B\to A}
  \DeltaRow{Retraktion}{a\in A\vdash r(a)=a}
  \DeltaRow{Produkt über dem Rückzug}{%
    x,y\in B\vdash x\diamond y=r(x)\star r(y)}
  \DeltaRow{\textbf{Neues Symbol}}{}
  \DeltaPrem{Inflation}{%
    \mathsf{Infl}_{r}\bigl((A,\star),(B,\diamond)\bigr)}
}

Für die nachfolgende Anwendung ist nicht die Gleichheit zweier Elemente
entscheidend, sondern die Gleichheit ihrer Links- und Rechtstranslationen.

\FormulaDefDeltaK[Translationelle Ununterscheidbarkeit]{%
  \begin{aligned}[t]
    &\SemigroupAlg{B}{\diamond}\dsep x,y\in B\vdash{}\\[-2pt]
    &x\mathrel{\equiv^{\mathrm{tr}}_{B,\diamond}}y
      \coloneqq
      \forall z\in B\,\bigl(
        x\diamond z=y\diamond z\land
        z\diamond x=z\diamond y\bigr)
  \end{aligned}%
}{SemigroupTranslationIndistinguishableDef}{%
  \DeltaRow{Mengen}{B\dsep x\dsep y\dsep z}
  \DeltaRow{Funktionen}{\diamond}
  \DeltaRow{Relationsoperatoren}{%
    \equiv^{\mathrm{tr}}_{B,\diamond}}
}

\FormulaThmDeltaK[Zugriff auf die gemeinsamen Translationen]{%
  \begin{aligned}[t]
    &\SemigroupAlg{B}{\diamond}\dsep
      x,y\in B\dsep
      x\mathrel{\equiv^{\mathrm{tr}}_{B,\diamond}}y\dsep z\in B
      \vdash{}\\[-2pt]
    &x\diamond z=y\diamond z\land z\diamond x=z\diamond y
  \end{aligned}%
}{SemigroupTranslationIndistinguishableActions}{%
  \DeltaRow{Mengen}{B\dsep x\dsep y\dsep z}
  \DeltaRow{Funktionen}{\diamond}
}
\begin{tabproofwide}
  \proofstepwidestar[1]{\SemigroupAlg{B}{\diamond}}{\rA}
  \proofstepwidestar[2]{x\in B}{\rA}
  \proofstepwidestar[3]{y\in B}{\rA}
  \proofstepwidestar[4]{%
    x\mathrel{\equiv^{\mathrm{tr}}_{B,\diamond}}y}{\rA}
  \proofstepwidestar[5]{z\in B}{\rA}
  \proofstepwidestar[1,2,3,4]{%
    \forall u\in B\,\bigl(
      x\diamond u=y\diamond u\land
      u\diamond x=u\diamond y\bigr)}{%
    \FormulaRefAuto{SemigroupTranslationIndistinguishableDef}[def]{1,2,3,4}}
  \proofstepwidestar[1,2,3,4,5]{%
    x\diamond z=y\diamond z\land z\diamond x=z\diamond y}{%
    \rRE{5,\rUE{6}}}
\end{tabproofwide}

Zunächst halten wir die für die Anwendung benötigte Reflexivität des
Translationsprofils fest.

\FormulaThmDeltaK[Reflexivität der translationellen Ununterscheidbarkeit]{%
  \SemigroupAlg{B}{\diamond}\dsep x\in B\vdash
  x\mathrel{\equiv^{\mathrm{tr}}_{B,\diamond}}x%
}{SemigroupTranslationIndistinguishableReflexive}{%
  \DeltaRow{Mengen}{B\dsep x\dsep z}
  \DeltaRow{Funktionen}{\diamond}
  \DeltaRow{Relationsoperatoren}{\equiv^{\mathrm{tr}}_{B,\diamond}}
}
\begin{tabproofwide}
  \proofstepwidestar[1]{\SemigroupAlg{B}{\diamond}}{\rA}
  \proofstepwidestar[2]{x\in B}{\rA}
  \proofstepwidestar[3]{z\in B}{\rA}
  \proofstepwidestar[3]{x\diamond z=x\diamond z}{\rII}
  \proofstepwidestar[3]{z\diamond x=z\diamond x}{\rII}
  \proofstepwidestar[3]{%
    x\diamond z=x\diamond z\land z\diamond x=z\diamond x}{\rAI{4,5}}
  \proofstepwidestar[]{%
    \forall z\in B\,\bigl(
      x\diamond z=x\diamond z\land z\diamond x=z\diamond x\bigr)}{%
    \rUI{\rRI{3,6}}}
  \proofstepwidestar[1,2]{%
    x\mathrel{\equiv^{\mathrm{tr}}_{B,\diamond}}x}{%
    \FormulaRefAuto{SemigroupTranslationIndistinguishableDef}[def]{1,2,2,7}}
\end{tabproofwide}

Der folgende Satz ist das algebraische Bindeglied zur späteren
Reservefamilie.  Er verlangt nicht, dass die Bijektion die einzelnen
Elemente festhält.  Es genügt, sie nur innerhalb ihrer Translationsprofile
zu verschieben und sämtliche wirklichen Produktwerte festzulassen.

\FormulaThmDeltaK[Isomorphiekriterium durch Fixierung der Produktwerte]{%
  \begin{aligned}[t]
    &\SemigroupAlg{A}{\star}\dsep\SemigroupAlg{B}{\diamond}\dsep
      A\subseteq B\dsep F\colon A\bij B\dsep{}\\[-2pt]
    &\forall x,y\in A\,(x\star y=x\diamond y)\dsep{}\\[-2pt]
    &\forall x\in A\,
      \bigl(x\mathrel{\equiv^{\mathrm{tr}}_{B,\diamond}}F(x)\bigr)
      \dsep{}\\[-2pt]
    &\forall x,y\in A\,F(x\star y)=x\star y
      \vdash \SemigroupIso{F}{A,\star}{B,\diamond}
  \end{aligned}%
}{SemigroupFixedProductsIsomorphismCriterion}{%
  \DeltaRow{Mengen}{A\dsep B\dsep x\dsep y}
  \DeltaRow{Funktionen}{F\dsep\star\dsep\diamond}
}
\begin{tabproofwide}
  \proofstepwidestar[1]{\SemigroupAlg{A}{\star}}{\rA}
  \proofstepwidestar[2]{\SemigroupAlg{B}{\diamond}}{\rA}
  \proofstepwidestar[3]{A\subseteq B}{\rA}
  \proofstepwidestar[4]{F\colon A\bij B}{\rA}
  \proofstepwidestar[5]{\forall u,v\in A\,(u\star v=u\diamond v)}{\rA}
  \proofstepwidestar[6]{\forall u\in A\,
    \bigl(u\mathrel{\equiv^{\mathrm{tr}}_{B,\diamond}}F(u)\bigr)}{\rA}
  \proofstepwidestar[7]{\forall u,v\in A\,F(u\star v)=u\star v}{\rA}
  \proofstepwidestar[4]{F\colon A\to B}{%
    \FormulaRefAuto{Bijektive Funktion}[def]{4}}
  \proofstepwidestar[9]{x\in A}{\rA}
  \proofstepwidestar[10]{y\in A}{\rA}
  \proofstepwidestar[3,9]{x\in B}{%
    \FormulaRefAuto{A\subseteq B\dsep x\in A\vdash x\in B}[thm]{3,9}}
  \proofstepwidestar[3,10]{y\in B}{%
    \FormulaRefAuto{A\subseteq B\dsep x\in A\vdash x\in B}[thm]{3,10}}
  \proofstepwidestar[4,9]{F(x)\in B}{%
    \FormulaRefAuto{F\colon A\to B\dsep x\in A\vdash F(x)\in B}[thm]{8,9}}
  \proofstepwidestar[4,10]{F(y)\in B}{%
    \FormulaRefAuto{F\colon A\to B\dsep x\in A\vdash F(x)\in B}[thm]{8,10}}
  \proofstepwidestar[6,9]{%
    x\mathrel{\equiv^{\mathrm{tr}}_{B,\diamond}}F(x)}{\rRE{9,\rUE{6}}}
  \proofstepwidestar[6,10]{%
    y\mathrel{\equiv^{\mathrm{tr}}_{B,\diamond}}F(y)}{\rRE{10,\rUE{6}}}
  \proofstepwidestar[2,3,4,6,9,10]{%
    x\diamond F(y)=F(x)\diamond F(y)}{%
    \rAEa{\FormulaRefAuto{SemigroupTranslationIndistinguishableActions}[thm]{%
      2,11,13,15,14}}}
  \proofstepwidestar[2,3,4,6,9,10]{%
    F(x)\diamond F(y)=x\diamond F(y)}{%
    \FormulaRefAuto{a=b\vdash b=a}[thm]{17}}
  \proofstepwidestar[2,3,4,6,9,10]{x\diamond y=x\diamond F(y)}{%
    \rAEb{\FormulaRefAuto{SemigroupTranslationIndistinguishableActions}[thm]{%
      2,12,14,16,11}}}
  \proofstepwidestar[2,3,4,6,9,10]{x\diamond F(y)=x\diamond y}{%
    \FormulaRefAuto{a=b\vdash b=a}[thm]{19}}
  \proofstepwidestar[5,9,10]{x\star y=x\diamond y}{%
    \rRE{10,\rUE{\rRE{9,\rUE{5}}}}}
  \proofstepwidestar[5,9,10]{x\diamond y=x\star y}{%
    \FormulaRefAuto{a=b\vdash b=a}[thm]{21}}
  \proofstepwidestar[7,9,10]{F(x\star y)=x\star y}{%
    \rRE{10,\rUE{\rRE{9,\rUE{7}}}}}
  \proofstepwidestar[7,9,10]{x\star y=F(x\star y)}{%
    \FormulaRefAuto{a=b\vdash b=a}[thm]{23}}
  \proofstepwidestar[2,3,4,5,6,7,9,10]{%
    F(x\star y)=F(x)\diamond F(y)}{%
    \rChain{23,21,19,17}}
  \proofstepwidestar[1,2,3,4,5,6,7]{%
    \SemigroupHom{F}{A,\star}{B,\diamond}}{%
    \FormulaRefAuto{SemigroupHomDef}[def]{1,2,8,25}}
  \proofstepwidestar[1,2,3,4,5,6,7]{%
    \SemigroupIso{F}{A,\star}{B,\diamond}}{%
    \FormulaRefAuto{SemigroupIsoDef}[def]{26,4}}
\end{tabproofwide}

\FormulaThmDeltaK[Gleiche Translationsprofile in der Potenzhalbgruppe]{%
  \begin{aligned}[t]
    &\SemigroupAlg{B}{\diamond}\dsep X,Y\in\PowNE{B}\dsep{}\\[-2pt]
    &\forall x\in X\,\exists y\in Y\,
      \bigl(x\mathrel{\equiv^{\mathrm{tr}}_{B,\diamond}}y\bigr)
      \dsep{}\\[-2pt]
    &\forall y\in Y\,\exists x\in X\,
      \bigl(y\mathrel{\equiv^{\mathrm{tr}}_{B,\diamond}}x\bigr)
      \vdash
      X\mathrel{\equiv^{\mathrm{tr}}_{\PowNE{B},
        \diamond_{\mathcal P}}}Y
  \end{aligned}%
}{PowerSemigroupTranslationSupportCriterion}{%
  \DeltaRow{Mengen}{B\dsep X\dsep Y\dsep Z\dsep
    x\dsep y\dsep z\dsep w}
  \DeltaRow{Funktionen}{\diamond\dsep\diamond_{\mathcal P}}
}
\begin{tabproofsplitwide}
  \proofpartwideindR[Rechte Mengenprodukte]{%
    \begin{aligned}[t]
      &\SemigroupAlg{B}{\diamond}\dsep X,Y,Z\in\PowNE{B}\dsep{}\\[-2pt]
      &\forall x\in X\,\exists y\in Y\,
        \bigl(x\mathrel{\equiv^{\mathrm{tr}}_{B,\diamond}}y\bigr)
        \dsep{}\\[-2pt]
      &\forall y\in Y\,\exists x\in X\,
        \bigl(y\mathrel{\equiv^{\mathrm{tr}}_{B,\diamond}}x\bigr)
        \vdash X\diamond_{\mathcal P}Z=Y\diamond_{\mathcal P}Z
    \end{aligned}%
  }{%
    \begin{aligned}[t]
      &\SemigroupAlg{B}{\diamond}\dsep X,Y,Z\in\PowNE{B}\dsep{}\\[-2pt]
      &\forall x\in X\,\exists y\in Y\,
        \bigl(x\mathrel{\equiv^{\mathrm{tr}}_{B,\diamond}}y\bigr)
        \dsep{}\\[-2pt]
      &\forall y\in Y\,\exists x\in X\,
        \bigl(y\mathrel{\equiv^{\mathrm{tr}}_{B,\diamond}}x\bigr)
        \vdash X\diamond_{\mathcal P}Z=Y\diamond_{\mathcal P}Z
    \end{aligned}%
  }[PowerSemigroupTranslationSupportRightProduct]
    \proofstepwidestar[1]{\SemigroupAlg{B}{\diamond}}{\rA}
    \proofstepwidestar[2]{X\in\PowNE{B}}{\rA}
    \proofstepwidestar[3]{Y\in\PowNE{B}}{\rA}
    \proofstepwidestar[4]{Z\in\PowNE{B}}{\rA}
    \proofstepwidestar[5]{\forall u\in X\,\exists v\in Y\,
      \bigl(u\mathrel{\equiv^{\mathrm{tr}}_{B,\diamond}}v\bigr)}{\rA}
    \proofstepwidestar[6]{\forall v\in Y\,\exists u\in X\,
      \bigl(v\mathrel{\equiv^{\mathrm{tr}}_{B,\diamond}}u\bigr)}{\rA}

    \proofstepwidestar[7]{w\in X\diamond_{\mathcal P}Z}{\rA}
    \proofstepwidestar[1,2,4,7]{%
      \exists x\in X\,\exists z\in Z\,(w=x\diamond z)}{%
      \FormulaRefAuto{PowerSemigroupProductWitnesses}[thm]{1,2,4,7}}
    \proofstepwidestar[8]{x\in X}{\rA}
    \proofstepwidestar[8]{z\in Z}{\rA}
    \proofstepwidestar[8]{w=x\diamond z}{\rA}
    \proofstepwidestar[5,8]{\exists y\in Y\,
      \bigl(x\mathrel{\equiv^{\mathrm{tr}}_{B,\diamond}}y\bigr)}{%
      \rRE{9,\rUE{5}}}
    \proofstepwidestar[12]{y\in Y}{\rA}
    \proofstepwidestar[12]{%
      x\mathrel{\equiv^{\mathrm{tr}}_{B,\diamond}}y}{\rA}
    \proofstepwidestar[2,8]{x\in B}{%
      \FormulaRefAuto{PowerSemigroupCarrierElement}[thm]{2,9}}
    \proofstepwidestar[3,12]{y\in B}{%
      \FormulaRefAuto{PowerSemigroupCarrierElement}[thm]{3,13}}
    \proofstepwidestar[4,8]{z\in B}{%
      \FormulaRefAuto{PowerSemigroupCarrierElement}[thm]{4,10}}
    \proofstepwidestar[1,2,3,4,8,12]{x\diamond z=y\diamond z}{%
      \rAEa{\FormulaRefAuto{SemigroupTranslationIndistinguishableActions}[thm]{%
        1,15,16,14,17}}}
    \proofstepwidestar[1,2,3,4,8,12]{w=y\diamond z}{\rChain{11,18}}
    \proofstepwidestar[1,3,4,8,12]{%
      y\diamond z\in Y\diamond_{\mathcal P}Z}{%
      \FormulaRefAuto{PowerSemigroupProductContains}[thm]{1,3,4,13,10}}
    \proofstepwidestar[1,2,3,4,8,12]{%
      w\in Y\diamond_{\mathcal P}Z}{\rIE{19,20}}
    \proofstepwidestar[1,2,3,4,5,8]{%
      w\in Y\diamond_{\mathcal P}Z}{\rEEStar{12\text{--}21}}
    \proofstepwidestar[1,2,3,4,5,7]{%
      w\in Y\diamond_{\mathcal P}Z}{\rEEStar{8\text{--}22}}
    \proofstepwidestar[1,2,3,4,5]{%
      \forall w\bigl(w\in X\diamond_{\mathcal P}Z\rightarrow
        w\in Y\diamond_{\mathcal P}Z\bigr)}{\rUI{\rRI{7,23}}}
    \proofstepwidestar[1,2,3,4,5]{%
      X\diamond_{\mathcal P}Z\subseteq Y\diamond_{\mathcal P}Z}{%
      \FormulaRefAuto{A\subseteq B\coloneq
        \forall x\,(x\in A\rightarrow x\in B)}[def]{24}}

    \proofstepwidestar[26]{w\in Y\diamond_{\mathcal P}Z}{\rA}
    \proofstepwidestar[1,3,4,26]{%
      \exists y\in Y\,\exists z\in Z\,(w=y\diamond z)}{%
      \FormulaRefAuto{PowerSemigroupProductWitnesses}[thm]{1,3,4,26}}
    \proofstepwidestar[27]{y\in Y}{\rA}
    \proofstepwidestar[27]{z\in Z}{\rA}
    \proofstepwidestar[27]{w=y\diamond z}{\rA}
    \proofstepwidestar[6,27]{\exists x\in X\,
      \bigl(y\mathrel{\equiv^{\mathrm{tr}}_{B,\diamond}}x\bigr)}{%
      \rRE{28,\rUE{6}}}
    \proofstepwidestar[31]{x\in X}{\rA}
    \proofstepwidestar[31]{%
      y\mathrel{\equiv^{\mathrm{tr}}_{B,\diamond}}x}{\rA}
    \proofstepwidestar[3,27]{y\in B}{%
      \FormulaRefAuto{PowerSemigroupCarrierElement}[thm]{3,28}}
    \proofstepwidestar[2,31]{x\in B}{%
      \FormulaRefAuto{PowerSemigroupCarrierElement}[thm]{2,32}}
    \proofstepwidestar[4,27]{z\in B}{%
      \FormulaRefAuto{PowerSemigroupCarrierElement}[thm]{4,29}}
    \proofstepwidestar[1,2,3,4,27,31]{y\diamond z=x\diamond z}{%
      \rAEa{\FormulaRefAuto{SemigroupTranslationIndistinguishableActions}[thm]{%
        1,34,35,33,36}}}
    \proofstepwidestar[1,2,3,4,27,31]{w=x\diamond z}{\rChain{30,37}}
    \proofstepwidestar[1,2,4,27,31]{%
      x\diamond z\in X\diamond_{\mathcal P}Z}{%
      \FormulaRefAuto{PowerSemigroupProductContains}[thm]{1,2,4,32,29}}
    \proofstepwidestar[1,2,3,4,27,31]{%
      w\in X\diamond_{\mathcal P}Z}{\rIE{38,39}}
    \proofstepwidestar[1,2,3,4,6,27]{%
      w\in X\diamond_{\mathcal P}Z}{\rEEStar{31\text{--}40}}
    \proofstepwidestar[1,2,3,4,6,26]{%
      w\in X\diamond_{\mathcal P}Z}{\rEEStar{27\text{--}41}}
    \proofstepwidestar[1,2,3,4,6]{%
      \forall w\bigl(w\in Y\diamond_{\mathcal P}Z\rightarrow
        w\in X\diamond_{\mathcal P}Z\bigr)}{\rUI{\rRI{26,42}}}
    \proofstepwidestar[1,2,3,4,6]{%
      Y\diamond_{\mathcal P}Z\subseteq X\diamond_{\mathcal P}Z}{%
      \FormulaRefAuto{A\subseteq B\coloneq
        \forall x\,(x\in A\rightarrow x\in B)}[def]{43}}
    \proofstepwidestar[1,2,3,4,5,6]{%
      X\diamond_{\mathcal P}Z=Y\diamond_{\mathcal P}Z}{%
      \FormulaRefAuto{A\subseteq B\dsep B\subseteq A\vdash A=B}[thm]{25,44}}
  \closeproofpartwideindR

  \proofpartwideindR[Linke Mengenprodukte]{%
    \begin{aligned}[t]
      &\SemigroupAlg{B}{\diamond}\dsep X,Y,Z\in\PowNE{B}\dsep{}\\[-2pt]
      &\forall x\in X\,\exists y\in Y\,
        \bigl(x\mathrel{\equiv^{\mathrm{tr}}_{B,\diamond}}y\bigr)
        \dsep{}\\[-2pt]
      &\forall y\in Y\,\exists x\in X\,
        \bigl(y\mathrel{\equiv^{\mathrm{tr}}_{B,\diamond}}x\bigr)
        \vdash Z\diamond_{\mathcal P}X=Z\diamond_{\mathcal P}Y
    \end{aligned}%
  }{%
    \begin{aligned}[t]
      &\SemigroupAlg{B}{\diamond}\dsep X,Y,Z\in\PowNE{B}\dsep{}\\[-2pt]
      &\forall x\in X\,\exists y\in Y\,
        \bigl(x\mathrel{\equiv^{\mathrm{tr}}_{B,\diamond}}y\bigr)
        \dsep{}\\[-2pt]
      &\forall y\in Y\,\exists x\in X\,
        \bigl(y\mathrel{\equiv^{\mathrm{tr}}_{B,\diamond}}x\bigr)
        \vdash Z\diamond_{\mathcal P}X=Z\diamond_{\mathcal P}Y
    \end{aligned}%
  }[PowerSemigroupTranslationSupportLeftProduct]
    \proofstepwidestar[1]{\SemigroupAlg{B}{\diamond}}{\rA}
    \proofstepwidestar[2]{X\in\PowNE{B}}{\rA}
    \proofstepwidestar[3]{Y\in\PowNE{B}}{\rA}
    \proofstepwidestar[4]{Z\in\PowNE{B}}{\rA}
    \proofstepwidestar[5]{\forall u\in X\,\exists v\in Y\,
      \bigl(u\mathrel{\equiv^{\mathrm{tr}}_{B,\diamond}}v\bigr)}{\rA}
    \proofstepwidestar[6]{\forall v\in Y\,\exists u\in X\,
      \bigl(v\mathrel{\equiv^{\mathrm{tr}}_{B,\diamond}}u\bigr)}{\rA}

    \proofstepwidestar[7]{w\in Z\diamond_{\mathcal P}X}{\rA}
    \proofstepwidestar[1,2,4,7]{%
      \exists z\in Z\,\exists x\in X\,(w=z\diamond x)}{%
      \FormulaRefAuto{PowerSemigroupProductWitnesses}[thm]{1,4,2,7}}
    \proofstepwidestar[8]{z\in Z}{\rA}
    \proofstepwidestar[8]{x\in X}{\rA}
    \proofstepwidestar[8]{w=z\diamond x}{\rA}
    \proofstepwidestar[5,8]{\exists y\in Y\,
      \bigl(x\mathrel{\equiv^{\mathrm{tr}}_{B,\diamond}}y\bigr)}{%
      \rRE{10,\rUE{5}}}
    \proofstepwidestar[12]{y\in Y}{\rA}
    \proofstepwidestar[12]{%
      x\mathrel{\equiv^{\mathrm{tr}}_{B,\diamond}}y}{\rA}
    \proofstepwidestar[2,8]{x\in B}{%
      \FormulaRefAuto{PowerSemigroupCarrierElement}[thm]{2,10}}
    \proofstepwidestar[3,12]{y\in B}{%
      \FormulaRefAuto{PowerSemigroupCarrierElement}[thm]{3,13}}
    \proofstepwidestar[4,8]{z\in B}{%
      \FormulaRefAuto{PowerSemigroupCarrierElement}[thm]{4,9}}
    \proofstepwidestar[1,2,3,4,8,12]{z\diamond x=z\diamond y}{%
      \rAEb{\FormulaRefAuto{SemigroupTranslationIndistinguishableActions}[thm]{%
        1,15,16,14,17}}}
    \proofstepwidestar[1,2,3,4,8,12]{w=z\diamond y}{\rChain{11,18}}
    \proofstepwidestar[1,3,4,8,12]{%
      z\diamond y\in Z\diamond_{\mathcal P}Y}{%
      \FormulaRefAuto{PowerSemigroupProductContains}[thm]{1,4,3,9,13}}
    \proofstepwidestar[1,2,3,4,8,12]{%
      w\in Z\diamond_{\mathcal P}Y}{\rIE{19,20}}
    \proofstepwidestar[1,2,3,4,5,8]{%
      w\in Z\diamond_{\mathcal P}Y}{\rEEStar{12\text{--}21}}
    \proofstepwidestar[1,2,3,4,5,7]{%
      w\in Z\diamond_{\mathcal P}Y}{\rEEStar{8\text{--}22}}
    \proofstepwidestar[1,2,3,4,5]{%
      \forall w\bigl(w\in Z\diamond_{\mathcal P}X\rightarrow
        w\in Z\diamond_{\mathcal P}Y\bigr)}{\rUI{\rRI{7,23}}}
    \proofstepwidestar[1,2,3,4,5]{%
      Z\diamond_{\mathcal P}X\subseteq Z\diamond_{\mathcal P}Y}{%
      \FormulaRefAuto{A\subseteq B\coloneq
        \forall x\,(x\in A\rightarrow x\in B)}[def]{24}}

    \proofstepwidestar[26]{w\in Z\diamond_{\mathcal P}Y}{\rA}
    \proofstepwidestar[1,3,4,26]{%
      \exists z\in Z\,\exists y\in Y\,(w=z\diamond y)}{%
      \FormulaRefAuto{PowerSemigroupProductWitnesses}[thm]{1,4,3,26}}
    \proofstepwidestar[27]{z\in Z}{\rA}
    \proofstepwidestar[27]{y\in Y}{\rA}
    \proofstepwidestar[27]{w=z\diamond y}{\rA}
    \proofstepwidestar[6,27]{\exists x\in X\,
      \bigl(y\mathrel{\equiv^{\mathrm{tr}}_{B,\diamond}}x\bigr)}{%
      \rRE{29,\rUE{6}}}
    \proofstepwidestar[31]{x\in X}{\rA}
    \proofstepwidestar[31]{%
      y\mathrel{\equiv^{\mathrm{tr}}_{B,\diamond}}x}{\rA}
    \proofstepwidestar[3,27]{y\in B}{%
      \FormulaRefAuto{PowerSemigroupCarrierElement}[thm]{3,29}}
    \proofstepwidestar[2,31]{x\in B}{%
      \FormulaRefAuto{PowerSemigroupCarrierElement}[thm]{2,32}}
    \proofstepwidestar[4,27]{z\in B}{%
      \FormulaRefAuto{PowerSemigroupCarrierElement}[thm]{4,28}}
    \proofstepwidestar[1,2,3,4,27,31]{z\diamond y=z\diamond x}{%
      \rAEb{\FormulaRefAuto{SemigroupTranslationIndistinguishableActions}[thm]{%
        1,34,35,33,36}}}
    \proofstepwidestar[1,2,3,4,27,31]{w=z\diamond x}{\rChain{30,37}}
    \proofstepwidestar[1,2,4,27,31]{%
      z\diamond x\in Z\diamond_{\mathcal P}X}{%
      \FormulaRefAuto{PowerSemigroupProductContains}[thm]{1,4,2,28,32}}
    \proofstepwidestar[1,2,3,4,27,31]{%
      w\in Z\diamond_{\mathcal P}X}{\rIE{38,39}}
    \proofstepwidestar[1,2,3,4,6,27]{%
      w\in Z\diamond_{\mathcal P}X}{\rEEStar{31\text{--}40}}
    \proofstepwidestar[1,2,3,4,6,26]{%
      w\in Z\diamond_{\mathcal P}X}{\rEEStar{27\text{--}41}}
    \proofstepwidestar[1,2,3,4,6]{%
      \forall w\bigl(w\in Z\diamond_{\mathcal P}Y\rightarrow
        w\in Z\diamond_{\mathcal P}X\bigr)}{\rUI{\rRI{26,42}}}
    \proofstepwidestar[1,2,3,4,6]{%
      Z\diamond_{\mathcal P}Y\subseteq Z\diamond_{\mathcal P}X}{%
      \FormulaRefAuto{A\subseteq B\coloneq
        \forall x\,(x\in A\rightarrow x\in B)}[def]{43}}
    \proofstepwidestar[1,2,3,4,5,6]{%
      Z\diamond_{\mathcal P}X=Z\diamond_{\mathcal P}Y}{%
      \FormulaRefAuto{A\subseteq B\dsep B\subseteq A\vdash A=B}[thm]{25,44}}
  \closeproofpartwideindR

  \proofpartwide{Schluss}
    \proofstepwidestar[1]{\SemigroupAlg{B}{\diamond}}{\rA}
    \proofstepwidestar[2]{X\in\PowNE{B}}{\rA}
    \proofstepwidestar[3]{Y\in\PowNE{B}}{\rA}
    \proofstepwidestar[4]{\forall x\in X\,\exists y\in Y\,
      \bigl(x\mathrel{\equiv^{\mathrm{tr}}_{B,\diamond}}y\bigr)}{\rA}
    \proofstepwidestar[5]{\forall y\in Y\,\exists x\in X\,
      \bigl(y\mathrel{\equiv^{\mathrm{tr}}_{B,\diamond}}x\bigr)}{\rA}
    \proofstepwidestar[1]{%
      \SemigroupAlg{\PowNE{B}}{\diamond_{\mathcal P}}}{%
      \FormulaRefAuto{NonemptyPowerSemigroup}[thm]{1}}
    \proofstepwidestar[7]{Z\in\PowNE{B}}{\rA}
    \proofstepwidestarbreak[1,2,3,4,5,7]{%
      X\diamond_{\mathcal P}Z=Y\diamond_{\mathcal P}Z}{%
      \FormulaRefAuto{PowerSemigroupTranslationSupportRightProduct}[thm]{%
        1,2,3,7,4,5}}
    \proofstepwidestarbreak[1,2,3,4,5,7]{%
      Z\diamond_{\mathcal P}X=Z\diamond_{\mathcal P}Y}{%
      \FormulaRefAuto{PowerSemigroupTranslationSupportLeftProduct}[thm]{%
        1,2,3,7,4,5}}
    \proofstepwidestar[1,2,3,4,5,7]{%
      \begin{aligned}[t]
        &X\diamond_{\mathcal P}Z=Y\diamond_{\mathcal P}Z\land{}\\[-2pt]
        &Z\diamond_{\mathcal P}X=Z\diamond_{\mathcal P}Y
      \end{aligned}}{\rAI{8,9}}
    \proofstepwidestarbreak[1,2,3,4,5]{%
      \forall Z\in\PowNE{B}\,\left(
        \begin{aligned}[t]
          &X\diamond_{\mathcal P}Z=Y\diamond_{\mathcal P}Z\land{}\\[-2pt]
          &Z\diamond_{\mathcal P}X=Z\diamond_{\mathcal P}Y
        \end{aligned}\right)}{%
      \rUI{\rRI{7,10}}}
    \proofstepwidestarbreak[1,2,3,4,5]{%
      X\mathrel{\equiv^{\mathrm{tr}}_{\PowNE{B},
        \diamond_{\mathcal P}}}Y}{%
      \FormulaRefAuto{SemigroupTranslationIndistinguishableDef}[def]{%
        6,2,3,11}}
  \closeproofpartwide
\end{tabproofsplitwide}

\newcommand{\MogStar}{\mathbin{\ast}}
\newcommand{\MogDiamond}{\mathbin{\diamond}}
\newcommand{\MogTableProofHeading}{%
  \par
  \ifdim\dimexpr\pagegoal-\pagetotal\relax<6\baselineskip
    \newpage
  \fi
  \medskip
  \noindent\textit{Tabellarische Fassung.}\par\nobreak
  \smallskip\nobreak
}
\newcommand{\MogProofStepWideStarKeep}[3][]{%
  \stepcounter{proofstepnr}%
  #1 & \theproofstepnr
     & \multicolumn{3}{@{}>{$}l<{$}@{}}{#2} & \\*[-.2ex]%
     & & \multicolumn{4}{@{}r@{}}{%
       \makebox[0pt][r]{\ensuremath{#3}}}\\%
}

\subsection{Ein unendliches Gegenbeispiel zur Umkehrung}

Die folgende Konstruktion geht auf
\href{https://doi.org/10.1007/BF02389140}%
  {Mogiljanskaja, \emph{Non-isomorphic Semigroups with Isomorphic
  Semigroups of Subsets}, 1973, S.~330--333}
zurück. Ihre Idee ist einfach: Ein zusätzliches Element zerstört die
Isomorphie der Grundhalbgruppen, wird auf der Ebene der Teilmengen aber von
einer unendlichen Reserve aufgenommen. Im Folgenden werden nur die für diesen
Mechanismus wesentlichen Aussagen ausgeführt.

Aus Band~21 übernehmen wir die disjunkten Schichten
\[
  L=\{a_i\mid i\in\mathbb N\},
  \qquad
  D=\{b_{ij}\mid i,j\in\mathbb N\}.
\]
Die Elemente werden durch ihre Indizes eindeutig bestimmt; siehe
\FormulaRefAuto{MogiljanskajaAIndexUnique}[thm],
\FormulaRefAuto{MogiljanskajaBIndexUnique}[thm] und
\FormulaRefAuto{MogiljanskajaLayersDisjoint}[thm].

Jeder Beweis dieses Gegenbeispiels erscheint unmittelbar nacheinander in zwei
Darstellungsformen. Zunächst wird er metasprachlich und auf den tragenden
Gedanken konzentriert geführt; darunter folgen dieselben Argumente in
Tabellenform. Die Kürzel \textup{(M1)} bis \textup{(M17)} verweisen auf die
ausführlicheren Einzelsätze im Anhang.

\FormulaDefDeltaK[Die beiden Träger und ihre Multiplikationen]{%
  \begin{gathered}
    \begin{aligned}
      \mathbf0&\coloneqq(2,0),&
      k&\coloneqq(3,0),\\[-2pt]
      S&\coloneqq L\cup D\cup\{\mathbf0\},&
      T&\coloneqq S\cup\{k\}
    \end{aligned}\\[3pt]
    \begin{aligned}[t]
      &\MogStar\colon S\times S\to S,\\[-2pt]
      &\forall i,j\in\mathbb N\,
        \bigl(a_i\MogStar a_j\coloneqq b_{ij}\bigr),\\[-2pt]
      &\forall x,y\in S\,
        \bigl((x\notin L\lor y\notin L)
          \to x\MogStar y\coloneqq\mathbf0\bigr),\\[2pt]
      &\MogDiamond\colon T\times T\to T,\\[-2pt]
      &\forall i,j\in\mathbb N\,
        \bigl(a_i\MogDiamond a_j\coloneqq b_{ij}\bigr),\\[-2pt]
      &\forall x,y\in T\,
        \bigl((x\notin L\lor y\notin L)
          \to x\MogDiamond y\coloneqq\mathbf0\bigr)
    \end{aligned}
  \end{gathered}
}{MogiljanskajaPairDef}{%
  \DeltaRow{Mengen}{L\dsep D\dsep S\dsep T\dsep\mathbf0\dsep k}
  \DeltaRow{Mengen}{a_i\dsep a_j\dsep b_{ij}\dsep i\dsep j\dsep x\dsep y}
  \DeltaRow{Funktionen}{\MogStar\dsep\MogDiamond}
}

\begin{proof}
Die verschiedenen ersten Tags halten alle Trägerschichten auseinander.
Liegen beide Faktoren in \(L\),
so besitzen sie eindeutige Indizes \(i,j\), und das vorgeschriebene Ergebnis
\(b_{ij}\) liegt in \(D\). Liegt mindestens ein Faktor außerhalb von \(L\),
so greift eindeutig der Nullfall; dabei gehört \(\mathbf0\) sowohl zu \(S\)
als auch zu \(T\). Damit definieren die beiden Fallvorschriften tatsächlich
Abbildungen \(S\times S\to S\) beziehungsweise \(T\times T\to T\). Die
Einzelheiten sind in \textup{(M1)} festgehalten.
\end{proof}

\MogTableProofHeading
\begin{tabproofwide}
  \proofstepwidestar[]{%
    \MogStar\colon S\times S\to S\dsep
    \MogDiamond\colon T\times T\to T}{%
    (M1) Tags, Indizes und Falltrennung}
\end{tabproofwide}

\FormulaThmDeltaK[Grundlegende Eigenschaften des Mogiljanskaja-Paares]{%
  \begin{aligned}[t]
    &\SemigroupAlg{S}{\MogStar}\land
      \SemigroupAlg{T}{\MogDiamond}\land{}\\[-2pt]
    &\neg\Finite{S}\land\neg\Finite{T}\land{}\\[-2pt]
    &\neg\exists f\,
      \SemigroupIso{f}{S,\MogStar}{T,\MogDiamond}
  \end{aligned}
}{MogiljanskajaBasePairProperties}{%
  \DeltaRow{Mengen}{L\dsep D\dsep S\dsep T\dsep\mathbf0\dsep k}
  \DeltaRow{Mengen}{s\dsep a_i\dsep a_m\dsep a_n\dsep a_0\dsep a_1}
  \DeltaRow{Mengen}{b_{ij}\dsep b_{mn}\dsep b_{i0}\dsep b_{i1}
    \dsep i\dsep m\dsep n}
  \DeltaRow{Funktionen}{a\dsep f\dsep\MogStar\dsep\MogDiamond}
}
\begin{proof}
Nach der Definition ist jedes erste Produkt entweder ein Element von \(D\)
oder \(\mathbf0\). Beide Bereiche liegen außerhalb von \(L\). Multipliziert
man ein solches Ergebnis noch mit einem dritten Faktor, so greift deshalb
stets der Nullfall. Für beide Operationen werden somit beide Klammerungen zu
\(\mathbf0\). Damit sind sie assoziativ; dies ist der Inhalt von
\textup{(M2)}.

Die Elemente \(a_i\) bilden eine injektiv indizierte Kopie von
\(\mathbb N\) in \(L\subseteq S\subseteq T\). Also können weder \(S\) noch
\(T\) endlich sein; siehe \textup{(M3)}.

Für die Nichtisomorphie ist der Produktstatus des neuen Punktes entscheidend.
Jedes Element von \(S\setminus L\) ist ein Produkt, während \(k\) in \(T\)
niemals als Produkt auftritt; das zeigt \textup{(M4)}. Angenommen, es gäbe
einen Isomorphismus
\(f\colon(S,\MogStar)\to(T,\MogDiamond)\). Seine Surjektivität liefert
zunächst nur ein \(s\in S\) mit \(f(s)=k\); sie sagt noch nicht, in welcher
Schicht dieses Urbild liegt.

Wegen der disjunkten Zerlegung
\(S=L\cup D\cup\{\mathbf0\}\) bleiben neben \(s\in L\) genau die beiden
Möglichkeiten \(s\in D\) und \(s=\mathbf0\). Im ersten Fall gibt es
\(m,n\in\mathbb N\) mit \(s=b_{mn}=a_m\MogStar a_n\), im zweiten Fall ist
\(s=\mathbf0=\mathbf0\MogStar\mathbf0\). Da ein Isomorphismus Produkte auf
Produkte abbildet, ergäbe sich entsprechend
\[
  \begin{aligned}
    s\in D&\quad\Longrightarrow\quad
      k=f(s)=f(a_m)\MogDiamond f(a_n),\\
    s=\mathbf0&\quad\Longrightarrow\quad
      k=f(s)=f(\mathbf0)\MogDiamond f(\mathbf0).
  \end{aligned}
\]
Die rechten Seiten sind Produkte zweier Elemente aus \(T\). Beides ist
unmöglich: Jedes Produkt in \(T\) liegt in
\(D\cup\{\mathbf0\}\), während \(k\notin D\cup\{\mathbf0\}\) gilt. Das
Urbild \(s\) liegt daher weder in \(D\) noch in \(\{\mathbf0\}\), sondern
in \(L\). Erst jetzt folgt \(s=a_i\) für ein \(i\in\mathbb N\), also
\(f(a_i)=k\). Dann wären aber
\[
  \begin{aligned}
    f(b_{i0})
      &=f(a_i\MogStar a_0)
        =f(a_i)\MogDiamond f(a_0)
        =k\MogDiamond f(a_0)=\mathbf0,\\
    f(b_{i1})
      &=f(a_i\MogStar a_1)
        =f(a_i)\MogDiamond f(a_1)
        =k\MogDiamond f(a_1)=\mathbf0.
  \end{aligned}
\]
Die verschiedenen Elemente \(b_{i0}\) und \(b_{i1}\) hätten dasselbe Bild,
im Widerspruch zur Injektivität von \(f\). Dies ist \textup{(M5)}.
\end{proof}

\MogTableProofHeading
\begin{tabproofwide}
  \MogProofStepWideStarKeep[]{%
    \SemigroupAlg{S}{\MogStar}\land
    \SemigroupAlg{T}{\MogDiamond}}{%
    \text{\normalfont (M2): Produktbereich und Nullabsorption}}
  \MogProofStepWideStarKeep[]{%
    \neg\Finite{S}\land\neg\Finite{T}}{%
    \text{\normalfont (M3): injizierte unendliche Schicht}}
  \MogProofStepWideStarKeep[]{%
    \SemigroupIso{f}{S,\MogStar}{T,\MogDiamond}
      \vdash\exists s\in S\,(f(s)=k)}{%
    \text{\normalfont Surjektivität und }k\in T}
  \MogProofStepWideStarKeep[]{%
    \begin{aligned}[t]
      &\SemigroupIso{f}{S,\MogStar}{T,\MogDiamond}
        \dsep s\in S\dsep f(s)=k\dsep s\notin L\\[-2pt]
      &\vdash\exists p\in S\,\exists q\in S\,(s=p\MogStar q)
    \end{aligned}}{%
    \text{\normalfont (M4): }s\in D\text{\normalfont{ oder }}s=\mathbf0}
  \MogProofStepWideStarKeep[4]{%
    \begin{aligned}[t]
      &\SemigroupIso{f}{S,\MogStar}{T,\MogDiamond}
        \dsep s\in S\dsep f(s)=k\dsep s\notin L\\[-2pt]
      &\vdash\exists u\in T\,\exists v\in T\,(k=u\MogDiamond v)
    \end{aligned}}{%
    \text{\normalfont Homomorphie; wähle }u=f(p)\text{\normalfont{ und }}v=f(q)}
  \MogProofStepWideStarKeep[]{%
    \neg\exists u\in T\,\exists v\in T\,(k=u\MogDiamond v)}{%
    \text{\normalfont (M4): }k\text{\normalfont{ liegt nicht im Produktbereich}}}
  \MogProofStepWideStarKeep[5,6]{%
    \begin{aligned}[t]
      &\SemigroupIso{f}{S,\MogStar}{T,\MogDiamond}
        \dsep s\in S\dsep f(s)=k\\[-2pt]
      &\vdash s\in L
    \end{aligned}}{%
    \text{\normalfont Ausschluss der Alternative }s\notin L}
  \MogProofStepWideStarKeep[3,7]{%
    \SemigroupIso{f}{S,\MogStar}{T,\MogDiamond}\vdash
      \exists i\in\mathbb N\,(f(a_i)=k)}{%
    \text{\normalfont Surjektivität und Beschreibung von }L}
  \MogProofStepWideStarKeep[8]{%
    \begin{aligned}[t]
      \SemigroupIso{f}{S,\MogStar}{T,\MogDiamond}\vdash
      \exists i\in\mathbb N\,\bigl(&f(b_{i0})=\mathbf0\land{}\\[-2pt]
        &f(b_{i1})=\mathbf0\bigr)
    \end{aligned}}{%
    \text{\normalfont (M5): Homomorphie und Multiplikationsvorschrift}}
  \MogProofStepWideStarKeep[]{i\in\mathbb N\vdash b_{i0}\neq b_{i1}}{%
    \text{\normalfont Indexeindeutigkeit und }0\neq1}
  \MogProofStepWideStarKeep[9,10]{%
    \neg\exists f\,\SemigroupIso{f}{S,\MogStar}{T,\MogDiamond}}{%
    \text{\normalfont Widerspruch zur Injektivität von }f}
  \MogProofStepWideStarKeep[1,2,11]{%
    \begin{aligned}[t]
      &\SemigroupAlg{S}{\MogStar}\land
        \SemigroupAlg{T}{\MogDiamond}\land{}\\[-2pt]
      &\neg\Finite{S}\land\neg\Finite{T}\land{}\\[-2pt]
      &\neg\exists f\,
        \SemigroupIso{f}{S,\MogStar}{T,\MogDiamond}
    \end{aligned}}{%
    \text{\normalfont Bündelung}}
\end{tabproofwide}

\FormulaDefDeltaK[Rechteck- und Nullanteil eines Mengenprodukts]{%
  \begin{aligned}[t]
    R(X,Y)&\coloneqq
      \bigl\{b_{ij}\bigm|
        i,j\in\mathbb N,\,
        a_i\in X\cap L,\,
        a_j\in Y\cap L\bigr\},\\[-2pt]
    \varepsilon(X,Y)&\coloneqq
      \begin{cases}
        \varnothing,&X\subseteq L\land Y\subseteq L,\\
        \{\mathbf0\},&\text{sonst}
      \end{cases}
  \end{aligned}
}{MogiljanskajaProductPartsDef}{%
  \DeltaRow{Mengen}{L\dsep D\dsep X\dsep Y\dsep
    a_i\dsep a_j\dsep b_{ij}\dsep i\dsep j
    \dsep R(X,Y)\dsep\varepsilon(X,Y)}
}

\FormulaThmDeltaK[Direkte Formel für die beiden Mengenprodukte]{%
  \begin{aligned}[t]
    X,Y\in\PowNE{S}
      &\vdash X\MogStar_{\mathcal P}Y
        =R(X,Y)\cup\varepsilon(X,Y),\\[-2pt]
    X,Y\in\PowNE{T}
      &\vdash X\MogDiamond_{\mathcal P}Y
        =R(X,Y)\cup\varepsilon(X,Y)
  \end{aligned}
}{MogiljanskajaPowerProductFormula}{%
  \DeltaRow{Mengen}{L\dsep D\dsep S\dsep T\dsep\mathbf0
    \dsep X\dsep Y\dsep R(X,Y)\dsep\varepsilon(X,Y)}
  \DeltaRow{Funktionen}{\MogStar_{\mathcal P}
    \dsep\MogDiamond_{\mathcal P}}
}
\begin{proof}
In einem Mengenprodukt gibt es nur zwei Arten von Ergebnissen. Treffen zwei
Elemente aus \(L\) aufeinander, so entsteht das durch ihre Indizes bestimmte
Element \(b_{ij}\); die Gesamtheit dieser Ergebnisse ist genau das Rechteck
\(R(X,Y)\). Sobald mindestens ein Faktor außerhalb von \(L\) liegt, entsteht
dagegen \(\mathbf0\).

Da \(X\) und \(Y\) nichtleer sind, existiert ein solches gemischtes Paar
genau dann, wenn nicht zugleich \(X\subseteq L\) und \(Y\subseteq L\) gilt.
Genau diesen Fall zeichnet \(\varepsilon(X,Y)=\{\mathbf0\}\) aus. Damit ist
das Mengenprodukt in beiden Trägern die Vereinigung aus Rechteck- und
Nullanteil. Eine elementweise Ausführung steht in \textup{(M6)}.
\end{proof}

\MogTableProofHeading
\begin{tabproofwide}
  \proofstepwidestarbreak[]{%
    \begin{aligned}[t]
      X,Y\in\PowNE{S}
        &\vdash X\MogStar_{\mathcal P}Y
          =R(X,Y)\cup\varepsilon(X,Y),\\[-2pt]
      X,Y\in\PowNE{T}
        &\vdash X\MogDiamond_{\mathcal P}Y
          =R(X,Y)\cup\varepsilon(X,Y)
    \end{aligned}}{%
    \text{\normalfont(M6) Paarweise Produktklassifikation}}
\end{tabproofwide}

Für die Reserve verwenden wir die in Band~21 definierten Mengen und
Bijektionen
\[
  \begin{gathered}
    c_0=b_{00},\qquad c_1=b_{11},\qquad c_2=b_{10},\\
    U=\{b_{\Succ(\Succ(n)),j}\mid n,j\in\mathbb N\},\qquad
    V=U\cup\{c_2\},\\
    P=\{c_0,c_1\},\qquad D'=P\cup V,\\
    \sigma\colon U\bij V,\qquad
    \theta\colon U\bij D.
  \end{gathered}
\]
Dies sind \FormulaRefAuto{MogiljanskajaDPrimeDef}[def],
\FormulaRefAuto{MogiljanskajaSigmaBijective}[thm] und
\FormulaRefAuto{MogiljanskajaThetaBijective}[thm]. Nach
\FormulaRefAuto{MogiljanskajaReserveMarkerFacts}[thm] und
\FormulaRefAuto{MogiljanskajaReserveCoreFacts}[thm] gilt insbesondere
\[
  P\cap U=\varnothing,\qquad
  c_2\notin P\cup U,\qquad
  D'\subseteq D,\qquad
  b_{01}\notin D'.
\]

\FormulaDefDeltaK[Reserve und lokale Potenzmengenabbildung]{%
  \begin{aligned}[t]
    \mathcal C&\coloneqq\CoreFamily{P}{V},&
    \mathcal C_0&\coloneqq\CoreFamily{P}{U},\\[-2pt]
    \mathcal C_1&\coloneqq\CoreFamily{P\cup\{c_2\}}{U},&
    \mathcal E&\coloneqq\CoreFamily{\{k\}}{D},\\[-2pt]
    \kappa_0&\coloneqq\CoreTransport{P}{P}{\sigma},&
    \kappa_1&\coloneqq
      \CoreTransport{P\cup\{c_2\}}{\{k\}}{\theta},\\[-2pt]
    \psi&\coloneqq
      \CaseFun{\mathcal C_0}{\mathcal C_1}{\kappa_0}{\kappa_1},\\[-2pt]
    \mathcal O&\coloneqq\powerset(D)\setminus\mathcal C,&
    \eta&\coloneqq
      \CaseFun{\mathcal O}{\mathcal C}{\Id_{\mathcal O}}{\psi}
  \end{aligned}
}{MogiljanskajaReserveMapsDef}{%
  \DeltaRow{Mengen}{D\dsep U\dsep V\dsep P\dsep c_2\dsep k
    \dsep\mathcal C\dsep\mathcal C_0\dsep\mathcal C_1
    \dsep\mathcal E\dsep\mathcal O}
  \DeltaRow{Funktionen}{\sigma\dsep\theta\dsep
    \kappa_0\dsep\kappa_1\dsep\psi\dsep\eta}
}

\FormulaThmDeltaK[Eigenschaften der Reserveabbildung]{%
  \begin{aligned}[t]
    &\psi\colon\mathcal C\bij\mathcal C\cup\mathcal E,\\[-2pt]
    &\eta\colon\powerset(D)\bij\powerset(D\cup\{k\})
      \dsep\eta(\varnothing)=\varnothing,\\[-2pt]
    &K\in\powerset(D)\dsep K\notin\mathcal C
      \vdash\eta(K)=K,\\[-2pt]
    &X,Y\in\powerset(T)\vdash R(X,Y)\notin\mathcal C
  \end{aligned}
}{MogiljanskajaReserveMapProperties}{%
  \DeltaRow{Mengen}{D\dsep D'\dsep P\dsep U\dsep V\dsep k
    \dsep\mathcal C\dsep\mathcal E\dsep\mathcal O
    \dsep K\dsep X\dsep Y\dsep R(X,Y)}
  \DeltaRow{Funktionen}{\psi\dsep\eta}
}
\begin{proof}
Die Reservefamilie \(\mathcal C\) wird in zwei disjunkte Teile zerlegt. Der
erste Teil wird mit \(\sigma\) wieder in \(\mathcal C\) transportiert, der
zweite mit \(\theta\) in die neue Familie \(\mathcal E\) der
\(k\)-haltigen Mengen. Die dafür nötigen Disjunktheits- und
Bijektivitätsvoraussetzungen stehen in \textup{(M7)}; der allgemeine
Absorptionssatz liefert dann die Bijektion
\(\psi\colon\mathcal C\bij\mathcal C\cup\mathcal E\), wie in
\textup{(M8)} ausgeführt.

Anschaulich besteht \(\mathcal E\) genau aus den Teilmengen von
\(D\cup\{k\}\), die den neuen Punkt \(k\) enthalten. Alle anderen
Zielmengen sind Teilmengen von \(D\) und liegen daher in
\(\mathcal O\cup\mathcal C\). Auf \(\mathcal O\) wirkt \(\eta\) als
Identität, auf \(\mathcal C\) wie \(\psi\). Diese disjunkte Verklebung macht
\(\eta\) bijektiv; siehe \textup{(M9)}. Weil der feste Kern
\(P=\{c_0,c_1\}\) nichtleer ist, gehört \(\varnothing\) nicht zur Reserve
\(\mathcal C\). Daher fixiert \(\eta\) die leere Menge und überhaupt jede
Menge außerhalb von \(\mathcal C\); dies ist \textup{(M10)}.

Schließlich kann kein Produktrechteck in der Reserve liegen. Ein Element von
\(\mathcal C\) müsste die beiden festen Ecken \(b_{00}\) und \(b_{11}\)
enthalten. Für ein Rechteck erzwingen diese beiden Ecken aber auch die
gemischte Ecke \(b_{01}\). Zugleich wäre das Rechteck als Element von
\(\mathcal C\) eine Teilmenge von \(D'\), obwohl \(b_{01}\notin D'\) gilt.
Dieser Widerspruch ist in \textup{(M11)} formalisiert.
\end{proof}

\MogTableProofHeading
\begin{tabproofwide}
  \proofstepwidestar[]{%
    \psi\colon\mathcal C\bij\mathcal C\cup\mathcal E}{%
    (M7)--(M8) Kerntransport und Disjunktheit}
  \proofstepwidestar[]{%
    \eta\colon\powerset(D)\bij\powerset(D\cup\{k\})}{%
    (M9) disjunkte Zielzerlegung}
  \proofstepwidestar[]{%
    \eta(\varnothing)=\varnothing}{%
    (M10) geschützter unveränderter Teil}
  \proofstepwidestar[]{%
    K\in\powerset(D)\dsep K\notin\mathcal C
      \vdash\eta(K)=K}{%
    (M10) Identitätszweig}
  \proofstepwidestar[]{%
    X,Y\in\powerset(T)\vdash R(X,Y)\notin\mathcal C}{%
    (M11) erzwungene Ecke \(b_{01}\)}
\end{tabproofwide}

\FormulaDefDeltaK[Die globale Potenzmengenabbildung]{%
  \begin{aligned}[t]
    A_*&\coloneqq L\cup\{\mathbf0\},\\[-2pt]
    \Phi&\coloneqq
      \LayerPowerMap{A_*}{D}{D\cup\{k\}}{\eta}
      \restriction_{\PowNE{S}}
  \end{aligned}
}{MogiljanskajaPhiDef}{%
  \DeltaRow{Mengen}{A_*\dsep L\dsep D\dsep S\dsep T
    \dsep\mathbf0\dsep k}
  \DeltaRow{Funktionen}{\eta\dsep\Phi}
}

Für \(X\in\PowNE{S}\) lautet die Grundgleichung
\[
  \Phi(X)=(X\cap A_*)\cup\eta(X\cap D).
\]

\FormulaThmDeltaK[Isomorphie der Potenzhalbgruppen]{%
  \SemigroupIso{\Phi}
    {\PowNE{S},\MogStar_{\mathcal P}}
    {\PowNE{T},\MogDiamond_{\mathcal P}}%
}{MogiljanskajaPowerSemigroupIsomorphism}{%
  \DeltaRow{Mengen}{A_*\dsep L\dsep D\dsep S\dsep T
    \dsep X\dsep Y\dsep Z}
  \DeltaRow{Funktionen}{\eta\dsep\Phi
    \dsep\MogStar_{\mathcal P}\dsep\MogDiamond_{\mathcal P}}
}
\begin{proof}
Die Träger zerfallen in die geschützte Schicht
\(A_*=L\cup\{\mathbf0\}\) und die jeweils bewegliche Schicht:
\[
  S=A_*\cup D,
  \qquad
  T=A_*\cup(D\cup\{k\}).
\]
Auf \(A_*\) ändert \(\Phi\) nichts; auf \(D\) verwendet sie die bijektive
und nulltreue Abbildung \(\eta\). Daher ist \(\Phi\) eine Bijektion zwischen
den nichtleeren Potenzmengen. Dies ist \textup{(M12)}.

Die Abbildung verändert also nur solche \(D\)-Anteile, die als Mengen in
\(\mathcal C\) liegen. Der \(L\)-Anteil einer Menge und die Frage, ob die
ganze Menge in \(L\) liegt, bleiben unverändert; siehe \textup{(M13)}.
Folglich ändern sich weder das
Produktrechteck \(R(X,Y)\) noch der Nullanteil \(\varepsilon(X,Y)\). Mit der
Produktformel aus \textup{(M6)} ergibt sich
\[
  \Phi(X)\MogDiamond_{\mathcal P}\Phi(Y)
    =R(X,Y)\cup\varepsilon(X,Y)
    =X\MogStar_{\mathcal P}Y.
\]
Das ist \textup{(M14)}.

Setzt man \(Z=X\MogStar_{\mathcal P}Y\), so ist sein \(D\)-Anteil gerade
\(R(X,Y)\) und sein geschützter Anteil gerade \(\varepsilon(X,Y)\). Nach
\textup{(M11)} liegt das Rechteck außerhalb der Reserve und wird daher von
\(\eta\) fixiert. Somit gilt \(\Phi(Z)=Z\); dies ist \textup{(M15)}. Beide
Seiten der Homomorphiegleichung sind also dieselbe Menge:
\[
  \Phi(X\MogStar_{\mathcal P}Y)
    =X\MogStar_{\mathcal P}Y
    =\Phi(X)\MogDiamond_{\mathcal P}\Phi(Y).
\]
Zusammen mit der Bijektivität von \(\Phi\) folgt die behauptete Isomorphie,
wie in \textup{(M16)} zusammengefasst.
\end{proof}

\MogTableProofHeading
\begin{tabproofwide}
  \proofstepwidestar[]{%
    \Phi\colon\PowNE{S}\bij\PowNE{T}}{%
    (M12) nulltreue Schichtfortsetzung}
  \proofstepwidestarbreak[]{%
    \begin{aligned}[t]
      \Phi(X)\cap L&=X\cap L,\\[-2pt]
      X\subseteq L&\Longleftrightarrow\Phi(X)\subseteq L
    \end{aligned}}{%
    \text{\normalfont(M13) Invarianz der geschützten Schicht}}
  \proofstepwidestarbreak[]{%
    \Phi(X)\MogDiamond_{\mathcal P}\Phi(Y)
      =X\MogStar_{\mathcal P}Y}{%
    \text{\normalfont(M14) Invarianz von }R\text{\normalfont{ und }}
      \varepsilon}
  \proofstepwidestarbreak[]{%
    \Phi(X\MogStar_{\mathcal P}Y)
      =X\MogStar_{\mathcal P}Y}{%
    \text{\normalfont(M15) Produktrechtecke sind geschützt}}
  \proofstepwidestarbreak[3,4]{%
    \Phi(X\MogStar_{\mathcal P}Y)
      =\Phi(X)\MogDiamond_{\mathcal P}\Phi(Y)}{%
    \text{\normalfont(M16) Vergleich der Fixpunktgleichungen}}
  \proofstepwidestarbreak[1,5]{%
    \SemigroupIso{\Phi}
      {\PowNE{S},\MogStar_{\mathcal P}}
      {\PowNE{T},\MogDiamond_{\mathcal P}}}{%
    \text{\normalfont(M16) Bijektivität und Homomorphie}}
\end{tabproofwide}

\FormulaThmDeltaK[Das Mogiljanskaja-Paar ist ein unendliches
Gegenbeispiel]{%
  \begin{aligned}[t]
    &\SemigroupAlg{S}{\MogStar}\land
      \SemigroupAlg{T}{\MogDiamond}\land{}\\[-2pt]
    &\neg\Finite{S}\land\neg\Finite{T}\land{}\\[-2pt]
    &\neg\exists f\,
      \SemigroupIso{f}{S,\MogStar}{T,\MogDiamond}\land{}\\[-2pt]
    &\exists F\,\SemigroupIso{F}
      {\PowNE{S},\MogStar_{\mathcal P}}
      {\PowNE{T},\MogDiamond_{\mathcal P}}
  \end{aligned}
}{MogiljanskajaPairCounterexample}{%
  \DeltaRow{Mengen}{S\dsep T}
  \DeltaRow{Funktionen}{f\dsep F\dsep\Phi
    \dsep\MogStar\dsep\MogDiamond
    \dsep\MogStar_{\mathcal P}\dsep\MogDiamond_{\mathcal P}}
}
\begin{proof}
Die grundlegenden Eigenschaften wurden in
\FormulaRefAuto{MogiljanskajaBasePairProperties}[thm] bewiesen: \(S\) und
\(T\) sind unendliche Halbgruppen, aber nicht isomorph. Für die letzte
Aussage wählen wir ausdrücklich \(F=\Phi\). Nach
\FormulaRefAuto{MogiljanskajaPowerSemigroupIsomorphism}[thm] ist diese
Abbildung ein Isomorphismus ihrer Potenzhalbgruppen. Damit sind alle vier
Bestandteile des behaupteten Gegenbeispiels gleichzeitig erfüllt; der
explizite Schlusszeuge wird in \textup{(M17)} nochmals zusammengefasst.
\end{proof}

\MogTableProofHeading
\begin{tabproofwide}
  \proofstepwidestar[]{%
    \begin{aligned}[t]
      &\SemigroupAlg{S}{\MogStar}\land
        \SemigroupAlg{T}{\MogDiamond}\land{}\\[-2pt]
      &\neg\Finite{S}\land\neg\Finite{T}\land{}\\[-2pt]
      &\neg\exists f\,
        \SemigroupIso{f}{S,\MogStar}{T,\MogDiamond}
    \end{aligned}}{%
    \FormulaRefAuto{MogiljanskajaBasePairProperties}[thm]}
  \proofstepwidestarbreak[]{%
    \exists F\,\SemigroupIso{F}
      {\PowNE{S},\MogStar_{\mathcal P}}
      {\PowNE{T},\MogDiamond_{\mathcal P}}}{%
    \text{\normalfont(M17) Zeuge }F=\Phi}
  \proofstepwidestar[1,2]{%
    \begin{aligned}[t]
      &\SemigroupAlg{S}{\MogStar}\land
        \SemigroupAlg{T}{\MogDiamond}\land{}\\[-2pt]
      &\neg\Finite{S}\land\neg\Finite{T}\land{}\\[-2pt]
      &\neg\exists f\,
        \SemigroupIso{f}{S,\MogStar}{T,\MogDiamond}\land{}\\[-2pt]
      &\exists F\,\SemigroupIso{F}
        {\PowNE{S},\MogStar_{\mathcal P}}
        {\PowNE{T},\MogDiamond_{\mathcal P}}
    \end{aligned}}{%
    \text{\normalfont Bündelung}}
\end{tabproofwide}

Damit ist das Mogiljanskaja-Paar ein Gegenbeispiel zur Umkehrung von
\FormulaRefAuto{SemigroupIsoImpliesPowerSemigroupIso}[thm].
Die Umkehrung ist also ohne Endlichkeitsvoraussetzung falsch.

\chapter{Formaler Anhang zum Mogiljanskaja-Gegenbeispiel}

In diesem Anhang wird jedes der Argumente \textup{(M1)} bis \textup{(M17)} als
eigener Satz formuliert. Auf den ausführlichen Beweis in gewöhnlicher
mathematischer Sprache folgt jeweils eine Tabelle mit derselben Argumentfolge
und denselben Abhängigkeiten. So lassen sich beide Darstellungsformen direkt
vergleichen. Die Reihenfolge stimmt mit der Argumentfolge des Haupttextes
überein; insbesondere bleibt der dort verwendete direkte Weg über \(R\) und
\(\varepsilon\) erhalten.

\section{Träger und Grundhalbgruppen}

\newcommand{\MogRule}[3]{\mathsf R_{\mathrm M}(#1,#2,#3)}

\FormulaDefDeltaK[Die Fallrelation der beiden Multiplikationen]{%
  \begin{aligned}[t]
    \MogRule{x}{y}{z}\ \leftrightarrow{}&
      \exists i,j\in\mathbb N\,
        (x=a_i\land y=a_j\land z=b_{ij})\ \lor{}\\[-2pt]
      &((x\notin L\lor y\notin L)\land z=\mathbf0)
  \end{aligned}
}{MogiljanskajaRuleDef}{%
  \DeltaRow{Mengen}{L\dsep\mathbf0\dsep x\dsep y\dsep z
    \dsep a_i\dsep a_j\dsep b_{ij}\dsep i\dsep j}
  \DeltaRow{Prädikatensymbole}{\mathsf R_{\mathrm M}}
}

\noindent\textit{Metasprachliche Lesart.}
Die Relation fasst die beiden Multiplikationsfälle in einem Satz zusammen:
Liegen beide Faktoren in \(L\), so bestimmen ihre eindeutigen Indizes
\(i,j\) das Ergebnis \(b_{ij}\). Sobald wenigstens ein Faktor außerhalb
von \(L\) liegt, ist das Ergebnis \(\mathbf0\). Sie ist damit die
relationale Schreibweise genau der Fallunterscheidung, durch die
\(\MogStar\) und \(\MogDiamond\) festgelegt sind.

\newpage
\FormulaThmDeltaK[Argument (M1): Wohldefiniertheit der Konstruktion]{%
  \begin{aligned}[t]
    &L\cap D=L\cap\{\mathbf0\}=D\cap\{\mathbf0\}
      =L\cap\{k\}=D\cap\{k\}=\{\mathbf0\}\cap\{k\}
      =\varnothing,\\[-2pt]
    &s\in L\vdash\exists!i\in\mathbb N\,(s=a_i),\\[-2pt]
    &x,y\in S\vdash\exists!z\in S\,\MogRule{x}{y}{z},\\[-2pt]
    &x,y\in T\vdash\exists!z\in T\,\MogRule{x}{y}{z},\\[-2pt]
    &\MogStar\colon S\times S\to S\dsep
      \MogDiamond\colon T\times T\to T,\\[-2pt]
    &x,y\in S\vdash\MogRule{x}{y}{x\MogStar y},\\[-2pt]
    &x,y\in T\vdash\MogRule{x}{y}{x\MogDiamond y}
  \end{aligned}
}{MogiljanskajaArgumentM1WellDefined}{%
  \DeltaRow{Mengen}{L\dsep D\dsep S\dsep T\dsep\mathbf0\dsep k
    \dsep s\dsep x\dsep y\dsep z\dsep a_i\dsep a_j\dsep b_{ij}
    \dsep i\dsep j}
  \DeltaRow{Funktionen}{a\dsep\MogStar\dsep\MogDiamond}
}
\begin{proof}[Metasprachliche Fassung]
Die verschiedenen ersten Koordinaten wirken hier wie unverwechselbare
Etiketten: Die Elemente von \(L\) tragen Tag \(0\), diejenigen von \(D\)
Tag \(1\), während \(\mathbf0\) und \(k\) die Tags \(2\) beziehungsweise
\(3\) besitzen. Daher sind \(L,D,\{\mathbf0\}\) und \(\{k\}\) paarweise
disjunkt; für \(L\cap D=\varnothing\) siehe auch
\FormulaRefAuto{MogiljanskajaLayersDisjoint}[thm]. Jedes \(s\in L\) hat
nach \FormulaRefAuto{MogiljanskajaLayerLMembership}[thm] die Form \(a_i\).
Der Index \(i\) ist nach
\FormulaRefAuto{MogiljanskajaAIndexUnique}[thm] eindeutig.

Seien nun \(x,y\in S\), beziehungsweise \(x,y\in T\). Liegen beide
Elemente in \(L\), so gibt es eindeutige Indizes \(i,j\) mit
\(x=a_i\) und \(y=a_j\); die Fallrelation schreibt dann eindeutig
\(z=b_{ij}\) vor. Dieses Element liegt nach
\FormulaRefAuto{MogiljanskajaBElementsDef}[def] in \(D\), also sowohl in
\(S\) als auch in \(T\). Liegt wenigstens einer der beiden Faktoren
außerhalb von \(L\), so schreibt die Relation eindeutig \(z=\mathbf0\)
vor. Die beiden Fälle schließen einander aus und
\(\mathbf0\in S\subseteq T\). Somit gibt es in beiden Trägern zu jedem
Faktorenpaar genau ein passendes \(z\).

Die in \FormulaRefAuto{MogiljanskajaPairDef}[def] festgelegten Operationen
wählen jeweils genau dieses eindeutige \(z\). Daher sind
\(\MogStar\colon S\times S\to S\) und
\(\MogDiamond\colon T\times T\to T\) wohldefiniert und erfüllen die
Fallrelation \FormulaRefAuto{MogiljanskajaRuleDef}[def].
\end{proof}

\MogTableProofHeading
\begin{tabproofwide}
  \proofstepwidestar[]{L\cap D=\varnothing}{%
    \FormulaRefAuto{MogiljanskajaLayersDisjoint}[thm]}
  \proofstepwidestar[]{\mathbf0\notin L\cup D}{%
    \text{\normalfont Tag 2 ist weder 0 noch 1}}
  \proofstepwidestarbreak[]{k\notin L\cup D\cup\{\mathbf0\}}{%
    \text{\normalfont Tag 3 ist von 0, 1 und 2 verschieden}}
  \proofstepwidestarbreak[1,2,3]{%
    \begin{aligned}[t]
      L\cap D&=L\cap\{\mathbf0\}=D\cap\{\mathbf0\}\\[-2pt]
        &=L\cap\{k\}=D\cap\{k\}
          =\{\mathbf0\}\cap\{k\}=\varnothing
    \end{aligned}}{%
    \text{\normalfont paarweise Disjunktheit}}
  \proofstepwidestarbreak[]{s\in L\leftrightarrow
    \exists i\in\mathbb N\,(s=a_i)}{%
    \FormulaRefAuto{MogiljanskajaLayerLMembership}[thm]}
  \proofstepwidestarbreak[5]{s\in L\vdash
    \exists!i\in\mathbb N\,(s=a_i)}{%
    \FormulaRefAuto{MogiljanskajaAIndexUnique}[thm]}
  \proofstepwidestarbreak[]{i,j\in\mathbb N\vdash
    b_{ij}\in D\subseteq S\subseteq T}{%
    \FormulaRefAuto{MogiljanskajaBElementsDef}[def],
    \FormulaRefAuto{MogiljanskajaPairDef}[def]}
  \proofstepwidestar[]{\mathbf0\in S\subseteq T}{%
    \FormulaRefAuto{MogiljanskajaPairDef}[def]}
  \proofstepwidestar[]{i,j,m,n\in\mathbb N\dsep b_{ij}=b_{mn}
    \vdash i=m\land j=n}{%
    \FormulaRefAuto{MogiljanskajaBIndexUnique}[thm]}
  \proofstepwidestarbreak[4,6,7,8,9]{x,y\in S\vdash
    \exists!z\in S\,\MogRule{x}{y}{z}}{%
    \FormulaRefAuto{MogiljanskajaRuleDef}[def]}
  \proofstepwidestarbreak[4,6,7,8,9]{x,y\in T\vdash
    \exists!z\in T\,\MogRule{x}{y}{z}}{%
    \FormulaRefAuto{MogiljanskajaRuleDef}[def]}
  \proofstepwidestarbreak[]{\MogStar\colon S\times S\to S\dsep
    \MogDiamond\colon T\times T\to T}{%
    \FormulaRefAuto{MogiljanskajaPairDef}[def]}
  \proofstepwidestarbreak[10,12]{x,y\in S\vdash
    \MogRule{x}{y}{x\MogStar y}}{%
    \FormulaRefAuto{MogiljanskajaPairDef}[def]}
  \proofstepwidestarbreak[11,12]{x,y\in T\vdash
    \MogRule{x}{y}{x\MogDiamond y}}{%
    \FormulaRefAuto{MogiljanskajaPairDef}[def]}
\end{tabproofwide}

\FormulaThmDeltaK[Argument (M2): Produktbereich und Nullabsorption]{%
  \begin{aligned}[t]
    x,y\in S&\vdash x\MogStar y\in D\cup\{\mathbf0\},\\[-2pt]
    u\in D\cup\{\mathbf0\}\dsep z\in S
      &\vdash u\MogStar z=\mathbf0=z\MogStar u,\\[-2pt]
    x,y,z\in S&\vdash
      (x\MogStar y)\MogStar z=\mathbf0
        =x\MogStar(y\MogStar z),\\[2pt]
    x,y\in T&\vdash x\MogDiamond y\in D\cup\{\mathbf0\},\\[-2pt]
    u\in D\cup\{\mathbf0\}\dsep z\in T
      &\vdash u\MogDiamond z=\mathbf0=z\MogDiamond u,\\[-2pt]
    x,y,z\in T&\vdash
      (x\MogDiamond y)\MogDiamond z=\mathbf0
        =x\MogDiamond(y\MogDiamond z),\\[2pt]
    &\SemigroupAlg{S}{\MogStar}\land
      \SemigroupAlg{T}{\MogDiamond}
  \end{aligned}
}{MogiljanskajaArgumentM2Semigroups}{%
  \DeltaRow{Mengen}{L\dsep D\dsep S\dsep T\dsep\mathbf0
    \dsep x\dsep y\dsep z\dsep u}
  \DeltaRow{Funktionen}{\MogStar\dsep\MogDiamond}
}
\begin{proof}[Metasprachliche Fassung]
Nach \FormulaRefAuto{MogiljanskajaPairDef}[def] kann ein Produkt nur auf
zwei Arten entstehen: Zwei Elemente der Schicht \(L\) liefern ein Element
\(b_{ij}\in D\), jedes andere Faktorenpaar liefert \(\mathbf0\). Deshalb
liegen sowohl \(x\MogStar y\) als auch \(x\MogDiamond y\) stets in
\(D\cup\{\mathbf0\}\).

Nach \FormulaRefAuto{MogiljanskajaArgumentM1WellDefined}[thm] ist
\(D\cup\{\mathbf0\}\) disjunkt zu \(L\). Sobald also einer der Faktoren
in \(D\cup\{\mathbf0\}\) liegt, greift stets der zweite
Multiplikationsfall, und das Ergebnis ist \(\mathbf0\). Für drei Elemente
\(x,y,z\) liegt das zuerst gebildete Produkt daher in
\(D\cup\{\mathbf0\}\); bei der zweiten Multiplikation entsteht auf beiden
Klammerungsseiten Null:
\[
  (x\MogStar y)\MogStar z=\mathbf0
    =x\MogStar(y\MogStar z)
\]
und ebenso
\[
  (x\MogDiamond y)\MogDiamond z=\mathbf0
    =x\MogDiamond(y\MogDiamond z).
\]
Beide Operationen sind somit assoziativ. Zusammen mit ihrer
Wohldefiniertheit aus
\FormulaRefAuto{MogiljanskajaArgumentM1WellDefined}[thm] folgt nach
\FormulaRefAuto{SemigroupDef}[def], dass \((S,\MogStar)\) und
\((T,\MogDiamond)\) Halbgruppen sind.
\end{proof}

\MogTableProofHeading
\begin{tabproofwide}
  \proofstepwidestar[]{x,y\in S\vdash
    x\MogStar y\in D\cup\{\mathbf0\}}{%
    \FormulaRefAuto{MogiljanskajaPairDef}[def]}
  \proofstepwidestar[]{x,y\in T\vdash
    x\MogDiamond y\in D\cup\{\mathbf0\}}{%
    \FormulaRefAuto{MogiljanskajaPairDef}[def]}
  \proofstepwidestar[]{(D\cup\{\mathbf0\})\cap L=\varnothing}{%
    \FormulaRefAuto{MogiljanskajaArgumentM1WellDefined}[thm]}
  \proofstepwidestarbreak[3]{u\in D\cup\{\mathbf0\}\dsep z\in S
    \vdash u\MogStar z=\mathbf0=z\MogStar u}{%
    \FormulaRefAuto{MogiljanskajaPairDef}[def]}
  \proofstepwidestarbreak[3]{u\in D\cup\{\mathbf0\}\dsep z\in T
    \vdash u\MogDiamond z=\mathbf0=z\MogDiamond u}{%
    \FormulaRefAuto{MogiljanskajaPairDef}[def]}
  \proofstepwidestarbreak[1,4]{x,y,z\in S\vdash
    (x\MogStar y)\MogStar z=\mathbf0
      =x\MogStar(y\MogStar z)}{%
    \text{\normalfont zweimal Nullabsorption}}
  \proofstepwidestarbreak[2,5]{x,y,z\in T\vdash
    (x\MogDiamond y)\MogDiamond z=\mathbf0
      =x\MogDiamond(y\MogDiamond z)}{%
    \text{\normalfont zweimal Nullabsorption}}
  \proofstepwidestar[6,7]{%
    \SemigroupAlg{S}{\MogStar}\land
    \SemigroupAlg{T}{\MogDiamond}}{%
    \FormulaRefAuto{SemigroupDef}[def]}
\end{tabproofwide}

\FormulaThmDeltaK[Argument (M3): Unendlichkeit beider Träger]{%
  \neg\Finite{S}\land\neg\Finite{T}%
}{MogiljanskajaArgumentM3Infinite}{%
  \DeltaRow{Mengen}{L\dsep S\dsep T}
  \DeltaRow{Funktionen}{a\dsep H}
}
\begin{proof}[Metasprachliche Fassung]
Die Abbildung \(i\mapsto a_i\) ist nach
\FormulaRefAuto{MogiljanskajaAElementsDef}[def] und
\FormulaRefAuto{MogiljanskajaAIndexUnique}[thm] eine Injektion
\(\mathbb N\inj L\). Wegen \(L\subseteq S\subseteq T\), siehe
\FormulaRefAuto{MogiljanskajaPairDef}[def], liefert die Erweiterung des
Zielbereichs gemäß \FormulaRefAuto{InjectiveCodomainExtension}[thm]
zugleich Injektionen \(\mathbb N\inj S\) und \(\mathbb N\inj T\). Eine
endliche Menge kann nach \FormulaRefAuto{FiniteNoNatInjection}[thm] keine
solche Injektion aufnehmen. Folglich sind weder \(S\) noch \(T\) endlich.
\end{proof}

\MogTableProofHeading
\begin{tabproofwide}
  \proofstepwidestar[]{a\colon\mathbb N\inj L}{%
    \FormulaRefAuto{MogiljanskajaAElementsDef}[def],
    \FormulaRefAuto{MogiljanskajaAIndexUnique}[thm]}
  \proofstepwidestar[]{L\subseteq S\subseteq T}{%
    \FormulaRefAuto{MogiljanskajaPairDef}[def]}
  \proofstepwidestar[1,2]{a\colon\mathbb N\inj S}{%
    \FormulaRefAuto{InjectiveCodomainExtension}[thm]}
  \proofstepwidestar[1,2]{a\colon\mathbb N\inj T}{%
    \FormulaRefAuto{InjectiveCodomainExtension}[thm]}
  \proofstepwidestar[]{\Finite{S}\vdash
    \neg\exists H\colon\mathbb N\inj S}{%
    \FormulaRefAuto{FiniteNoNatInjection}[thm]}
  \proofstepwidestar[3,5]{\neg\Finite{S}}{%
    Widerspruch mit dem Zeugen \(a\)}
  \proofstepwidestar[]{\Finite{T}\vdash
    \neg\exists H\colon\mathbb N\inj T}{%
    \FormulaRefAuto{FiniteNoNatInjection}[thm]}
  \proofstepwidestar[4,7]{\neg\Finite{T}}{%
    Widerspruch mit dem Zeugen \(a\)}
  \proofstepwidestar[6,8]{\neg\Finite{S}\land\neg\Finite{T}}{\rAI{6,8}}
\end{tabproofwide}

\FormulaThmDeltaK[Argument (M4): Produktstatus der ausgezeichneten Elemente]{%
  \begin{aligned}[t]
    s\in S\setminus L
      &\vdash\exists x\in S\,\exists y\in S\,(s=x\MogStar y),\\[-2pt]
    x,y\in T&\vdash x\MogDiamond y\neq k
  \end{aligned}
}{MogiljanskajaArgumentM4ProductStatus}{%
  \DeltaRow{Mengen}{L\dsep D\dsep S\dsep T\dsep\mathbf0\dsep k
    \dsep s\dsep x\dsep y\dsep a_i\dsep a_j\dsep b_{ij}
    \dsep i\dsep j}
  \DeltaRow{Funktionen}{\MogStar\dsep\MogDiamond}
}
\begin{proof}[Metasprachliche Fassung]
Sei \(s\in S\setminus L\). Wegen der disjunkten Zerlegung
\(S=L\cup D\cup\{\mathbf0\}\) aus
\FormulaRefAuto{MogiljanskajaPairDef}[def] und
\FormulaRefAuto{MogiljanskajaArgumentM1WellDefined}[thm] liegt \(s\)
entweder in \(D\) oder es gilt \(s=\mathbf0\). Im ersten Fall besitzt
\(s\) nach \FormulaRefAuto{MogiljanskajaLayerDMembership}[thm] die Form
\(b_{ij}\); dann ist
\[
  s=b_{ij}=a_i\MogStar a_j.
\]
Im zweiten Fall genügt \(\mathbf0=\mathbf0\MogStar\mathbf0\). Damit ist
jedes Element von \(S\setminus L\) ein Produkt zweier Elemente aus \(S\).

Dagegen liegt jedes Produkt \(x\MogDiamond y\) nach
\FormulaRefAuto{MogiljanskajaArgumentM2Semigroups}[thm] in
\(D\cup\{\mathbf0\}\). Das Element \(k\) liegt nach
\FormulaRefAuto{MogiljanskajaArgumentM1WellDefined}[thm] außerhalb dieser
Menge. Daher kann \(k\) kein Produkt in \(T\) sein.
\end{proof}

\MogTableProofHeading
\begin{tabproofwide}
  \proofstepwidestar[]{s\in S\setminus L\vdash
    s\in D\lor s=\mathbf0}{%
    \FormulaRefAuto{MogiljanskajaPairDef}[def],
    \FormulaRefAuto{MogiljanskajaArgumentM1WellDefined}[thm]}
  \proofstepwidestar[]{s\in D\vdash
    \exists i,j\in\mathbb N\,(s=b_{ij})}{%
    \FormulaRefAuto{MogiljanskajaLayerDMembership}[thm]}
  \proofstepwidestar[2]{s\in D\vdash
    \exists x\in S\,\exists y\in S\,(s=x\MogStar y)}{%
    Wähle \(x=a_i\) und \(y=a_j\)}
  \proofstepwidestar[]{\mathbf0=\mathbf0\MogStar\mathbf0}{%
    \FormulaRefAuto{MogiljanskajaPairDef}[def]}
  \proofstepwidestarbreak[1,3,4]{s\in S\setminus L\vdash
    \exists x\in S\,\exists y\in S\,(s=x\MogStar y)}{%
    \text{\normalfont Fallunterscheidung}}
  \proofstepwidestar[]{x,y\in T\vdash
    x\MogDiamond y\in D\cup\{\mathbf0\}}{%
    \FormulaRefAuto{MogiljanskajaArgumentM2Semigroups}[thm]}
  \proofstepwidestar[]{k\notin D\cup\{\mathbf0\}}{%
    \FormulaRefAuto{MogiljanskajaArgumentM1WellDefined}[thm]}
  \proofstepwidestar[6,7]{x,y\in T\vdash x\MogDiamond y\neq k}{%
    \text{\normalfont Ausschluss aus dem Produktbereich}}
\end{tabproofwide}

\FormulaThmDeltaK[Argument (M5): Kollaps eines angenommenen Isomorphismus]{%
  \neg\exists f\,\SemigroupIso{f}{S,\MogStar}{T,\MogDiamond}%
}{MogiljanskajaArgumentM5NotIsomorphic}{%
  \DeltaRow{Mengen}{L\dsep D\dsep S\dsep T\dsep\mathbf0\dsep k}
  \DeltaRow{Mengen}{s\dsep p\dsep q\dsep u\dsep v\dsep x\dsep y}
  \DeltaRow{Mengen}{a_i\dsep a_m\dsep a_n\dsep a_0\dsep a_1}
  \DeltaRow{Mengen}{b_{mn}\dsep b_{i0}\dsep b_{i1}\dsep i\dsep m\dsep n}
  \DeltaRow{Funktionen}{f\dsep\MogStar\dsep\MogDiamond}
}
\begin{proof}[Metasprachliche Fassung]
Angenommen, es gäbe einen Halbgruppenisomorphismus
\(f\colon(S,\MogStar)\to(T,\MogDiamond)\). Da \(f\) surjektiv ist, besitzt
\(k\) zunächst ein Urbild \(s\in S\). Die Surjektivität allein besagt noch
nicht, dass dieses Urbild in \(L\) liegt. Nach der Trägerdefinition und der
Disjunktheit aus \FormulaRefAuto{MogiljanskajaArgumentM1WellDefined}[thm]
zerfällt \(S\) in die drei Schichten \(L\), \(D\) und
\(\{\mathbf0\}\).

Angenommen, \(s\notin L\). Dann gilt entweder \(s\in D\) oder
\(s=\mathbf0\). Im ersten Fall hat \(s\) nach
\FormulaRefAuto{MogiljanskajaLayerDMembership}[thm] die Form \(b_{mn}\),
also
\[
  s=b_{mn}=a_m\MogStar a_n.
\]
Im zweiten Fall gilt
\[
  s=\mathbf0=\mathbf0\MogStar\mathbf0.
\]
In beiden Fällen ist \(s\) somit ein Produkt in \(S\), wie es
\FormulaRefAuto{MogiljanskajaArgumentM4ProductStatus}[thm] zusammenfasst.
Für geeignete Produktzeugen \(p,q\in S\) gilt also
\(s=p\MogStar q\). Die Homomorphie von \(f\) würde dann
\[
  k=f(s)=f(p\MogStar q)=f(p)\MogDiamond f(q)
\]
liefern. Da \(f(p),f(q)\in T\) gilt, wäre \(k\) damit ein Produkt in \(T\),
obwohl dies nach demselben Satz unmöglich ist. Die Alternativen \(s\in D\)
und \(s=\mathbf0\) sind daher ausgeschlossen; es muss \(s\in L\) gelten.

Nach \FormulaRefAuto{MogiljanskajaLayerLMembership}[thm] ist somit
\(s=a_i\) für einen Index \(i\), also \(f(a_i)=k\). Mit der
Homomorphieeigenschaft aus \FormulaRefAuto{SemigroupIsoDef}[def] und der
Multiplikationsvorschrift aus \FormulaRefAuto{MogiljanskajaPairDef}[def]
folgt
\[
  f(b_{i0})
    =f(a_i\MogStar a_0)
    =k\MogDiamond f(a_0)
    =\mathbf0,
  \qquad
  f(b_{i1})
    =f(a_i\MogStar a_1)
    =k\MogDiamond f(a_1)
    =\mathbf0.
\]
Anschaulich kollabiert also die ganze von \(a_i\) ausgehende Produktzeile
auf Null, weil \(k\) außerhalb von \(L\) liegt.

Wären \(b_{i0}\) und \(b_{i1}\) gleich, so ergäbe die Indexeindeutigkeit
aus \FormulaRefAuto{MogiljanskajaBIndexUnique}[thm] die Gleichung
\(0=1\). Dies widerspricht
\FormulaRefAuto{PeanoZeroNeOne}[thm]. Damit bildet \(f\) zwei verschiedene
Elemente auf dasselbe Element ab, obwohl ein Isomorphismus injektiv sein
muss. Ein solcher Isomorphismus existiert also nicht.
\end{proof}

\MogTableProofHeading
\begin{tabproofwide}
  \proofstepwidestar[1]{\exists f\,
    \SemigroupIso{f}{S,\MogStar}{T,\MogDiamond}}{\rA}
  \proofstepwidestar[2]{\SemigroupIso{f}{S,\MogStar}{T,\MogDiamond}}{\rA}
  \proofstepwidestar[2]{f\colon S\bij T}{%
    \FormulaRefAuto{SemigroupIsoDef}[def]{2}}
  \proofstepwidestar[2]{x,y\in S\vdash
    f(x\MogStar y)=f(x)\MogDiamond f(y)}{%
    \FormulaRefAuto{SemigroupIsoDef}[def]{2}}
  \proofstepwidestar[3]{\exists s\in S\,(f(s)=k)}{%
    Surjektivität und \(k\in T\)}
  \proofstepwidestar[5]{s\in S\dsep f(s)=k}{\rA}
  \proofstepwidestar[6]{s\notin L\vdash
    s\in D\lor s=\mathbf0}{%
    \FormulaRefAuto{MogiljanskajaPairDef}[def],
    \FormulaRefAuto{MogiljanskajaArgumentM1WellDefined}[thm]}
  \MogProofStepWideStarKeep[7]{s\notin L\vdash
    \exists p\in S\,\exists q\in S\,(s=p\MogStar q)}{%
    \FormulaRefAuto{MogiljanskajaArgumentM4ProductStatus}[thm];
    \text{\normalfont Zeugen im }D\text{\normalfont{-Fall: }}a_m,a_n
    \dsep\text{\normalfont im Nullfall: }\mathbf0,\mathbf0}
  \MogProofStepWideStarKeep[3,4,6,8]{s\notin L\vdash
    \exists u\in T\,\exists v\in T\,(k=u\MogDiamond v)}{%
    \text{\normalfont Wähle }u=f(p)\text{\normalfont{ und }}v=f(q)}
  \proofstepwidestar[]{%
    \neg\exists u\in T\,\exists v\in T\,(k=u\MogDiamond v)}{%
    \FormulaRefAuto{MogiljanskajaArgumentM4ProductStatus}[thm]}
  \proofstepwidestar[9,10]{s\in L}{%
    Ausschluss der Alternative \(s\notin L\)}
  \proofstepwidestar[11]{\exists i\in\mathbb N\,(s=a_i)}{%
    \FormulaRefAuto{MogiljanskajaLayerLMembership}[thm]{11}}
  \proofstepwidestar[12]{i\in\mathbb N\dsep s=a_i}{\rA}
  \proofstepwidestar[4,6,13]{f(b_{i0})=\mathbf0}{%
    \ensuremath{b_{i0}=a_i\MogStar a_0\dsep
      k\MogDiamond f(a_0)=\mathbf0}}
  \proofstepwidestar[4,6,13]{f(b_{i1})=\mathbf0}{%
    \ensuremath{b_{i1}=a_i\MogStar a_1\dsep
      k\MogDiamond f(a_1)=\mathbf0}}
  \proofstepwidestar[13]{b_{i0}\neq b_{i1}}{%
    \FormulaRefAuto{MogiljanskajaBIndexUnique}[thm],
    \FormulaRefAuto{PeanoZeroNeOne}[thm]}
  \proofstepwidestar[3,14,15,16]{\bot}{%
    Injektivität von \(f\)}
  \proofstepwidestar[2,6]{\bot}{\rEE{12,13,17}}
  \proofstepwidestar[2]{\bot}{\rEE{5,6,18}}
  \proofstepwidestar[1]{\bot}{\rEE{1,2,19}}
  \proofstepwidestar[]{\neg\exists f\,
    \SemigroupIso{f}{S,\MogStar}{T,\MogDiamond}}{\rCI{1,20}}
\end{tabproofwide}

\section{Mengenprodukte und Reserve}

\FormulaThmDeltaK[Argument (M6): Elementweise Produktklassifikation]{%
  \begin{aligned}[t]
    X,Y\in\PowNE{S}&\vdash
      \bigl(z\in X\MogStar_{\mathcal P}Y
        \leftrightarrow z\in R(X,Y)\cup\varepsilon(X,Y)\bigr),\\[-2pt]
    X,Y\in\PowNE{T}&\vdash
      \bigl(z\in X\MogDiamond_{\mathcal P}Y
        \leftrightarrow z\in R(X,Y)\cup\varepsilon(X,Y)\bigr),\\[2pt]
    X,Y\in\PowNE{S}&\vdash
      X\MogStar_{\mathcal P}Y=R(X,Y)\cup\varepsilon(X,Y),\\[-2pt]
    X,Y\in\PowNE{T}&\vdash
      X\MogDiamond_{\mathcal P}Y=R(X,Y)\cup\varepsilon(X,Y)
  \end{aligned}
}{MogiljanskajaArgumentM6ProductClassification}{%
  \DeltaRow{Mengen}{L\dsep S\dsep T\dsep X\dsep Y
    \dsep R(X,Y)\dsep\varepsilon(X,Y)}
  \DeltaRow{Funktionen}{\MogStar_{\mathcal P}
    \dsep\MogDiamond_{\mathcal P}}
}
\begin{proof}[Metasprachliche Fassung]
Betrachten wir zunächst \(X,Y\in\PowNE{S}\). Nach
\FormulaRefAuto{MogiljanskajaArgumentM2Semigroups}[thm] und
\FormulaRefAuto{PowerSemigroupProductWitnesses}[thm] stammt jedes
\(z\in X\MogStar_{\mathcal P}Y\) von einem Faktorenpaar \(x\in X\),
\(y\in Y\). Liegen beide Faktoren in \(L\), so haben sie eindeutig die
Form \(x=a_i\), \(y=a_j\), und ihr Produkt ist \(b_{ij}\in R(X,Y)\).
Liegt wenigstens ein Faktor außerhalb von \(L\), so ist das Produkt nach
\FormulaRefAuto{MogiljanskajaPairDef}[def] gleich \(\mathbf0\); zugleich
gilt dann nach \FormulaRefAuto{MogiljanskajaProductPartsDef}[def]
\(\varepsilon(X,Y)=\{\mathbf0\}\). Damit liegt jedes Produkt in
\(R(X,Y)\cup\varepsilon(X,Y)\).

Umgekehrt besitzt jedes \(b_{ij}\in R(X,Y)\) die Produktzeugen
\(a_i\in X\) und \(a_j\in Y\), liegt also in
\(X\MogStar_{\mathcal P}Y\). Ist \(\mathbf0\in\varepsilon(X,Y)\), so
sind nicht beide Mengen Teilmengen von \(L\). In mindestens einer von ihnen
gibt es daher einen Faktor außerhalb von \(L\); aus der anderen,
nichtleeren Menge wählen wir einen beliebigen Faktor. Ihr Produkt ist
\(\mathbf0\). Somit gilt elementweise
\[
  z\in X\MogStar_{\mathcal P}Y
  \quad\Longleftrightarrow\quad
  z\in R(X,Y)\cup\varepsilon(X,Y),
\]
und durch Extensionalität folgt die behauptete Mengengleichheit.

Für \(X,Y\in\PowNE{T}\) ist das Argument wortgleich: Die Operation
\(\MogDiamond\) verwendet dieselbe Unterscheidung zwischen zwei
\(L\)-Faktoren und einem Faktorenpaar mit mindestens einem Element außerhalb
von \(L\). Daher erhält man ebenso
\[
  z\in X\MogDiamond_{\mathcal P}Y
  \quad\Longleftrightarrow\quad
  z\in R(X,Y)\cup\varepsilon(X,Y)
\]
und damit
\(X\MogDiamond_{\mathcal P}Y=R(X,Y)\cup\varepsilon(X,Y)\).
\end{proof}

\MogTableProofHeading
\begin{tabproofwide}
  \proofstepwidestarbreak[]{X,Y\in\PowNE{S}\dsep
    z\in X\MogStar_{\mathcal P}Y\vdash
    \exists x\in X\,\exists y\in Y\,(z=x\MogStar y)}{%
    \FormulaRefAuto{PowerSemigroupProductWitnesses}[thm]}
  \proofstepwidestarbreak[]{x\in X\cap L\dsep y\in Y\cap L\vdash
    x\MogStar y\in R(X,Y)}{%
    \FormulaRefAuto{MogiljanskajaLayerLMembership}[thm],
    \FormulaRefAuto{MogiljanskajaProductPartsDef}[def]}
  \proofstepwidestarbreak[]{x\in X\dsep y\in Y\dsep
    (x\notin L\lor y\notin L)\vdash
    x\MogStar y=\mathbf0\in\varepsilon(X,Y)}{%
    \FormulaRefAuto{MogiljanskajaPairDef}[def],
    \FormulaRefAuto{MogiljanskajaProductPartsDef}[def]}
  \proofstepwidestarbreak[1,2,3]{X,Y\in\PowNE{S}\dsep
    z\in X\MogStar_{\mathcal P}Y\vdash
    z\in R(X,Y)\cup\varepsilon(X,Y)}{%
    \text{\normalfont Fallunterscheidung für die Produktzeugen}}
  \proofstepwidestarbreak[]{z\in R(X,Y)\vdash
    z\in X\MogStar_{\mathcal P}Y}{%
    \FormulaRefAuto{MogiljanskajaProductPartsDef}[def],
    \FormulaRefAuto{MogiljanskajaPairDef}[def]}
  \proofstepwidestarbreak[]{%
    \begin{aligned}[t]
      X\neq\varnothing\dsep Y\neq\varnothing\dsep
        &\neg(X\subseteq L\land Y\subseteq L)\vdash{}\\[-2pt]
      &\exists x\in X\,\exists y\in Y\,(x\notin L\lor y\notin L)
    \end{aligned}}{%
    \text{\normalfont Fälle }X\nsubseteq L\text{\normalfont{ und }}
      Y\nsubseteq L\text{\normalfont{; der andere Faktor ist nichtleer}}}
  \proofstepwidestarbreak[6]{X,Y\in\PowNE{S}\dsep
    z\in\varepsilon(X,Y)\vdash
    z\in X\MogStar_{\mathcal P}Y}{%
    \FormulaRefAuto{MogiljanskajaProductPartsDef}[def],
    \FormulaRefAuto{MogiljanskajaPairDef}[def]}
  \proofstepwidestarbreak[5,7]{X,Y\in\PowNE{S}\dsep
    z\in R(X,Y)\cup\varepsilon(X,Y)\vdash
    z\in X\MogStar_{\mathcal P}Y}{%
    \text{\normalfont Vereinigungsfall}}
  \proofstepwidestarbreak[4,8]{X,Y\in\PowNE{S}\vdash
    \bigl(z\in X\MogStar_{\mathcal P}Y
      \leftrightarrow z\in R(X,Y)\cup\varepsilon(X,Y)\bigr)}{%
    \text{\normalfont Äquivalenzeinführung}}
  \proofstepwidestarbreak[]{X,Y\in\PowNE{T}\dsep
    z\in X\MogDiamond_{\mathcal P}Y\vdash
    \exists x\in X\,\exists y\in Y\,(z=x\MogDiamond y)}{%
    \FormulaRefAuto{PowerSemigroupProductWitnesses}[thm]}
  \proofstepwidestarbreak[]{x\in X\cap L\dsep y\in Y\cap L\vdash
    x\MogDiamond y\in R(X,Y)}{%
    \FormulaRefAuto{MogiljanskajaLayerLMembership}[thm],
    \FormulaRefAuto{MogiljanskajaProductPartsDef}[def]}
  \proofstepwidestarbreak[]{x\in X\dsep y\in Y\dsep
    (x\notin L\lor y\notin L)\vdash
    x\MogDiamond y=\mathbf0\in\varepsilon(X,Y)}{%
    \FormulaRefAuto{MogiljanskajaPairDef}[def],
    \FormulaRefAuto{MogiljanskajaProductPartsDef}[def]}
  \proofstepwidestarbreak[10,11,12]{X,Y\in\PowNE{T}\dsep
    z\in X\MogDiamond_{\mathcal P}Y\vdash
    z\in R(X,Y)\cup\varepsilon(X,Y)}{%
    \text{\normalfont Fallunterscheidung für die Produktzeugen}}
  \proofstepwidestarbreak[]{z\in R(X,Y)\vdash
    z\in X\MogDiamond_{\mathcal P}Y}{%
    \FormulaRefAuto{MogiljanskajaProductPartsDef}[def],
    \FormulaRefAuto{MogiljanskajaPairDef}[def]}
  \proofstepwidestarbreak[6]{X,Y\in\PowNE{T}\dsep
    z\in\varepsilon(X,Y)\vdash
    z\in X\MogDiamond_{\mathcal P}Y}{%
    \FormulaRefAuto{MogiljanskajaProductPartsDef}[def],
    \FormulaRefAuto{MogiljanskajaPairDef}[def]}
  \proofstepwidestarbreak[14,15]{X,Y\in\PowNE{T}\dsep
    z\in R(X,Y)\cup\varepsilon(X,Y)\vdash
    z\in X\MogDiamond_{\mathcal P}Y}{%
    \text{\normalfont Vereinigungsfall}}
  \proofstepwidestarbreak[13,16]{X,Y\in\PowNE{T}\vdash
    \bigl(z\in X\MogDiamond_{\mathcal P}Y
      \leftrightarrow z\in R(X,Y)\cup\varepsilon(X,Y)\bigr)}{%
    \text{\normalfont Äquivalenzeinführung}}
  \proofstepwidestarbreak[9]{X,Y\in\PowNE{S}\vdash
    X\MogStar_{\mathcal P}Y=R(X,Y)\cup\varepsilon(X,Y)}{%
    \text{\normalfont Extensionalität}}
  \proofstepwidestarbreak[17]{X,Y\in\PowNE{T}\vdash
    X\MogDiamond_{\mathcal P}Y=R(X,Y)\cup\varepsilon(X,Y)}{%
    \text{\normalfont Extensionalität}}
\end{tabproofwide}

\FormulaThmDeltaK[Argument (M7): Voraussetzungen der Reserveabsorption]{%
  \begin{aligned}[t]
    &P\cap U=\varnothing\dsep c_2\notin P\cup U\dsep
      \{k\}\cap D=\varnothing,\\[-2pt]
    &\sigma\colon U\bij V\dsep\theta\colon U\bij D
  \end{aligned}
}{MogiljanskajaArgumentM7AbsorptionHypotheses}{%
  \DeltaRow{Mengen}{D\dsep U\dsep V\dsep P\dsep c_2\dsep k}
  \DeltaRow{Funktionen}{\sigma\dsep\theta}
}
\begin{proof}[Metasprachliche Fassung]
Nach \FormulaRefAuto{MogiljanskajaReserveCoreFacts}[thm] sind der feste
Kern \(P\) und die Reserve \(U\) disjunkt, und \(c_2\) gehört nicht zu
\(P\). Zugleich liegt \(c_2\) nach
\FormulaRefAuto{MogiljanskajaReserveMarkerFacts}[thm] auch nicht in
\(U\). Damit liegt \(c_2\) außerhalb der gesamten Menge \(P\cup U\).
Das Element \(k\) gehört wegen seines eigenen Tags nicht zu \(D\), wie
bereits \FormulaRefAuto{MogiljanskajaArgumentM1WellDefined}[thm] zeigt;
daher ist auch \(\{k\}\cap D=\varnothing\).

Schließlich stehen die benötigten Bijektionen bereits zur Verfügung:
\(\sigma\colon U\bij V\) folgt aus
\FormulaRefAuto{MogiljanskajaSigmaBijective}[thm] und
\(\theta\colon U\bij D\) aus
\FormulaRefAuto{MogiljanskajaThetaBijective}[thm].
\end{proof}

\MogTableProofHeading
\begin{tabproofwide}
  \proofstepwidestar[]{P\cap U=\varnothing\dsep c_2\notin P}{%
    \FormulaRefAuto{MogiljanskajaReserveCoreFacts}[thm]}
  \proofstepwidestar[]{c_2\notin U}{%
    \FormulaRefAuto{MogiljanskajaReserveMarkerFacts}[thm]}
  \proofstepwidestar[1,2]{c_2\notin P\cup U}{%
    \text{\normalfont Vereinigungskriterium}}
  \proofstepwidestar[]{k\notin D}{%
    \FormulaRefAuto{MogiljanskajaArgumentM1WellDefined}[thm]}
  \proofstepwidestar[4]{\{k\}\cap D=\varnothing}{%
    \text{\normalfont Einermengenkriterium}}
  \proofstepwidestar[]{\sigma\colon U\bij V}{%
    \FormulaRefAuto{MogiljanskajaSigmaBijective}[thm]}
  \proofstepwidestar[]{\theta\colon U\bij D}{%
    \FormulaRefAuto{MogiljanskajaThetaBijective}[thm]}
  \proofstepwidestarbreak[1,3,5,6,7]{%
    \begin{aligned}[t]
      &P\cap U=\varnothing\dsep c_2\notin P\cup U\dsep
        \{k\}\cap D=\varnothing,\\[-2pt]
      &\sigma\colon U\bij V\dsep\theta\colon U\bij D
    \end{aligned}}{%
    \text{\normalfont Bündelung}}
\end{tabproofwide}

\pagebreak[2]
\FormulaThmDeltaK[Argument (M8): Disjunktheit und Absorption der Kernfamilien]{%
  \mathcal C\subseteq\powerset(D)\dsep
  \mathcal C\cap\mathcal E=\varnothing\dsep
  \psi\colon\mathcal C\bij\mathcal C\cup\mathcal E%
}{MogiljanskajaArgumentM8PsiBijective}{%
  \DeltaRow{Mengen}{D\dsep D'\dsep U\dsep V\dsep P\dsep c_2\dsep k
    \dsep\mathcal C\dsep\mathcal C_0\dsep\mathcal C_1
    \dsep\mathcal E\dsep H}
  \DeltaRow{Funktionen}{\sigma\dsep\theta\dsep
    \kappa_0\dsep\kappa_1\dsep\psi}
}
\begin{proof}[Metasprachliche Fassung]
Nach \FormulaRefAuto{MogiljanskajaDPrimeDef}[def] gilt
\(V=U\cup\{c_2\}\) und \(P\cup V=D'\); außerdem ist
\(D'\subseteq D\) nach
\FormulaRefAuto{MogiljanskajaReserveMarkerFacts}[thm]. Jedes Element von
\(\mathcal C=\CoreFamily{P}{V}\) ist also eine Teilmenge von \(D'\) und
damit von \(D\), vgl.
\FormulaRefAuto{CoreFamilySubsetPowerset}[thm]. Dagegen enthält jede
Menge aus \(\mathcal E=\CoreFamily{\{k\}}{D}\) den Punkt \(k\), nach
\FormulaRefAuto{CoreFamilyMembership}[thm]. Da \(k\notin D\) nach
\FormulaRefAuto{MogiljanskajaArgumentM7AbsorptionHypotheses}[thm], kann
eine solche Menge nicht zugleich zu \(\mathcal C\) gehören. Somit sind
\(\mathcal C\) und \(\mathcal E\) disjunkt.

Anschaulich zerlegt der Marker \(c_2\) die Familie \(\mathcal C\) in
die Mengen ohne \(c_2\), also \(\mathcal C_0\), und die Mengen mit
\(c_2\), also \(\mathcal C_1\). Auf dem ersten Teil transportiert
\(\kappa_0\) die freien \(U\)-Koordinaten mittels \(\sigma\), auf dem
zweiten transportiert \(\kappa_1\) sie mittels \(\theta\) und ersetzt
den Kern \(P\cup\{c_2\}\) durch \(\{k\}\). Die dafür nötigen
Disjunktheits- und Bijektivitätsvoraussetzungen sind genau die Aussagen
von \FormulaRefAuto{MogiljanskajaArgumentM7AbsorptionHypotheses}[thm].
Mit den Definitionen aus
\FormulaRefAuto{MogiljanskajaReserveMapsDef}[def] liefert daher
\FormulaRefAuto{CoreFamilyAbsorptionBijective}[thm] unmittelbar
\[
  \psi\colon\mathcal C\bij\mathcal C\cup\mathcal E.
\]
\end{proof}

\MogTableProofHeading
\begin{tabproofwide}
  \proofstepwidestar[]{V=U\cup\{c_2\}\dsep P\cup V=D'\subseteq D}{%
    \FormulaRefAuto{MogiljanskajaDPrimeDef}[def],
    \FormulaRefAuto{MogiljanskajaReserveMarkerFacts}[thm]}
  \proofstepwidestar[]{\mathcal C=\CoreFamily{P}{V}}{%
    \FormulaRefAuto{MogiljanskajaReserveMapsDef}[def]}
  \proofstepwidestar[1,2]{\mathcal C\subseteq\powerset(D')}{%
    \FormulaRefAuto{CoreFamilySubsetPowerset}[thm]}
  \proofstepwidestar[1,3]{\mathcal C\subseteq\powerset(D)}{%
    \text{\normalfont Monotonie der Potenzmenge}}
  \proofstepwidestar[]{%
    \begin{aligned}[t]
      \mathcal C_0&=\CoreFamily{P}{U},&
      \mathcal C_1&=\CoreFamily{P\cup\{c_2\}}{U},\\[-2pt]
      \kappa_0&=\CoreTransport{P}{P}{\sigma},&
      \kappa_1&=\CoreTransport{P\cup\{c_2\}}{\{k\}}{\theta},\\[-2pt]
      \mathcal E&=\CoreFamily{\{k\}}{D},&
      \psi&=\CaseFun{\mathcal C_0}{\mathcal C_1}{\kappa_0}{\kappa_1}
    \end{aligned}}{%
    \FormulaRefAuto{MogiljanskajaReserveMapsDef}[def]}
  \proofstepwidestar[5]{H\in\mathcal E\vdash k\in H}{%
    \FormulaRefAuto{CoreFamilyMembership}[thm]}
  \proofstepwidestar[]{k\notin D}{%
    \FormulaRefAuto{MogiljanskajaArgumentM7AbsorptionHypotheses}[thm]}
  \proofstepwidestar[4,6,7]{\mathcal C\cap\mathcal E=\varnothing}{%
    Kern \(k\) liegt außerhalb von \(D\)}
  \proofstepwidestar[]{%
    \begin{aligned}[t]
      &P\cap U=\varnothing\dsep c_2\notin P\cup U\dsep
        \{k\}\cap D=\varnothing,\\[-2pt]
      &\sigma\colon U\bij V\dsep\theta\colon U\bij D
    \end{aligned}}{%
    \FormulaRefAuto{MogiljanskajaArgumentM7AbsorptionHypotheses}[thm]}
  \proofstepwidestar[1,5,8,9]{\psi\colon\mathcal C\bij
    \mathcal C\cup\mathcal E}{%
    \FormulaRefAuto{CoreFamilyAbsorptionBijective}[thm]}
\end{tabproofwide}

\FormulaThmDeltaK[Argument (M9): Zielzerlegung und Bijektivität von \(\eta\)]{%
  \begin{aligned}[t]
    &\powerset(D\cup\{k\})
      =\mathcal O\cup\mathcal C\cup\mathcal E,\\[-2pt]
    &\mathcal O\cap\mathcal C=\mathcal O\cap\mathcal E
      =\mathcal C\cap\mathcal E=\varnothing,\\[-2pt]
    &\eta\colon\powerset(D)\bij\powerset(D\cup\{k\})
  \end{aligned}
}{MogiljanskajaArgumentM9EtaBijective}{%
  \DeltaRow{Mengen}{D\dsep P\dsep k\dsep\mathcal O\dsep\mathcal C
    \dsep\mathcal E\dsep H}
  \DeltaRow{Funktionen}{\psi\dsep\eta}
}
\begin{proof}[Metasprachliche Fassung]
Nach \FormulaRefAuto{CoreFamilyBooleanInterval}[thm] und
\FormulaRefAuto{MogiljanskajaReserveMapsDef}[def] gilt
\[
  \mathcal E=\mathfrak I[\{k\},D\cup\{k\}].
\]
Weil \(k\notin D\),
lässt sich jede Teilmenge von \(D\cup\{k\}\) anschaulich danach
sortieren, ob sie den neuen Punkt \(k\) enthält: Ohne \(k\) ist sie eine
Teilmenge von \(D\), mit \(k\) gehört sie zu \(\mathcal E\). Daher
\[
  \powerset(D\cup\{k\})=\powerset(D)\cup\mathcal E,
\]
und die beiden Teile sind disjunkt. Andererseits zerlegt die Definition
\(\mathcal O=\powerset(D)\setminus\mathcal C\) die Potenzmenge von \(D\)
in die disjunkten Teile \(\mathcal O\) und \(\mathcal C\). Zusammen mit
\(\mathcal C\cap\mathcal E=\varnothing\) aus
\FormulaRefAuto{MogiljanskajaArgumentM8PsiBijective}[thm] ergibt dies die
behauptete paarweise disjunkte Zerlegung
\[
  \powerset(D\cup\{k\})=\mathcal O\cup\mathcal C\cup\mathcal E.
\]

Es bleibt nur die Abbildung zu lesen: Auf \(\mathcal O\) ist \(\eta\)
die Identität, während sie auf \(\mathcal C\) mit \(\psi\) bijektiv auf
\(\mathcal C\cup\mathcal E\) abbildet. Dabei ist
\(\psi\colon\mathcal C\bij\mathcal C\cup\mathcal E\) durch
\FormulaRefAuto{MogiljanskajaArgumentM8PsiBijective}[thm] gegeben. Ferner
ist \(P\neq\varnothing\), da \(c_0\in P\) nach
\FormulaRefAuto{MogiljanskajaDPrimeDef}[def]; deshalb gehört die leere
Menge nach \FormulaRefAuto{CoreFamilyMembership}[thm] nicht zu
\(\mathcal C\). Alle Voraussetzungen von
\FormulaRefAuto{ProtectedPowerSetExtension}[thm] sind somit erfüllt, und
mit \FormulaRefAuto{MogiljanskajaReserveMapsDef}[def] folgt
\(\eta\colon\powerset(D)\bij\powerset(D\cup\{k\})\).
\end{proof}

\MogTableProofHeading
\begin{tabproofwide}
  \proofstepwidestarbreak[]{\mathcal E
    =\mathfrak I[\{k\},D\cup\{k\}]}{%
    \FormulaRefAuto{CoreFamilyBooleanInterval}[thm],
    \FormulaRefAuto{MogiljanskajaReserveMapsDef}[def]}
  \proofstepwidestar[]{k\notin D}{%
    \FormulaRefAuto{MogiljanskajaArgumentM7AbsorptionHypotheses}[thm]}
  \proofstepwidestar[]{\mathcal C\subseteq\powerset(D)}{%
    \FormulaRefAuto{MogiljanskajaArgumentM8PsiBijective}[thm]}
  \proofstepwidestar[]{P\neq\varnothing}{%
    \FormulaRefAuto{MogiljanskajaDPrimeDef}[def]}
  \proofstepwidestarbreak[4]{\varnothing\notin\mathcal C}{%
    \FormulaRefAuto{CoreFamilyMembership}[thm],
    \FormulaRefAuto{MogiljanskajaReserveMapsDef}[def]}
  \proofstepwidestar[]{\psi\colon\mathcal C\bij
    \mathcal C\cup\mathcal E}{%
    \FormulaRefAuto{MogiljanskajaArgumentM8PsiBijective}[thm]}
  \proofstepwidestarbreak[]{\mathcal O=\powerset(D)\setminus\mathcal C\dsep
    \eta=\CaseFun{\mathcal O}{\mathcal C}{\Id_{\mathcal O}}{\psi}}{%
    \FormulaRefAuto{MogiljanskajaReserveMapsDef}[def]}
  \proofstepwidestarbreak[1,2,3,5,6,7]{%
    \eta\colon\powerset(D)\bij\powerset(D\cup\{k\})}{%
    \FormulaRefAuto{ProtectedPowerSetExtension}[thm]}
  \proofstepwidestarbreak[1]{H\in\mathcal E\leftrightarrow
    H\subseteq D\cup\{k\}\land k\in H}{%
    \FormulaRefAuto{CoreFamilyMembership}[thm]}
  \proofstepwidestarbreak[]{H\subseteq D\cup\{k\}\dsep k\notin H
    \vdash H\subseteq D}{%
    \text{\normalfont Abspaltung des einzigen neuen Punktes}}
  \proofstepwidestarbreak[9,10]{\powerset(D\cup\{k\})
    =\powerset(D)\cup\mathcal E}{%
    \text{\normalfont Fall }k\in H
      \text{\normalfont{ oder }}k\notin H}
  \proofstepwidestarbreak[]{\powerset(D)=\mathcal O\cup\mathcal C\dsep
    \mathcal O\cap\mathcal C=\varnothing}{%
    \FormulaRefAuto{MogiljanskajaReserveMapsDef}[def]}
  \proofstepwidestarbreak[11,12]{\powerset(D\cup\{k\})
    =\mathcal O\cup\mathcal C\cup\mathcal E}{%
    \text{\normalfont Einsetzen}}
  \proofstepwidestar[2,9,12]{\mathcal O\cap\mathcal E=\varnothing}{%
    Mengen aus \(\mathcal O\) enthalten \(k\) nicht}
  \proofstepwidestar[]{\mathcal C\cap\mathcal E=\varnothing}{%
    \FormulaRefAuto{MogiljanskajaArgumentM8PsiBijective}[thm]}
  \proofstepwidestarbreak[12,14,15]{%
    \mathcal O\cap\mathcal C=\mathcal O\cap\mathcal E
      =\mathcal C\cap\mathcal E=\varnothing}{%
    \text{\normalfont paarweise Disjunktheit}}
\end{tabproofwide}

\FormulaThmDeltaK[Argument (M10): Nulltreue und Fixierung außerhalb der Reserve]{%
  \eta(\varnothing)=\varnothing\dsep
  \forall K\in\powerset(D)\,
    \bigl(K\notin\mathcal C\rightarrow\eta(K)=K\bigr)%
}{MogiljanskajaArgumentM10EtaFixedPart}{%
  \DeltaRow{Mengen}{D\dsep P\dsep c_0\dsep\mathcal C
    \dsep\mathcal O\dsep K}
  \DeltaRow{Funktionen}{\eta}
}
\begin{proof}[Metasprachliche Fassung]
Wegen \(c_0\in P\), siehe
\FormulaRefAuto{MogiljanskajaDPrimeDef}[def], ist der feste Kern \(P\)
nicht leer. Eine Menge aus \(\mathcal C=\CoreFamily{P}{V}\) muss aber
den ganzen Kern \(P\) enthalten, nach
\FormulaRefAuto{CoreFamilyMembership}[thm]. Folglich liegt
\(\varnothing\) nicht in \(\mathcal C\), sondern im Komplement
\(\mathcal O=\powerset(D)\setminus\mathcal C\). Auf diesem Teil ist
\(\eta\) definitionsgemäß die Identität; daher gilt
\(\eta(\varnothing)=\varnothing\).

Dasselbe Argument gilt ohne Besonderheit der leeren Menge: Ist
\(K\in\powerset(D)\) und \(K\notin\mathcal C\), so ist
\(K\in\mathcal O\). Der Identitätszweig der in
\FormulaRefAuto{MogiljanskajaReserveMapsDef}[def] festgelegten
Fallfunktion liefert dann \(\eta(K)=K\). Da \(K\) beliebig war, folgt
die behauptete Allaussage.
\end{proof}

\MogTableProofHeading
\begin{tabproofwide}
  \proofstepwidestar[]{c_0\in P}{%
    \FormulaRefAuto{MogiljanskajaDPrimeDef}[def]}
  \proofstepwidestar[1]{P\neq\varnothing}{%
    Nichtleerheitszeuge \(c_0\)}
  \proofstepwidestar[2]{\varnothing\notin\mathcal C}{%
    \FormulaRefAuto{CoreFamilyMembership}[thm],
    \FormulaRefAuto{MogiljanskajaReserveMapsDef}[def]}
  \proofstepwidestar[3]{\varnothing\in\mathcal O}{%
    \FormulaRefAuto{MogiljanskajaReserveMapsDef}[def]}
  \proofstepwidestar[4]{\eta(\varnothing)=\varnothing}{%
    \FormulaRefAuto{MogiljanskajaReserveMapsDef}[def]}
  \proofstepwidestar[]{K\in\powerset(D)\dsep K\notin\mathcal C
    \vdash K\in\mathcal O}{%
    \FormulaRefAuto{MogiljanskajaReserveMapsDef}[def]}
  \proofstepwidestar[6]{K\in\powerset(D)\dsep K\notin\mathcal C
    \vdash\eta(K)=K}{%
    \FormulaRefAuto{MogiljanskajaReserveMapsDef}[def]}
  \proofstepwidestarbreak[5,7]{\eta(\varnothing)=\varnothing\dsep
    \forall K\in\powerset(D)\,
      (K\notin\mathcal C\rightarrow\eta(K)=K)}{%
    \text{\normalfont Generalisierung}}
\end{tabproofwide}

\FormulaThmDeltaK[Argument (M11): Kein Produktrechteck liegt in der Reserve]{%
  X,Y\in\powerset(T)\vdash R(X,Y)\notin\mathcal C%
}{MogiljanskajaArgumentM11RectangleProtected}{%
  \DeltaRow{Mengen}{L\dsep D'\dsep T\dsep P\dsep\mathcal C
    \dsep X\dsep Y\dsep R(X,Y)\dsep a_0\dsep a_1
    \dsep b_{00}\dsep b_{11}\dsep b_{01}}
}
\begin{proof}[Metasprachliche Fassung]
Seien \(X\) und \(Y\) Teilmengen von \(T\). Nehmen wir zum Widerspruch an,
dass \(R(X,Y)\in\mathcal C\) gilt. Das Elementkriterium
\FormulaRefAuto{CoreFamilyMembership}[thm] gibt zusammen mit
\FormulaRefAuto{MogiljanskajaReserveMapsDef}[def] und
\FormulaRefAuto{MogiljanskajaDPrimeDef}[def]
\[
  P\subseteq R(X,Y)\subseteq D'.
\]
Da \(P=\{b_{00},b_{11}\}\), liegen insbesondere \(b_{00}\) und
\(b_{11}\) im Rechteck \(R(X,Y)\). Nach der Definition von \(R\) in
\FormulaRefAuto{MogiljanskajaProductPartsDef}[def] und der Eindeutigkeit
der Indizes aus \FormulaRefAuto{MogiljanskajaBIndexUnique}[thm] bedeutet
dies, dass sowohl \(a_0\) als auch \(a_1\) in \(X\cap L\) und in
\(Y\cap L\) liegen. Das Rechteck enthält daher auch die gemischte Ecke
\(b_{01}\).

Aus \(R(X,Y)\subseteq D'\) würde nun \(b_{01}\in D'\) folgen. Nach
\FormulaRefAuto{MogiljanskajaReserveMarkerFacts}[thm] gilt andererseits
\[
  b_{01}\notin D'.
\]
Das ist der gesuchte Widerspruch. Also kann \(R(X,Y)\) nicht in
\(\mathcal C\) liegen.
\end{proof}

\MogTableProofHeading
\begin{tabproofwide}
  \proofstepwidestar[1]{X,Y\in\powerset(T)}{\rA}
  \proofstepwidestar[2]{R(X,Y)\in\mathcal C}{\rA}
  \proofstepwidestarbreak[2]{P\subseteq R(X,Y)\dsep R(X,Y)\subseteq D'}{%
    \FormulaRefAuto{CoreFamilyMembership}[thm]{2},
    \FormulaRefAuto{MogiljanskajaReserveMapsDef}[def],
    \FormulaRefAuto{MogiljanskajaDPrimeDef}[def]}
  \proofstepwidestar[3]{b_{00},b_{11}\in R(X,Y)}{%
    \FormulaRefAuto{MogiljanskajaDPrimeDef}[def]{3}}
  \proofstepwidestar[4]{a_0\in X\cap L\dsep a_0\in Y\cap L}{%
    \FormulaRefAuto{MogiljanskajaProductPartsDef}[def],
    \FormulaRefAuto{MogiljanskajaBIndexUnique}[thm]}
  \proofstepwidestar[4]{a_1\in X\cap L\dsep a_1\in Y\cap L}{%
    \FormulaRefAuto{MogiljanskajaProductPartsDef}[def],
    \FormulaRefAuto{MogiljanskajaBIndexUnique}[thm]}
  \proofstepwidestar[5,6]{b_{01}\in R(X,Y)}{%
    \FormulaRefAuto{MogiljanskajaProductPartsDef}[def]}
  \proofstepwidestar[3,7]{b_{01}\in D'}{\rIE{7,3}}
  \proofstepwidestar[]{b_{01}\notin D'}{%
    \FormulaRefAuto{MogiljanskajaReserveMarkerFacts}[thm]}
  \proofstepwidestar[2]{\bot}{\rBI{8,9}}
  \proofstepwidestar[1]{R(X,Y)\notin\mathcal C}{\rCI{2,10}}
  \proofstepwidestar[]{X,Y\in\powerset(T)\vdash
    R(X,Y)\notin\mathcal C}{\rRI{1,11}}
\end{tabproofwide}

\section{Globale Abbildung und Isomorphieschluss}

\FormulaThmDeltaK[Argument (M12): Schichtzerlegung und globale Bijektion]{%
  \begin{aligned}[t]
    &A_*\cap D=\varnothing\dsep
      A_*\cap(D\cup\{k\})=\varnothing,\\[-2pt]
    &S=A_*\cup D\dsep T=A_*\cup(D\cup\{k\}),\\[-2pt]
    &\Phi\colon\PowNE{S}\bij\PowNE{T}
  \end{aligned}
}{MogiljanskajaArgumentM12PhiBijective}{%
  \DeltaRow{Mengen}{A_*\dsep L\dsep D\dsep S\dsep T
    \dsep\mathbf0\dsep k}
  \DeltaRow{Funktionen}{\eta\dsep\Phi}
}
\begin{proof}[Metasprachliche Fassung]
Anschaulich besteht der Träger aus einer festen Schicht
\(A_*=L\cup\{\mathbf0\}\) und einer beweglichen Schicht \(D\). Beim
Übergang von \(S\) zu \(T\) kommt nur in der zweiten Schicht das Element
\(k\) hinzu. Aus \FormulaRefAuto{MogiljanskajaPhiDef}[def] und der
paarweisen Disjunktheit aus
\FormulaRefAuto{MogiljanskajaArgumentM1WellDefined}[thm] folgen
\[
  A_*\cap D=\varnothing,
  \qquad
  A_*\cap(D\cup\{k\})=\varnothing.
\]
Die Trägerdefinition in \FormulaRefAuto{MogiljanskajaPairDef}[def] ergibt
zugleich
\[
  S=A_*\cup D,
  \qquad
  T=A_*\cup(D\cup\{k\}).
\]

Nach \FormulaRefAuto{MogiljanskajaArgumentM9EtaBijective}[thm] bildet
\(\eta\) die Potenzmenge von \(D\) bijektiv auf die Potenzmenge von
\(D\cup\{k\}\) ab, und nach
\FormulaRefAuto{MogiljanskajaArgumentM10EtaFixedPart}[thm] bleibt die leere
Menge fest. Damit sind alle Voraussetzungen von
\FormulaRefAuto{LayerPowerMapNonemptyBijective}[thm] erfüllt: Die feste
Schicht wird unverändert übernommen, während \(\eta\) die zweite Schicht
bijektiv erweitert. Wegen \FormulaRefAuto{MogiljanskajaPhiDef}[def] ist die
dortige Schichtabbildung, eingeschränkt auf die nichtleeren Teilmengen,
genau \(\Phi\). Also gilt \(\Phi\colon\PowNE{S}\bij\PowNE{T}\).
\end{proof}

\MogTableProofHeading
\begin{tabproofwide}
  \proofstepwidestar[]{A_*=L\cup\{\mathbf0\}}{%
    \FormulaRefAuto{MogiljanskajaPhiDef}[def]}
  \proofstepwidestar[1]{A_*\cap D=\varnothing}{%
    \FormulaRefAuto{MogiljanskajaArgumentM1WellDefined}[thm]}
  \proofstepwidestar[1]{A_*\cap(D\cup\{k\})=\varnothing}{%
    \FormulaRefAuto{MogiljanskajaArgumentM1WellDefined}[thm]}
  \proofstepwidestarbreak[1]{S=A_*\cup D\dsep
    T=A_*\cup(D\cup\{k\})}{%
    \FormulaRefAuto{MogiljanskajaPairDef}[def]}
  \proofstepwidestarbreak[]{\eta\colon\powerset(D)\bij
    \powerset(D\cup\{k\})}{%
    \FormulaRefAuto{MogiljanskajaArgumentM9EtaBijective}[thm]}
  \proofstepwidestar[]{\eta(\varnothing)=\varnothing}{%
    \FormulaRefAuto{MogiljanskajaArgumentM10EtaFixedPart}[thm]}
  \proofstepwidestar[2,3,4,5,6]{\Phi\colon\PowNE{S}\bij\PowNE{T}}{%
    \FormulaRefAuto{LayerPowerMapNonemptyBijective}[thm],
    \FormulaRefAuto{MogiljanskajaPhiDef}[def]}
\end{tabproofwide}

\FormulaThmDeltaK[Argument (M13): \(L\)-Invarianz von \(\Phi\)]{%
  \begin{aligned}[t]
    X\in\PowNE{S}\vdash{}&
      \Phi(X)=(X\cap A_*)\cup\eta(X\cap D),\\[-2pt]
    X\in\PowNE{S}\vdash{}&
      \Phi(X)\cap L=X\cap L,\\[-2pt]
    X\in\PowNE{S}\vdash{}&
      \bigl(X\subseteq L\leftrightarrow\Phi(X)\subseteq L\bigr)
  \end{aligned}
}{MogiljanskajaArgumentM13PhiInvariants}{%
  \DeltaRow{Mengen}{A_*\dsep L\dsep D\dsep S\dsep T\dsep X}
  \DeltaRow{Funktionen}{\eta\dsep\Phi}
}
\begin{proof}[Metasprachliche Fassung]
Die Abbildung \(\Phi\) verändert ausschließlich den Anteil einer Menge in
\(D\). Aus der Schichtzerlegung in
\FormulaRefAuto{MogiljanskajaArgumentM12PhiBijective}[thm], der
Bijektivität von \(\eta\) aus
\FormulaRefAuto{MogiljanskajaArgumentM9EtaBijective}[thm] und
\FormulaRefAuto{LayerPowerMapValue}[thm] folgt daher unmittelbar
\[
  \Phi(X)=(X\cap A_*)\cup\eta(X\cap D).
\]
Der Satz \FormulaRefAuto{LayerPowerMapStablePart}[thm] sagt zugleich, dass
der feste Anteil wirklich fest bleibt:
\(\Phi(X)\cap A_*=X\cap A_*\). Da \(L\subseteq A_*\), liefert ein weiterer
Schnitt mit \(L\)
\[
  \Phi(X)\cap L=X\cap L.
\]

Auch die Frage, ob eine Menge vollständig in \(L\) liegt, kann durch die
Verschiebung in der disjunkten \(D\)-Schicht nicht geändert werden. Formal
folgt dies aus \FormulaRefAuto{LayerPowerMapPreservesSubsets}[thm]; dabei
wird außerdem \(\eta(\varnothing)=\varnothing\) aus
\FormulaRefAuto{MogiljanskajaArgumentM10EtaFixedPart}[thm] verwendet. Somit
ist
\[
  X\subseteq L\quad\Longleftrightarrow\quad\Phi(X)\subseteq L.
\]
\end{proof}

\MogTableProofHeading
\begin{tabproofwide}
  \proofstepwidestarbreak[]{A_*\cap D=\varnothing\dsep
    A_*\cap(D\cup\{k\})=\varnothing\dsep
    S=A_*\cup D}{%
    \FormulaRefAuto{MogiljanskajaArgumentM12PhiBijective}[thm]}
  \proofstepwidestarbreak[]{\eta\colon\powerset(D)\bij
    \powerset(D\cup\{k\})}{%
    \FormulaRefAuto{MogiljanskajaArgumentM9EtaBijective}[thm]}
  \proofstepwidestar[]{\eta(\varnothing)=\varnothing}{%
    \FormulaRefAuto{MogiljanskajaArgumentM10EtaFixedPart}[thm]}
  \proofstepwidestar[1]{X\in\PowNE{S}\vdash X\subseteq A_*\cup D}{%
    \text{\normalfont Trägerzerlegung}}
  \proofstepwidestarbreak[1,2,4]{X\in\PowNE{S}\vdash
    \Phi(X)=(X\cap A_*)\cup\eta(X\cap D)}{%
    \FormulaRefAuto{LayerPowerMapValue}[thm],
    \FormulaRefAuto{MogiljanskajaPhiDef}[def]}
  \proofstepwidestarbreak[1,2,4]{X\in\PowNE{S}\vdash
    \Phi(X)\cap A_*=X\cap A_*}{%
    \FormulaRefAuto{LayerPowerMapStablePart}[thm]}
  \proofstepwidestar[]{L\subseteq A_*}{%
    \FormulaRefAuto{MogiljanskajaPhiDef}[def]}
  \proofstepwidestarbreak[6,7]{X\in\PowNE{S}\vdash
    \Phi(X)\cap L=X\cap L}{%
    \text{\normalfont Schnitt mit der Teilschicht }L}
  \proofstepwidestarbreak[1,2,3,4,7]{X\in\PowNE{S}\vdash
    \bigl(X\subseteq L\leftrightarrow\Phi(X)\subseteq L\bigr)}{%
    \FormulaRefAuto{LayerPowerMapPreservesSubsets}[thm]}
\end{tabproofwide}

\FormulaThmDeltaK[Argument (M14): Invarianz der Produktanteile]{%
  \begin{aligned}[t]
    X,Y\in\PowNE{S}\vdash{}&
      R(\Phi(X),\Phi(Y))=R(X,Y),\\[-2pt]
    X,Y\in\PowNE{S}\vdash{}&
      \varepsilon(\Phi(X),\Phi(Y))=\varepsilon(X,Y),\\[-2pt]
    X,Y\in\PowNE{S}\vdash{}&
      \Phi(X)\MogDiamond_{\mathcal P}\Phi(Y)
        =R(X,Y)\cup\varepsilon(X,Y)
        =X\MogStar_{\mathcal P}Y
  \end{aligned}
}{MogiljanskajaArgumentM14ProductPartsInvariant}{%
  \DeltaRow{Mengen}{L\dsep S\dsep T\dsep X\dsep Y
    \dsep R(X,Y)\dsep\varepsilon(X,Y)}
  \DeltaRow{Funktionen}{\Phi\dsep\MogStar_{\mathcal P}
    \dsep\MogDiamond_{\mathcal P}}
}
\begin{proof}[Metasprachliche Fassung]
\leavevmode\par
Der Rechteckanteil \(R(X,Y)\) sieht nach
\FormulaRefAuto{MogiljanskajaProductPartsDef}[def] nur die beiden Schnitte
\(X\cap L\) und \(Y\cap L\). Diese Schnitte bleiben nach
\FormulaRefAuto{MogiljanskajaArgumentM13PhiInvariants}[thm] unter \(\Phi\)
unverändert; daher ist
\[
  R(\Phi(X),\Phi(Y))=R(X,Y).
\]
Der Nullanteil
\(\varepsilon(X,Y)\) hängt nur davon ab, ob beide Mengen vollständig in
\(L\) liegen. Auch diese Eigenschaft bewahrt \(\Phi\), also gilt ebenso
\[
  \varepsilon(\Phi(X),\Phi(Y))=\varepsilon(X,Y).
\]

Nach \FormulaRefAuto{MogiljanskajaArgumentM12PhiBijective}[thm] sind
\(\Phi(X)\) und \(\Phi(Y)\) nichtleere Teilmengen von \(T\). Wendet man
die direkte Produktformel aus
\FormulaRefAuto{MogiljanskajaArgumentM6ProductClassification}[thm] einmal
in \(T\) und einmal in \(S\) an, so erhält man die gemeinsame Beschreibung
\[
  \Phi(X)\MogDiamond_{\mathcal P}\Phi(Y)
  =R(X,Y)\cup\varepsilon(X,Y)
  =X\MogStar_{\mathcal P}Y.
\]
\end{proof}

\MogTableProofHeading
\begin{tabproofwide}
  \proofstepwidestarbreak[]{X,Y\in\PowNE{S}\vdash
    \Phi(X)\cap L=X\cap L\dsep
    \Phi(Y)\cap L=Y\cap L}{%
    \FormulaRefAuto{MogiljanskajaArgumentM13PhiInvariants}[thm]}
  \proofstepwidestarbreak[1]{X,Y\in\PowNE{S}\vdash
    R(\Phi(X),\Phi(Y))=R(X,Y)}{%
    \FormulaRefAuto{MogiljanskajaProductPartsDef}[def]}
  \proofstepwidestarbreak[]{%
    \begin{aligned}[t]
      X,Y\in\PowNE{S}&\vdash{}\\[-2pt]
      X\subseteq L&\leftrightarrow\Phi(X)\subseteq L,\\[-2pt]
      Y\subseteq L&\leftrightarrow\Phi(Y)\subseteq L
    \end{aligned}}{%
    \FormulaRefAuto{MogiljanskajaArgumentM13PhiInvariants}[thm]}
  \proofstepwidestarbreak[3]{X,Y\in\PowNE{S}\vdash
    \varepsilon(\Phi(X),\Phi(Y))=\varepsilon(X,Y)}{%
    \FormulaRefAuto{MogiljanskajaProductPartsDef}[def]}
  \proofstepwidestarbreak[]{X,Y\in\PowNE{S}\vdash
    \Phi(X),\Phi(Y)\in\PowNE{T}}{%
    \FormulaRefAuto{MogiljanskajaArgumentM12PhiBijective}[thm]}
  \proofstepwidestarbreak[2,4,5]{X,Y\in\PowNE{S}\vdash
    \Phi(X)\MogDiamond_{\mathcal P}\Phi(Y)
      =R(X,Y)\cup\varepsilon(X,Y)}{%
    \FormulaRefAuto{MogiljanskajaArgumentM6ProductClassification}[thm]}
  \proofstepwidestarbreak[]{X,Y\in\PowNE{S}\vdash
    X\MogStar_{\mathcal P}Y
      =R(X,Y)\cup\varepsilon(X,Y)}{%
    \FormulaRefAuto{MogiljanskajaArgumentM6ProductClassification}[thm]}
  \proofstepwidestarbreak[6,7]{X,Y\in\PowNE{S}\vdash
    \Phi(X)\MogDiamond_{\mathcal P}\Phi(Y)
      =R(X,Y)\cup\varepsilon(X,Y)
      =X\MogStar_{\mathcal P}Y}{%
    \text{\normalfont Gleichheitskette}}
\end{tabproofwide}

\FormulaThmDeltaK[Argument (M15): Fixierung jedes ersten Mengenprodukts]{%
  X,Y\in\PowNE{S}\vdash
  \Phi(X\MogStar_{\mathcal P}Y)=X\MogStar_{\mathcal P}Y%
}{MogiljanskajaArgumentM15ProductFixed}{%
  \DeltaRow{Mengen}{A_*\dsep D\dsep S\dsep X\dsep Y\dsep Z
    \dsep R(X,Y)\dsep\varepsilon(X,Y)\dsep\mathcal C}
  \DeltaRow{Funktionen}{\eta\dsep\Phi\dsep\MogStar_{\mathcal P}}
}
\begin{proof}[Metasprachliche Fassung]
Seien \(X,Y\in\PowNE{S}\) und \(Z=X\MogStar_{\mathcal P}Y\). Da
\((S,\MogStar)\) nach
\FormulaRefAuto{MogiljanskajaArgumentM2Semigroups}[thm] eine Halbgruppe ist,
ist \(Z\) nach \FormulaRefAuto{PowerSemigroupProductClosure}[thm] wieder eine
nichtleere Teilmenge von \(S\). Die Produktklassifikation
\FormulaRefAuto{MogiljanskajaArgumentM6ProductClassification}[thm] liefert
\[
  Z=R(X,Y)\cup\varepsilon(X,Y).
\]
Hier liegt \(R(X,Y)\) in \(D\), während
\(\varepsilon(X,Y)\subseteq\{\mathbf0\}\subseteq A_*\) gilt. Wegen der
disjunkten Schichten aus
\FormulaRefAuto{MogiljanskajaArgumentM12PhiBijective}[thm] sind also
\[
  Z\cap D=R(X,Y),
  \qquad
  Z\cap A_*=\varepsilon(X,Y).
\]

Anschaulich ist nun entscheidend, dass ein Produktrechteck nie in den für
die Verschiebung reservierten Bereich fällt. Wegen \(S\subseteq T\) liefert
\FormulaRefAuto{MogiljanskajaArgumentM11RectangleProtected}[thm]
\[
  R(X,Y)\notin\mathcal C.
\]
Deshalb greift der Identitätszweig aus
\FormulaRefAuto{MogiljanskajaArgumentM10EtaFixedPart}[thm], also
\[
  \eta(R(X,Y))=R(X,Y).
\]
Setzt man dies in die Grundgleichung aus
\FormulaRefAuto{MogiljanskajaArgumentM13PhiInvariants}[thm] ein, erhält man
\[
  \Phi(Z)
  =(Z\cap A_*)\cup\eta(Z\cap D)
  =\varepsilon(X,Y)\cup R(X,Y)
  =Z.
\]
Damit ist
\(\Phi(X\MogStar_{\mathcal P}Y)=X\MogStar_{\mathcal P}Y\).
\end{proof}

\MogTableProofHeading
\begin{tabproofwide}
  \proofstepwidestar[]{X,Y\in\PowNE{S}\dsep
    Z=X\MogStar_{\mathcal P}Y}{\rA}
  \proofstepwidestar[1]{X,Y\in\powerset(T)}{%
    \FormulaRefAuto{MogiljanskajaPairDef}[def]}
  \proofstepwidestar[1]{Z\in\PowNE{S}}{%
    \FormulaRefAuto{MogiljanskajaArgumentM2Semigroups}[thm],
    \FormulaRefAuto{NonemptyPowerSemigroup}[thm]}
  \proofstepwidestar[1]{Z=R(X,Y)\cup\varepsilon(X,Y)}{%
    \FormulaRefAuto{MogiljanskajaArgumentM6ProductClassification}[thm]{1}}
  \proofstepwidestar[]{R(X,Y)\subseteq D\dsep
    \varepsilon(X,Y)\subseteq\{\mathbf0\}\subseteq A_*}{%
    \FormulaRefAuto{MogiljanskajaProductPartsDef}[def],
    \FormulaRefAuto{MogiljanskajaPhiDef}[def]}
  \proofstepwidestar[4,5]{Z\cap D=R(X,Y)\dsep
    Z\cap A_*=\varepsilon(X,Y)}{%
    \FormulaRefAuto{MogiljanskajaArgumentM12PhiBijective}[thm]}
  \proofstepwidestar[2]{R(X,Y)\notin\mathcal C}{%
    \FormulaRefAuto{MogiljanskajaArgumentM11RectangleProtected}[thm]}
  \proofstepwidestar[5,7]{\eta(R(X,Y))=R(X,Y)}{%
    \FormulaRefAuto{MogiljanskajaArgumentM10EtaFixedPart}[thm]}
  \proofstepwidestar[3,6]{\Phi(Z)
    =(Z\cap A_*)\cup\eta(Z\cap D)}{%
    \FormulaRefAuto{MogiljanskajaArgumentM13PhiInvariants}[thm]}
  \proofstepwidestar[4,6,8,9]{\Phi(Z)
    =\varepsilon(X,Y)\cup R(X,Y)=Z}{%
    \text{\normalfont Einsetzen}}
  \proofstepwidestar[1,10]{\Phi(X\MogStar_{\mathcal P}Y)
    =X\MogStar_{\mathcal P}Y}{%
    Definition von \(Z\)}
\end{tabproofwide}

\FormulaThmDeltaK[Argument (M16): Direkter Homomorphie- und Isomorphieschluss]{%
  \begin{aligned}[t]
    &\SemigroupHom{\Phi}
      {\PowNE{S},\MogStar_{\mathcal P}}
      {\PowNE{T},\MogDiamond_{\mathcal P}}\land{}\\[-2pt]
    &\SemigroupIso{\Phi}
      {\PowNE{S},\MogStar_{\mathcal P}}
      {\PowNE{T},\MogDiamond_{\mathcal P}}
  \end{aligned}
}{MogiljanskajaArgumentM16PowerIsomorphism}{%
  \DeltaRow{Mengen}{S\dsep T\dsep X\dsep Y}
  \DeltaRow{Funktionen}{\Phi\dsep\MogStar_{\mathcal P}
    \dsep\MogDiamond_{\mathcal P}}
}
\begin{proof}[Metasprachliche Fassung]
Die beiden vorangegangenen Argumente treffen sich in derselben Menge. Nach
\FormulaRefAuto{MogiljanskajaArgumentM15ProductFixed}[thm] gilt
\[
  \Phi(X\MogStar_{\mathcal P}Y)=X\MogStar_{\mathcal P}Y,
\]
während \FormulaRefAuto{MogiljanskajaArgumentM14ProductPartsInvariant}[thm]
\[
  \Phi(X)\MogDiamond_{\mathcal P}\Phi(Y)
  =X\MogStar_{\mathcal P}Y
\]
liefert. Folglich bewahrt \(\Phi\) die Multiplikation.

Nach \FormulaRefAuto{MogiljanskajaArgumentM2Semigroups}[thm] und
\FormulaRefAuto{NonemptyPowerSemigroup}[thm] sind die beiden
Potenzhalbgruppen tatsächlich Halbgruppen. Zusammen mit der Bijektivität von
\(\Phi\) aus \FormulaRefAuto{MogiljanskajaArgumentM12PhiBijective}[thm]
zeigt die erhaltene Produktgleichung gemäß
\FormulaRefAuto{SemigroupHomDef}[def] zunächst die Homomorphie und gemäß
\FormulaRefAuto{SemigroupIsoDef}[def] anschließend die behauptete
Isomorphie.
\end{proof}

\MogTableProofHeading
\begin{tabproofwide}
  \proofstepwidestar[]{X,Y\in\PowNE{S}\vdash
    \Phi(X\MogStar_{\mathcal P}Y)
      =X\MogStar_{\mathcal P}Y}{%
    \FormulaRefAuto{MogiljanskajaArgumentM15ProductFixed}[thm]}
  \proofstepwidestar[]{X,Y\in\PowNE{S}\vdash
    \Phi(X)\MogDiamond_{\mathcal P}\Phi(Y)
      =X\MogStar_{\mathcal P}Y}{%
    \FormulaRefAuto{MogiljanskajaArgumentM14ProductPartsInvariant}[thm]}
  \proofstepwidestarbreak[1,2]{X,Y\in\PowNE{S}\vdash
    \Phi(X\MogStar_{\mathcal P}Y)
      =\Phi(X)\MogDiamond_{\mathcal P}\Phi(Y)}{%
    \text{\normalfont gemeinsame rechte Seite}}
  \proofstepwidestar[]{%
    \SemigroupAlg{\PowNE{S}}{\MogStar_{\mathcal P}}\land
    \SemigroupAlg{\PowNE{T}}{\MogDiamond_{\mathcal P}}}{%
    \FormulaRefAuto{MogiljanskajaArgumentM2Semigroups}[thm],
    \FormulaRefAuto{NonemptyPowerSemigroup}[thm]}
  \proofstepwidestar[]{\Phi\colon\PowNE{S}\bij\PowNE{T}}{%
    \FormulaRefAuto{MogiljanskajaArgumentM12PhiBijective}[thm]}
  \proofstepwidestar[3,4,5]{\SemigroupHom{\Phi}
    {\PowNE{S},\MogStar_{\mathcal P}}
    {\PowNE{T},\MogDiamond_{\mathcal P}}}{%
    \FormulaRefAuto{SemigroupHomDef}[def]}
  \proofstepwidestar[5,6]{\SemigroupIso{\Phi}
    {\PowNE{S},\MogStar_{\mathcal P}}
    {\PowNE{T},\MogDiamond_{\mathcal P}}}{%
    \FormulaRefAuto{SemigroupIsoDef}[def]}
\end{tabproofwide}

\FormulaThmDeltaK[Argument (M17): Expliziter Schlusszeuge]{%
  \begin{aligned}[t]
    &\SemigroupAlg{S}{\MogStar}\land
      \SemigroupAlg{T}{\MogDiamond}\land{}\\[-2pt]
    &\neg\Finite{S}\land\neg\Finite{T}\land{}\\[-2pt]
    &\neg\exists f\,
      \SemigroupIso{f}{S,\MogStar}{T,\MogDiamond}\land{}\\[-2pt]
    &\SemigroupIso{\Phi}
      {\PowNE{S},\MogStar_{\mathcal P}}
      {\PowNE{T},\MogDiamond_{\mathcal P}}
  \end{aligned}
}{MogiljanskajaArgumentM17CounterexampleWitness}{%
  \DeltaRow{Mengen}{S\dsep T}
  \DeltaRow{Funktionen}{f\dsep\Phi\dsep\MogStar\dsep\MogDiamond
    \dsep\MogStar_{\mathcal P}\dsep\MogDiamond_{\mathcal P}}
}
\begin{proof}[Metasprachliche Fassung]
Nach \FormulaRefAuto{MogiljanskajaArgumentM2Semigroups}[thm] sind beide
Grundstrukturen Halbgruppen. Dass ihre Träger unendlich sind, folgt aus
\FormulaRefAuto{MogiljanskajaArgumentM3Infinite}[thm]. Ihre Nichtisomorphie
ist \FormulaRefAuto{MogiljanskajaArgumentM5NotIsomorphic}[thm]. Auf der
Ebene der nichtleeren Teilmengen liefert dagegen
\FormulaRefAuto{MogiljanskajaArgumentM16PowerIsomorphism}[thm] mit der
explizit konstruierten Abbildung \(\Phi\) einen Isomorphismus. Damit sind
alle vier Aussagen des Satzes erfüllt.
\end{proof}

\MogTableProofHeading
\begin{tabproofwide}
  \proofstepwidestar[]{\SemigroupAlg{S}{\MogStar}\land
    \SemigroupAlg{T}{\MogDiamond}}{%
    \FormulaRefAuto{MogiljanskajaArgumentM2Semigroups}[thm]}
  \proofstepwidestar[]{\neg\Finite{S}\land\neg\Finite{T}}{%
    \FormulaRefAuto{MogiljanskajaArgumentM3Infinite}[thm]}
  \proofstepwidestar[]{\neg\exists f\,
    \SemigroupIso{f}{S,\MogStar}{T,\MogDiamond}}{%
    \FormulaRefAuto{MogiljanskajaArgumentM5NotIsomorphic}[thm]}
  \proofstepwidestar[]{\SemigroupIso{\Phi}
    {\PowNE{S},\MogStar_{\mathcal P}}
    {\PowNE{T},\MogDiamond_{\mathcal P}}}{%
    \FormulaRefAuto{MogiljanskajaArgumentM16PowerIsomorphism}[thm]}
  \proofstepwidestar[1,2,3,4]{%
    \begin{aligned}[t]
      &\SemigroupAlg{S}{\MogStar}\land
        \SemigroupAlg{T}{\MogDiamond}\land{}\\[-2pt]
      &\neg\Finite{S}\land\neg\Finite{T}\land{}\\[-2pt]
      &\neg\exists f\,
        \SemigroupIso{f}{S,\MogStar}{T,\MogDiamond}\land{}\\[-2pt]
      &\SemigroupIso{\Phi}
        {\PowNE{S},\MogStar_{\mathcal P}}
        {\PowNE{T},\MogDiamond_{\mathcal P}}
    \end{aligned}}{%
    \text{\normalfont Bündelung}}
\end{tabproofwide}

\end{document}
